% !TeX program = pdflatex
% Tom II: Procesy stochastyczne; rachunek Ito w tomie III.
% Tom II: rozdzialy 1--14. Audyt calosciowy matematyczny, dydaktyczny i jezykowy.
\documentclass[11pt,a4paper,oneside]{report}
\usepackage[T1]{fontenc}
\usepackage[utf8]{inputenc}
\usepackage[polish]{babel}
\usepackage{lmodern}
\usepackage{amsmath,amssymb,amsthm,mathtools}
\usepackage{mathrsfs}
\usepackage{microtype}
\usepackage{geometry}
\usepackage{enumitem}
\usepackage{xcolor}
\usepackage{tikz}
\usetikzlibrary{arrows.meta,positioning,calc,shapes.geometric,decorations.pathmorphing}
\usepackage[unicode]{hyperref}
% Odsylacze miedzytomowe sa opcjonalne: brak .aux nie blokuje budowy tomu.
% Po zbudowaniu innych tomow ponowna kompilacja aktywuje odnosniki xr-hyper.
\usepackage{xr-hyper}
\IfFileExists{probability_volume_i.aux}{%
  \externaldocument[I-]{probability_volume_i}[probability_volume_i.pdf]%
}{}
\IfFileExists{probability_volume_iii.aux}{%
  \externaldocument[III-]{probability_volume_iii}[probability_volume_iii.pdf]%
}{}
\IfFileExists{probability_volume_iv.aux}{%
  \externaldocument[IV-]{probability_volume_iv}[probability_volume_iv.pdf]%
}{}

\geometry{margin=2.65cm}
\hypersetup{colorlinks=true,linkcolor=blue!45!black,urlcolor=blue!45!black}
\setlength{\parindent}{0pt}
\setlength{\parskip}{4pt}
\setlength{\emergencystretch}{2em}
\setlist[enumerate]{topsep=3pt,itemsep=2pt}
\setlist[itemize]{topsep=3pt,itemsep=2pt}
\numberwithin{equation}{chapter}
\theoremstyle{plain}
\newtheorem{twierdzenie}{Twierdzenie}[chapter]
\newtheorem{lemat}[twierdzenie]{Lemat}
\newtheorem{wniosek}[twierdzenie]{Wniosek}
\newtheorem{stwierdzenie}[twierdzenie]{Stwierdzenie}
\theoremstyle{definition}
\newtheorem{definicja}{Definicja}[chapter]
\newtheorem{przyklad}{Przykład}[chapter]
\newtheorem{zadanie}{Zadanie rozwiązane}[chapter]
\theoremstyle{remark}
\newtheorem{uwaga}{Uwaga}[chapter]
\renewcommand{\proofname}{Dowód}
\newcommand{\R}{\mathbb{R}}
\newcommand{\N}{\mathbb{N}}
\newcommand{\Z}{\mathbb{Z}}
\newcommand{\Q}{\mathbb{Q}}
\newcommand{\Prob}{\mathbb{P}}
\newcommand{\E}{\mathbb{E}}
\newcommand{\Var}{\operatorname{Var}}
\newcommand{\Cov}{\operatorname{Cov}}
\newcommand{\Law}{\operatorname{Law}}
\newcommand{\calF}{\mathcal{F}}
\newcommand{\calE}{\mathcal{E}}
\newcommand{\calH}{\mathcal{H}}
\newcommand{\calN}{\mathcal{N}}
\newcommand{\calB}{\mathcal{B}}
\newcommand{\calP}{\mathcal{P}}
\newcommand{\calG}{\mathcal{G}}
\newcommand{\calD}{\mathcal{D}}
\newcommand{\calC}{\mathcal{C}}
\newcommand{\ind}[1]{\mathbf{1}_{#1}}
\newcommand{\dd}{\mathrm{d}}
\newcommand{\norm}[1]{\lVert #1\rVert}
\newcommand{\abs}[1]{\lvert #1\rvert}
\newcommand{\iid}{\mathrel{\overset{\mathrm{iid}}{\sim}}}
\newcommand{\Normal}{\mathcal{N}}
\DeclareMathOperator{\dist}{dist}
\DeclareMathOperator{\supp}{supp}
\DeclareMathOperator{\spanop}{span}

% BEGIN GLOBAL CODE LISTINGS (generated from latex/pl/code-listings.tex)
\usepackage{xcolor}
\usepackage{listings}
\definecolor{codebg}{RGB}{248,249,251}
\definecolor{codeborder}{RGB}{180,191,203}
\definecolor{codecomment}{RGB}{61,113,79}
\definecolor{codekeyword}{RGB}{41,78,132}
\lstdefinelanguage{Julia}{
  morekeywords={abstract,baremodule,begin,break,catch,const,continue,do,else,elseif,
    end,export,false,finally,for,function,global,if,import,in,isa,let,local,macro,
    module,mutable,primitive,quote,return,struct,true,try,type,using,where,while,
    missing,nothing,println,print},
  sensitive=true,morecomment=[l]{\#},morecomment=[n]{\#=}{=\#},
  morestring=[b]",
}
\lstdefinestyle{skrypt}{
  backgroundcolor=\color{codebg},frame=single,rulecolor=\color{codeborder},
  framesep=5pt,frameround=tttt,framerule=.5pt,
  basicstyle=\ttfamily\footnotesize,keywordstyle=\color{codekeyword}\bfseries,
  commentstyle=\color{codecomment}\itshape,stringstyle=\color{purple!65!black},
  numbers=left,numberstyle=\tiny\color{gray!75!black},numbersep=8pt,
  stepnumber=1,breaklines=true,breakatwhitespace=false,
  columns=fullflexible,keepspaces=true,showstringspaces=false,tabsize=4,
  linewidth=.90\linewidth,captionpos=b,aboveskip=11pt,belowskip=12pt,
  xleftmargin=12pt,xrightmargin=2pt,
}
\lstdefinestyle{pythonpde}{style=skrypt,language=Python}
\lstdefinestyle{python}{style=skrypt,language=Python}
\lstdefinestyle{julia}{style=skrypt,language=Julia}
\lstset{style=skrypt}
\renewcommand{\lstlistingname}{Kod}
\renewcommand{\lstlistlistingname}{Spis kodów}
% END GLOBAL CODE LISTINGS
\begin{document}
\hypersetup{pageanchor=false}
\begin{titlepage}
\centering
\vspace*{3.4cm}
{\Large Podstawy teorii prawdopodobieństwa i statystyki\par}
\vspace{1.2cm}
{\Huge\bfseries Tom II\\[0.2cm] Procesy stochastyczne\par}
\vspace{0.8cm}
{\LARGE Od procesów Markowa do martyngałów i wariacji kwadratowej\par}
\vspace{0.6cm}
{\Large Rozdziały 1--14\par}
\vfill
{\large Wersja po całościowym audycie matematycznym i językowym\par}
\vspace{1.2cm}
\end{titlepage}
\hypersetup{pageanchor=true}
\tableofcontents
\chapter*{Wprowadzenie do tomu II}
\addcontentsline{toc}{chapter}{Wprowadzenie do tomu II}
Tom I dotyczył zmiennych losowych, ich rozkładów, zbieżności, warunkowej
wartości oczekiwanej, łańcuchów Markowa i martyngałów w czasie dyskretnym.
Tom II bada całe \emph{rodziny} zmiennych losowych, indeksowane czasem,
w szczególności czasem ciągłym. Jest to zmiana perspektywy: oprócz rozkładu
w danej chwili interesuje nas wspólna zależność wielu chwil, a także
regularność całej trajektorii.

W pierwszej części odpowiemy na trzy odrębne pytania. (1) Czym dokładnie
jest proces i co znaczy znać jego rozkład? (2) Czy zadanie zgodnych
rozkładów skończeniewymiarowych wyznacza rzeczywisty proces, a jeśli tak,
kiedy ma on ciągłe trajektorie? (3) Jak formalnie opisać narastającą
w czasie informację i podejmować decyzje bez znajomości przyszłości?
Rozdziały zostały ułożone tak, by konstrukcja procesu Wienera,
miary Wienera i regularności trajektorii opierała się na ścisłych
fundamentach probabilistycznych. Całka stochastyczna Itô,
równania stochastyczne i systematyczna teoria całek po trajektoriach
będą rozwinięte w tomie III, poświęconym analizie stochastycznej.
Tom IV obejmie statystykę matematyczną i wnioskowanie bayesowskie.

Zakładamy znajomość teorii miary i całki Lebesgue'a, twierdzenia
Radona--Nikodyma, warunkowej wartości oczekiwanej, lematu
Borela--Cantellego oraz podstaw zbieżności zmiennych losowych z tomu I.
Jako znane wyniki z teorii miary traktujemy twierdzenie Carathéodory'ego
o rozszerzaniu miar, regularność borelowskich miar probabilistycznych
na przestrzeniach polskich oraz twierdzenie o zwartości iloczynu
przeliczalnej rodziny niepustych przestrzeni zwartych.
\emph{Nie} zakładamy twierdzenia Kołmogorowa o istnieniu procesu ani jego
kryterium ciągłości: oba zostaną udowodnione.

Równości i nierówności między zmiennymi losowymi rozumiemy
jako zachodzące prawie na pewno, o ile nie zaznaczono inaczej.
W rozdziałach 1--3 podkreślimy, że pojęcie modyfikacji nie wymaga
równości procesów poza jednym zbiorem zerowym dla wszystkich chwil.
To odróżnia je od nierozróżnialności.


\chapter{Procesy stochastyczne: od rozkładów do trajektorii}
\label{ch:procesy}
\paragraph{Odsyłacze międzytomowe.} Przypomnienie konstrukcji probabilistycznych znajduje się w tomie I, rozdziale~\ref{I-ch:probability}; statystyczny problem danych zależnych omawia tom IV, rozdział~\ref{IV-ch:time-series}.


\section{Dlaczego pojedyncza zmienna losowa już nie wystarcza?}
Rozkład zmiennej $X$ odpowiada na pytanie, jakie wartości może przyjmować
jeden pomiar. Jeżeli mierzymy temperaturę co godzinę, pozycję cząstki
przez kilka sekund albo liczbę klientów w sklepie w ciągu dnia, to
kolejne pomiary są ze sobą powiązane. Sama znajomość rozkładów
$X_0,X_1,\ldots$ nie wystarcza do odpowiedzi na pytania o tę zależność.

Weźmy dwie zmienne o jednakowym rozkładzie Bernoulliego z parametrem
$1/2$. Mogą być niezależne, a mogą być wręcz identyczne: $X_1=X_0$.
W obu przypadkach $\Prob(X_0=1)=\Prob(X_1=1)=1/2$, lecz
$\Prob(X_0=X_1)$ wynosi odpowiednio $1/2$ albo $1$.
Zrozumienie procesu wymaga więc danych o \emph{wspólnych} rozkładach
pomiarów, a potem także informacji o zachowaniu trajektorii.

\begin{definicja}[Proces stochastyczny]
Niech $(\Omega,\calF,\Prob)$ będzie przestrzenią probabilistyczną,
$T\ne\varnothing$ dowolnym zbiorem indeksów, a $(E,\calE)$
przestrzenią mierzalną. \emph{Procesem stochastycznym o przestrzeni
stanów $E$ i zbiorze czasów $T$} nazywamy rodzinę
$X=(X_t)_{t\in T}$, w której każde
$X_t:(\Omega,\calF)\to(E,\calE)$ jest odwzorowaniem mierzalnym.
Dla ustalonego $t$ zmienna $X_t$ opisuje stan w chwili $t$;
dla ustalonego $\omega\in\Omega$ funkcja
$X(\omega):T\to E$, $t\mapsto X_t(\omega)$, jest
\emph{trajektorią} lub \emph{realizacją} procesu.
\end{definicja}

Typowe przypadki to $T=\N_0:=\{0,1,2,\ldots\}$ (czas dyskretny)
oraz $T=[0,\infty)$ (czas ciągły). Przestrzenią stanów może być
$\R^d$, zbiór skończony, przestrzeń funkcji czy przestrzeń miar.
Ważne, że sam proces jest rodziną funkcji mierzalnych; nie
wymagamy jeszcze mierzalności \emph{wspólnej} odwzorowania
$(t,\omega)\mapsto X_t(\omega)$, ani żadnej ciągłości trajektorii.

\begin{center}
\begin{tikzpicture}[>=Latex,scale=0.95]
\node[draw,rounded corners,minimum width=2.9cm,minimum height=2.5cm,fill=blue!3] (a) at (0,0) {};
\node[above=3pt] at (a.north) {losowy wynik $\omega$};
\node at (-0.8,0.35) {$\omega_1$};
\node at (0.6,0.35) {$\omega_2$};
\node at (-0.3,-0.55) {$\omega_3$};
\node[draw,rounded corners,minimum width=5.8cm,minimum height=2.5cm,fill=green!3] (b) at (7.2,0) {};
\node[above=3pt] at (b.north) {trajektoria $t\mapsto X_t(\omega)$};
\draw[->] (4.8,-0.8)--(9.6,-0.8) node[right] {$t$};
\draw[->] (4.85,-0.85)--(4.85,0.8) node[above] {$x$};
\draw[thick,blue!70!black] (5,-0.35)--(5.7,0.1)--(6.5,-0.15)--(7.25,0.58)--(8.05,0.25)--(8.8,0.65)--(9.3,0.3);
\draw[->,thick] (1.7,0.3) -- (4.25,0.3) node[midway,above] {$\omega\mapsto X(\omega)$};
\end{tikzpicture}
\par\smallskip\small Schemat: jedna realizacja losowości wyznacza całą funkcję czasu, nie tylko jeden pomiar.
\end{center}

\begin{przyklad}[Losowa prosta]
Niech $Z$ będzie rzeczywistą zmienną losową i $X_t=tZ$ dla $t\ge0$.
W chwili $t$ rozkład $X_t$ jest obrazem rozkładu $Z$ przez
odwzorowanie $z\mapsto tz$. Każda trajektoria jest prostą o losowym
nachyleniu. Zmienność po czasie jest całkowicie zdeterminowana przez
jedną liczbę $Z(\omega)$: znajomość $X_1$ wyznacza wszystkie $X_t$.
Jeśli $\E Z=0$ i $\Var(Z)=\sigma^2<\infty$, to
$\Cov(X_s,X_t)=st\sigma^2$. Przyrosty
$X_t-X_s=(t-s)Z$ są zwykle silnie zależne; nie tworzą procesu o
niezależnych przyrostach.
\end{przyklad}

\begin{przyklad}[Dwa procesy o tych samych rozkładach pojedynczych]
Niech $\xi_0,\xi_1,\ldots$ będą niezależnymi zmiennymi Bernoulliego
z parametrem $1/2$, a $Z$ kolejną taką zmienną. Definiujemy
$X_n=\xi_n$ i $Y_n=Z$ dla każdego $n\in\N_0$.
Wówczas $\Law(X_n)=\Law(Y_n)$ dla każdego $n$, ale
$\Prob(X_0=X_1)=1/2$, podczas gdy $\Prob(Y_0=Y_1)=1$.
Dla $n\ge1$ mamy również
$\Prob(X_0=\cdots=X_n)=2^{-n}$,
a $\Prob(Y_0=\cdots=Y_n)=1$.
To pierwszy powód, aby badać rozkłady \emph{skończeniewymiarowe}.
\end{przyklad}

\section{Rozkłady skończeniewymiarowe i prawo procesu}
\begin{definicja}[Rozkład skończeniewymiarowy]
Niech $X=(X_t)_{t\in T}$ będzie procesem o stanach w $(E,\calE)$.
Dla każdego $m\ge1$ oraz $t_1,\ldots,t_m\in T$ definiujemy
rozkład wektora $(X_{t_1},\ldots,X_{t_m})$ na
$(E^m,\calE^{\otimes m})$ przez
\[
\mu_{t_1,\ldots,t_m}(A)
:=\Prob\big((X_{t_1},\ldots,X_{t_m})\in A\big),
\qquad A\in\calE^{\otimes m}.
\]
Rodzinę wszystkich miar $\mu_{t_1,\ldots,t_m}$ nazywamy rodziną
\emph{rozkładów skończeniewymiarowych} procesu.
\end{definicja}

Na przykład dla $A_0,\ldots,A_n\in\calE$ mamy
\[
\mu_{t_0,\ldots,t_n}(A_0\times\cdots\times A_n)
=\Prob(X_{t_0}\in A_0,\ldots,X_{t_n}\in A_n).
\]
Przecinki wewnątrz $\Prob(\ldots)$ oznaczają \emph{przecięcie zdarzeń}.
W szczególności $\Prob(X_s\in A,X_t\in B)$ nie jest na ogół
iloczynem $\Prob(X_s\in A)\Prob(X_t\in B)$; taki iloczyn
wymaga niezależności zmiennych $X_s$ i $X_t$.

\begin{twierdzenie}[Zgodność rzeczywistych rozkładów procesu]
Niech $X=(X_t)_{t\in T}$ będzie procesem o stanach w $(E,\calE)$.
Wówczas dla dowolnych $m\ge1$, $t_1,\ldots,t_m\in T$,
permutacji $\pi$ zbioru $\{1,\ldots,m\}$ oraz każdego
$A\in\calE^{\otimes m}$ zachodzi
\[
\mu_{t_{\pi(1)},\ldots,t_{\pi(m)}}(A)
=\mu_{t_1,\ldots,t_m}(\Pi_\pi^{-1}(A)),
\quad
\Pi_\pi(x_1,\ldots,x_m):=(x_{\pi(1)},\ldots,x_{\pi(m)}).
\]
Ponadto dla $k<m$ i $B\in\calE^{\otimes k}$ zachodzi
$\mu_{t_1,\ldots,t_m}(B\times E^{m-k})=\mu_{t_1,\ldots,t_k}(B)$.
Jeśli $t_i=t_j$, rozkład $\mu_{t_1,\ldots,t_m}$ jest skupiony na
zbiorze $\{(x_1,\ldots,x_m):x_i=x_j\}$, o ile diagonalny zbiór jest mierzalny
(np. gdy $E$ jest przestrzenią polską z borelowskim $\sigma$-ciałem).
\end{twierdzenie}
\begin{proof}
Równość dotycząca permutacji wynika wprost z równości zmiennych
wektorowych
$(X_{t_{\pi(1)}},\ldots,X_{t_{\pi(m)}})
=\Pi_\pi(X_{t_1},\ldots,X_{t_m})$.
Własność marginalizacji otrzymujemy z tożsamości zdarzeń
$\{(X_{t_1},\ldots,X_{t_m})\in B\times E^{m-k}\}
=\{(X_{t_1},\ldots,X_{t_k})\in B\}$.
Jeżeli $t_i=t_j$, to $X_{t_i}(\omega)=X_{t_j}(\omega)$
dla każdego $\omega$, więc wektor należy do wskazanej diagonali
z prawdopodobieństwem $1$.
\end{proof}

W drugą stronę powstaje podstawowe pytanie: \emph{jeśli zadamy
rodzinę miar o tych własnościach, czy istnieje proces mający właśnie
te rozkłady?} Odpowiemy na nie w rozdziale~\ref{ch:konstrukcja}.

\begin{definicja}[Prawo procesu na przestrzeni iloczynowej]
Niech $(E,\calE)$ będzie przestrzenią mierzalną, a
$\calE^{\otimes T}$ najmniejszym $\sigma$-ciałem na $E^T$,
które czyni wszystkie rzuty $\pi_t:E^T\to E$, $\pi_t(x)=x(t)$,
mierzalnymi. Dla procesu $X=(X_t)_{t\in T}$ jego
\emph{prawem na przestrzeni iloczynowej} nazywamy
$\Law(X):=\Prob\circ\mathbf X^{-1}$, gdzie
$\mathbf X:\Omega\to E^T$ jest dane wzorem
$\mathbf X(\omega)=(X_t(\omega))_{t\in T}$.
\end{definicja}

\begin{stwierdzenie}[Mierzalność i jednoznaczność prawa]
Dla dowolnego procesu $X$ odwzorowanie
$\mathbf X:(\Omega,\calF)\to(E^T,\calE^{\otimes T})$ jest
mierzalne. Jeśli procesy $X,Y$ o tym samym zbiorze czasów $T$
i przestrzeni stanów $(E,\calE)$ mają jednakowe rozkłady
skończeniewymiarowe, to $\Law(X)=\Law(Y)$ na
$(E^T,\calE^{\otimes T})$.
\end{stwierdzenie}
\begin{proof}
Przeciwobraz zbioru $\pi_t^{-1}(A)$ przez $\mathbf X$ jest
$\{X_t\in A\}\in\calF$. Zbiory te generują $\calE^{\otimes T}$,
więc $\mathbf X$ jest mierzalne. Dla drugiej części rozważmy
klasę cylindrów $\pi_{t_1,\ldots,t_m}^{-1}(A)$,
$A\in\calE^{\otimes m}$. Tworzą one $\pi$-układ, gdyż przecięcie
dwóch takich cylindrów jest cylindrem zależnym od sumy ich
skończonych zbiorów indeksów. Obie miary praw procesów
są równe na cylindrach na mocy założenia o rozkładach
skończeniewymiarowych. Ponieważ cylindry generują
$\calE^{\otimes T}$, twierdzenie o jednoznaczności miar
na generującym $\pi$-układzie kończy dowód.
\end{proof}

\begin{uwaga}[Co jest, a co nie jest ustalone przez rozkłady?]
Rozkłady skończeniewymiarowe wyznaczają prawo procesu na
\emph{cylindrycznym} $\sigma$-ciele $\calE^{\otimes T}$.
Nie znaczy to, że dowolny zbiór trajektorii w $E^T$ jest mierzalny
ani że własność ``trajektoria jest ciągła'' da się badać
wyłącznie przez dowolną, nieprzeliczalną operację przecięcia
zdarzeń. Aby poprawnie mówić o ciągłych trajektoriach, będziemy
konstruować modyfikacje ciągłe i korzystać z mierzalnej przestrzeni
$C([0,T];\R^d)$.
\end{uwaga}


\begin{zadanie}[Prawdopodobieństwo dwóch pomiarów i maksimum trajektorii]
Niech $Z$ ma rozkład $\Normal(0,1)$, a $X_t=tZ$ dla
$t\in[0,2]$. Obliczyć $\Prob(X_1>0,X_2>1)$ oraz
$\Prob(\sup_{0\le t\le2}X_t>1)$ i wyjaśnić, czy w obu pytaniach
wystarczy rozkład pojedynczego pomiaru.

\emph{Rozwiązanie.} Ponieważ $X_1=Z$, $X_2=2Z$,
zdarzenie $\{X_1>0,X_2>1\}$ jest równe $\{Z>1/2\}$.
Oznaczając dystrybuantę standardowego rozkładu normalnego przez
$\Phi$, otrzymujemy $1-\Phi(1/2)$.
Dla ustalonego wyniku $\omega$ funkcja $t\mapsto tZ(\omega)$
rośnie, gdy $Z(\omega)\ge0$, a maleje, gdy $Z(\omega)<0$.
Wobec tego $\sup_{0\le t\le2}X_t=\max(0,2Z)$ i również
$\Prob(\sup_{t\le2}X_t>1)=1-\Phi(1/2)$.
Zgodność liczb jest szczególną cechą tego modelu, a nie ogólnym
prawem: w drugim pytaniu występuje funkcjonał \emph{całej}
trajektorii. Dla procesu o ciągłych trajektoriach maksimum można
też wyrazić jako supremum po przeliczalnym zbiorze chwil
wymiernych, co zapewnia mierzalność zdarzenia.
\end{zadanie}

\section{Modyfikacje i nierozróżnialność}
W czasie dyskretnym można zsumować przeliczalnie wiele zbiorów
zerowych. W czasie ciągłym jest nieprzeliczalnie wiele chwil;
nie wolno zatem automatycznie zamieniać zdań ``dla każdego $t$
prawie na pewno'' oraz ``prawie na pewno dla każdego $t$''.

\begin{definicja}[Modyfikacja i nierozróżnialność]
Niech $X,Y$ będą procesami na tej samej przestrzeni probabilistycznej,
o wspólnym zbiorze indeksów $T$ i przestrzeni stanów $E$.
Mówimy, że $Y$ jest \emph{modyfikacją} (wersją) $X$, jeśli
dla każdego $t\in T$ zdarzenie $\{X_t=Y_t\}$ jest mierzalne
i zachodzi $\Prob(X_t=Y_t)=1$.
Warunek mierzalności jest automatyczny, gdy $E$ jest przestrzenią
metryczną z borelowskim $\sigma$-ciałem: diagonala
$\{(x,y)\in E^2:x=y\}$ jest wtedy zbiorem domkniętym.
Mówimy, że $X$ i $Y$ są \emph{nierozróżnialne}, jeżeli istnieje
$N\in\calF$ z $\Prob(N)=0$ takie, że dla wszystkich
$\omega\notin N$ i wszystkich $t\in T$ zachodzi
$X_t(\omega)=Y_t(\omega)$.
\end{definicja}

\begin{stwierdzenie}[Czas przeliczalny]
Jeśli $T$ jest przeliczalny, a $X$ i $Y$ są swoimi
modyfikacjami, to są nierozróżnialne.
\end{stwierdzenie}
\begin{proof}
Dla każdego $t\in T$ zbiór
$N_t=\{\omega:X_t(\omega)\ne Y_t(\omega)\}$ jest mierzalny
i ma miarę zero (przyjmujemy mierzalną diagonalę w $E\times E$).
Wtedy $N=\bigcup_{t\in T}N_t$ jest przeliczalną sumą
zbiorów zerowych. Poza $N$ zachodzi $X_t=Y_t$ dla
wszystkich $t$ równocześnie.
\end{proof}

\begin{przyklad}[Modyfikacja, która nie jest nierozróżnialna]
Niech $\Omega=[0,1]$ z miarą Lebesgue'a i niech
$U(\omega)=\omega$. Dla $t\in[0,1]$ definiujemy
$X_t(\omega)=0$ oraz $Y_t(\omega)=\ind{\{U(\omega)=t\}}$.
Dla ustalonego $t$ zdarzenie $\{U=t\}$ jest pojedynczym
punktem i ma prawdopodobieństwo $0$, a zatem
$\Prob(X_t=Y_t)=1$. Natomiast dla \emph{każdego}
$\omega\in[0,1]$ w chwili $t=U(\omega)$ mamy
$Y_t(\omega)=1\ne X_t(\omega)$. Nie istnieje więc
nawet jeden wynik $\omega$, dla którego trajektorie byłyby
wszędzie równe. Trajektorie $Y$ mają losowy ``kolec''.
\end{przyklad}

\begin{twierdzenie}[Ciągłe modyfikacje są nierozróżnialne]
Niech $T=[0,T_0]$ dla $T_0>0$, a $E$ będzie przestrzenią
metryczną. Niech procesy $X,Y$ mają trajektorie ciągłe na $T$
poza ewentualnymi, osobnymi mierzalnymi zbiorami zerowymi.
Jeżeli $X$ i $Y$ są modyfikacjami, to są nierozróżnialne.
\end{twierdzenie}
\begin{proof}
Ustalmy przeliczalny zbiór gęsty
$D=(\Q\cap[0,T_0])\cup\{T_0\}$.
Dla każdego $t\in D$ mamy $\Prob(X_t=Y_t)=1$,
więc poza jednym mierzalnym zbiorem zerowym $N_1$
zachodzi $X_t(\omega)=Y_t(\omega)$ dla wszystkich $t\in D$.
Do $N_1$ dodajmy zbiory zerowe, na których którakolwiek z
trajektorii może nie być ciągła, otrzymując $N$.
Dla $\omega\notin N$ i dowolnego $s\in[0,T_0]$
wybierzmy $t_n\in D$ z $t_n\to s$.
Wówczas przez ciągłość obu trajektorii
$X_s(\omega)=\lim_n X_{t_n}(\omega)
=\lim_n Y_{t_n}(\omega)=Y_s(\omega)$.
Zbiór $N$ nie zależy od $s$, co jest kluczowym elementem dowodu.
\end{proof}

\section{Przestrzeń trajektorii ciągłych}
\begin{definicja}[{Przestrzeń $C([0,T];\R^d)$}]
Dla $T>0$ i $d\in\N$ oznaczamy przez
$C([0,T];\R^d)$ zbiór wszystkich funkcji ciągłych
$x:[0,T]\to\R^d$ i wyposażamy go w metrykę
$\norm{x-y}_\infty:=\sup_{t\in[0,T]}\norm{x(t)-y(t)}$.
Jego borelowskie $\sigma$-ciało oznaczamy przez
$\calB(C([0,T];\R^d))$.
\end{definicja}

\begin{twierdzenie}[Mierzalność współrzędnych i generowanie $\sigma$-ciała]
Niech $T>0$, $d\ge1$ oraz
$D=(\Q\cap[0,T])\cup\{T\}$.
Dla każdego $t\in[0,T]$ ewaluacja
$e_t:C([0,T];\R^d)\to\R^d$, $e_t(x)=x(t)$,
jest ciągła. Ponadto
\[
\calB(C([0,T];\R^d))=\sigma(e_t:t\in D).
\]
W konsekwencji, jeśli proces $X=(X_t)_{t\in[0,T]}$
ma ciągłe trajektorie dla każdego $\omega$, to odwzorowanie
$\omega\mapsto(t\mapsto X_t(\omega))$ z $\Omega$ do
$C([0,T];\R^d)$ jest mierzalne.
\end{twierdzenie}
\begin{proof}
Ponieważ $\norm{x(t)-y(t)}\le\norm{x-y}_\infty$,
ewaluacja jest $1$-Lipschitzowska, a zatem ciągła.
Stąd $\sigma(e_t:t\in D)\subseteq\calB(C)$.

Dla dowolnego $f\in C$ i $r>0$, z ciągłości funkcji
$t\mapsto\norm{x(t)-f(t)}$ i gęstości $D$ wynika
$\norm{x-f}_\infty=\sup_{t\in D}\norm{x(t)-f(t)}$.
Otwartą kulę można opisać następująco:
\[
B(f,r)=\bigcup_{k\ge1:\,1/k<r}
\bigcap_{t\in D}
\left\{x:\norm{x(t)-f(t)}\le r-\frac1k\right\}.
\]
Każdy zbiór po prawej stronie jest mierzalny względem
$\sigma(e_t:t\in D)$. Aby przejść od kul do wszystkich
zbiorów otwartych, użyjemy separowalności $C$.
Funkcje łamane, których węzły mają racjonalne chwile
(z dołączeniem $0,T$), a wartości w węzłach należą do
$\Q^d$, tworzą zbiór przeliczalny. Jest on gęsty w
$C$: jednostajna ciągłość pozwala przybliżyć funkcję
odpowiednio drobną interpolacją liniową, a wartości
w skończonej liczbie węzłów przybliżyć liczbami wymiernymi.
Wobec tego $C$ ma przeliczalną bazę otwartych kul,
więc każdy zbiór otwarty jest przeliczalną sumą takich kul.
Dowodzi to odwrotnej inkluzji $\calB(C)\subseteq\sigma(e_t:t\in D)$.
Ostatni wniosek wynika z mierzalności $e_t\circ X=X_t$
dla $t\in D$.
\end{proof}

\begin{uwaga}[Dlaczego wymagamy wspólnego zbioru pełnej miary?]
Jeśli trajektorie $X$ są ciągłe tylko prawie na pewno,
to aby traktować $X$ jako zmienną losową o wartościach w $C$,
wybieramy mierzalny zbiór $\Omega_0$ miary $1$, na którym
ciągłość zachodzi \emph{jednocześnie dla wszystkich $t$},
i poza nim przypisujemy dowolną ustaloną trajektorię ciągłą.
Tak zmieniony proces jest modyfikacją poprzedniego, a jego
odwzorowanie do $C$ jest mierzalne z poprzedniego twierdzenia.
Przypadek $C([0,\infty);\R^d)$ można potraktować przy pomocy
metryki jednostajnej zbieżności na przedziałach zwartych.
\end{uwaga}



\begin{uwaga}[Prawo na przestrzeni trajektorii ciągłych]
Jeśli proces $X$ ma trajektorie ciągłe, traktujemy go jako
zmienną losową o wartościach w $C([0,T];\R^d)$. Wówczas
rozkłady skończeniewymiarowe wyznaczają jego prawo
\emph{jako miarę na przestrzeni trajektorii ciągłych}:
równość miar na cylindrach określonych chwilami wymiernymi
przechodzi na całe borelowskie $\sigma$-ciało, ponieważ
te cylindry je generują. Jest to mocniejszy kontekst niż
kanoniczna przestrzeń $E^T$, gdzie sama ciągłość trajektorii
nie musi odpowiadać mierzalnemu zbiorowi cylindrycznemu.
\end{uwaga}

\section{Trajektorie ze skokami: przestrzeń funkcji c\`adl\`ag}
Między procesem o trajektoriach ciągłych a procesem
zupełnie nieregularnym istnieje ważna klasa pośrednia:
trajektorie mogą mieć skoki, ale nie powinny oscylować
w dowolny sposób bezpośrednio przed każdą chwilą.
Taki warunek będzie podstawowy dla procesów Poissona,
łańcuchów Markowa w czasie ciągłym i procesów Lévy'ego.

\begin{definicja}[Trajektoria c\`adl\`ag]
Niech $(E,d)$ będzie przestrzenią metryczną, a $T>0$.
Funkcja $f:[0,T]\to E$ jest \emph{c\`adl\`ag}, jeżeli:
\begin{enumerate}
\item jest prawostronnie ciągła w każdym $t\in[0,T)$,
tzn. $\lim_{s\downarrow t}f(s)=f(t)$;
\item w każdej chwili $t\in(0,T]$ istnieje
granica lewostronna
$f(t-):=\lim_{s\uparrow t}f(s)\in E$.
\end{enumerate}
Przez $D([0,T];E)$ oznaczamy zbiór wszystkich takich funkcji.
Gdy $E=\R^d$, wielkość $\Delta f(t)=f(t)-f(t-)$
(nieokreślona w $0$) nazywamy \emph{skokiem} w chwili $t$.
\end{definicja}

\begin{przyklad}[Losowy skok, ale uporządkowana trajektoria]
Dla $U\sim\operatorname{U}([0,1])$ proces
$X_t=\ind{\{U\le t\}}$ ma trajektorie c\`adl\`ag.
Przed chwilą $U$ przyjmuje wartość $0$,
a w chwili $U$ i później wartość $1$.
Stąd $X_U=1$, $X_{U-}=0$ oraz $\Delta X(U)=1$.
Jest to funkcja prawostronnie ciągła z jednym skokiem.
Jej trajektorie nie są ciągłe, chociaż -- jak wykażemy
w dalszym przykładzie -- proces jest stochastycznie ciągły.
To rozróżnienie motywuje używanie przestrzeni $D([0,T];E)$
zamiast $C([0,T];E)$ przy procesach skokowych.
\end{przyklad}

\begin{uwaga}[Topologia Skorochoda]
Przy zbieżności procesów skokowych zwykła metryka
jednostajna na $D([0,T];\R^d)$ bywa zbyt silna:
przesunięcie chwili pojedynczego skoku o $1/n$
może nie zmniejszać odległości jednostajnej.
Dlatego w bardziej zaawansowanej teorii wprowadza się
topologię Skorochoda $J_1$, dopuszczającą niewielkie
przesunięcia czasu. Na razie potrzebujemy tylko
samego pojęcia trajektorii c\`adl\`ag.
\end{uwaga}

\section{Przyrosty, stacjonarność i regularność w czasie}
\begin{definicja}[Przyrosty]
Niech $X=(X_t)_{t\ge0}$ będzie procesem o wartościach w $\R^d$.
Mówimy, że ma \emph{niezależne przyrosty}, jeżeli dla
każdych $m\ge1$ i $0\le t_0<t_1<\cdots<t_m$
zmienne $X_{t_1}-X_{t_0},\ldots,X_{t_m}-X_{t_{m-1}}$
są niezależne. Proces ma \emph{stacjonarne przyrosty}, jeśli
dla każdych $0\le s<t$ rozkład $X_t-X_s$ zależy tylko od $t-s$.
Mówimy, że $X$ jest \emph{stochastycznie ciągły w chwili $t$},
jeśli dla każdego $\varepsilon>0$ zachodzi
$\lim_{s\to t}\Prob(\norm{X_s-X_t}>\varepsilon)=0$
(granica jest brana w zbiorze indeksów).
\end{definicja}


\begin{twierdzenie}[Rozkład procesu wyznaczony przez niezależne przyrosty]
Niech $X=(X_t)_{t\ge0}$ będzie rzeczywistym procesem
o niezależnych przyrostach i niech $X_0=0$ prawie na pewno.
Dla $0=t_0<t_1<\cdots<t_m$ oznaczmy
$\Delta_j=X_{t_j}-X_{t_{j-1}}$ i
$\varphi_j(u)=\E e^{iu\Delta_j}$, $u\in\R$.
Wówczas dla dowolnych $u_1,\ldots,u_m\in\R$ zachodzi
\[
\E\exp\left(i\sum_{k=1}^m u_k X_{t_k}\right)
=\prod_{j=1}^m\varphi_j\left(\sum_{k=j}^m u_k\right).
\]
W szczególności rozkłady wszystkich przyrostów
$X_t-X_s$, $0\le s<t$, łącznie z niezależnością
przyrostów i warunkiem $X_0=0$, wyznaczają
wszystkie rozkłady skończeniewymiarowe procesu.
\end{twierdzenie}
\begin{proof}
Dla $k=1,\ldots,m$ mamy
$X_{t_k}=\Delta_1+\cdots+\Delta_k$.
Stąd przez przestawienie dwóch \emph{skończonych} sum
$\sum_{k=1}^m u_k X_{t_k}
=\sum_{j=1}^m(\sum_{k=j}^m u_k)\Delta_j$.
Korzystając z niezależności $\Delta_1,\ldots,\Delta_m$,
możemy rozłożyć wartość oczekiwaną iloczynu
wykładników na iloczyn wartości oczekiwanych,
co daje wzór. Funkcja charakterystyczna wyznacza
jednoznacznie rozkład wektora losowego (twierdzenie
o jednoznaczności funkcji charakterystycznej z tomu I),
a dowolną skończoną rodzinę chwil można uporządkować
i dołączyć chwilę $0$. Dowodzi to ostatniego zdania.
\end{proof}

\begin{przyklad}[Dwa pierwsze kroki spaceru losowego]
Niech $\xi_1,\xi_2$ będą niezależne, przyjmujące
wartości $-1,1$ z prawdopodobieństwem $1/2$,
a $S_0=0$, $S_1=\xi_1$, $S_2=\xi_1+\xi_2$.
Prawdopodobieństwo ścieżki $(0,1,0)$ wynosi
$\Prob(\xi_1=1,\xi_2=-1)=1/4$.
Z kolei $\Prob(S_2=0)=1/2$, bo wynik $0$ otrzymujemy
przy ścieżkach przyrostów $(1,-1)$ albo $(-1,1)$.
Stąd $\Prob(S_1=1\mid S_2=0)=(1/4)/(1/2)=1/2$.
Wniosek: nawet gdy pojedyncze przyrosty są niezależne,
położenia $S_1,S_2$ nie są niezależne; drugie położenie
zawiera informację o pierwszym.
\end{przyklad}

\begin{uwaga}
Stacjonarne przyrosty nie oznaczają stacjonarności samego procesu.
Proces Wienera zaczynający w zerze będzie miał stacjonarne
przyrosty, lecz wariancja $X_t$ będzie rosnąć jak $t$.
Stochastyczna ciągłość mówi jedynie o zbieżności w prawdopodobieństwie
przy zbliżaniu czasu, nie o ciągłości \emph{każdej trajektorii}.
Trzeba też odróżnić warunek dla każdej ustalonej chwili $t$
od jednego zdarzenia pełnej miary, na którym wszystkie trajektorie
mają żądaną regularność.
\end{uwaga}

\begin{definicja}[Stacjonarność ścisła i w drugim rzędzie]
Niech $T=\R$ lub $T=[0,\infty)$, a $X_t\in\R^d$.
Proces jest \emph{ściśle stacjonarny}, jeżeli dla każdego
$m\ge1$, $t_1,\ldots,t_m\in T$ i $h$ takich, że wszystkie
$t_i+h\in T$, zachodzi
$(X_{t_1+h},\ldots,X_{t_m+h})\overset{d}= (X_{t_1},\ldots,X_{t_m})$.
Dla procesu rzeczywistego z $\E X_t^2<\infty$ dla wszystkich $t$
mówimy o \emph{stacjonarności w drugim rzędzie}, gdy
$\E X_t=m$ nie zależy od $t$, a $\Cov(X_s,X_t)=c(t-s)$
dla pewnej funkcji $c$ i wszystkich $s,t\in T$.
\end{definicja}

\begin{twierdzenie}
Każdy ściśle stacjonarny proces rzeczywisty,
dla którego $\E X_t^2<\infty$ dla wszystkich $t$,
jest stacjonarny w drugim rzędzie.
\end{twierdzenie}
\begin{proof}
Równość rozkładów $X_t$ i $X_{t+h}$ daje równość średnich
oraz drugich momentów. Równość rozkładów par
$(X_s,X_t)$ i $(X_{s+h},X_{t+h})$ daje równość
$\E[X_sX_t]$ i $\E[X_{s+h}X_{t+h}]$, więc i kowariancji.
Wybierając przesunięcie, które przenosi jedną z chwil na $0$
(gdy $T=[0,\infty)$, używamy dozwolonego przesunięcia
obu nieujemnych chwil), otrzymujemy zależność kowariancji
wyłącznie od różnicy czasów.
\end{proof}

\begin{przyklad}[Stacjonarny proces z gładkimi trajektoriami]
Niech $Z_1,Z_2$ będą niezależnymi zmiennymi standardowo normalnymi
i niech $X_t=Z_1\cos t+Z_2\sin t$ dla $t\in\R$.
Mamy $\E X_t=0$ i
\[
\Cov(X_s,X_t)=\cos s\cos t+\sin s\sin t=\cos(t-s).
\]
Dla dowolnego $m$ wektor $(X_{t_1},\ldots,X_{t_m})$
jest gaussowski, więc wyznaczają go średnie i macierz kowariancji.
Po przesunięciu wszystkich chwil o $h$ te dane nie zmieniają się,
zatem proces jest ściśle stacjonarny. Każda trajektoria jest przy
tym funkcją gładką; stacjonarność nie oznacza ``poszarpania'' ścieżek.
\end{przyklad}

\begin{przyklad}[Ciągłość stochastyczna bez ciągłych trajektorii]
Niech $U$ ma rozkład jednostajny na $[0,1]$ i
$X_t=\ind{\{U\le t\}}$ dla $t\in[0,1]$.
Dla ustalonego $t$ i $s$ zachodzi
$\Prob(\abs{X_s-X_t}>0)=\Prob(U\in(\min(s,t),\max(s,t)])
=\abs{s-t}$, zatem $X$ jest stochastycznie ciągły.
Niemniej dla każdego $\omega$ z $U(\omega)\in(0,1)$
trajektoria ma skok z $0$ do $1$ w chwili $U(\omega)$.
Co więcej, nie istnieje ciągła modyfikacja: gdyby $Y$ była ciągłą
modyfikacją, to poza jednym zbiorem zerowym mielibyśmy
$Y_q=X_q$ dla wszystkich $q\in\Q\cap[0,1]$ jednocześnie.
Dla każdego $\omega$ z tego zbioru pełnej miary oraz
$U(\omega)\in(0,1)$ wartości $X_q(\omega)$ na wymiernych
po obu stronach $U(\omega)$ miałyby różne granice jednostronne.
Ciągła funkcja $Y(\omega)$ nie może mieć takiego zachowania.
\end{przyklad}

\section{Procesy gaussowskie jako zapowiedź konstrukcji}
\begin{definicja}[Proces gaussowski]
Proces rzeczywisty $X=(X_t)_{t\in T}$ nazywamy
\emph{gaussowskim}, jeżeli dla każdego $m\ge1$
i $t_1,\ldots,t_m\in T$ wektor
$(X_{t_1},\ldots,X_{t_m})$ ma wielowymiarowy rozkład
normalny (dopuszczamy rozkłady zdegenerowane).
Jego funkcjami średniej i kowariancji są odpowiednio
$m(t)=\E X_t$ i $K(s,t)=\Cov(X_s,X_t)$.
\end{definicja}

\begin{stwierdzenie}[Dodatnia półokreśloność kowariancji]
Jeśli $X=(X_t)_{t\in T}$ jest procesem rzeczywistym o
skończonych drugich momentach, to funkcja
$K(s,t)=\Cov(X_s,X_t)$ jest symetryczna i dla każdych
$m\ge1$, $t_1,\ldots,t_m\in T$, $a_1,\ldots,a_m\in\R$
zachodzi
$\sum_{i,j=1}^m a_i a_j K(t_i,t_j)\ge0$.
\end{stwierdzenie}
\begin{proof}
Symetria wynika z symetrii kowariancji. Po rozwinięciu
wariancji kombinacji liniowej otrzymujemy
\[
\sum_{i,j=1}^m a_i a_j K(t_i,t_j)
=\Var\left(\sum_{i=1}^m a_i X_{t_i}\right)\ge0.
\]
Każda macierz $[K(t_i,t_j)]_{i,j}$ musi więc być dodatnio
półokreślona, co stanowi istotne ograniczenie przy próbie
zadania procesu przez samą funkcję kowariancji.
\end{proof}

\begin{zadanie}[Czy dwa procesy mają ten sam rozkład?]
Niech $U,V$ będą niezależnymi zmiennymi Bernoulliego z parametrem
$1/2$. Na przedziale $[0,1]$ definiujemy
$X_t=U$ dla wszystkich $t$ oraz
$Y_t=U$ dla $t<1/2$ i $Y_t=V$ dla $t\ge1/2$.
Obliczymy rozkłady jednowymiarowe, prawdopodobieństwo
$\Prob(X_{1/4}=X_{3/4})$ oraz jego odpowiednik dla $Y$.

Dla każdego $t$ obie zmienne $X_t$ i $Y_t$ mają rozkład
Bernoulliego z parametrem $1/2$, więc same rozkłady
jednowymiarowe nie rozróżniają tych procesów. Natomiast
$X_{1/4}=X_{3/4}=U$, zatem
$\Prob(X_{1/4}=X_{3/4})=1$.
Dla drugiego procesu
$Y_{1/4}=U$, $Y_{3/4}=V$; stąd przez niezależność
\[
\Prob(Y_{1/4}=Y_{3/4})
=\Prob(U=V=0)+\Prob(U=V=1)=\tfrac14+\tfrac14=\tfrac12.
\]
Ponadto $X$ ma ciągłe trajektorie, a $Y$ ma skok w chwili
$1/2$ z prawdopodobieństwem $\Prob(U\ne V)=1/2$.
Zadanie pokazuje równocześnie rolę rozkładów łącznych
i regularności trajektorii.
\end{zadanie}

\begin{uwaga}[Granice pierwszego rozdziału]
Dotąd startowaliśmy od \emph{istniejącego} procesu i wyznaczaliśmy
jego prawo. Jeśli chcemy w modelu narzucić rozkłady, np.
$(X_{t_1},\ldots,X_{t_m})\sim\Normal(0,[\min(t_i,t_j)]_{i,j})$,
musimy wykazać istnienie procesu o takich rozkładach.
Dalej będziemy jeszcze wymagać, by jego trajektorie były ciągłe.
Te dwa kroki nie mogą być zlane w jeden.
\end{uwaga}


\begin{przyklad}[Ta sama kowariancja nie oznacza tego samego prawa]
Niech $Z_1,Z_2$ będą niezależnymi zmiennymi standardowo
normalnymi, a $\xi_1,\xi_2$ niezależnymi symetrycznymi
zmiennymi o wartościach $-1,1$. Definiujemy na $\R$
\[
X_t=Z_1\cos t+Z_2\sin t,
\qquad Y_t=\xi_1\cos t+\xi_2\sin t.
\]
Oba procesy mają średnią zero i identyczną funkcję
kowariancji $K(s,t)=\cos(t-s)$, ponieważ
$\E Z_i^2=\E\xi_i^2=1$ oraz wszystkie mieszane
kowariancje są zerowe. Jednak $X_0\sim\Normal(0,1)$,
a $Y_0=\xi_1$ przyjmuje tylko wartości $-1,1$.
Zatem procesy mają inne rozkłady już w chwili $0$.
Funkcja kowariancji wyznacza rozkłady skończeniewymiarowe
\emph{w klasie procesów gaussowskich}, ale nie w klasie
wszystkich procesów o skończonych drugich momentach.
\end{przyklad}


\chapter{Konstrukcja procesów i regularność trajektorii}
\label{ch:konstrukcja}

\section{Problem istnienia: od tabel rozkładów do całego procesu}
Załóżmy, że chcemy opisać losowy ruch bez podawania jawnego wzoru
$X_t(\omega)$. Możemy zamiast tego ustalić, jak mają wyglądać
rozkłady wszystkich wektorów
$(X_{t_1},\ldots,X_{t_m})$. Nie każda dowolna lista rozkładów
może jednak pochodzić od jednej rodziny zmiennych losowych.
Jeżeli np. zadamy $\Prob(X_1=0)=1$, a jednocześnie
$\Prob(X_1=1,X_2=0)=1$, otrzymamy sprzeczność.
Rozkłady muszą dawać zgodne odpowiedzi na pytania o te same
pomiarowe współrzędne.

W rozdziale~\ref{ch:procesy} wykazaliśmy, że rzeczywiste
rozkłady procesu spełniają relacje permutacji i marginalizacji.
Teraz udowodnimy twierdzenie odwrotne. Jego zasadnicza trudność
nie leży w sprawdzeniu skończenie wielu zmiennych, lecz w
zbudowaniu miary na przestrzeni funkcji indeksowanych
nieprzeliczalnym zbiorem chwil.

\begin{definicja}[Zgodny układ rozkładów]
Niech $T$ będzie niepustym zbiorem, a $E$ niepustą
przestrzenią polską, tj. ośrodkową przestrzenią topologiczną,
której topologię można zadać metryką zupełną.
Wyposażamy $E$ w borelowskie $\sigma$-ciało $\calB(E)$. Dla każdej skończonej
\emph{listy} $\mathbf t=(t_1,\ldots,t_m)$ indeksów z $T$
(niekoniecznie parami różnych) niech $\mu_{\mathbf t}$
będzie miarą probabilistyczną na $E^m$.
Rodzinę $(\mu_{\mathbf t})$ nazywamy \emph{zgodną}, jeżeli:
\begin{enumerate}
\item dla każdej permutacji $\pi$ indeksów $1,\ldots,m$,
$\mu_{t_{\pi(1)},\ldots,t_{\pi(m)}}=(\Pi_\pi)_\#\mu_{t_1,\ldots,t_m}$;
\item dla $k<m$ rzut $E^m\to E^k$ na pierwsze $k$ współrzędnych
przenosi $\mu_{t_1,\ldots,t_m}$ w $\mu_{t_1,\ldots,t_k}$;
\item jeżeli $t_i=t_j$, to $\mu_{t_1,\ldots,t_m}$ nadaje
prawdopodobieństwo $1$ zbiorowi $\{x\in E^m:x_i=x_j\}$.
\end{enumerate}
Symbol $(f)_\#\mu$ oznacza miarę obrazową
$((f)_\#\mu)(A)=\mu(f^{-1}(A))$.
\end{definicja}

Warunek dotyczący powtórzonych indeksów nie jest dodatkową
``własnością niezależności''. Wyraża po prostu to, że
$X_t$ oglądane dwa razy w tej samej chwili musi być tą samą
zmienną. Równoważnie można zadawać rozkłady tylko dla
skończonych \emph{zbiorów} czasów i wymagać zgodności
względem wszystkich rzutowań; wtedy problem powtórzeń znika.

\begin{lemat}[Algebra cylindrów]
Niech $E\ne\varnothing$ będzie przestrzenią polską, a $T$
niepustym zbiorem. Niech $\cal A$ będzie rodziną podzbiorów
$E^T$ postaci $\pi_J^{-1}(B)$, gdzie
$J\subseteq T$ jest skończony, $\pi_J:E^T\to E^J$
jest rzutowaniem, a $B\in\calB(E^J)$.
Wówczas $\cal A$ jest algebrą zbiorów i
$\sigma(\cal A)=\calB(E)^{\otimes T}$.
\end{lemat}
\begin{proof}
Zbiory $\varnothing,E^T$ należą do $\cal A$.
Dopełnienie $\pi_J^{-1}(B)$ jest równe
$\pi_J^{-1}(E^J\setminus B)$. Jeśli
$A=\pi_J^{-1}(B)$, $A'=\pi_K^{-1}(B')$, to
$A\cup A'$ jest cylindrem zależnym od skończonego zbioru
$J\cup K$: jego baza jest sumą przeciwobrazów $B,B'$ przez
odpowiednie rzuty $E^{J\cup K}\to E^J,E^K$.
Dowodzi to domknięcia na skończone sumy i dopełnienia.
Dla skończonego $J$ zachodzi
$\mathcal B(E^J)=\mathcal B(E)^{\otimes J}$.
Istotnie, rzuty są ciągłe, co daje jedną inkluzję,
a przeliczalna baza topologii $E$ daje przeliczalną
bazę prostokątów otwartych w $E^J$, co daje drugą inkluzję.
Każdy cylinder jest więc mierzalny względem produktowego
$\sigma$-ciała, a cylindry zależne od jednej współrzędnej
generują to $\sigma$-ciało. Stąd równość.
\end{proof}


\begin{przyklad}[Dlaczego nie wystarczy zgodność pojedynczych rozkładów?]
Niech $T=\{0,1\}$, $E=\{0,1\}$ i zażądajmy, aby
$X_0$ oraz $X_1$ miały symetryczny rozkład Bernoulliego.
Te wymagania dopuszczają zarówno niezależność
($\Prob(X_0=X_1)=1/2$), jak i równość $X_0=X_1$ prawie
na pewno. Trzeba więc wybrać także \emph{wspólne} prawo pary.
Jeśli rozszerzymy zbiór indeksów o kolejne chwile,
ten wybór musi być zgodny z prawami trójek, czwórek itd. Twierdzenie Kołmogorowa mówi, że przy
zgodności \emph{całej} rodziny problem jest rozwiązywalny;
nie mówi natomiast niczego o gładkości lub ciągłości
otrzymanych trajektorii.
\end{przyklad}

\begin{twierdzenie}[Twierdzenie Kołmogorowa o istnieniu procesu]
\label{thm:kol-ext}
Niech $T$ będzie dowolnym niepustym zbiorem, a $(E,d)$
niepustą przestrzenią polską z borelowskim $\sigma$-ciałem.
Niech dla każdej skończonej listy $\mathbf t=(t_1,\ldots,t_m)$
indeksów z $T$ dana będzie miara probabilistyczna
$\mu_{\mathbf t}$ na $E^m$, przy czym rodzina tych miar
jest zgodna w sensie powyższej definicji.
Wówczas istnieją przestrzeń probabilistyczna
$(\Omega,\calF,\Prob)$ i proces
$(X_t)_{t\in T}$ o stanach w $E$ takie, że
\[
\Law(X_{t_1},\ldots,X_{t_m})=\mu_{t_1,\ldots,t_m}
\quad\text{dla wszystkich }m\ge1\text{ i }t_1,\ldots,t_m\in T.
\]
Można wybrać $\Omega=E^T$,
$\calF=\calB(E)^{\otimes T}$ i $X_t(\omega)=\omega(t)$.
Prawo procesu na tej przestrzeni jest jednoznacznie wyznaczone
przez rodzinę $\mu_{\mathbf t}$.
\end{twierdzenie}

\begin{proof}
\emph{Krok 1: miara na cylindrach.}
Niech $\cal A$ będzie algebrą z poprzedniego lematu.
Dla cylindra $A=\pi_J^{-1}(B)$ definiujemy
$\mu_0(A)=\mu_J(B)$, gdzie dla skończonego $J$
przez $\mu_J$ rozumiemy rozkład we współrzędnych $J$
w dowolnej ustalonej kolejności, utożsamiony z miarą na $E^J$.
Dla $J=\varnothing$ przyjmujemy, że $E^J$ jest jednopunktowe
i $\mu_J$ jest jedyną miarą probabilistyczną na tej przestrzeni. Relacje permutacji i
marginalizacji dowodzą, że wynik nie zależy od kolejności
ani od dołączenia sztucznych współrzędnych. W szczególności
wartość $\mu_0(A)$ nie zależy od wybranej prezentacji
cylindra: jeżeli $A=\pi_J^{-1}(B)=\pi_K^{-1}(B')$,
obie prezentacje przenosimy na wspólny skończony
zbiór współrzędnych $L=J\cup K$.
Rzut $\pi_L:E^T\to E^L$ jest surjekcją, gdyż $E$ jest
niepuste, więc równość cylindrów oznacza równość ich
baz po przeniesieniu na $E^L$. Zgodność rozkładów daje
równość ich miar pod $\mu_L$.
Stąd $\mu_0(E^T)=1$ i $\mu_0$ jest skończenie addytywna:
dla rozłącznych cylindrów $A,B$ wyrażamy oba na
wspólnym skończonym zbiorze współrzędnych i stosujemy
addytywność miary skończeniewymiarowej.

\emph{Krok 2: dlaczego potrzebujemy ciągłości z góry.}
Aby rozszerzyć $\mu_0$ na $\sigma(\cal A)$ za pomocą
twierdzenia Carathéodory'ego, musimy pokazać przeliczalną
addytywność $\mu_0$ na $\cal A$. W przypadku skończonej
miary skończenie addytywnej wystarczy następujący warunek:
\[
A_n\in\cal A,\quad A_{n+1}\subseteq A_n,\quad
\bigcap_{n\ge1}A_n=\varnothing
\quad\Longrightarrow\quad \mu_0(A_n)\longrightarrow0.
\tag{$*$}
\]
Rzeczywiście, jeżeli $A=\bigsqcup_{k\ge1}B_k$,
przy $A,B_k\in\cal A$, to zbiory
$R_n=A\setminus\bigcup_{k=1}^n B_k$ należą do $\cal A$,
maleją do $\varnothing$, a skończona addytywność daje
$\mu_0(A)=\sum_{k=1}^n\mu_0(B_k)+\mu_0(R_n)$.
Przejście do granicy daje przeliczalną addytywność.

\emph{Krok 3: dowód warunku $(*)$ przy użyciu zwartości.}
Przypuśćmy, że $A_n\downarrow\varnothing$ i że
$\mu_0(A_n)\not\to0$. Ponieważ $\mu_0$ jest nieujemna i skończenie addytywna,
jest monotoniczna, więc ciąg $\mu_0(A_n)$ jest nierosnący. Istnieje $\varepsilon>0$ takie, że
$\mu_0(A_n)\ge\varepsilon$ dla wszystkich $n$.
Każdy $A_n$ zależy od pewnego skończonego zbioru
współrzędnych $J_n\subseteq T$. Ich suma
$S:=\bigcup_{n\ge1}J_n$ jest co najwyżej przeliczalna.
Jeśli $S=\varnothing$, wszystkie cylindry $A_n$ są
równe $\varnothing$ albo $E^T$, więc sprzeczność jest
natychmiastowa. W przeciwnym razie wypiszmy elementy
$S$ jako $s_1,s_2,\ldots$ (dla skończonego $S$
lista kończy się po skończenie wielu wyrazach). Ponieważ rozkłady
jednowymiarowe są miarami borelowskimi na przestrzeni
polskiej, są ciasne na mocy twierdzenia o ciasności
skończonych miar na przestrzeniach polskich
(por. skrypt \emph{Teoria miary i całki}, rozdział
\emph{Miary borelowskie, regularność i miary Radona}): dla każdego $j$ istnieje niepusty
zwarty $K_j\subseteq E$ taki, że
\[
\mu_{s_j}(E\setminus K_j)<\varepsilon2^{-j-3}.
\tag{1}
\]
Dla każdego $n$ zapiszmy $A_n=\pi_{J_n}^{-1}(B_n)$.
Skończony produkt przestrzeni polskich jest polski
(metryka będąca skończoną sumą metryk zupełnych jest
zupełna, a produkt skończenie wielu przeliczalnych
zbiorów gęstych jest przeliczalny i gęsty).
Wewnętrzna regularność względem zbiorów zwartych
skończonej miary borelowskiej $\mu_{J_n}$ na tej
przestrzeni daje niepusty zwarty
$C_n\subseteq B_n$ taki, że
\[
\mu_{J_n}(B_n\setminus C_n)<\varepsilon2^{-n-3}.
\tag{2}
\]
Zbiory $C_n$ są niepuste, ponieważ
$\mu_{J_n}(B_n)=\mu_0(A_n)\ge\varepsilon$.

Ustalmy $N$ i niech $I_N=\bigcup_{n=1}^N J_n$.
W skończonym produkcie $E^{I_N}$ rozważmy jednoczesne
warunki $\pi_{J_n}(x)\in C_n$ dla $1\le n\le N$
(tutaj $\pi_{J_n}$ oznacza rzut z $E^{I_N}$ na $E^{J_n}$)
oraz $x(s_j)\in K_j$ dla wszystkich $s_j\in I_N$.
Ich wspólna miara pod $\mu_{I_N}$ jest co najmniej
\[
\begin{aligned}
\mu_0(A_N)
&-\sum_{n=1}^N \mu_0\big(A_n\setminus\pi_{J_n}^{-1}(C_n)\big)
-\sum_{\{j:s_j\in I_N\}}\mu_{s_j}(E\setminus K_j)\\
&>\varepsilon-\frac\varepsilon8-\frac\varepsilon8
=\frac{3\varepsilon}{4}>0.
\end{aligned}
\tag{3}
\]
Użyliśmy tutaj $A_N\subseteq A_n$ dla $n\le N$ oraz
nierówności na miarę sumy, wyłącznie dla miary
$\mu_{I_N}$ (zgodność pozwala utożsamić jej wartości
na odpowiednich bazach cylindrów z wartościami $\mu_0$).
Nie korzystamy z przeliczalnej addytywności $\mu_0$,
której dopiero dowodzimy. Zatem istnieje wybór
współrzędnych na $I_N$, spełniający wszystkie te warunki.

Produkt $K:=\prod_{s_j\in S}K_j$ jest niepustą przestrzenią
zwartą (dla skończonego $S$ jest to produkt skończony).
Dla każdego $n$ zbiór
$F_n:=\{x\in K:\pi_{J_n}(x)\in C_n\}$
jest domknięty w $K$, ponieważ $C_n$ jest zwarty,
więc domknięty w $E^{J_n}$. Oszacowanie (3) pokazuje,
że dla każdego $N$ przecięcie $F_1\cap\cdots\cap F_N$
nie jest puste: znalezione wartości na skończonym $I_N$
można dopełnić dowolnymi punktami $K_j$ dla
pozostałych współrzędnych. Z własności zwartości
i skończonych przecięć dostajemy
$\bigcap_{n\ge1}F_n\ne\varnothing$.
Wybierzmy $x$ z tego przecięcia. Wówczas
$\pi_{J_n}(x)\in C_n\subseteq B_n$ dla każdego $n$.
Rozszerzając $x$ do elementu $E^T$ poprzez wybór
stałego punktu $e_0\in E$ poza $S$, otrzymujemy
$\omega\in\bigcap_{n\ge1}A_n$, sprzeczność.
Wobec tego warunek $(*)$ jest spełniony.

\emph{Krok 4: rozszerzenie i proces kanoniczny.}
Na mocy kroków 2--3 i twierdzenia Carathéodory'ego istnieje
miara probabilistyczna $\Prob$ na
$\calF=\sigma(\cal A)=\calB(E)^{\otimes T}$,
której ograniczenie do $\cal A$ jest równe $\mu_0$.
Odwzorowania współrzędne $X_t(\omega)=\omega(t)$ są
mierzalne z definicji produktowego $\sigma$-ciała,
a ich rozkłady dla list parami różnych indeksów są zadanymi
$\mu_{\mathbf t}$ z definicji $\mu_0$.
Dla dowolnej listy $\mathbf t=(t_1,\ldots,t_m)$ wybierzmy
listę $\mathbf s=(s_1,\ldots,s_k)$ wszystkich jej różnych
indeksów. Niech $r:E^m\to E^k$ wybiera po jednej
współrzędnej dla każdego $s_i$, a $j:E^k\to E^m$
odtwarza kolejność i powtórzenia z listy $\mathbf t$.
Z permutacji i marginalizacji wynika
$r_\#\mu_{\mathbf t}=\mu_{\mathbf s}$.
Na zbiorze, na którym wszystkie powtórzone współrzędne
są równe, zachodzi $j\circ r=\operatorname{id}$.
Zbiór ten jest mierzalny i ma miarę $1$ pod $\mu_{\mathbf t}$
na mocy trzeciego warunku zgodności. W konsekwencji
$\mu_{\mathbf t}=j_\#\mu_{\mathbf s}$.
Wektor kanonicznych współrzędnych dla listy $\mathbf t$
jest obrazem wektora dla $\mathbf s$ przez $j$, więc
również ma wymagany rozkład.
Jednoznaczność prawa wynika z jednoznaczności miar,
które zgadzają się na generującej algebrze cylindrów.
\end{proof}

\begin{uwaga}[Rola założenia o przestrzeni polskiej]
Założenie polskości $E$ pojawia się w kroku 3: gwarantuje,
że miary skończeniewymiarowe można od wewnątrz
przybliżać zbiorami zwartymi, a marginesy są ciasne.
Bez odpowiednich założeń o przestrzeni stanów
rozszerzenie zgodnej rodziny rozkładów może zawieść.
Twierdzenie nie zapewnia jeszcze ciągłości trajektorii.
Jeżeli zbiór indeksów jest wyposażony w $\sigma$-ciało,
nie gwarantuje też mierzalności wspólnej odwzorowania
$(t,\omega)\mapsto X_t(\omega)$ względem odpowiedniego
produktowego $\sigma$-ciała. Otrzymaliśmy tylko
mierzalność każdej współrzędnej $X_t$ z osobna.
\end{uwaga}

\section{Przykłady konstrukcji z zadanych rozkładów}
\begin{przyklad}[Nieskończona rodzina niezależnych zmiennych]
Niech $T$ będzie dowolnym zbiorem, a $\nu$ miarą
probabilistyczną na przestrzeni polskiej $E$.
Dla $t_1,\ldots,t_m$ parami różnych zadajmy
$\mu_{t_1,\ldots,t_m}=\nu^{\otimes m}$,
a dla powtórzonych indeksów utożsammy odpowiednie
współrzędne. Jest to rodzina zgodna, ponieważ
permutacja współrzędnych nie zmienia miary produktowej,
a rzut usuwa niepotrzebne czynniki.
Twierdzenie Kołmogorowa tworzy proces o niezależnych
$X_t$ i identycznych rozkładach $\nu$.

Jeżeli $T=[0,1]$ i $\nu$ jest rozkładem Bernoulliego z
parametrem $1/2$, to dla różnych $s,t$
$\Prob(|X_t-X_s|=1)=1/2$. Prawa strona nie dąży do
zera przy $s\to t$. Taki proces nie jest stochastycznie
ciągły i nie może mieć ciągłej modyfikacji: ciągłość
trajektorii prawie na pewno implikowałaby zbieżność
w prawdopodobieństwie, a modyfikacja zachowuje
rozkłady par $(X_s,X_t)$, a więc także rozkłady przyrostów. Samo
istnienie procesu jest więc znacznie słabsze niż
regularność jego trajektorii.
\end{przyklad}

\begin{przyklad}[Proces gaussowski z zadanej kowariancji]
Niech $T$ będzie zbiorem, $m:T\to\R$ dowolną funkcją,
a $K:T\times T\to\R$ funkcją symetryczną taką, że
dla wszystkich $k\ge1$, $t_1,\ldots,t_k\in T$
i $a_1,\ldots,a_k\in\R$ zachodzi
$\sum_{i,j=1}^k a_ia_jK(t_i,t_j)\ge0$.
Dla listy czasów $t_1,\ldots,t_k$ zadawajmy
$\mu_{t_1,\ldots,t_k}=\Normal(m_{\mathbf t},K_{\mathbf t})$,
gdzie $m_{\mathbf t}=(m(t_i))_{i=1}^k$,
$K_{\mathbf t}=[K(t_i,t_j)]_{i,j=1}^k$.
Nawet gdy $K_{\mathbf t}$ jest osobliwa, rozkład normalny
istnieje: można zapisać $m_{\mathbf t}+AZ$, gdzie
$Z\sim\Normal(0,I_k)$ i $AA^\top=K_{\mathbf t}$.
Macierz $A$ istnieje na mocy twierdzenia spektralnego
dla rzeczywistych macierzy symetrycznych: jeśli
$K_{\mathbf t}=Q\Lambda Q^\top$, gdzie $Q$ jest ortogonalna,
a $\Lambda$ diagonalna o nieujemnych wyrazach, można
wziąć $A=Q\Lambda^{1/2}$.
Rozkład ten jest jednoznacznie wyznaczony przez średnią
i kowariancję, gdyż jego funkcja charakterystyczna w punkcie
$u\in\R^k$ ma postać
$\exp(iu^\top m_{\mathbf t}-u^\top K_{\mathbf t}u/2)$
(por. rozdział o procesach gaussowskich).
Rzut i permutacja wektora gaussowskiego zachowują jego
gaussowskość i odpowiednio przekształcają średnią
oraz kowariancję. Z jednoznaczności rozkładu wynika
zatem zgodność z permutacją i marginalizacją. Jeżeli $t_i=t_j$, to różnica
odpowiednich współrzędnych ma średnią $m(t_i)-m(t_j)=0$
i wariancję $K(t_i,t_i)+K(t_j,t_j)-2K(t_i,t_j)=0$,
więc współrzędne te są równe prawie na pewno. Zatem rodzina jest zgodna
i istnieje proces gaussowski o średniej $m$ i kowariancji $K$.
\end{przyklad}


\begin{przyklad}[Konstrukcja łańcucha Markowa z macierzy przejścia]
Niech $E=\{0,1\}$, $T=\N_0$, $\mu=(\mu_0,\mu_1)$
będzie rozkładem początkowym, a
$P=(p_{ij})_{i,j\in E}$ macierzą stochastyczną:
$p_{ij}\ge0$ i $\sum_{j\in E}p_{ij}=1$ dla każdego $i$.
Zadajmy dla $n\ge0$ i $x_0,\ldots,x_n\in E$
\[
\mu_{0,1,\ldots,n}(\{(x_0,\ldots,x_n)\})
:=\mu_{x_0}\prod_{k=1}^n p_{x_{k-1}x_k}.
\tag{7}
\]
Sumując po $x_n$ i korzystając z
$\sum_{x_n\in E}p_{x_{n-1}x_n}=1$, widzimy,
że ostatnia współrzędna może zostać usunięta
bez zmiany prawa wcześniejszych współrzędnych.
Gdy pomijamy chwilę pośrednią $k$, pojawia się suma
$\sum_j p_{ij}p_{j\ell}=(P^2)_{i\ell}$,
a po pominięciu większej liczby chwil
odpowiednie potęgi $P$.
Dokładniej, dla $t_1,\ldots,t_m\in\N_0$ takich, że
$0\le t_1<\cdots<t_m$, i $x_1,\ldots,x_m\in E$ określamy
\[
\mu_{t_1,\ldots,t_m}(\{(x_1,\ldots,x_m)\})
=(\mu P^{t_1})_{x_1}
\prod_{k=2}^m (P^{t_k-t_{k-1}})_{x_{k-1}x_k}.
\]
Każda potęga $P^r$, $r\in\N_0$, jest macierzą stochastyczną:
iloczyn dwóch macierzy stochastycznych jest nieujemny
i zachowuje wektor kolumnowy jedynek. Sumowanie
powyższego wzoru po kolejnych współrzędnych od końca
daje zatem $1$, więc jest to miara probabilistyczna.
Równość $P^{a+b}=P^aP^b$ pokazuje zgodność
z dowolnymi rzutowaniami; powtórzenie chwili
oznacza z definicji powtórzenie tej samej współrzędnej.
Twierdzenie~\ref{thm:kol-ext} daje więc proces
$(X_n)_{n\ge0}$ o tych rozkładach.

Dla zdarzenia o dodatnim prawdopodobieństwie
$\{X_0=x_0,\ldots,X_n=x_n\}$ obliczamy z (7)
\[
\begin{aligned}
&\Prob(X_{n+1}=j\mid X_0=x_0,\ldots,X_n=x_n)\\
&\quad=\frac{\mu_{x_0}\,p_{x_0x_1}\cdots
p_{x_{n-1}x_n}\,p_{x_nj}}
{\mu_{x_0}\,p_{x_0x_1}\cdots p_{x_{n-1}x_n}}
=p_{x_nj}.
\end{aligned}
\]
Tak otrzymujemy własność Markowa, a nie tylko
jej formalną deklarację.

Dla konkretnego modelu weźmy
$\mu=(1,0)$ i
$P=\begin{pmatrix}3/4&1/4\\1/2&1/2\end{pmatrix}$.
Wtedy $\Prob(X_0=0,X_1=0,X_2=1)
=1\cdot(3/4)\cdot(1/4)=3/16$, podczas gdy
$\Prob(X_2=1)=(P^2)_{01}
=3/4\cdot1/4+1/4\cdot1/2=5/16$.
Różnica polega na tym, że w drugim pytaniu
sumujemy po wszystkich możliwych stanach w chwili $1$.
\end{przyklad}

\section{Kiedy proces może mieć trajektorie ciągłe?}
Rozkłady skończeniewymiarowe odpowiadają na pytania o
skończoną liczbę chwil, ale własność ciągłości wymaga
jednoczesnej kontroli \emph{nieskończenie wielu} przyrostów.
Narzędziem będzie oszacowanie momentów przyrostów
$\E\norm{X_t-X_s}^p$ o wykładniku potęgowym ściśle
większym niż $1$ w jednowymiarowym parametrze czasu.

\begin{definicja}[Ciągłość Höldera]
Niech $T>0$ oraz $\gamma>0$.
Funkcję $f:[0,T]\to\R^d$ nazywamy $\gamma$-h"olderowską,
jeśli istnieje stała $C_f<\infty$ taka, że
$\norm{f(t)-f(s)}\le C_f\abs{t-s}^{\gamma}$
dla wszystkich $s,t\in[0,T]$.
Gdy $\gamma>1$, funkcja o tej własności musi być stała:
dzieląc przedział $[s,t]$ na $n$ równych części,
to ten sam warunek daje
$\norm{f(t)-f(s)}\le C_f\abs{t-s}^{\gamma}n^{1-\gamma}\to0$.
\end{definicja}

\begin{lemat}[Kontrola przyrostów na siatkach dyadycznych]
\label{lem:dyadyczny}
Niech $T>0$, $d\ge1$ i $0<\gamma<\infty$.
Niech $D_n:=\{kT2^{-n}:k=0,\ldots,2^n\}$,
$D:=\bigcup_{n\ge0}D_n$. Jeżeli $f:D\to\R^d$
spełnia dla pewnego $N\ge0$
\[
\norm{f((k+1)T2^{-n})-f(kT2^{-n})}
\le 2^{-n\gamma}
\quad(n\ge N,\ 0\le k<2^n),
\]
to $f$ jest jednostajnie ciągła na $D$,
przedłuża się jednoznacznie do ciągłej funkcji
$\bar f:[0,T]\to\R^d$, a dla pewnej stałej
$C_{\gamma,T,N,f}<\infty$ zachodzi
$\norm{\bar f(t)-\bar f(s)}\le C_{\gamma,T,N,f}|t-s|^\gamma$.
\end{lemat}
\begin{proof}
Dla $u\in D$ i $n\ge N$ oznaczmy przez
$u_n=T2^{-n}\lfloor 2^n u/T\rfloor$ największy punkt
siatki $D_n$ nieprzekraczający $u$ (w szczególności
dla $u=T$ mamy $u_n=T$).
Dla dostatecznie dużego $m$ mamy $u_m=u$,
a kolejne przybliżenia $u_j,u_{j+1}$ są albo równe,
albo sąsiadują na siatce $D_{j+1}$. Z założenia
\[
\norm{f(u)-f(u_n)}
\le\sum_{j=n}^{m-1}2^{-(j+1)\gamma}
\le\frac{2^{-(n+1)\gamma}}{1-2^{-\gamma}}.
\tag{4}
\]
Jeśli $u,v\in D$ i $\abs{u-v}\le T2^{-n}$,
to ich lewostronne przybliżenia $u_n,v_n$
są identyczne lub sąsiadują na siatce $D_n$.
Wraz z (4) otrzymujemy
\[
\norm{f(u)-f(v)}\le
\left(1+\frac{2^{1-\gamma}}{1-2^{-\gamma}}\right)
2^{-n\gamma}.
\tag{5}
\]
Gdy $0<\abs{u-v}\le T2^{-N}$, wybierzmy $n\ge N$
tak, aby $T2^{-(n+1)}<\abs{u-v}\le T2^{-n}$.
Stąd $2^{-n\gamma}<(2\abs{u-v}/T)^\gamma$.
Z (5) dostajemy oszacowanie h"olderowskie na parach
wystarczająco bliskich. Dla par odległych o co najmniej
$T2^{-N}$ wykorzystujemy ograniczoność funkcji $f$ na
$D$, wynikającą z (4) oraz skończoności siatki $D_N$.
W ten sposób oszacowanie obowiązuje dla wszystkich par
w $D$ z pewną skończoną stałą. Ponieważ $\R^d$ jest
zupełna, $f$ ma jednoznaczne ciągłe przedłużenie
na domknięcie $D$, tj. na $[0,T]$; przejście do
odpowiednich granic przenosi oszacowanie na wszystkie $s,t$.
\end{proof}

\begin{twierdzenie}[Kryterium ciągłości Kołmogorowa]
\label{thm:kol-cont}
Niech $T>0$, $d\ge1$ i niech
$X=(X_t)_{t\in[0,T]}$ będzie procesem o wartościach w $\R^d$
na przestrzeni probabilistycznej $(\Omega,\calF,\Prob)$.
Załóżmy, że istnieją stałe $p>0$, $\beta>0$ i $C<\infty$
takie, że dla wszystkich $s,t\in[0,T]$ zachodzi
\[
\E\norm{X_t-X_s}^{p}\le C\abs{t-s}^{1+\beta}.
\tag{6}
\]
Wówczas istnieje modyfikacja $\widetilde X$ procesu $X$,
której trajektorie są ciągłe. Co więcej, można wybrać
$\widetilde X$ tak, że z prawdopodobieństwem $1$ jej
trajektorie są $\gamma$-h"olderowskie dla
\emph{każdego} $0<\gamma<\beta/p$.
\end{twierdzenie}
\begin{proof}
\emph{Krok 1: duże przyrosty na pojedynczej siatce.}
Ustalmy $0<\gamma_0<\beta/p$ i siatki
$D_n=\{kT2^{-n}:0\le k\le2^n\}$.
Niech $A_n$ będzie zdarzeniem, że istnieje sąsiednia
para punktów $D_n$, dla której przyrost przekracza
$2^{-n\gamma_0}$. Z nierówności Markowa, (6) i oszacowania
miary sumy mamy
\[
\begin{aligned}
\Prob(A_n)
&\le\sum_{k=0}^{2^n-1}
\frac{\E\norm{X_{(k+1)T2^{-n}}-X_{kT2^{-n}}}^{p}}
{2^{-np\gamma_0}}\\
&\le C T^{1+\beta}\,2^{-n(\beta-p\gamma_0)}.
\end{aligned}
\]
Ponieważ $\beta-p\gamma_0>0$, szereg
$\sum_{n\ge0}\Prob(A_n)$ jest zbieżny.

\emph{Krok 2: kontrola wszystkich dostatecznie drobnych siatek.}
Z pierwszego lematu Borela--Cantellego wynika, że z
prawdopodobieństwem $1$ zachodzi tylko skończenie wiele
zdarzeń $A_n$. Zatem na pewnym mierzalnym zbiorze
$\Omega_0$ o mierze $1$, dla każdego $\omega\in\Omega_0$
istnieje liczba całkowita $N(\omega)$ taka, że
wszystkie przyrosty sąsiednie na każdej siatce
$D_n$, $n\ge N(\omega)$, mają normę nie większą niż
$2^{-n\gamma_0}$. Dla takiego $\omega$ lemat
\ref{lem:dyadyczny} zastosowany do funkcji
$t\mapsto X_t(\omega)$ na $D:=\bigcup_nD_n$
daje jednoznaczne przedłużenie
$t\mapsto\widetilde X_t(\omega)$, które jest
$\gamma_0$-h"olderowskie.
Dla $\omega\notin\Omega_0$ ustawiamy
$\widetilde X_t(\omega)=0$ dla wszystkich $t$.
Przy ustalonym $t$ wybierzmy deterministyczny ciąg
$t_m\in D$ z $t_m\to t$. Na $\Omega_0$ mamy
$\widetilde X_t=\lim_m X_{t_m}$, a poza $\Omega_0$
zdefiniowaliśmy wartość stałą. Dokładniej, $\widetilde X_t$ jest granicą punktową na
całym $\Omega$ zmiennych mierzalnych
$\chi_{\Omega_0}X_{t_m}$: na $\Omega_0$ zachodzi zbieżność
do wartości przedłużenia, a poza nim wszystkie wyrazy
są zerowe. Nie wymaga to zupełności przestrzeni probabilistycznej.

\emph{Krok 3: przedłużenie jest modyfikacją.}
Dla ustalonego $t$ z (6) wynika
$\E\norm{X_{t_m}-X_t}^{p}\le C|t_m-t|^{1+\beta}\to0$.
Nierówność Markowa daje zbieżność $X_{t_m}\to X_t$
w prawdopodobieństwie, również gdy $0<p<1$.
Z drugiej strony $X_{t_m}\to\widetilde X_t$
prawie na pewno na $\Omega_0$, zatem także w
prawdopodobieństwie. Jednoznaczność granicy w
prawdopodobieństwie daje $\Prob(\widetilde X_t=X_t)=1$.
Ważne, że argument dotyczy osobno każdej ustalonej chwili;
nie zakładamy z góry nierozróżnialności.

\emph{Krok 4: wszystkie wykładniki mniejsze od $\beta/p$.}
Wybierzmy rosnący ciąg liczb
$0<\gamma_j<\beta/p$ taki, że
$\gamma_j\uparrow\beta/p$. Dla każdego $j$
powtórzmy kroki 1--2 z wykładnikiem $\gamma_j$.
Każdy z tych argumentów daje zbiór pełnej miary,
na którym ograniczenie procesu do \emph{tego samego}
zbioru $D$ ma ciągłe przedłużenie.
Przecinamy te przeliczalnie liczne mierzalne zbiory
pełnej miary również ze zbiorem $\Omega_0$ z kroku 2.
Na otrzymanym zbiorze jednoznaczność ciągłego
przedłużenia z $D$ zapewnia, że już skonstruowana
modyfikacja $\widetilde X$ ma trajektorie
$\gamma_j$-h"olderowskie dla wszystkich $j$.
Jeśli $0<\gamma<\beta/p$, wybieramy $j$ z
$\gamma<\gamma_j$ i wykorzystujemy
$|t-s|^{\gamma_j}\le T^{\gamma_j-\gamma}|t-s|^\gamma$.
To kończy dowód.
\end{proof}

\begin{uwaga}[Znaczenie wykładnika $1+\beta$]
Przy siatce dyadycznej liczba sąsiednich par jest rzędu
$2^n$, czyli rośnie odwrotnie proporcjonalnie do długości
kroku. Moment rzędu $p$ pojedynczego przyrostu jest
ograniczony z góry przez $CT^{1+\beta}2^{-n(1+\beta)}$. Po zsumowaniu po wszystkich parach
zyskujemy więc czynnik $2^{-n\beta}$.
Właśnie ta \emph{nadwyżka wykładnika ponad $1$}
pozwala zastosować Borela--Cantellego i otrzymać
jedno wspólne zdarzenie pełnej miary, kontrolujące
nieskończenie wiele chwil.
\end{uwaga}

\begin{przyklad}[Jak wybrać moment w kryterium Kołmogorowa?]
Załóżmy, że dla pewnego $p=4$ i $s,t\in[0,T]$
zachodzi $\E|X_t-X_s|^4\le C|t-s|^2$.
W zapisie twierdzenia mamy $1+\beta=2$, zatem $\beta=1$
i $\beta/p=1/4$. Otrzymujemy modyfikację o trajektoriach
h"olderowskich każdego rzędu $\gamma<1/4$.
Nie twierdzimy, że wykładnik $1/4$ jest optymalny:
może się okazać, że dysponujemy mocniejszymi
oszacowaniami dla wyższych momentów, prowadzącymi do
większych wykładników.
\end{przyklad}

\section{Konstrukcja procesu Wienera: pierwszy kontakt}
Rozdziały poświęcone ruchowi Browna i mierze Wienera
rozwiną tę konstrukcję, ale teraz możemy ją już wykonać
matematycznie poprawnie, korzystając z dwóch udowodnionych
twierdzeń. Zauważmy, że nie potrzebujemy żadnej całki
stochastycznej.

\begin{twierdzenie}[Istnienie standardowego procesu Wienera]
\label{thm:wiener-exists}
Istnieją przestrzeń probabilistyczna i proces rzeczywisty
$W=(W_t)_{t\ge0}$ takie, że:
\begin{enumerate}
\item $W_0=0$ prawie na pewno;
\item $W$ ma niezależne przyrosty;
\item dla $0\le s<t$ zachodzi
$W_t-W_s\sim\Normal(0,t-s)$;
\item trajektorie $t\mapsto W_t$ są ciągłe prawie na pewno.
\end{enumerate}
Można ponadto wybrać wspólny zbiór pełnej miary,
na którym trajektorie są $\gamma$-h"olderowskie
na każdym przedziale $[0,T]$ dla każdego $0<\gamma<1/2$.
\end{twierdzenie}
\begin{proof}
\emph{Krok 1: zgodne rozkłady gaussowskie.}
Dla $s,t\ge0$ niech $K(s,t)=\min(s,t)$.
Jest to funkcja dodatnio półokreślona, ponieważ
\[
\sum_{i,j=1}^m a_ia_j\min(t_i,t_j)
=\int_0^\infty
\left(\sum_{i=1}^m a_i\chi_{[0,t_i]}(u)\right)^2
\,\mathrm d\lambda_1(u)\ge0.
\]
Z przykładu o procesach gaussowskich i twierdzenia
\ref{thm:kol-ext} istnieje proces gaussowski $X$
o zerowej średniej i kowariancji $K$.
Ponieważ $K(0,0)=0$, $X_0=0$ prawie na pewno.

\emph{Krok 2: niezależne przyrosty i ich rozkłady.}
Dla $0\le s<t$ wariancja przyrostu wynosi
\[
\Var(X_t-X_s)=K(t,t)+K(s,s)-2K(s,t)=t-s.
\]
Jest to zmienna gaussowska o średniej $0$.
Dla dwóch rozłącznych przedziałów $(a,b]$ i $(c,d]$,
gdzie $a<b\le c<d$, obliczamy
\[
\Cov(X_b-X_a,X_d-X_c)
=K(b,d)-K(b,c)-K(a,d)+K(a,c)
=b-b-a+a=0.
\]
Wektor wszystkich przyrostów jest gaussowski, ponieważ
jest liniowym obrazem skończonego wektora gaussowskiego.
Jego macierz kowariancji jest diagonalna dla
parami rozłącznych przedziałów, więc funkcja charakterystyczna
wektora jest iloczynem funkcji charakterystycznych jego
współrzędnych. Z jednoznaczności funkcji charakterystycznych
wynika, że te przyrosty są \emph{łącznie niezależne}; sama parami zerowa kowariancja
nie wystarczyłaby dla dowolnych rozkładów, lecz wystarcza
dla wektora gaussowskiego.

\emph{Krok 3: ciągła modyfikacja.}
Jeśli $Z\sim\Normal(0,1)$, to dla każdego $p>0$
stała $c_p:=\E|Z|^p$ jest skończona, gdyż
$|z|^p\exp(-z^2/2)$ jest całkowalne na $\R$
(w nieskończoności wykładniczy zanik dominuje wzrost
potęgowy, a w otoczeniu zera funkcja jest ograniczona).
Z rozkładu przyrostów wynika
\[
\E|X_t-X_s|^p=c_p|t-s|^{p/2}.
\]
Dla $p>2$ wykładnik $p/2$ jest większy niż $1$.
Kryterium \ref{thm:kol-cont} na $[0,T]$ daje
ciągłą modyfikację o trajektoriach h"olderowskich
każdego rzędu $\gamma<1/2-1/p$.
Aby otrzymać jeden proces na $[0,\infty)$, stosujemy
konstrukcję na przedziałach $[0,n]$, $n\in\N$.
Na wspólnym zbiorze pełnej miary przedłużenia
zgadzają się na wspólnych punktach wymiernych,
bo wszystkie są modyfikacjami tego samego $X$.
Z ciągłości zgadzają się więc na całych częściach wspólnych
przedziałów. Oznaczmy ten mierzalny zbiór pełnej miary przez $G$,
a skonstruowane modyfikacje na $[0,n]$ przez $X^{(n)}$.
Dla $\omega\in G$ oraz $t\ge0$ definiujemy
$W_t(\omega)=X_t^{(n)}(\omega)$ dla dowolnego całkowitego
$n\ge\max(1,t)$, a dla $\omega\notin G$ kładziemy
$W_t(\omega)=0$. Zgodność na częściach wspólnych
zapewnia niezależność wyboru $n$ i ciągłość wszystkich
trajektorii. Dla ustalonego $t$ można wybrać deterministyczne
$n$, więc $W_t$ jest mierzalna; ponieważ $G$ ma miarę $1$
i $X_t^{(n)}=X_t$ prawie na pewno, $W$ jest modyfikacją $X$.

\emph{Krok 4: zachowanie rozkładów i wykładniki.}
Modyfikacja nie zmienia żadnego rozkładu
skończeniewymiarowego: dla skończonej liczby chwil
przecięcie zbiorów, na których $W_t=X_t$, ma miarę $1$.
Stąd $W$ zachowuje wszystkie własności z kroków 1--2.
Dla $n\ge1$ i $k\ge1$ wybierzmy $p_k=k+2$ i zastosujmy
kryterium ciągłości do ograniczenia $X$ na $[0,n]$.
Otrzymana modyfikacja jest nierozróżnialna z ograniczeniem
$W$ do $[0,n]$, na mocy twierdzenia o ciągłych
modyfikacjach z pierwszego rozdziału. Dla każdej pary
$(n,k)$ istnieje więc mierzalny zbiór pełnej miary,
na którym trajektoria $W$ jest $\gamma$-h"olderowska
na $[0,n]$ dla wszystkich $0<\gamma<1/2-1/p_k$.
Przecięcie tych przeliczalnie licznych zbiorów ma miarę $1$.
Dla dowolnych $T>0$ i $0<\gamma<1/2$ wybieramy
$n\ge T$ oraz $k$ tak duże, że $\gamma<1/2-1/p_k$,
i ograniczamy odpowiednie oszacowanie do $[0,T]$.
\end{proof}

\begin{uwaga}[Co dokładnie skonstruowaliśmy?]
Proces Wienera jest tutaj \emph{wnioskiem} z twierdzenia
o istnieniu rozkładów i kryterium ciągłości.
Jego pełną teorię -- zasadę odbicia, czasy trafienia,
własność Markowa, nieróżniczkowalność, wariację
kwadratową, most Browna i miarę Wienera --
przedstawimy w kolejnych rozdziałach tomu II.
Żadna z tych dodatkowych własności nie została
jeszcze użyta w konstrukcji.
\end{uwaga}

\begin{center}
\begin{tikzpicture}[>=Latex,scale=1]
\draw[->] (0,0)--(10.4,0) node[right] {$t$};
\draw[->] (0,-1.7)--(0,1.9) node[above] {$W_t$};
\draw[gray!40,step=0.5] (0,-1.5) grid (10,1.5);
\draw[very thick,blue!65!black]
(0,0)--(0.45,0.35)--(0.85,0.06)--(1.1,0.43)--(1.45,0.11)--
(1.8,0.57)--(2.15,0.30)--(2.35,0.74)--(2.65,0.28)--(3,0.46)--
(3.3,-0.05)--(3.62,0.37)--(3.85,0.04)--(4.1,-0.50)--(4.32,-0.22)--
(4.58,-0.84)--(4.82,-0.50)--(5.10,-0.95)--(5.34,-0.63)--(5.62,-1.12)--
(5.88,-0.61)--(6.12,-0.74)--(6.36,-0.25)--(6.60,-0.55)--(6.92,0.05)--
(7.20,-0.30)--(7.46,0.19)--(7.74,-0.12)--(8.02,0.30)--(8.28,0.02)--
(8.54,0.42)--(8.81,0.05)--(9.1,0.53)--(9.42,0.31)--(9.75,0.81);
\end{tikzpicture}
\par\smallskip\small Poglądowa łamana, a nie numerycznie wygenerowana realizacja procesu Wienera.
\end{center}

\section{Zadanie rozwiązane: dlaczego ciągłość wymaga osobnego dowodu?}
\begin{zadanie}[Proces o jednym losowym skoku]
Niech $U\sim\operatorname{U}([0,1])$ i
$X_t=\chi_{\{U\le t\}}$ dla $t\in[0,1]$.
\emph{(a)} Obliczyć $\Prob(X_s=0,X_t=1)$ dla $0\le s<t\le1$.
\emph{(b)} Wykazać stochastyczną ciągłość.
\emph{(c)} Sprawdzić, czy kryterium Kołmogorowa może
zostać zastosowane do procesu $X$.
\emph{(d)} Ustalić, czy istnieje ciągła modyfikacja.

\emph{Rozwiązanie.}
(a) Zdarzenie $\{X_s=0,X_t=1\}$ jest dokładnie
$\{s<U\le t\}$, a zatem ma prawdopodobieństwo $t-s$.
Podobnie $\Prob(X_s=1,X_t=0)=0$, co odzwierciedla
niemalejące trajektorie.

(b) Dla $s<t$ mamy $|X_t-X_s|\in\{0,1\}$ i
$\Prob(|X_t-X_s|=1)=t-s$. W konsekwencji dla każdego
$0<\varepsilon<1$ zachodzi
$\Prob(|X_t-X_s|>\varepsilon)=|t-s|\to0$.
Dla $\varepsilon\ge1$ prawdopodobieństwo wynosi $0$.
To dowodzi ciągłości stochastycznej.

(c) Dla każdego $p>0$ zachodzi
$\E|X_t-X_s|^p=|t-s|$, ponieważ zmienna
$|X_t-X_s|$ przyjmuje tylko wartości $0,1$.
Nie istnieją więc $C<\infty$ i $\beta>0$ takie, że
$|t-s|\le C|t-s|^{1+\beta}$ dla wszystkich $s,t$:
dzieląc przez $|t-s|>0$ i posyłając $t\to s$,
dostalibyśmy $1\le0$. Założenie kryterium
Kołmogorowa nie jest spełnione.

(d) Dla każdej realizacji z $U\in(0,1)$ funkcja
$t\mapsto X_t$ ma skok. Sam ten fakt nie wystarcza,
by odrzucić możliwość \emph{innej} modyfikacji.
Jednak każda modyfikacja musi zgadzać się z $X$
we wszystkich chwilach wymiernych poza jednym zbiorem
zerowym. Po obu stronach $U$ wartości w chwilach
wymiernych są odpowiednio równe $0$ i $1$; nie
mogą więc być wartościami funkcji ciągłej.
Ciągła modyfikacja nie istnieje.
\end{zadanie}


\chapter{Filtracje i obserwacja procesów w czasie ciągłym}
\label{ch:filtracje}

\section{Informacja: co wolno wiedzieć w chwili \texorpdfstring{$t$}{t}?}
Proces $X=(X_t)_{t\ge0}$ opisuje, co dzieje się z układem.
Do zdefiniowania decyzji zależnej od przeszłości potrzebujemy jednak
osobnego pojęcia \emph{informacji dostępnej w danej chwili}.
Obserwator widzący ruch cząstki do czasu $t$ może obliczyć jej
aktualne położenie, maksymalną dotychczasową odległość lub moment
pierwszego przekroczenia progu. Nie powinien natomiast wiedzieć,
czy cząstka przekroczy próg \emph{jutro}.

Matematycznym modelem tego rozróżnienia jest rosnąca rodzina
$\sigma$-ciał: im późniejsza chwila, tym więcej zdarzeń umiemy
rozpoznać na podstawie dostępnej informacji.

\begin{definicja}[Filtracja i przestrzeń filtrowana]
Niech $(\Omega,\calF,\Prob)$ będzie przestrzenią
probabilistyczną, a $T=[0,\infty)$ lub $T=[0,T_0]$
przy $T_0>0$. \emph{Filtracją} nazywamy rodzinę
$\mathbb F=(\calF_t)_{t\in T}$ pod-$\sigma$-ciał $\calF$
spełniającą $\calF_s\subseteq\calF_t$ dla wszystkich
$s,t\in T$ z $s\le t$. Układ
$(\Omega,\calF,\mathbb F,\Prob)$ nazywamy
\emph{przestrzenią probabilistyczną z filtracją}.
\end{definicja}

\begin{definicja}[Filtracja naturalna i adaptacja]
Niech $X=(X_t)_{t\in T}$ będzie procesem o stanach w
$(E,\calE)$. Jego \emph{filtracją naturalną} jest
$\calF_t^X:=\sigma(X_s:0\le s\le t)$.
Mówimy, że $X$ jest \emph{adaptowany} do filtracji
$\mathbb F$, jeżeli $X_t$ jest $\calF_t/\calE$-mierzalny
dla każdego $t\in T$.
\end{definicja}

\begin{stwierdzenie}
Dla każdego procesu $X$ z czasem $T=[0,\infty)$
jego filtracja naturalna $\mathbb F^X$ jest filtracją,
$X$ jest do niej adaptowany, a jeżeli $\mathbb F$
jest dowolną filtracją, do której $X$ jest adaptowany,
to $\calF_t^X\subseteq\calF_t$ dla każdego $t\ge0$.
\end{stwierdzenie}
\begin{proof}
Jeżeli $s\le t$, to rodzina zmiennych
$\{X_u:0\le u\le s\}$ jest podrodziną
$\{X_u:0\le u\le t\}$, a więc
$\calF_s^X\subseteq\calF_t^X$.
Z definicji $X_t$ jest mierzalna względem
$\calF_t^X$. Dla ostatniej własności
adaptacja do $\mathbb F$ oznacza, że wszystkie
$X_u$, $u\le t$, są $\calF_u$-mierzalne,
a przez monotoniczność filtracji również
$\calF_t$-mierzalne. Zatem generowane przez nie
$\sigma$-ciało zawiera się w $\calF_t$.
\end{proof}

\begin{przyklad}[Losowe obserwacje w chwilach całkowitych]
Niech $\xi_1,\xi_2,\ldots$ będą niezależnymi rzutami symetryczną
monetą, o wartościach $\pm1$, i niech
$S_n=\sum_{k=1}^n\xi_k$, $S_0=0$.
Możemy rozważać $X_t=S_{\lfloor t\rfloor}$ dla $t\ge0$.
Filtracja $\calF_t:=\sigma(\xi_1,\ldots,\xi_{\lfloor t\rfloor})$
(z trywialnym $\calF_t$ dla $t<1$) opisuje poznawanie
wyników kolejnych rzutów. Dla $1\le t<2$ widzimy
pierwszy rzut, ale nie drugi. Proces $X$ jest adaptowany.
Zmienna $\xi_2$ nie jest $\calF_1$-mierzalna:
gdyby była, zdarzenie $\{\xi_2=1\}$ należałoby do
$\sigma(\xi_1)$, choć jest od niej niezależne i ma
prawdopodobieństwo $1/2$. Zdarzenie niezależne od
samego siebie miałoby prawdopodobieństwo $0$ lub $1$,
co daje sprzeczność.
\end{przyklad}

\begin{center}
\begin{tikzpicture}[>=Latex]
\draw[->] (0,0)--(12,0) node[right] {$t$};
\foreach \x/\lab in {1/0,4/1,7/2,10/3}{
  \draw (\x,0.12)--(\x,-0.12) node[below] {$\lab$};
}
\draw[very thick,blue!70!black] (1,1.1)--(4,1.1);
\draw[very thick,blue!70!black] (4,1.55)--(7,1.55);
\draw[very thick,blue!70!black] (7,0.68)--(10,0.68);
\draw[blue!70!black,fill=white,thick] (4,1.1) circle (2pt) (7,1.55) circle (2pt);
\fill[blue!70!black] (1,1.1) circle (2pt) (4,1.55) circle (2pt) (7,0.68) circle (2pt);
\node[left,blue!65!black] at (1,1.12) {$X_t$};
\node[draw,rounded corners,fill=green!8,align=center] at (2.5,2.45) {$\calF_0$\\bez rzutów};
\node[draw,rounded corners,fill=green!8,align=center] at (5.5,2.45) {$\calF_1$\\po I rzucie};
\node[draw,rounded corners,fill=green!8,align=center] at (8.5,2.45) {$\calF_2$\\po II rzucie};
\draw[->,green!45!black] (3.55,2.45)--(4.4,2.45);
\draw[->,green!45!black] (6.55,2.45)--(7.4,2.45);
\node[below=18pt,align=center] at (5.9,-0.2) {informacja rośnie, a przyszły przyrost\\nie jest znany przed wykonaniem rzutu};
\end{tikzpicture}
\par\smallskip\small Rozróżnienie: wykres trajektorii oraz zbiory informacji $\calF_t$.
\end{center}

\section{Zupełność i prawostronna ciągłość filtracji}
\begin{definicja}[Warunki zwykłe]
Niech $\mathbb F=(\calF_t)_{t\ge0}$ będzie filtracją
na $(\Omega,\calF,\Prob)$. Mówimy, że filtracja jest
\emph{zupełna}, jeśli $\calF_0$ zawiera wszystkie
podzbiory zbiorów $N\in\calF$ z $\Prob(N)=0$
(przyjmujemy, że przestrzeń bazowa jest odpowiednio uzupełniona).
Mówimy, że jest \emph{prawostronnie ciągła}, jeśli
dla każdego $t\ge0$ zachodzi
\[
\calF_t=\calF_{t+}:=\bigcap_{u>t}\calF_u.
\]
Filtracja spełnia \emph{warunki zwykłe}, jeśli jest
zupełna i prawostronnie ciągła.
\end{definicja}

Skąd prawostronna ciągłość? Jeśli informacja, że pewne
zdarzenie zaszło, jest dostępna w każdej chwili
$t+\varepsilon$ dla dowolnie małego $\varepsilon>0$,
to naturalnie chcielibyśmy traktować ją jako dostępną
już w chwili $t$. To założenie ma realne konsekwencje:
pomaga wykazać, że czasy pierwszego wejścia do
zbiorów otwartych są czasami zatrzymania.
Nie jest jednak automatyczną własnością każdej filtracji.

\begin{stwierdzenie}[Prawostronnie ciągłe powiększenie]
Niech $\mathbb F^0=(\calF_t^0)_{t\ge0}$ będzie
filtracją na pełnej przestrzeni probabilistycznej
$(\Omega,\calF,\Prob)$, a $\calN$ rodziną podzbiorów
zdarzeń zerowych z $\calF$. Niech
$\calH_t:=\sigma(\calF_t^0\cup\calN)$ i
$\calG_t:=\bigcap_{u>t}\calH_u$.
Wówczas $(\calG_t)_{t\ge0}$ jest zupełną,
prawostronnie ciągłą filtracją oraz
$\calF_t^0\subseteq\calG_t$ dla każdego $t$.
\end{stwierdzenie}
\begin{proof}
Rodzina $\calH_t$ jest rosnąca, dlatego również
$\calG_t$ jest rosnąca. Każde $\calG_t$ zawiera
$\calH_t$, ponieważ $\calH_t\subseteq\calH_u$
dla $u>t$, a zatem zawiera $\calF_t^0$ i $\calN$.
Dla $t\ge0$ obliczamy
\[
\bigcap_{s>t}\calG_s
=\bigcap_{s>t}\bigcap_{u>s}\calH_u
=\bigcap_{u>t}\calH_u=\calG_t.
\]
Środkowa równość wynika stąd, że dla każdego $u>t$
można wybrać $s$ z $t<s<u$, a warunki indeksowane
przez wszystkie takie pary $(s,u)$ są więc dokładnie
warunkami indeksowanymi przez $u>t$.
To dowodzi prawostronnej ciągłości i zupełności.
\end{proof}

\begin{uwaga}
Uzupełnienie i prawostronne domknięcie filtracji
nie zmienia samego procesu $X_t$ jako rodziny funkcji.
Może jednak zmienić, jakie zdarzenia uważa się za
obserwowalne w chwili $t$. Dlatego przy twierdzeniach
odwołujących się do warunków zwykłych trzeba je
wymienić jawnie, a nie zakładać milcząco.
\end{uwaga}


\begin{przyklad}[Filtracja naturalna nie musi być prawostronnie ciągła]
Niech $Z$ będzie niedegenerowaną zmienną Bernoulliego,
a $X_0=0$, $X_t=Z$ dla każdego $t>0$.
Wówczas $\calF_0^X=\{\varnothing,\Omega\}$, lecz
$\calF_t^X=\sigma(Z)$ dla każdego $t>0$, więc
\[
\calF_{0+}^X=\bigcap_{u>0}\calF_u^X
=\sigma(Z)\ne\calF_0^X.
\]
Wynik losowania $Z$ zostaje ujawniony natychmiast
\emph{po} chwili zerowej, choć w chwili $0$ nie jest
jeszcze obserwowalny. Prawostronne domknięcie filtracji
przenosi tę informację do chwili $0$.
Nie twierdzimy tutaj, że $X$ ma prawostronnie ciągłe
trajektorie w zerze: właśnie skok przy $t=0$ umożliwia
tak prosty kontrprzykład.
\end{przyklad}

\section{Mierzalność wspólna, progresywna i przewidywalna}
\begin{definicja}[Mierzalność procesu]
Niech $E$ będzie przestrzenią polską,
$X:[0,\infty)\times\Omega\to E$ oraz
$\mathbb F=(\calF_t)_{t\ge0}$ filtracją.
\begin{enumerate}
\item $X$ jest \emph{wspólnie mierzalny}, jeśli
$(t,\omega)\mapsto X_t(\omega)$ jest mierzalny względem
$\calB([0,\infty))\otimes\calF$.
\item $X$ jest \emph{progresywnie mierzalny}, jeśli dla każdego
$T>0$ jego ograniczenie do $[0,T]\times\Omega$
jest mierzalne względem
$\calB([0,T])\otimes\calF_T$.
\item \emph{Przewidywalne $\sigma$-ciało} $\calP$ na
$[0,\infty)\times\Omega$ jest $\sigma$-ciałem generowanym
przez wszystkie rzeczywiste procesy adaptowane,
których trajektorie są lewostronnie ciągłe na $(0,\infty)$
(ich wartości w $0$ są $\calF_0$-mierzalne).
Proces $E$-wartościowy nazywamy \emph{przewidywalnym},
jeśli jest mierzalny względem $\calP$ i $\calB(E)$.
\end{enumerate}
\end{definicja}

Znaczenie przewidywalności jest operacyjne: decyzja
realizowana w krótkim przedziale \emph{po chwili $t$}
powinna być wyznaczona informacją dostępną \emph{do}
chwili $t$, a nie skokiem, który nastąpi dopiero w tym
przedziale. Dlatego procesy proste w rachunku Itô
będą zbudowane z wartości znanych na lewych końcach
przedziałów.

\begin{twierdzenie}[Regularne trajektorie dają mierzalność progresywną]
\label{thm:prog}
Niech $E$ będzie przestrzenią polską,
$\mathbb F=(\calF_t)_{t\ge0}$ filtracją, a
$X=(X_t)_{t\ge0}$ procesem adaptowanym.
Jeśli każda trajektoria $t\mapsto X_t(\omega)$
jest prawostronnie ciągła, to $X$ jest progresywnie
mierzalny. To samo zachodzi przy lewostronnej
ciągłości na $(0,\infty)$ oraz dowolnej wartości w
$t=0$. W szczególności proces adaptowany o ciągłych
trajektoriach jest progresywnie mierzalny i przewidywalny.
\end{twierdzenie}
\begin{proof}
Ustalmy $T>0$. Dla każdej liczby naturalnej $n$
rozpatrzmy podział $[0,T]$ o węzłach
$t_k=kT/n$, $k=0,\ldots,n$. Jeśli trajektorie są
prawostronnie ciągłe, zdefiniujmy funkcję schodkową
\[
X^{(n)}_s(\omega)=
\begin{cases}
X_{t_k}(\omega),&t_{k-1}<s\le t_k,\quad 1\le k\le n,\\
X_0(\omega),&s=0.
\end{cases}
\]
Każda funkcja $X^{(n)}$ jest mierzalna względem
$\calB([0,T])\otimes\calF_T$, ponieważ jest sumą
skończenie wielu funkcji o stałej po czasie wartości
$X_{t_k}$ (w przypadku ogólnego $E$ oznacza to
funkcję sklejoną na skończonym podziale mierzalnym).
Dla każdego $s<T$ prawy koniec odpowiedniej komórki
zbiega do $s$ od strony prawej, toteż przez
prawostronną ciągłość $X^{(n)}_s\to X_s$;
dla $s=T$ równość zachodzi dokładnie, a dla $s=0$
równość jest w definicji. Granica punktowa funkcji
mierzalnych o wartościach w przestrzeni metrycznej
jest mierzalna, więc $X$ jest progresywnie mierzalny.

Jeżeli trajektorie są lewostronnie ciągłe na $(0,T]$,
zamiast prawych końców używamy lewych końców
komórek: $X^{(n)}_s=X_{t_{k-1}}$ dla
$t_{k-1}\le s<t_k$, a w chwili $T$ ustawiamy
$X^{(n)}_T=X_T$. Dla $s>0$ czasy aproksymacji
zbiegają do $s$ z lewej, a w chwili $0$ mamy
$X^{(n)}_0=X_0$. To dowodzi drugiej wersji.
Proces ciągły adaptowany jest z definicji
przewidywalny, ponieważ jest lewostronnie ciągły;
pierwsza część dowodu daje jego progresywną mierzalność.
\end{proof}

\begin{uwaga}[Dlaczego lewa aproksymacja wymaga ostrożności?]
W powyższym dowodzie lewostronnej ciągłości lewy koniec
komórki dla $s\in(t_{k-1},t_k)$ zbiega do $s$ z lewej.
Dla samego punktu siatki $s=t_k$ wartość
aproksymacji może być równa $X_s$, ale nie ma to
znaczenia dla granicy. Procesy jedynie adaptowane,
bez regularności trajektorii, nie muszą być wspólnie,
a tym bardziej progresywnie mierzalne.
\end{uwaga}

\begin{stwierdzenie}[Proste strategie przewidywalne]
Niech $0=t_0<t_1<\cdots<t_m<\infty$, a
$H_0,\ldots,H_{m-1}$ będą rzeczywistymi zmiennymi
losowymi, przy czym $H_k$ jest $\calF_{t_k}$-mierzalna.
Wtedy proces
\[
H_t=H_{-1}\ind{\{0\}}(t)
+\sum_{k=0}^{m-1}H_k\ind{(t_k,t_{k+1}]}(t),
\qquad 0\le t\le t_m,
\]
gdzie $H_{-1}$ jest dowolną $\calF_0$-mierzalną
zmienną losową, jest przewidywalny.
\end{stwierdzenie}
\begin{proof}
Dla $t\in(t_k,t_{k+1}]$ wartość $H_t=H_k$ jest
$\calF_{t_k}\subseteq\calF_t$-mierzalna.
W punkcie $t_k$ ($k\ge1$) wartość z poprzedniego
przedziału wynosi $H_{k-1}$ i też jest mierzalna
względem $\calF_{t_k}$. Proces jest więc adaptowany.
Przy każdym $t>0$ jego wartość jest granicą wartości
z czasów $s\uparrow t$, ponieważ przedziały mają
postać $(t_k,t_{k+1}]$. Trajektorie są lewostronnie
ciągłe, a więc $H$ jest przewidywalny z definicji.
\end{proof}

\begin{uwaga}
Zapis $H_{-1}$ nie oznacza informacji w chwili ujemnej;
służy wyłącznie do oddzielnego określenia wartości w $0$.
Przy konstrukcji całki Itô wybierzemy zwykle $H_0=0$
i pominiemy wkład pojedynczej chwili $0$.
\end{uwaga}

\section{Czasy zatrzymania}
\begin{definicja}[Czas zatrzymania]
Niech $\mathbb F=(\calF_t)_{t\ge0}$ będzie filtracją.
Zmienną losową $\tau:\Omega\to[0,\infty]$
nazywamy \emph{czasem zatrzymania względem $\mathbb F$},
jeżeli dla każdego $t\ge0$ zachodzi
$\{\tau\le t\}\in\calF_t$.
Jeśli dodatkowo $\Prob(\tau<\infty)=1$, to
mówimy, że $\tau$ jest prawie na pewno skończony.
\end{definicja}

Warunek mówi, że w chwili $t$ umiemy rozstrzygnąć,
czy zatrzymanie nastąpiło \emph{najpóźniej do $t$}.
Wymóg ten wyklucza strategie, które wybierają czas
na podstawie wyniku z przyszłości. Nie wymaga natomiast,
aby w chwili $0$ było wiadomo, jaka będzie wartość $\tau$.

\begin{stwierdzenie}[Działania na czasach zatrzymania]
Niech $\sigma,\tau$ będą czasami zatrzymania dla filtracji
$\mathbb F=(\calF_t)_{t\ge0}$. Wtedy
$\sigma\wedge\tau:=\min(\sigma,\tau)$,
$\sigma\vee\tau:=\max(\sigma,\tau)$ oraz
$\tau\wedge a$ dla każdej stałej $a\ge0$ są czasami
zatrzymania. Jeśli $\mathbb F$ jest prawostronnie ciągła,
to również $\inf_{n\ge1}\tau_n$ jest czasem zatrzymania
dla dowolnego ciągu czasów zatrzymania $(\tau_n)_{n\ge1}$.
\end{stwierdzenie}
\begin{proof}
Dla wszystkich $t\ge0$ mamy
$\{\sigma\wedge\tau\le t\}
=\{\sigma\le t\}\cup\{\tau\le t\}$ oraz
$\{\sigma\vee\tau\le t\}
=\{\sigma\le t\}\cap\{\tau\le t\}$.
Oba zdarzenia są w $\calF_t$.
Dla $\tau\wedge a$ stosujemy pierwszy przypadek,
bo stała $a$ jest czasem zatrzymania.

Niech $\rho=\inf_n\tau_n$. Dla dowolnego $u>0$
$\{\rho<u\}=\bigcup_{n\ge1}\{\tau_n<u\}$:
jeśli infimum jest mniejsze od $u$, to pewien
składnik jest mniejszy od $u$.
Każde $\{\tau_n<u\}$ należy do $\calF_u$,
bo jest przeliczalną sumą zbiorów
$\{\tau_n\le q\}$ po wymiernych $q<u$.
Stąd $\{\rho<u\}\in\calF_u$.
Ponieważ
$\{\rho\le t\}=\bigcap_{k\ge1}\{\rho<t+1/k\}$,
to zdarzenie po lewej należy do
$\bigcap_{k\ge1}\calF_{t+1/k}=\calF_t$
przy prawostronnej ciągłości filtracji.
\end{proof}

\begin{twierdzenie}[Czas pierwszego trafienia w zbiór domknięty]
\label{thm:hitting-closed}
Niech $(E,d)$ będzie przestrzenią metryczną,
$F\subseteq E$ niepustym zbiorem domkniętym,
$\mathbb F=(\calF_t)_{t\ge0}$ filtracją, a
$X=(X_t)_{t\ge0}$ procesem $E$-wartościowym,
adaptowanym, którego wszystkie trajektorie są ciągłe.
Definiujemy $\tau_F(\omega)=\inf\{t\ge0:X_t(\omega)\in F\}$,
przy $\inf\varnothing=\infty$.
Wówczas $\tau_F$ jest czasem zatrzymania.
\end{twierdzenie}
\begin{proof}
Dla każdego $t\ge0$ funkcja
$u\mapsto d(X_u(\omega),F)$ jest ciągła na
przedziale zwartym $[0,t]$; odległość od zbioru
domkniętego jest funkcją $1$-Lipschitzowską.
Wobec tego
\[
\{\tau_F\le t\}
=\left\{\inf_{0\le u\le t}d(X_u,F)=0\right\}
=\left\{\inf_{u\in D_t}d(X_u,F)=0\right\},
\]
gdzie $D_t=(\Q\cap[0,t])\cup\{0,t\}$ jest przeliczalny
i gęsty w $[0,t]$ (dla $t=0$ zawiera punkt $0$).
Pierwsza równość wymaga uzasadnienia w obie strony.
Jeżeli proces trafia do $F$ przed $t$ lub w chwili $t$,
infimum odległości wynosi $0$. Jeśli infimum wynosi $0$,
to na mocy zwartości $[0,t]$ jest osiągane w pewnej
chwili $u\in[0,t]$; domkniętość $F$ implikuje
$X_u\in F$, więc $\tau_F\le t$.
Druga równość wynika z ciągłości funkcji odległości
i gęstości $D_t$.
Wszystkie $X_u$ dla $u\in D_t$ są
$\calF_u\subseteq\calF_t$-mierzalne,
więc infimum przeliczalnej rodziny ich odległości
jest $\calF_t$-mierzalne. Zatem
$\{\tau_F\le t\}\in\calF_t$.
\end{proof}

\begin{twierdzenie}[Czas wejścia do zbioru otwartego]
Niech $(E,d)$ będzie przestrzenią metryczną,
$O\subseteq E$ otwartym zbiorem,
$\mathbb F=(\calF_t)_{t\ge0}$ prawostronnie ciągłą filtracją,
a $X=(X_t)_{t\ge0}$ adaptowanym procesem o ciągłych
trajektoriach. Wtedy
$\tau_O:=\inf\{t\ge0:X_t\in O\}$ (z $\inf\varnothing=\infty$)
jest czasem zatrzymania.
\end{twierdzenie}
\begin{proof}
Dla każdego $t>0$ mamy
$\{\tau_O<t\}=\bigcup_{q\in\Q\cap[0,t)}\{X_q\in O\}$.
Jeżeli $\tau_O<t$, istnieje chwila $u<t$ z $X_u\in O$
(z definicji infimum). Ponieważ $O$ jest otwarty,
a trajektoria ciągła, istnieje otoczenie $u$,
w którym trajektoria nadal należy do $O$; zawiera
ono wymierną chwilę $q<t$. Implikacja odwrotna
jest oczywista. Zatem $\{\tau_O<t\}\in\calF_t$.
Dla każdego $t\ge0$ zachodzi
\[
\{\tau_O\le t\}=\bigcap_{n\ge1}\{\tau_O<t+1/n\}
\in\bigcap_{n\ge1}\calF_{t+1/n}=\calF_t,
\]
co dowodzi tezy. Założenie prawostronnej ciągłości
pojawia się dokładnie w ostatnim przejściu.
\end{proof}

\begin{uwaga}[Wersja dla procesu ciągłego tylko prawie na pewno]
Jeśli ciągłość trajektorii zachodzi na mierzalnym
zbiorze pełnej miary $\Omega_0$, to zbiory opisujące
czas trafienia mogą różnić się od zdarzeń mierzalnych
przez podzbiory $\Omega\setminus\Omega_0$. Wystarczy
pracować z ciągłą modyfikacją o trajektoriach ciągłych
wszędzie (zmieniając proces poza $\Omega_0$) albo
założyć zupełność filtracji; wówczas powyższe
argumenty zachowują ważność po uwzględnieniu
zbiorów zerowych.
\end{uwaga}

\section{Informacja dostępna w losowej chwili}
\begin{definicja}[$\sigma$-ciało w czasie zatrzymania]
Niech $\tau$ będzie czasem zatrzymania dla filtracji
$\mathbb F=(\calF_t)_{t\ge0}$ na $(\Omega,\calF,\Prob)$.
Definiujemy
\[
\calF_\tau:=\big\{A\in\calF:\ 
A\cap\{\tau\le t\}\in\calF_t\quad
\text{dla każdego }t\ge0\big\}.
\]
Interpretujemy $\calF_\tau$ jako informację dostępną
w chwili losowej $\tau$.
\end{definicja}

\begin{twierdzenie}[Podstawowe własności $\calF_\tau$]
Niech $\mathbb F$ będzie filtracją i $\tau$ jej czasem
zatrzymania. Wtedy $\calF_\tau$ jest $\sigma$-ciałem.
Jeśli $\sigma$ jest innym czasem zatrzymania i
$\sigma\le\tau$ punktowo, to
$\calF_\sigma\subseteq\calF_\tau$.
Dla każdego $A\in\calF$ zachodzi również równoważność
\[
A\in\calF_\tau
\quad\Longleftrightarrow\quad
A\cap\{\tau=t\}\in\calF_t
\ \text{dla każdego }t\in\N_0,
\]
\emph{jeżeli} $\tau$ przyjmuje wartości w $\N_0\cup\{\infty\}$.
Nie jest wtedy potrzebny osobny warunek na zdarzeniu
$\{\tau=\infty\}$, ponieważ definicja zawiera tylko
przecięcia ze zdarzeniami $\{\tau\le t\}$ dla skończonych $t$.
\end{twierdzenie}
\begin{proof}
Własność $\Omega\in\calF_\tau$ wynika z
$\{\tau\le t\}\in\calF_t$.
Jeśli $A\in\calF_\tau$, to
$A^c\cap\{\tau\le t\}
=\{\tau\le t\}\setminus(A\cap\{\tau\le t\})\in\calF_t$.
Dla ciągu $(A_n)$ elementów $\calF_\tau$ mamy
$(\bigcup_nA_n)\cap\{\tau\le t\}
=\bigcup_n(A_n\cap\{\tau\le t\})\in\calF_t$.
Tak otrzymujemy $\sigma$-ciało.

Jeśli $\sigma\le\tau$ i $A\in\calF_\sigma$, to
$A\cap\{\tau\le t\}
=(A\cap\{\sigma\le t\})\cap\{\tau\le t\}\in\calF_t$,
ponieważ $\{\tau\le t\}\subseteq\{\sigma\le t\}$.
Stąd $A\in\calF_\tau$.

Dla ostatniej równoważności, gdy $\tau$ jest
$\N_0\cup\{\infty\}$-wartościowy, zauważmy, że
$\{\tau=k\}=\{\tau\le k\}\setminus\{\tau\le k-1\}$
(z odpowiednią konwencją dla $k=0$), więc jest w $\calF_k$.
Jeśli $A\in\calF_\tau$, to
$A\cap\{\tau=k\}
=(A\cap\{\tau\le k\})\cap\{\tau=k\}\in\calF_k$.
W drugą stronę, jeśli każde $A\cap\{\tau=k\}$
jest w $\calF_k$, to dla dowolnego $t\ge0$
\[
A\cap\{\tau\le t\}
=\bigcup_{k=0}^{\lfloor t\rfloor}
(A\cap\{\tau=k\})\in\calF_t.
\]
Skończona suma nie obejmuje $\{\tau=\infty\}$,
ponieważ ta część nie należy do $\{\tau\le t\}$
dla skończonego $t$. Nie trzeba narzucać osobnego
warunku na $A\cap\{\tau=\infty\}$.
\end{proof}

\begin{uwaga}[Kiedy $\tau=\infty$ z dodatnim prawdopodobieństwem]
Definicja $\calF_\tau$ nie wymaga $\tau<\infty$
prawie na pewno. Jeżeli jednak chcemy badać zmienną
$X_\tau$, musimy osobno określić, co znaczy $X_\infty$,
albo ograniczyć się do zdarzenia $\{\tau<\infty\}$.
W literaturze często przyjmuje się $\calF=\calF_\infty$,
gdzie $\calF_\infty=\sigma(\bigcup_t\calF_t)$
(po odpowiednim uzupełnieniu). Bez takiego ustalenia
nie należy utożsamiać $\calF_\infty$ z całym $\calF$.
\end{uwaga}

\begin{twierdzenie}[Mierzalność obserwacji w chwili zatrzymania]
\label{thm:sampled}
Niech $E$ będzie przestrzenią polską,
$X=(X_t)_{t\ge0}$ progresywnie mierzalnym procesem
$E$-wartościowym, a $\tau$ czasem zatrzymania
względem $\mathbb F$. Załóżmy $\tau<\infty$
prawie na pewno i że $\calF_0$ zawiera wszystkie
zbiory zerowe (aby można było dowolnie określić
$X_\tau$ na $\{\tau=\infty\}$).
Wtedy zmienna $X_\tau(\omega):=X_{\tau(\omega)}(\omega)$
na $\{\tau<\infty\}$, dowolnie uzupełniona w
$\{\tau=\infty\}$, jest $\calF_\tau$-mierzalna.
\end{twierdzenie}
\begin{proof}
Ustalmy $t\ge0$. Z własności czasu zatrzymania
zmienna $\tau\wedge t$ jest $\calF_t$-mierzalna.
Progresywna mierzalność $X$ oznacza, że odwzorowanie
$(s,\omega)\mapsto X_s(\omega)$ na $[0,t]\times\Omega$
jest mierzalne względem
$\calB([0,t])\otimes\calF_t$.
Odwzorowanie
$\omega\mapsto(\tau(\omega)\wedge t,\omega)$
z $\Omega$ do $[0,t]\times\Omega$ jest
$\calF_t$-mierzalne, stąd
$X_{\tau\wedge t}$ jest $\calF_t$-mierzalna.
Na zdarzeniu $\{\tau\le t\}$ mamy
$X_\tau=X_{\tau\wedge t}$, a więc dla każdego
$B\in\calB(E)$ zachodzi
\[
\{X_\tau\in B\}\cap\{\tau\le t\}
=\{X_{\tau\wedge t}\in B\}\cap\{\tau\le t\}\in\calF_t.
\]
Z definicji oznacza to
$\{X_\tau\in B\}\in\calF_\tau$.
Dla $\tau=\infty$ na zbiorze zerowym uzupełnienie
jest mierzalne przy założonej zupełności filtracji.
\end{proof}

\section{Zatrzymywanie procesów i próbkowanie opcjonalne}
\begin{definicja}[Proces zatrzymany]
Niech $X=(X_t)_{t\ge0}$ będzie procesem $E$-wartościowym,
a $\tau$ czasem zatrzymania. \emph{Procesem zatrzymanym}
nazywamy $X^\tau=(X_{t\wedge\tau})_{t\ge0}$.
W szczególności dla $\tau=\infty$ mamy
$X_{t\wedge\tau}=X_t$ dla wszystkich skończonych $t$.
\end{definicja}

\begin{stwierdzenie}
Niech $X$ będzie progresywnie mierzalny i adaptowany do
filtracji $\mathbb F$, a $\tau$ czasem zatrzymania.
Wtedy proces $X^\tau$ jest adaptowany. Jeśli każda
trajektoria $X$ jest ciągła (odpowiednio prawostronnie
ciągła), to taka sama regularność zachodzi dla $X^\tau$.
\end{stwierdzenie}
\begin{proof}
Dla ustalonego $t$ zmienna $X_{\tau\wedge t}$
jest $\calF_t$-mierzalna na mocy argumentu z dowodu
twierdzenia~\ref{thm:sampled}, ponieważ $\tau\wedge t$
jest $\calF_t$-mierzalna, a proces jest progresywny.
To jest adaptacja. Dla ustalonego $\omega$ funkcja
$t\mapsto t\wedge\tau(\omega)$ jest ciągła.
Złożenie jej z ciągłą trajektorią $u\mapsto X_u(\omega)$
jest ciągłe. Jeśli ta ostatnia jest tylko prawostronnie
ciągła, to przy $t_n\downarrow t$ również
$t_n\wedge\tau\downarrow t\wedge\tau$, więc
zachodzi prawostronna ciągłość złożenia.
\end{proof}

W tomie I udowodniliśmy twierdzenie o zatrzymywaniu
martyngałów w czasie dyskretnym. W czasie ciągłym
przyda się jego odpowiednik dla czasów ograniczonych.
Przytoczymy go tu z dowodem, aby od początku pilnować
różnicy pomiędzy \emph{ograniczeniem czasu} a jedynie
jego prawie pewną skończonością.

\begin{lemat}[Jednostajna całkowalność warunkowych wartości oczekiwanych]
\label{lem:ui-cond}
Niech $Y\in L^1(\Omega,\calF,\Prob)$.
Rodzina $\{\E[Y\mid\calG]:\calG\subseteq\calF
\text{ jest pod-$\sigma$-ciałem}\}$ jest jednostajnie całkowalna.
\end{lemat}
\begin{proof}
Ustalmy $K>0$ i pod-$\sigma$-ciało $\calG$.
Rozpiszmy $Y=Y^{\le}+Y^{>}$, gdzie
$Y^{\le}=Y\ind{\{|Y|\le K/2\}}$ oraz
$Y^{>}=Y\ind{\{|Y|>K/2\}}$.
Dla $Z=\E[Y\mid\calG]$ mamy
$|\E[Y^{\le}\mid\calG]|\le K/2$ oraz
$|Z|\le K/2+\E[|Y^{>}|\mid\calG]$.
Na zdarzeniu $A=\{|Z|>K\}$ otrzymujemy
$\E[|Y^{>}|\mid\calG]>K/2$ i stąd
$|Z|\ind A\le2\E[|Y^{>}|\mid\calG]\ind A$.
Po scałkowaniu i wykorzystaniu własności
warunkowej wartości oczekiwanej dostajemy
\[
\E[|Z|\ind{\{|Z|>K\}}]
\le2\E[|Y|\ind{\{|Y|>K/2\}}]\xrightarrow[K\to\infty]{}0.
\]
Oszacowanie nie zależy od $\calG$, co dowodzi
jednostajnej całkowalności całej rodziny.
\end{proof}

\begin{lemat}[Próbkowanie opcjonalne na skończonej siatce]
\label{lem:disc-optional}
Niech $(M_k)_{k=0}^N$ będzie całkowalnym martyngałem
względem filtracji $(\calH_k)_{k=0}^N$,
a $\sigma,\tau$ czasami zatrzymania o wartościach w
$\{0,1,\ldots,N\}$, spełniającymi $\sigma\le\tau$.
Wtedy $\E[M_\tau\mid\calH_\sigma]=M_\sigma$
prawie na pewno.
\end{lemat}
\begin{proof}
Niech $A\in\calH_\sigma$. Rozpiszmy różnicę
$M_\tau-M_\sigma$ wzdłuż przyrostów:
\[
M_\tau-M_\sigma
=\sum_{k=1}^N\ind{\{\sigma<k\le\tau\}}
(M_k-M_{k-1}).
\]
Dla każdego $k$ zdarzenie
$A\cap\{\sigma<k\le\tau\}$ jest $\calH_{k-1}$-mierzalne.
Rzeczywiście,
$A\cap\{\sigma\le k-1\}\in\calH_{k-1}$
z definicji $\calH_\sigma$, a
$\{\tau\ge k\}=\{\tau\le k-1\}^c\in\calH_{k-1}$.
Własność martyngałowa daje więc
\[
\E[\ind{A\cap\{\sigma<k\le\tau\}}(M_k-M_{k-1})]
=\E\!\left[\ind{A\cap\{\sigma<k\le\tau\}}
\E[M_k-M_{k-1}\mid\calH_{k-1}]\right]=0.
\]
Sumujemy po skończenie wielu $k$ i otrzymujemy
$\E[\ind A M_\tau]=\E[\ind A M_\sigma]$.
Ponieważ $M_\sigma$ jest $\calH_\sigma$-mierzalny
(jest sumą $M_k\ind{\{\sigma=k\}}$),
charakterystyka warunkowej wartości oczekiwanej
daje tezę.
\end{proof}

\begin{twierdzenie}[Opcjonalne próbkowanie w ograniczonym horyzoncie]
\label{thm:optional-cont}
Niech $(\Omega,\calF,\Prob)$ będzie przestrzenią
probabilistyczną z filtracją $(\calF_t)_{0\le t\le T}$,
gdzie $T>0$. Niech $M=(M_t)_{0\le t\le T}$
będzie całkowalnym martyngałem względem tej filtracji
o trajektoriach prawostronnie ciągłych.
Jeśli $\sigma,\tau$ są czasami zatrzymania spełniającymi
punktowo $0\le\sigma\le\tau\le T$,
to $M_\sigma,M_\tau\in L^1$ oraz
\[
\E[M_\tau\mid\calF_\sigma]=M_\sigma
\quad\text{prawie na pewno}.
\]
\end{twierdzenie}
\begin{proof}
\emph{Krok 1: aproksymacja czasów.}
Dla każdego $n$ weźmy siatkę $D_n=\{kT2^{-n}:0\le k\le2^n\}$
i niech $\tau_n$ (odpowiednio $\sigma_n$) będzie
najmniejszym punktem siatki nie mniejszym od $\tau$
(odpowiednio $\sigma$). Wtedy
$\sigma_n\le\tau_n$, $\tau_n\downarrow\tau$,
$\sigma_n\downarrow\sigma$ i oba przybliżenia są
czasami zatrzymania o wartościach w $D_n$.
Istotnie, dla $kT2^{-n}\in D_n$ zachodzi
$\{\tau_n\le kT2^{-n}\}=\{\tau\le kT2^{-n}\}\in\calF_{kT2^{-n}}$,
podobnie dla $\sigma_n$.
Lemat~\ref{lem:disc-optional} daje
$\E[M_{\tau_n}\mid\calF_{\sigma_n}]=M_{\sigma_n}$.

\emph{Krok 2: jednostajna całkowalność.}
Stosując ten sam lemat do czasu $\tau_n$ i
stałego czasu końcowego $T$, otrzymujemy
$M_{\tau_n}=\E[M_T\mid\calF_{\tau_n}]$.
Z lematu~\ref{lem:ui-cond} rodzina
$(M_{\tau_n})$ jest jednostajnie całkowalna;
tak samo $(M_{\sigma_n})$.
Prawostronna ciągłość trajektorii wraz z
$\tau_n\downarrow\tau$, $\sigma_n\downarrow\sigma$
powoduje zbieżności
$M_{\tau_n}\to M_\tau$ i
$M_{\sigma_n}\to M_\sigma$ prawie na pewno.
Twierdzenie Vitaliego (zbieżność prawie na pewno i
jednostajna całkowalność implikują zbieżność w $L^1$)
daje obie zbieżności w $L^1$. W szczególności
zmienne graniczne są całkowalne.

\emph{Krok 3: przejście do warunkowania w losowej chwili.}
Niech $A\in\calF_\sigma$. Ponieważ
$\sigma\le\sigma_n$ punktowo, z poprzedniego
twierdzenia o informacjach w czasach zatrzymania
wynika $A\in\calF_{\sigma_n}$. Z równości
warunkowych wartości oczekiwanych z kroku 1 dostajemy
$\E[\ind A M_{\tau_n}]=\E[\ind A M_{\sigma_n}]$.
Przechodząc do granicy w $L^1$, otrzymujemy
$\E[\ind A M_\tau]=\E[\ind A M_\sigma]$.
Zmienna $M_\sigma$ jest $\calF_\sigma$-mierzalna,
ponieważ $M$ jest progresywnie mierzalny jako
proces adaptowany prawostronnie ciągły oraz na mocy
twierdzenia~\ref{thm:sampled}.
Wobec tego charakterystyka warunkowej wartości
oczekiwanej daje żądane równanie.
\end{proof}

\begin{uwaga}[Przykład pułapki: czas nieograniczony]
Twierdzenie wymaga $\tau\le T$ dla ustalonego $T$.
Dla symetrycznego spaceru losowego
$S_n=\xi_1+\cdots+\xi_n$, $\E S_n=0$,
czas $\tau=\inf\{n\ge0:S_n=1\}$ jest prawie
na pewno skończony, ale $S_\tau=1$ prawie na pewno.
Zatem $\E S_\tau=1\ne0=\E S_0$.
Nie ma sprzeczności: brak ograniczenia $\tau$
uniemożliwia bez dodatkowych hipotez przejście
od $S_{\tau\wedge n}$ do $S_\tau$ w $L^1$.
To samo ostrzeżenie będzie obowiązywać dla
martyngałów w czasie ciągłym.
\end{uwaga}


\begin{twierdzenie}[Dwa martyngały związane z procesem Wienera]
Niech $W=(W_t)_{t\ge0}$ będzie standardowym procesem
Wienera o ciągłych trajektoriach,
a $\calF_t^W=\sigma(W_s:0\le s\le t)$ jego filtracją
naturalną. Wtedy procesy
$M_t=W_t$ oraz $N_t=W_t^2-t$, $t\ge0$,
są całkowalnymi martyngałami względem $\mathbb F^W$.
\end{twierdzenie}
\begin{proof}
Ustalmy $0\le s<t$ i połóżmy
$Z=W_t-W_s$. Wiemy, że $Z\sim\Normal(0,t-s)$
i jest niezależne od wektora
$(W_{u_1},\ldots,W_{u_m})$ dla dowolnych
$u_1,\ldots,u_m\le s$: wynika to z gaussowskości
wspólnego wektora i zerowych kowariancji $Z$
z każdą wcześniejszą współrzędną.
Aby przejść od skończenie wielu współrzędnych do
całego $\sigma$-ciała $\calF_s^W$, stosujemy
twierdzenie $\pi$--$\lambda$: niezależność zachodzi
na $\pi$-układzie zdarzeń cylindrycznych
$\{(W_{u_1},\ldots,W_{u_m})\in B\}$, które generują
$\calF_s^W$, a klasa zdarzeń niezależnych od $Z$
jest $\lambda$-układem. Zatem $Z$ jest niezależne
od $\calF_s^W$.

Ponieważ $\E Z=0$ i $\E Z^2=t-s$, z własności
warunkowej wartości oczekiwanej dostajemy
$\E[Z\mid\calF_s^W]=0$ oraz
$\E[Z^2\mid\calF_s^W]=t-s$.
Z $W_t=W_s+Z$ otrzymujemy
\[
\E[W_t\mid\calF_s^W]=W_s,
\quad
\E[W_t^2-t\mid\calF_s^W]
=W_s^2+2W_s\E[Z\mid\calF_s^W]
+(t-s)-t=W_s^2-s.
\]
Całkowalność wynika ze skończonych momentów
rozkładu normalnego, a adaptacja z definicji
filtracji naturalnej. To są wszystkie warunki
bycia martyngałem.
\end{proof}

\begin{zadanie}[Wyjście procesu Wienera z przedziału]
Niech $W$ będzie procesem Wienera z ciągłymi trajektoriami,
$W_0=0$, a $a>0$. Względem jego filtracji naturalnej
zdefiniujmy
$\tau:=\inf\{t\ge0:|W_t|=a\}$.
Wykażemy, że $\tau<\infty$ prawie na pewno,
obliczymy $\E\tau$ oraz prawdopodobieństwa
wyjścia z przedziału przez jego prawy i lewy koniec.

Najpierw $\tau$ jest czasem zatrzymania,
ponieważ jest czasem trafienia w domknięty zbiór
$\{-a,a\}$ przez ciągły adaptowany proces.
Dla $T>0$ zastosujmy twierdzenie o próbkowaniu
opcjonalnym do ograniczonego czasu
$\tau\wedge T$ i obu martyngałów
$W_t$ i $W_t^2-t$. Otrzymujemy
\[
\E W_{\tau\wedge T}=0,
\qquad \E W_{\tau\wedge T}^2=\E(\tau\wedge T).
\tag{8}
\]
Z ciągłości trajektorii i definicji pierwszego
wyjścia wynika $|W_{\tau\wedge T}|\le a$
punktowo. Wobec tego z (8) mamy
$0\le\E(\tau\wedge T)\le a^2$.
Gdy $T\uparrow\infty$, twierdzenie o zbieżności
monotonicznej daje $\E\tau\le a^2<\infty$.
W szczególności $\tau<\infty$ prawie na pewno:
gdyby $\Prob(\tau=\infty)>0$, całka $\E\tau$
byłaby nieskończona.

Skoro $\tau$ jest skończony i trajektorie są ciągłe,
mamy $|W_\tau|=a$ prawie na pewno.
Ponadto $W_{\tau\wedge T}\to W_\tau$ prawie na pewno
i $|W_{\tau\wedge T}|\le a$, więc przez zbieżność
zdominowaną $\E W_\tau=0$ oraz
$\E W_{\tau\wedge T}^2\to\E W_\tau^2=a^2$.
Przechodząc do granicy w (8), dostajemy
$\E\tau=a^2$.

Na koniec oznaczmy $p=\Prob(W_\tau=a)$.
Ponieważ $W_\tau\in\{-a,a\}$ prawie na pewno,
$0=\E W_\tau=ap-a(1-p)$, skąd $p=1/2$.
Zatem cząstka wychodzi przez każdy z obu końców z
prawdopodobieństwem $1/2$, a oczekiwany czas wyjścia
z przedziału $(-a,a)$ wynosi $a^2$.

Ten przykład ilustruje właściwe użycie twierdzenia
o zatrzymywaniu: \emph{najpierw} zatrzymujemy proces
w chwili $\tau\wedge T$, gdzie czas jest ograniczony,
a dopiero potem uzasadniamy przejście $T\to\infty$
za pomocą oszacowania i odpowiednich twierdzeń
o zbieżności.
\end{zadanie}

\section{Zadanie rozwiązane: informacja, zatrzymanie i uczciwa gra}
\begin{zadanie}[Dwa rzuty monetą i chwila zatrzymania]
Niech $\Omega=\{HH,HT,TH,TT\}$, gdzie litera $H$
oznacza wynik $+1$, a $T$ wynik $-1$.
Wszystkim czterem wynikom przypisujemy prawdopodobieństwo $1/4$.
Niech $\xi_1,\xi_2$ będą odpowiednimi wynikami
rzutów, $\calF_0=\{\varnothing,\Omega\}$,
$\calF_1=\sigma(\xi_1)$, $\calF_2=\sigma(\xi_1,\xi_2)$,
$S_0=0$, $S_1=\xi_1$ i $S_2=\xi_1+\xi_2$.
Zdefiniujmy $\tau=1$ wtedy i tylko wtedy, gdy
$\xi_1=+1$, oraz $\tau=2$ w przeciwnym razie.

\emph{(a) Sprawdzić, że $\tau$ jest czasem zatrzymania.}
Dla $t<1$ zdarzenie $\{\tau\le t\}$ jest puste.
Dla $1\le t<2$ jest równe $\{\xi_1=+1\}$,
a więc należy do $\calF_1\subseteq\calF_t$
(przy stałym przedłużeniu filtracji między rzutami).
Dla $t\ge2$ jest równe $\Omega$. Stąd warunek czasu
zatrzymania jest spełniony.

\emph{(b) Opisać dokładnie $\calF_\tau$.}
Wyniki $HH$ i $HT$ są nierozróżnialne w chwili
zatrzymania: po pierwszym $H$ zatrzymujemy się,
więc drugiego rzutu nie obserwujemy. Gdy pierwszy
rzut da $T$, czekamy do drugiego rzutu i rozróżniamy
$TH$ oraz $TT$. Zatem atomami $\calF_\tau$ są
$\{HH,HT\}$, $\{TH\}$ i $\{TT\}$.
Można to sprawdzić z definicji: warunek
$A\cap\{\tau\le1\}\in\calF_1$ pozwala
albo wziąć obydwa wyniki $HH,HT$, albo żaden;
na $\{\tau=2\}$ nie ma dodatkowego ograniczenia,
bo $\calF_2=2^\Omega$.

\emph{(c) Obliczyć rozkład i średnią $S_\tau$.}
Dla $HH$ i $HT$ mamy $\tau=1$ oraz $S_\tau=1$.
Dla $TH$ mamy $\tau=2$ i $S_\tau=0$,
a dla $TT$ mamy $\tau=2$ i $S_\tau=-2$. Stąd
\[
\Prob(S_\tau=1)=\tfrac12,\qquad
\Prob(S_\tau=0)=\tfrac14,\qquad
\Prob(S_\tau=-2)=\tfrac14,
\]
co daje $\E S_\tau=\tfrac12\cdot1+
\tfrac14\cdot0+\tfrac14\cdot(-2)=0$.
Jest to konkretny przykład próbkowania opcjonalnego:
$S$ jest martyngałem, a $\tau\le2$.

\emph{(d) Co będzie, jeśli wybierzemy czas na podstawie przyszłości?}
Niech $\rho=1$ dla $\xi_2=+1$ oraz $\rho=2$
dla $\xi_2=-1$. Zdarzenie $\{\rho\le1\}
=\{\xi_2=+1\}$ nie należy do $\calF_1$,
więc $\rho$ nie jest czasem zatrzymania.
Mamy $S_\rho=\xi_1+\xi_2\ind{\{\xi_2=-1\}}$,
zatem $\E S_\rho=0-\tfrac12=-\tfrac12$.
Nie narusza to twierdzenia o próbkowaniu opcjonalnym:
strategia używa przyszłego wyniku, czyli jej czas nie
spełnia kluczowego założenia.
\end{zadanie}

\section{Co będzie potrzebne w rachunku Itô?}
Na zakończenie zbierzmy logiczne zależności.
Aby konstruować całkę stochastyczną względem procesu Wienera,
będziemy musieli odróżniać następujące własności:
\begin{itemize}
\item proces jest \emph{adaptowany}, a więc stan w chwili $t$
jest znany przy informacji $\calF_t$;
\item proces jest \emph{progresywnie mierzalny}, co umożliwia
mierzalne podstawianie losowych chwil i zwykłe całkowanie
funkcji $(t,\omega)$;
\item proces jest \emph{przewidywalny}, co pozwala zbudować
całkę z wartości współczynnika znanych przed
przyrostem $W_{t_{k+1}}-W_{t_k}$;
\item proces jest \emph{zatrzymany} w czasie $\tau$,
co pozwala lokalizować oszacowania i badać zjawiska
do losowej chwili, bez wiedzy o przyszłości.
\end{itemize}
W dalszych rozdziałach omówimy proces Wienera
wraz z jego naturalną filtracją, procesy Markowa
w czasie ciągłym, martyngały lokalne, miarę Wienera
na przestrzeni trajektorii i całkę Itô.


% ================================================================
% TOM II. ROZDZIAŁ 4
% ================================================================

\chapter{Proces Poissona: losowe zliczanie zdarzeń}
\label{ch:poisson-process}
\paragraph{Odsyłacze międzytomowe.} Estymację intensywności procesu Poissona i modele Poisson GLM omawia tom IV, rozdziały~\ref{IV-ch:conjugate-families} i~\ref{IV-ch:glm}.


\section{Motywacja: od losowych odstępów do procesu zliczającego}
W poprzednich rozdziałach skonstruowaliśmy procesy jako rodziny
zmiennych losowych, a następnie opisaliśmy informację dostępną w czasie.
Teraz zajmiemy się pierwszym ważnym modelem, w którym \emph{cała
trajektoria} ma prostą interpretację. Liczymy zgłoszenia do serwisu,
przybycia klientów, impulsy z detektora promieniowania albo błędy
transmisji. Interesują nas nie tylko liczby zdarzeń, lecz także czasy
ich wystąpienia, odstępy, zdarzenia po chwili obserwacji i reakcja na
zmianę intensywności.

Dwa pozornie różne opisy mogą prowadzić do tego samego procesu:
\begin{itemize}
\item odstępy między kolejnymi zgłoszeniami są niezależne i wykładnicze;
\item liczby zgłoszeń w rozłącznych przedziałach są niezależne, a liczba
w przedziale długości $t$ ma rozkład Poissona o parametrze $\lambda t$.
\end{itemize}
Udowodnimy równoważność tych opisów, zamiast przyjmować ją jako
nieuzasadnioną własność modelu.

\begin{definicja}[Czasy nadejścia i proces zliczający]
Niech $(\Omega,\calF,\Prob)$ będzie przestrzenią probabilistyczną,
a $E_1,E_2,\ldots$ dodatnimi zmiennymi losowymi. Niech
$T_0:=0$ i $T_n:=\sum_{k=1}^n E_k$, $n\ge1$.
Jeżeli $T_n\to\infty$ prawie na pewno, to dla $t\ge0$ określamy
\[
N_t:=\sum_{n=1}^\infty\ind{\{T_n\le t\}}
      =\max\{n\in\N_0:T_n\le t\}.
\]
Proces $N=(N_t)_{t\ge0}$ jest \emph{procesem zliczającym} zdarzenia
w chwilach $T_1,T_2,\ldots$. Warunek $T_n\to\infty$ oznacza, że na
żadnym ograniczonym przedziale nie występuje nieskończenie wiele
zdarzeń (brak eksplozji).
\end{definicja}

Każda trajektoria $N$ jest niemalejącą funkcją schodkową, prawostronnie
ciągłą i posiadającą granice lewostronne; jej skoki są równe jeden,
gdyż $E_k>0$. Dla przejrzystości wszystkie twierdzenia tego rozdziału
dotyczą czasu $[0,\infty)$.

\begin{center}
\begin{tikzpicture}[>=Latex,x=1.05cm,y=.62cm]
\draw[->] (-.3,0)--(8.1,0) node[right] {$t$};
\draw[->] (0,-.3)--(0,4.8) node[above] {$N_t$};
\foreach \j in {1,2,3,4} {\draw (-.09,\j)--(.09,\j) node[left=4pt] {$\j$};}
\draw[very thick,blue!70!black] (0,0)--(1.05,0)--(1.05,1)--(2.55,1)--(2.55,2)--(3.6,2)--(3.6,3)--(6.2,3)--(6.2,4)--(7.75,4);
\foreach \x/\j in {1.05/1,2.55/2,3.6/3,6.2/4} {
\fill[blue!70!black] (\x,\j) circle (2pt);
\draw[fill=white,blue!70!black] (\x,{\j-1}) circle (2pt);
\draw[densely dashed,gray] (\x,0)--(\x,\j);
\node[below=4pt] at (\x,0) {$T_{\j}$};}
\end{tikzpicture}
\par\smallskip\small Przykładowa trajektoria: $N_t$ liczy zdarzenia w $(0,t]$;
punkty pełne pokazują przyjętą prawostronną ciągłość.
\end{center}

\section{Konstrukcja ze zmiennych wykładniczych}
\begin{definicja}[Jednorodny proces Poissona]
Niech $\lambda>0$. \emph{Jednorodnym procesem Poissona o intensywności
$\lambda$} nazywamy proces zliczający skonstruowany z niezależnych
zmiennych $E_1,E_2,\ldots$ o wspólnym rozkładzie wykładniczym
$\operatorname{Exp}(\lambda)$, czyli
$\Prob(E_k>u)=e^{-\lambda u}$ dla $u\ge0$.
Czasy $T_n$ nazywamy chwilami nadejścia, a $E_n$ odstępami między
nadejściami.
\end{definicja}

Najpierw musimy sprawdzić, że konstrukcja ma sens dla wszystkich
$t\ge0$. Nie wolno poprzestać na intuicji, że ``czasy są dodatnie'':
suma nieskończenie wielu dodatnich liczb może być skończona.

\begin{lemat}[Brak eksplozji dla stałej intensywności]
Jeżeli $(E_n)_{n\ge1}$ są niezależnymi zmiennymi o rozkładzie
$\operatorname{Exp}(\lambda)$, $\lambda>0$, to
$T_n=\sum_{k=1}^n E_k\to\infty$ prawie na pewno.
\end{lemat}
\begin{proof}
Ustalmy $a=1/\lambda$. Wydarzenia $A_n=\{E_n\ge a\}$ są niezależne
i $\Prob(A_n)=e^{-1}>0$. Mamy więc
$\sum_n\Prob(A_n)=\infty$. Drugi lemat Borela--Cantellego daje
$\Prob(A_n\ \text{nieskończenie często})=1$. Na każdym takim
wyniku $\omega$ nieskończenie wiele składników $E_n(\omega)$ jest
nie mniejszych od $a$. Suma częściowa $T_n(\omega)$ jest niemalejąca
i przekracza dowolną ustaloną wielokrotność $a$, a zatem dąży do
nieskończoności.
\end{proof}

\begin{lemat}[Gęstość chwil nadejścia]
Dla $\lambda>0$, $n\ge1$ i niezależnych
$E_1,\ldots,E_n\sim\operatorname{Exp}(\lambda)$ wektor
$(T_1,\ldots,T_n)$ ma gęstość
\[
(s_1,\ldots,s_n)\longmapsto
\lambda^n e^{-\lambda s_n}\ind{\{0<s_1<\cdots<s_n\}}.
\]
W szczególności $T_n$ ma gęstość gamma
$f_{T_n}(t)=\lambda^n t^{n-1}e^{-\lambda t}/(n-1)!$ dla $t>0$.
\end{lemat}
\begin{proof}
Wektor $(E_1,\ldots,E_n)$ ma na dodatnim ortancie gęstość
$\lambda^n\exp(-\lambda\sum_{k=1}^n e_k)$.
Podstawiamy $s_k=e_1+\cdots+e_k$. Odwrotna transformacja ma postać
$e_1=s_1$, $e_k=s_k-s_{k-1}$ dla $k\ge2$ i trójkątny jakobian
równy $1$. Dodatniość wszystkich $e_k$ jest równoważna
$0<s_1<\cdots<s_n$, natomiast suma $e_k$ równa się $s_n$.
Otrzymujemy wskazaną gęstość wspólną. Dla $s_n=t$ całkujemy po
sympleksie $0<s_1<\cdots<s_{n-1}<t$ o objętości
$t^{n-1}/(n-1)!$, co daje gęstość $T_n$.
\end{proof}

\begin{twierdzenie}[Rozkład liczby nadejść i niezależność przyrostów]
\label{thm:poisson-fdd}
Niech $N$ będzie procesem Poissona o intensywności $\lambda>0$.
Wówczas $N_0=0$ oraz dla każdego $t\ge0$ i $n\in\N_0$
\[
\Prob(N_t=n)=e^{-\lambda t}\frac{(\lambda t)^n}{n!}.
\]
Jeżeli $0=t_0<t_1<\cdots<t_m$, to przyrosty
$N_{t_k}-N_{t_{k-1}}$, $k=1,\ldots,m$, są niezależne, a ich
rozkłady to odpowiednio
$\operatorname{Poisson}(\lambda(t_k-t_{k-1}))$.
W szczególności przyrosty są stacjonarne.
\end{twierdzenie}
\begin{proof}
Zdarzenie $\{N_t=n\}$ oznacza $\{T_n\le t<T_{n+1}\}$,
przy $n=0$ rozumiane jako $\{E_1>t\}$.
Dla $n\ge1$ po ustaleniu $0<s_1<\cdots<s_n<t$ prawdopodobieństwo,
że następny niezależny odstęp przekracza $t-s_n$, wynosi
$e^{-\lambda(t-s_n)}$. Wykorzystując poprzedni lemat, otrzymujemy
gęstość wspólną $n$ kolejnych chwil nadejścia \emph{i braku następnego
zdarzenia do chwili $t$}:
\[
\lambda^n e^{-\lambda s_n}e^{-\lambda(t-s_n)}
\ind{\{0<s_1<\cdots<s_n<t\}}
=\lambda^n e^{-\lambda t}\ind{\{0<s_1<\cdots<s_n<t\}}.
\]
Jej całka po sympleksie ma wartość
$e^{-\lambda t}\lambda^n t^n/n!$. Dla $n=0$ wzór wynika
bezpośrednio z rozkładu $E_1$.

Aby udowodnić \emph{niezależność}, nie wystarczy znać każdego
rozkładu brzegowego. Ustalmy liczby całkowite $n_1,\ldots,n_m\ge0$
i $n=\sum_k n_k$. Warunkowo względem $N_{t_m}=n$ chwile
$T_1,\ldots,T_n$ mają na uporządkowanym sympleksie stałą gęstość
$n!/t_m^n$: otrzymujemy ją przez podzielenie wyliczonej przed chwilą
gęstości wspólnej przez $\Prob(N_{t_m}=n)$.
Jest to prawo statystyk pozycyjnych $n$ niezależnych punktów
jednostajnych na $(0,t_m)$. Rzeczywiście, $n$ punktów
nieuporządkowanych ma stałą gęstość $t_m^{-n}$ na
$(0,t_m)^n$, a niemal każdy ich układ ma dokładnie $n!$
różnych uporządkowań. Po uporządkowaniu otrzymujemy więc
gęstość $n!/t_m^n$ na sympleksie. Z kolei przypisanie
$n$ niezależnych punktów do $m$ przedziałów ma
$\frac{n!}{n_1!\cdots n_m!}$ możliwych ustawień
z określonymi licznościami $n_1,\ldots,n_m$; prawdopodobieństwo
każdego ustawienia to iloczyn potęg ich względnych długości.
Dlatego liczby punktów w kolejnych
przedziałach $(t_{k-1},t_k]$ mają warunkowo rozkład wielomianowy:
\[
\Prob\bigl(N_{t_k}-N_{t_{k-1}}=n_k,\ 1\le k\le m
 \mid N_{t_m}=n\bigr)
=\frac{n!}{\prod_k n_k!}\prod_{k=1}^m
 \left(\frac{t_k-t_{k-1}}{t_m}\right)^{n_k}.
\]
Mnożąc tę liczbę przez $e^{-\lambda t_m}(\lambda t_m)^n/n!$
i rozkładając $e^{-\lambda t_m}=\prod_k e^{-\lambda(t_k-t_{k-1})}$,
uzyskujemy
\[
\Prob\bigl(N_{t_k}-N_{t_{k-1}}=n_k,\ 1\le k\le m\bigr)
=\prod_{k=1}^m e^{-\lambda(t_k-t_{k-1})}
  \frac{[\lambda(t_k-t_{k-1})]^{n_k}}{n_k!}.
\]
To dokładnie iloczyn rozkładów brzegowych. Dowód zakończony.
\end{proof}


\begin{twierdzenie}[Dwa równoważne opisy jednorodnego procesu Poissona]
\label{thm:poisson-reverse}
Niech $\lambda>0$, a $N=(N_t)_{t\ge0}$ będzie procesem
o niemalejących trajektoriach c\`adl\`ag, wartościach
w $\N_0$, $N_0=0$ i skokach o wielkości dokładnie $1$.
Załóżmy, że $N$ ma niezależne, stacjonarne przyrosty
oraz $N_t\sim\operatorname{Poisson}(\lambda t)$ dla każdego
$t\ge0$. Wtedy chwile nadejść
$T_k:=\inf\{t\ge0:N_t\ge k\}$, $k\ge1$,
są prawie na pewno skończone, a odstępy
$E_k:=T_k-T_{k-1}$, $T_0=0$, są niezależne i mają
wspólny rozkład $\operatorname{Exp}(\lambda)$.
\end{twierdzenie}
\begin{proof}
\emph{Krok 1: chwile $T_k$ istnieją i są skończone.}
Dla każdego $k\ge1$ oraz $t\ge0$, dzięki niemalejącym
trajektoriom i ich prawostronnej ciągłości,
$\{T_k\le t\}=\{N_t\ge k\}$.
W szczególności $\Prob(T_k>t)=\Prob(N_t<k)$.
Dla ustalonego $k$ prawa strona dąży do zera, gdy
$t\to\infty$: dla $N_t\sim\operatorname{Poisson}(\lambda t)$
na przykład nierówność Czebyszewa daje
$\Prob(N_t<k)\le\lambda t/(\lambda t-k)^2$
dla $\lambda t>k$.
Zatem każdy $T_k$ jest skończony prawie na pewno;
przecinamy te przeliczalnie liczne zdarzenia.

\emph{Krok 2: rozkłady wszystkich skończonych rodzin czasów.}
Prawo wektora $(T_1,\ldots,T_m)$ jest wyznaczone
przez prawdopodobieństwa
\[
\Prob(T_1\le t_1,\ldots,T_m\le t_m)
=\Prob(N_{t_1}\ge1,\ldots,N_{t_m}\ge m),
\qquad t_1,\ldots,t_m\ge0.
\]
Prawa strona zależy wyłącznie od rozkładów
skończeniewymiarowych $N$. Te zaś są jednoznacznie
wyznaczone przez $N_0=0$ oraz rozkłady niezależnych
przyrostów: dla uporządkowanych chwil $0=t_0<t_1<\cdots<t_r$
wektor wartości $N_{t_j}$ jest obrazem wektora
niezależnych przyrostów $N_{t_j}-N_{t_{j-1}}$ przez
odwzorowanie polegające na sumowaniu współrzędnych.

\emph{Krok 3: porównanie z wcześniej skonstruowanym modelem.}
W twierdzeniu~\ref{thm:poisson-fdd} skonstruowaliśmy
proces $\widehat N$ z niezależnych odstępów
$\widehat E_k\sim\operatorname{Exp}(\lambda)$ i wykazaliśmy,
że ma dokładnie takie same rozkłady skończeniewymiarowe jak $N$.
Krok 2 oznacza zatem, że dla każdego $m$
$(T_1,\ldots,T_m)$ i $(\widehat T_1,\ldots,\widehat T_m)$
mają te same rozkłady. Przekształcenie liniowe
$(t_1,\ldots,t_m)\mapsto(t_1,t_2-t_1,\ldots,t_m-t_{m-1})$
pokazuje identyczność praw wektorów
$(E_1,\ldots,E_m)$ i $(\widehat E_1,\ldots,\widehat E_m)$.
Drugi wektor ma rozkład produktowy
$\operatorname{Exp}(\lambda)^{\otimes m}$.
Ponieważ $m$ jest dowolne, wszystkie odstępy $E_k$
są niezależne i jednakowo wykładnicze, co kończy dowód.
\end{proof}

\begin{uwaga}
Kierunek odwrotny nie wymagał dodatkowego założenia o
mocnej własności Markowa w losowych chwilach $T_k$.
Najpierw porównaliśmy \emph{rozkłady skończeniewymiarowe},
a dopiero potem wyprowadziliśmy z nich wspólne prawo
czasów nadejść. To dobry przykład konstrukcyjnego
użycia twierdzenia o jednoznaczności prawa procesu.
\end{uwaga}

Wynika stąd $\E N_t=\Var(N_t)=\lambda t$,
a dla $0\le s\le t$ mamy
$\Cov(N_s,N_t)=\Var(N_s)=\lambda s$: zapisujemy
$N_t=N_s+(N_t-N_s)$ i korzystamy z niezależności przyrostu.

\begin{uwaga}[Co dokładnie znaczy niezależność przyrostów?]
Własność oznacza niezależność \emph{liczb zdarzeń} na rozłącznych
przedziałach, nie niezależność $N_s$ i $N_t$. Przeciwnie,
$N_t\ge N_s$ dla $t\ge s$ i zwykle
$\Cov(N_s,N_t)=\lambda s>0$. Co więcej, dla ustalonych
$t\ge0$ i $h>0$ mamy
$\Prob(N_{t+h}-N_t=0)=e^{-\lambda h}$,
$\Prob(N_{t+h}-N_t=1)=\lambda h+O(h^2)$
oraz $\Prob(N_{t+h}-N_t\ge2)=O(h^2)$ przy $h\downarrow0$.
To wyjaśnia pojęcie \emph{intensywności}: prawdopodobieństwo
jednego nadejścia na bardzo krótkim przedziale jest w pierwszym
przybliżeniu równe $\lambda h$.
\end{uwaga}

\section{Własność braku pamięci, filtracja i martyngał}
\begin{stwierdzenie}[Brak pamięci rozkładu wykładniczego]
Jeżeli $E\sim\operatorname{Exp}(\lambda)$, $\lambda>0$, to
$\Prob(E>s+t\mid E>s)=\Prob(E>t)$ dla wszystkich $s,t\ge0$.
Wśród rozkładów na $[0,\infty)$ o ciągłej funkcji przeżycia,
ściśle dodatniej dla argumentów skończonych i przyjmującej wartość
$1$ w zerze, ta własność charakteryzuje rozkład wykładniczy.
\end{stwierdzenie}
\begin{proof}
Pierwsza równość wynika ze wzoru
$e^{-\lambda(s+t)}/e^{-\lambda s}=e^{-\lambda t}$.
Dla odwrotnego kierunku oznaczmy $g(t)=\Prob(E>t)$.
Z własności warunkowania wynika $g(s+t)=g(s)g(t)$.
Funkcja $h(t):=-\log g(t)$ jest ciągła, spełnia równanie
$h(s+t)=h(s)+h(t)$ dla $s,t\ge0$ i $h(0)=0$.
Stąd najpierw dla argumentów wymiernych, a potem przez ciągłość
dla wszystkich $t\ge0$ otrzymujemy $h(t)=\lambda t$.
Ponieważ $g$ jest nierosnąca, $\lambda\ge0$; aby otrzymać
właściwy rozkład ze skończoną wartością $E$ prawie na pewno,
musimy mieć $\lim_{t\to\infty}g(t)=0$, a więc $\lambda>0$.
\end{proof}

\begin{twierdzenie}[Proces Poissona jest procesem Markowa]
Niech $N$ będzie procesem Poissona o intensywności $\lambda>0$,
a $\calF_t^N:=\sigma(N_s:0\le s\le t)$ jego naturalną
filtracją. Dla $0\le s<t$ i każdego $k\in\N_0$ zachodzi
\[
\Prob(N_t=N_s+k\mid\calF_s^N)
=e^{-\lambda(t-s)}\frac{[\lambda(t-s)]^k}{k!}
\quad\text{prawie na pewno}.
\]
W szczególności rozkład przyszłego stanu warunkowo względem
całej przeszłości zależy wyłącznie od aktualnego $N_s$.
\end{twierdzenie}
\begin{proof}
Z twierdzenia~\ref{thm:poisson-fdd} wynika niezależność
$N_t-N_s$ od wszystkich wektorów
$(N_{u_1},\ldots,N_{u_r})$ z $u_j\le s$:
możemy każdy $N_{u_j}$ zapisać jako sumę przyrostów na podziale
przedziału $[0,s]$ zawierającym wszystkie $u_j$.
Klasa zdarzeń cylindrycznych wyznaczonych przez takie wektory
jest $\pi$-układem generującym $\calF_s^N$.
Przez twierdzenie o jednoznaczności miar (argument
$\pi$--$\lambda$) niezależność rozciąga się na całe
$\calF_s^N$. Zatem rozkład warunkowy przyrostu jest jego
rozkładem bezwarunkowym, wyliczonym wcześniej.
Ponieważ $N_t=N_s+(N_t-N_s)$, dostajemy tezę.
\end{proof}

\begin{wniosek}[Kompensowany proces Poissona]
Dla procesu Poissona o intensywności $\lambda>0$ proces
$M_t:=N_t-\lambda t$ jest kwadratowalnym martyngałem względem
filtracji naturalnej $\mathbb F^N$, a
$\E(M_t-M_s)^2=\lambda(t-s)$ dla $0\le s\le t$.
\end{wniosek}
\begin{proof}
Dla $s\le t$ niezależność przyrostów daje
$\E[N_t-N_s\mid\calF_s^N]=\lambda(t-s)$,
więc $\E[M_t\mid\calF_s^N]=M_s$.
Całkowalność i kwadratowalność wynikają z momentów
rozkładu Poissona. Wreszcie
$\E(M_t-M_s)^2=\Var(N_t-N_s)=\lambda(t-s)$.
\end{proof}


\begin{center}
\begin{tikzpicture}[>=Latex,x=.95cm,y=.78cm]
\draw[->] (0,0)--(10,0) node[right] {$t$};
\foreach \x/\lab in {1/$T_1$,3/$T_2$,6/$T_3$,8/$T_4$}{
  \draw[red!65!black,very thick] (\x,-.16)--(\x,.16);
  \node[above=3pt] at (\x,.17) {\lab};
}
\draw[<->,blue!60!black] (1,-.52)--(3,-.52)
 node[midway,below] {$E_2=T_2-T_1$};
\draw[<->,blue!60!black] (3,-1.1)--(6,-1.1)
 node[midway,below] {$E_3=T_3-T_2$};
\node[align=left] at (5,1.05)
 {czasy nadejść: $T_1,T_2,\ldots$\qquad odstępy: $E_1,E_2,\ldots$};
\end{tikzpicture}
\par\smallskip\small Oś chwil nadejścia. Zliczenia na przedziale $(s,t]$
 zależą od liczby zaznaczonych punktów, a odstępy opisują
 czas pomiędzy kolejnymi punktami.
\end{center}

\section{Czasy nadejść i warunkowanie na liczbie zdarzeń}
\begin{twierdzenie}[Chwila $n$-tego zdarzenia]
Dla $n\ge1$ w procesie Poissona z intensywnością $\lambda>0$
czas $T_n$ ma rozkład gamma o parametrach $n$ i $\lambda$,
przy konwencji, że drugi parametr jest intensywnością. Zachodzi
\[
\Prob(T_n>t)=e^{-\lambda t}\sum_{k=0}^{n-1}
\frac{(\lambda t)^k}{k!},\qquad
\E T_n=\frac n\lambda,\qquad \Var(T_n)=\frac n{\lambda^2}.
\]
\end{twierdzenie}
\begin{proof}
Gęstość została wyprowadzona wcześniej. Zdarzenie
$\{T_n>t\}$ jest identyczne ze zdarzeniem $\{N_t\le n-1\}$;
stąd wzór na funkcję przeżycia. Ponieważ
$T_n=E_1+\cdots+E_n$, niezależność i momenty rozkładu
wykładniczego dają wskazane średnią i wariancję.
\end{proof}

\begin{twierdzenie}[Jednostajność chwil nadejść przy ustalonej liczbie]
\label{thm:uniform-poisson-times}
Niech $t>0$ i $n\ge1$. Warunkowo względem zdarzenia
$\{N_t=n\}$ wektor $(T_1,\ldots,T_n)$ ma taki sam rozkład,
jak uporządkowane statystyki pozycyjne $n$ niezależnych
zmiennych o rozkładzie jednostajnym na $(0,t)$.
W szczególności dla $0<s<t$ i $0\le k\le n$
\[
\Prob(N_s=k\mid N_t=n)
=\binom nk\left(\frac{s}{t}\right)^k
\left(1-\frac{s}{t}\right)^{n-k}.
\]
\end{twierdzenie}
\begin{proof}
W dowodzie twierdzenia~\ref{thm:poisson-fdd} otrzymaliśmy
warunkową gęstość $n!/t^n$ na sympleksie
$0<s_1<\cdots<s_n<t$. Dokładnie taka jest gęstość
statystyk pozycyjnych z jednostajnego losowania
(gęstość wspólna nieuporządkowanego wektora to $t^{-n}$,
a każdy porządek z $n!$ ma jednakowe prawdopodobieństwo).
Liczba wylosowanych punktów należących do $(0,s]$ jest sumą
$n$ niezależnych zmiennych Bernoulliego z parametrem $s/t$,
co daje wzór dwumianowy.
\end{proof}

Warto zauważyć zmianę perspektywy: bez warunkowania liczby
nadejść w rozłącznych przedziałach są niezależne, ale po warunku
$N_t=n$ przestają być niezależne, bo ich suma musi wynosić $n$.

\begin{zadanie}[Cztery zgłoszenia w godzinę]
Zgłoszenia pojawiają się zgodnie z procesem Poissona o
intensywności $\lambda=4$ na godzinę. Obliczyć: (a)
prawdopodobieństwo dokładnie trzech zgłoszeń w pierwszej godzinie;
(b) prawdopodobieństwo, że do chwili $15$ minut nadejdzie
przynajmniej jedno zgłoszenie; (c) prawdopodobieństwo,
że przy dokładnie czterech zgłoszeniach w godzinie dwa z nich
wystąpiły w pierwszych piętnastu minutach;
(d) oczekiwany czas czwartego zgłoszenia.

\emph{Rozwiązanie.} Przyjmujemy czas w godzinach.
(a) Z rozkładu Poissona dostajemy
$\Prob(N_1=3)=e^{-4}4^3/3!$.
(b) Na odcinku $[0,1/4]$ oczekiwana liczba zgłoszeń wynosi
$1$, więc $\Prob(N_{1/4}\ge1)=1-e^{-1}$.
(c) Z twierdzenia~\ref{thm:uniform-poisson-times} liczba
zgłoszeń w pierwszym kwadransie pod warunkiem $N_1=4$
ma rozkład $\operatorname{Bin}(4,1/4)$; prawdopodobieństwo
wynosi $\binom42(1/4)^2(3/4)^2=27/128$.
(d) Suma czterech niezależnych odstępów wykładniczych ma średnią
$4/\lambda=1$ godzinę. Średnia ta nie oznacza oczywiście, że
czwarte zgłoszenie przychodzi dokładnie po godzinie.
\end{zadanie}


\begin{zadanie}[Dokładne chwile nadejść po warunkowaniu]
Klienci przybywają według procesu Poissona o intensywności
$\lambda>0$. Wiemy, że w ciągu dwóch godzin przybyło dokładnie
dwóch klientów. Niech $T_1<T_2$ będą ich chwilami przybycia.
Obliczyć $\Prob(T_1<1<T_2\mid N_2=2)$,
$\E[T_1\mid N_2=2]$ oraz gęstość wspólną $(T_1,T_2)$
przy tym warunku.

\emph{Rozwiązanie.} Warunkowo względem $N_2=2$
dwie chwile są statystykami pozycyjnymi dwóch niezależnych
zmiennych o rozkładzie jednostajnym na $[0,2]$.
Ich uporządkowana gęstość wynosi
$f(s_1,s_2)=\frac{2!}{2^2}=\frac12$ na trójkącie
$0<s_1<s_2<2$ i jest zerowa poza nim. Zdarzenie
$T_1<1<T_2$ oznacza, że jedna chwila znajduje się
w $(0,1)$, a druga w $(1,2)$; dlatego ma warunkowe
prawdopodobieństwo $2(1/2)^2=1/2$.
Dystrybuanta minimum dwóch zmiennych jednostajnych daje
$\Prob(T_1>u\mid N_2=2)=(1-u/2)^2$ dla $0\le u\le2$.
Stąd, całkując funkcję przeżycia,
\[
\E[T_1\mid N_2=2]
=\int_0^2(1-u/2)^2\,\dd u=\frac23.
\]
Wyniki nie zależą od $\lambda$: po ustaleniu liczby nadejść
intensywność nie wpływa na względne położenie ich chwil.
\end{zadanie}

\section{Superpozycja i rozrzedzanie}
\begin{twierdzenie}[Superpozycja]
Niech $N^{(1)},\ldots,N^{(r)}$ będą niezależnymi procesami
Poissona o intensywnościach $\lambda_1,\ldots,\lambda_r>0$.
Wtedy $N_t:=\sum_{j=1}^rN_t^{(j)}$ jest procesem Poissona
o intensywności $\lambda:=\sum_j\lambda_j$.
\end{twierdzenie}
\begin{proof}
Trajektorie sumy są niemalejące i prawostronnie ciągłe.
W ustalonej chwili prawdopodobieństwo jednoczesnego skoku
dwóch niezależnych procesów wynosi zero: warunkowo względem
skończonej liczby skoków pierwszego procesu na przedziale
$[0,m]$, ich chwile mają rozkład z gęstością;
dla każdej deterministycznej chwili prawdopodobieństwo skoku
drugiego procesu jest zero. Sumując po chwilach i po $m\in\N$,
wykluczamy wspólne skoki prawie na pewno.

Dla rozłącznych przedziałów przyrosty sumy są niezależne,
ponieważ wszystkie przyrosty procesów składowych z tych
przedziałów tworzą niezależną rodzinę. Dla pojedynczego
przedziału długości $h$ funkcja tworząca prawdopodobieństwa
przyrostu wynosi
\[
\E z^{\sum_j (N_{t+h}^{(j)}-N_t^{(j)})}
=\prod_j e^{\lambda_j h(z-1)}
=e^{(\sum_j\lambda_j)h(z-1)},\qquad 0\le z\le1.
\]
Jest to funkcja tworząca rozkładu
$\operatorname{Poisson}(\lambda h)$; jednoznaczność
współczynników szeregu potęgowego daje taki właśnie rozkład.
Zatem suma spełnia charakterystykę z twierdzenia
\ref{thm:poisson-fdd}.
\end{proof}

\begin{twierdzenie}[Niezależne rozrzedzanie]
Niech $N$ będzie procesem Poissona o intensywności $\lambda>0$.
Każdy jego skok, niezależnie od innych skoków i od $N$,
oznaczmy jako typu $A$ z prawdopodobieństwem $p\in[0,1]$,
a jako typu $B$ z prawdopodobieństwem $1-p$.
Niech $N_t^A$ oraz $N_t^B$ zliczają odpowiednie typy do chwili $t$.
Wtedy $N^A$ i $N^B$ są \emph{niezależnymi} procesami Poissona
o intensywnościach $\lambda p$ i $\lambda(1-p)$
(z konwencją, że intensywność $0$ oznacza proces zerowy).
\end{twierdzenie}
\begin{proof}
Na ustalonym przedziale długości $h$, warunkowo względem
$N_{t+h}-N_t=n$, liczba oznaczeń $A$ ma rozkład
$\operatorname{Bin}(n,p)$, a liczba $B$ jest dopełnieniem do $n$.
Dla $|u|,|v|\le1$ wspólna funkcja tworząca wynosi
\[
\E\left[u^{N^A_{t+h}-N^A_t}v^{N^B_{t+h}-N^B_t}\right]
=\E[ (pu+(1-p)v)^{N_{t+h}-N_t}]
=e^{\lambda h p(u-1)}e^{\lambda h(1-p)(v-1)}.
\]
Iloczyn wskazuje jednocześnie rozkłady Poissona i
niezależność dwóch liczników na tym przedziale.
Ponieważ podstawowy proces ma niezależne przyrosty, a
oznaczenia na rozłącznych przedziałach są niezależne,
rozumowanie stosuje się łącznie do dowolnej skończonej
rodziny rozłącznych przedziałów. Tym samym oba powstałe
procesy są niezależne jako rodziny zmiennych losowych,
a nie tylko dla każdego ustalonego $t$.
\end{proof}

\begin{zadanie}[Dwa strumienie i klasyfikacja zgłoszeń]
Przychodzą niezależne zgłoszenia z działów $A$ i $B$
o intensywnościach odpowiednio $2$ i $3$ na godzinę.
Dodatkowo każde zgłoszenie, niezależnie od pozostałych, jest
pilne z prawdopodobieństwem $1/5$. Obliczyć
prawdopodobieństwo braku pilnych zgłoszeń w ciągu dwóch godzin;
oraz prawdopodobieństwo, że pierwsze zgłoszenie całego systemu
pochodzi z działu $A$.

\emph{Rozwiązanie.} Superpozycja daje intensywność $5$.
Rozrzedzenie pilnych zgłoszeń ma intensywność $5/5=1$;
w ciągu dwóch godzin liczba pilnych zgłoszeń jest więc
$\operatorname{Poisson}(2)$, a prawdopodobieństwo zera wynosi
$e^{-2}$.
Pierwsze czasy nadejść z działów $A$ i $B$ są niezależne,
$U\sim\operatorname{Exp}(2)$ i
$V\sim\operatorname{Exp}(3)$. Zatem
\[
\Prob(U<V)=\int_0^\infty \Prob(V>u)\,2e^{-2u}\,\dd u
=\int_0^\infty 2e^{-5u}\,\dd u=\frac25.
\]
Interpretacja: spośród wszystkich zdarzeń frakcja $2/5$
pochodzi średnio z działu $A$, ale wynikiem jest tu
konkretne prawdopodobieństwo typu \emph{pierwszego} zgłoszenia.
\end{zadanie}

\section{Złożony proces Poissona i kompensacja}
\begin{definicja}[Złożony proces Poissona]
Niech $N$ będzie procesem Poissona o intensywności $\lambda>0$,
a $(Y_k)_{k\ge1}$ ciągiem niezależnych rzeczywistych
zmiennych losowych o wspólnym rozkładzie, niezależnym od $N$.
Definiujemy $S_t=\sum_{k=1}^{N_t}Y_k$,
gdzie pusta suma wynosi zero. Proces $S$ nazywamy
\emph{złożonym procesem Poissona}. Skok $S$ w chwili $T_k$
ma wielkość $Y_k$.
\end{definicja}

Model ma szerokie zastosowania: wypadki występują w chwilach
Poissona, natomiast kwota szkody jest niezależnym losowym
\emph{znacznikiem} zdarzenia. Dla $Y_k=1$ odzyskujemy $N_t$.

\begin{twierdzenie}[Moment, wariancja i funkcja charakterystyczna]
Załóżmy w poprzedniej definicji $\E|Y_1|<\infty$.
Wtedy $\E S_t=\lambda t\E Y_1$.
Jeśli ponadto $\E Y_1^2<\infty$, to
$\Var(S_t)=\lambda t\E Y_1^2$.
Bez założeń momentowych zachodzi dla $u\in\R$
\[
\E e^{iuS_t}=\exp\!\bigl(\lambda t(\E e^{iuY_1}-1)\bigr).
\]
Przy $\E|Y_1|<\infty$ proces
$S_t-\lambda t\E Y_1$ jest martyngałem względem
filtracji generowanej przez zaznaczone zdarzenia do chwili $t$.
\end{twierdzenie}
\begin{proof}
Warunkując względem $N_t=n$ i używając niezależności znaczników,
uzyskujemy $\E[S_t\mid N_t=n]=n\E Y_1$,
$\Var(S_t\mid N_t=n)=n\Var(Y_1)$ i
$\E[e^{iuS_t}\mid N_t=n]=(\E e^{iuY_1})^n$.
Po uśrednieniu pierwsza równość jest oczywista. Ze wzoru
na wariancję całkowitą wynika
\begin{align*}
\Var(S_t)&=\E[N_t\Var(Y_1)]
 +\Var(N_t\E Y_1)\\
&=\lambda t\bigl(\Var(Y_1)+(\E Y_1)^2\bigr)
=\lambda t\E Y_1^2.
\end{align*}
Uśrednienie funkcji charakterystycznej wykorzystuje szereg
$\sum_{n\ge0}e^{-\lambda t}(\lambda t)^n z^n/n!
=e^{\lambda t(z-1)}$.
Przyrosty $S_t-S_s$ są niezależne od informacji do chwili $s$
i mają średnią $\lambda(t-s)\E Y_1$:
wynika to z niezależnych przyrostów procesu Poissona i z faktu,
że znaczniki przyszłych zdarzeń są niezależne od przeszłych.
Stąd właściwość martyngałowa po odjęciu deterministycznej średniej.
\end{proof}

\section{Niejednorodny proces Poissona}
Stała intensywność jest wygodna, ale liczba zgłoszeń w ciągu
dnia na ogół zależy od godziny. Następująca konstrukcja pokazuje,
że zmienna intensywność nie wymaga zmiany podstawowych idei.

\begin{definicja}[Proces Poissona z deterministyczną intensywnością]
Niech $a:[0,\infty)\to[0,\infty)$ będzie funkcją borelowską,
całkowalną na każdym przedziale $[0,t]$. Oznaczmy
$A(t)=\int_0^t a(s)\,\dd s$.
Jeżeli $\widetilde N$ jest standardowym procesem Poissona
o intensywności $1$, to proces $N_t=\widetilde N_{A(t)}$
nazywamy \emph{niejednorodnym procesem Poissona o intensywności $a$}.
\end{definicja}

\begin{stwierdzenie}
Dla procesu z poprzedniej definicji i $0\le s<t$ przyrost
$N_t-N_s$ ma rozkład
$\operatorname{Poisson}(\int_s^t a(u)\,\dd u)$, a przyrosty
na rozłącznych przedziałach są niezależne. Ponadto
$\Prob(N_t-N_s=0)=\exp(-\int_s^t a(u)\,\dd u)$.
\end{stwierdzenie}
\begin{proof}
Funkcja $A$ jest ciągła i niemalejąca, a
$A(t)-A(s)=\int_s^t a(u)\,\dd u$.
Stąd $N_t-N_s=\widetilde N_{A(t)}-\widetilde N_{A(s)}$.
Przedziały obrazowe $[A(s),A(t)]$ odpowiadające rozłącznym
przedziałom czasu mają rozłączne wnętrza; przedziały o
zerowej długości dają zerowe przyrosty. Zatem po zastosowaniu
niezależnych i stacjonarnych przyrostów $\widetilde N$
dostajemy wszystkie tezy.
\end{proof}

\begin{przyklad}[Intensywność sezonowa]
Niech $a(t)=2+t$ dla $0\le t\le2$ i $a(t)=4$ dla $t>2$.
W ciągu pierwszej godziny średnia liczba zdarzeń wynosi
$A(1)=\int_0^1(2+s)\,\dd s=5/2$,
a prawdopodobieństwo braku zgłoszeń jest $e^{-5/2}$.
Pomiędzy godziną pierwszą a drugą średnia przyrostu
wynosi $A(2)-A(1)=7/2$. Obie liczby zdarzeń są niezależne,
choć ich rozkłady nie są jednakowe: przyrosty pozostają
niezależne, ale przestają być stacjonarne.
\end{przyklad}

\section{Charakterystyka przez przyrosty i zdarzenia rzadkie}
Dotychczas budowaliśmy proces Poissona z czasów wykładniczych.
Można jednak rozpocząć od modelowania zliczeń w krótkich
przedziałach. Poniższy wynik uzasadnia tę drogę konstrukcji.

\begin{twierdzenie}[Charakterystyka procesu Poissona przez małe przyrosty]
\label{thm:poisson-infinitesimal}
Niech $N=(N_t)_{t\ge0}$ będzie procesem o wartościach
w $\N_0$, $N_0=0$, z niezależnymi i stacjonarnymi
nieujemnymi przyrostami. Załóżmy, że istnieje $\lambda>0$
takie, iż przy $h\downarrow0$
\[
\Prob(N_h=1)=\lambda h+o(h),\qquad
\Prob(N_h\ge2)=o(h).
\]
Wtedy dla każdego $t\ge0$ zachodzi
$N_t\sim\operatorname{Poisson}(\lambda t)$,
a rozkłady skończeniewymiarowe $N$ są takie same jak
rozkłady procesu Poissona o intensywności $\lambda$.
Jeżeli ponadto $N$ ma niemalejące trajektorie c\`adl\`ag,
to proces ten ma prawo procesu Poissona także jako
zmienna losowa w przestrzeni trajektorii c\`adl\`ag.
\end{twierdzenie}
\begin{proof}
Ustalmy $z\in[0,1]$ i funkcję tworzącą
$G_t(z)=\E z^{N_t}$, przyjmując $0^0=1$.
Z niezależnych i stacjonarnych przyrostów wynika
$G_{t+h}(z)=G_t(z)G_h(z)$. Ponieważ
$\Prob(N_h=0)=1-\lambda h+o(h)$, a wkład stanów
$N_h\ge2$ do $G_h(z)$ ma moduł nie większy niż
$\Prob(N_h\ge2)=o(h)$, otrzymujemy
\[
G_h(z)=1+\lambda h(z-1)+o(h).
\]
Zatem dla $t\ge0$ pochodna prawostronna spełnia
\[
\frac{\partial}{\partial t^+}G_t(z)
=\lim_{h\downarrow0}G_t(z)\frac{G_h(z)-1}{h}
=\lambda(z-1)G_t(z).
\]
Dodatkowo z $\Prob(N_h=0)\to1$ oraz niezależnych przyrostów
wynika ciągłość $t\mapsto G_t(z)$: z prawej na mocy
$G_{t+h}=G_tG_h$, a z lewej, dla $0<h\le t$,
$G_t=G_{t-h}G_h$ i $G_h\to1$; ponieważ
$0\le G_{t-h}\le1$, różnica dąży do zera.
Funkcja ciągła z ciągłą pochodną prawostronną zadaną
powyższym równaniem ma jego zwykłe rozwiązanie,
co można sprawdzić dla
$H_t=e^{-\lambda(z-1)t}G_t(z)$:
jej pochodna prawostronna znika w każdym punkcie,
a więc $H$ jest stała (wynika to z twierdzenia
o wartości średniej w wersji dla pochodnych jednostronnych,
lub z zastosowania argumentu o ekstremum do
$H_t\pm\varepsilon t$). Ponieważ $G_0(z)=1$,
\[
G_t(z)=e^{\lambda t(z-1)}.
\]
Współczynniki szeregu potęgowego są równe
$e^{-\lambda t}(\lambda t)^n/n!$, co daje rozkład Poissona.
Ponieważ przyrosty są niezależne i stacjonarne,
wszystkie rozkłady skończeniewymiarowe są wyznaczone
przez właśnie otrzymane rozkłady pojedynczych przyrostów.
Dla procesów c\`adl\`ag na przestrzeni $D([0,\infty),\N_0)$
rozkłady skończeniewymiarowe wyznaczają prawo procesu
(jego borelowską $\sigma$-algebrę generują współrzędne),
co kończy dowód.
\end{proof}

\begin{przyklad}[Rozrzedzony schemat Bernoulliego]
Ustalmy $t>0$ i rozważmy $m$ niezależnych prób
na siatce $t/m,2t/m,\ldots,t$, z prawdopodobieństwem
sukcesu $\lambda t/m$ w każdej próbie (dla dostatecznie
dużych $m$). Liczba sukcesów ma rozkład
$\operatorname{Bin}(m,\lambda t/m)$.
Dla ustalonego $k\ge0$ jej prawdopodobieństwo wynosi
\[
\binom mk\left(\frac{\lambda t}{m}\right)^k
\left(1-\frac{\lambda t}{m}\right)^{m-k}
\longrightarrow e^{-\lambda t}\frac{(\lambda t)^k}{k!}.
\]
Granica wynika z
$\binom mk/m^k\to1/k!$ i
$(1-\lambda t/m)^m\to e^{-\lambda t}$.
Dla wszystkich rozłącznych odcinków siatki liczby
sukcesów są niezależne; analogicznie otrzymujemy
zbieżność ich wspólnych rozkładów do niezależnych
przyrostów Poissona. Jest to matematyczne wyjaśnienie,
dlaczego rozkład Poissona naturalnie opisuje wiele
niezależnych, rzadkich zdarzeń.
\end{przyklad}

\section{Co wynika z modelu, a czego model nie mówi?}
Rozkład Poissona nie jest uniwersalnym opisem każdego
zliczania. Jeśli nadejście klienta zwiększa lub zmniejsza
szansę kolejnego nadejścia, niezależność przyrostów jest
nieadekwatna. Jeśli urządzenie ma okres refrakcji, to
bardzo krótkie odstępy między zdarzeniami nie są wykładnicze.
Na dalszych stronach uogólnimy ideę zliczania przez
\emph{intensywności zależne od stanu}. Otrzymamy łańcuchy
Markowa w czasie ciągłym, których czasy oczekiwania będą
nadal wykładnicze, ale ich parametry będą zmieniać się
po każdym skoku procesu.
% ================================================================
% TOM II. ROZDZIAŁ 5
% ================================================================

\chapter{Łańcuchy Markowa w czasie ciągłym}
\label{ch:ctmc}
\paragraph{Odsyłacze międzytomowe.} Łańcuchy w czasie dyskretnym opisuje tom I, rozdział~\ref{I-ch:markov}, a ich wykorzystanie w MCMC tom IV, rozdział~\ref{IV-ch:mcmc}.


\section{Dlaczego macierz przejścia nie wystarcza?}
W łańcuchu Markowa z czasu dyskretnego macierz $P=(p_{ij})$
opisuje jeden krok. W czasie ciągłym pojawiają się dwie odrębne
losowości: \emph{kiedy} następuje zmiana stanu oraz \emph{dokąd}
proces przechodzi. W modelu awarii i naprawy urządzenie może długo
pozostawać sprawne, a po awarii przebywać w stanie uszkodzenia
przez losowy czas. Sam ciąg odwiedzanych stanów nie opisuje
czasów trwania poszczególnych faz.

\begin{definicja}[Macierz intensywności]
Niech $E$ będzie skończonym lub przeliczalnym zbiorem stanów.
\emph{Konserwatywną macierzą intensywności} (macierzą $Q$)
nazywamy macierz $Q=(q_{ij})_{i,j\in E}$ spełniającą:
$q_{ij}\ge0$ dla $i\ne j$, $q_i:=-q_{ii}\in[0,\infty)$
oraz $\sum_{j\ne i}q_{ij}=q_i$ dla każdego $i\in E$.
Gdy $q_i>0$, definiujemy prawdopodobieństwa skoku
$r_{ij}=q_{ij}/q_i$ dla $j\ne i$ i $r_{ii}=0$.
Stan z $q_i=0$ jest pochłaniający; dla wygody
w macierzy łańcucha skoków przyjmujemy wówczas
$r_{ii}=1$ i $r_{ij}=0$ dla $j\ne i$.
\end{definicja}

Wyraz $q_{ij}$, $i\ne j$, jest natężeniem \emph{bezpośredniego}
przejścia z $i$ do $j$, a $q_i$ całkowitą intensywnością
opuszczenia stanu $i$. Zapis $q_{ii}=-q_i$ nie oznacza
ujemnego prawdopodobieństwa: elementy macierzy $Q$ są
intensywnościami, a nie prawdopodobieństwami.

\begin{przyklad}[Układ sprawny--uszkodzony]
Niech $E=\{0,1\}$, gdzie $0$ oznacza urządzenie sprawne,
a $1$ uszkodzone. Jeśli awarie występują z intensywnością
$\alpha>0$, a zakończenia napraw z intensywnością $\beta>0$,
to
\[
Q=\begin{pmatrix}-\alpha&\alpha\\\beta&-\beta\end{pmatrix}.
\]
Czas do awarii, rozpoczynając od stanu $0$, ma rozkład
$\operatorname{Exp}(\alpha)$; czas naprawy w stanie $1$
ma rozkład $\operatorname{Exp}(\beta)$. Po każdym przejściu
zaczyna się nowy, niezależny czas oczekiwania.
\end{przyklad}

\begin{center}
\begin{tikzpicture}[>=Latex,thick,node distance=2.4cm,
 st/.style={draw,circle,minimum size=11mm,fill=blue!5}]
\node[st] (zero) {$0$};
\node[st,right=of zero] (one) {$1$};
\node[st,right=of one] (two) {$2$};
\draw[->,bend left=22] (zero) to node[above] {$q_{01}=2$} (one);
\draw[->,bend left=22] (one) to node[below] {$q_{10}=1$} (zero);
\draw[->] (one) -- node[above] {$q_{12}=3$} (two);
\node[below=9mm of two] {pochłanianie};
\end{tikzpicture}
\par\smallskip\small Schemat: intensywności skoków dla trzech stanów.
W stanie $1$ łączna intensywność wynosi $1+3=4$.
\end{center}


\begin{uwaga}[Czym macierz intensywności różni się od macierzy przejścia?]
Jeżeli proces zaczyna w stanie $i$, to dla $j\ne i$
$\Prob_i(X_h=j)=q_{ij}h+O(h^2)$ przy $h\downarrow0$
\emph{w skończonym modelu}. Tymczasem dla dowolnego
ustalonego $t>0$ prawdopodobieństwo $p_{ij}(t)$ może zawierać
wkład wielu następujących po sobie zmian stanu.
Zatem $q_{ij}$ nie jest ani prawdopodobieństwem pojedynczego
skoku (nim jest $q_{ij}/q_i$), ani prawdopodobieństwem
znalezienia się w $j$ po czasie $t$ (nim jest $p_{ij}(t)$).
\end{uwaga}

\section{Konstrukcja przez łańcuch skoków i czasy przebywania}
\begin{twierdzenie}[Konstrukcja minimalnego procesu skokowego]
\label{thm:jump-construction}
Niech $E$ będzie zbiorem przeliczalnym, $Q$ konserwatywną
macierzą intensywności z $q_i<\infty$ dla każdego $i\in E$,
a $\mu$ rozkładem probabilistycznym na $E$.
Istnieje przestrzeń probabilistyczna, na której można
zdefiniować ciąg stanów $Y_0,Y_1,\ldots$, ciąg czasów skoków
$0=J_0<J_1<\cdots$ aż do ewentualnego pochłonięcia oraz
proces $X$ następująco:
\begin{itemize}
\item $Y_0\sim\mu$, a ze stanu $i$ o $q_i>0$ następny stan
ma rozkład $(r_{ij})_{j\in E}$;
\item przy wejściu do $i$ o $q_i>0$ losujemy czas przebywania
o rozkładzie $\operatorname{Exp}(q_i)$, warunkowo względem
bieżącego stanu niezależny od przeszłości i od wyboru
następnego stanu;
\item dla $J_n\le t<J_{n+1}$ ustalamy $X_t=Y_n$,
a po wejściu do stanu z $q_i=0$ proces pozostaje w nim na stałe.
\end{itemize}
Jeśli suma wszystkich kolejnych czasów przebywania jest skończona,
jej granicę $\zeta$ nazywamy \emph{czasem eksplozji}; dla $t\ge\zeta$
ustalamy $X_t=\partial$ w dodatkowym stanie cmentarnym
$\partial\notin E$. Tak otrzymany proces nazywamy
\emph{minimalnym procesem skokowym o macierzy $Q$}.
Na przedziale $[0,\zeta)$ jego trajektorie są schodkowe i
c\`adl\`ag, a czasy rzeczywistych skoków nie mają atomów
w skończonych chwilach.
\end{twierdzenie}
\begin{proof}
Dla każdego stanu $i$ z $q_i>0$ wiersz $(r_{ij})$ jest
rozkładem probabilistycznym. Używając niezależnych zmiennych
jednostajnych na $(0,1)$ oraz ustalonego wyliczenia stanów,
możemy metodą odwrotnej dystrybuanty wylosować następny stan
z tego rozkładu. Pozwala to rekurencyjnie skonstruować
$Y_0,Y_1,\ldots$; po pochłonięciu pozostawiamy
$Y_n$ w stanie pochłaniającym. Zmienne pomocnicze
możemy zrealizować na przeliczalnym iloczynie przestrzeni
probabilistycznych wraz z niezależnym stanem początkowym.

Niezależnie od całego łańcucha $Y$ losujemy ciąg
niezależnych zmiennych $Z_0,Z_1,\ldots$ o rozkładzie
$\operatorname{Exp}(1)$, a czas przebywania przy wizycie
w stanie $Y_n=i$ definiujemy jako $H_n=Z_n/q_i$,
jeśli $q_i>0$, oraz $H_n=\infty$, jeśli $q_i=0$.
Z równości
$\Prob(Z_n/q_i>t)=e^{-q_i t}$ wynika żądany rozkład
warunkowy. Ustalamy $J_0=0$ i $J_{n+1}=J_n+H_n$ w rozszerzonej
półprostej. W szczególności po pochłonięciu kolejne
czasy skoków są nieskończone. Definiujemy
$\zeta=\lim_n J_n$. Wszystkie $J_n$ i $\zeta$ są mierzalne.
Dla $i\in E$ i skończonego $t$ mamy
\[
\{X_t=i\}=\bigcup_{n\ge0}
\{Y_n=i,\ J_n\le t<J_{n+1}\},
\qquad \{X_t=\partial\}=\{\zeta\le t\}.
\]
Zatem $X_t$ jest zmienną losową na dyskretnej
przestrzeni $E\cup\{\partial\}$.

Jeśli $J_n\uparrow\zeta<\infty$, to przed chwilą $\zeta$
występuje nieskończenie wiele skoków. Nie należy wtedy
arbitralnie przedłużać trajektorii w $E$, bo może ona nie
mieć granicy. Dodanie stanu $\partial$ daje jednoznaczny
minimalny proces do końca czasu. Każdy skończony $H_n>0$, więc
kolejne skończone chwile $J_n$ są ściśle rosnące; na $[0,\zeta)$
na każdym przedziale domkniętym zawartym w $[0,\zeta)$
występuje tylko skończenie wiele skoków. Stąd wskazana
regularność trajektorii przed eksplozją.
Dla $n\ge1$ i ustalonego skończonego $t$ mamy
$\Prob(J_n=t)=0$: warunkowo względem ciągu stanów
$Y$, jeśli $J_n<\infty$, jest on sumą $n$ niezależnych
zmiennych wykładniczych o dodatnich parametrach, a więc
ma gęstość (przez splot gęstości). Możliwy jest natomiast
atom w nieskończoności.
\end{proof}

\begin{uwaga}[Dwie różne macierze]
Macierz $R=(r_{ij})$ opisuje \emph{dokąd} idziemy po skoku.
Macierz $Q$ opisuje łącznie wybór nowego stanu i tempo
opuszczania poprzedniego. Nawet jeśli dwa procesy mają taki sam
łańcuch skoków $Y_n$, lecz różne intensywności $q_i$,
mają inne rozkłady $X_t$ w czasie kalendarzowym.
\end{uwaga}


\begin{przyklad}[Dlaczego rozkład wykładniczy jest naturalny dla jednorodnego łańcucha?]
Przypuśćmy, że układ zaczyna w stanie $0$, czeka losowy
czas $H$, a następnie nieodwracalnie przechodzi do $1$.
Jeżeli $H$ ma rozkład jednostajny na $[0,2]$, to przy $0<s<s+h<2$
\[
\Prob(H>s+h\mid H>s)=\frac{2-s-h}{2-s}.
\]
To prawdopodobieństwo zależy od czasu $s$, który upłynął
od uruchomienia. Dla procesu, który \emph{jednorodnie}
zachowuje się w stanie $0$, oczekiwalibyśmy natomiast,
aby przy obecnym stanie $0$ szansa pozostania w nim przez
kolejne $h$ jednostek czasu nie zależała od wieku pobytu
w tym stanie. Gdy $H\sim\operatorname{Exp}(q)$, rzeczywiście
\[
\Prob(H>s+h\mid H>s)=e^{-qh}.
\]
Równanie braku pamięci wyjaśnia, dlaczego w konstrukcji
jednorodnego łańcucha Markowa wybieramy czasy wykładnicze.
Przy innych prawach czasu przebywania można otrzymać
\emph{proces pół-Markowa}, w którym do stanu należy
dodać informację o czasie od ostatniej zmiany.
\end{przyklad}

\section{Eksplozja i jej wykluczanie}
\begin{twierdzenie}[Ograniczone intensywności wykluczają eksplozję]
\label{thm:bounded-rates}
Niech proces minimalny z twierdzenia~\ref{thm:jump-construction}
ma $\sup_{i\in E}q_i\le\Lambda<\infty$.
Wtedy $\Prob(\zeta=\infty)=1$.
\end{twierdzenie}
\begin{proof}
Jeśli $\Lambda=0$, wszystkie stany są pochłaniające
i teza jest natychmiastowa. Załóżmy $\Lambda>0$.
Jeżeli proces wchodzi do stanu pochłaniającego, to
nie może mieć nieskończenie wielu rzeczywistych skoków
w czasie skończonym. W przeciwnym razie kolejne czasy
przebywania są równe $H_n=Z_n/q_{Y_n}$, gdzie
$Z_n$ są niezależne o rozkładzie $\operatorname{Exp}(1)$.
Ponieważ $q_{Y_n}\le\Lambda$, mamy
$H_n\ge Z_n/\Lambda$ prawie na pewno.
Z lematu o braku eksplozji dla niezależnych wykładniczych
z rozdziału 4 wynika $\sum_n Z_n/\Lambda=\infty$
prawie na pewno. Także $\sum_n H_n=\infty$.
\end{proof}

\begin{przyklad}[Dlaczego skończoność każdego $q_i$ nie wystarcza?]
Niech $E=\N_0$, $X_0=0$, a jedyne dozwolone przejście ma postać
$n\to n+1$ z intensywnością $q_{n,n+1}=\lambda_n>0$.
Czas eksplozji jest sumą niezależnych zmiennych
$H_n\sim\operatorname{Exp}(\lambda_n)$.
Dla $\lambda_n=(n+1)^2$ mamy
\[
\E\zeta=\sum_{n=0}^\infty\frac1{(n+1)^2}<\infty.
\]
Ponieważ $\zeta\ge0$, skończona wartość oczekiwana
implikuje $\zeta<\infty$ prawie na pewno.
Nieskończenie wiele skoków występuje więc przed pewną
skończoną, losową chwilą! To model eksplozji.
Natomiast dla $\lambda_n=n+1$ zbieżność $\sum1/\lambda_n$
nie zachodzi; poniższe twierdzenie pokazuje brak eksplozji.
\end{przyklad}

\begin{twierdzenie}[Kryterium eksplozji procesu czystych urodzin]
\label{thm:pure-birth-explosion}
Niech $E=\N_0$, $X_0=0$, a jedyną możliwą zmianą stanu
$n$ będzie przejście do $n+1$ z intensywnością
$\lambda_n\in(0,\infty)$. Wówczas
\[
\Prob(\zeta<\infty)=1
\quad\Longleftrightarrow\quad
\sum_{n=0}^\infty\frac1{\lambda_n}<\infty.
\]
Jeśli szereg jest rozbieżny, to $\Prob(\zeta=\infty)=1$.
\end{twierdzenie}
\begin{proof}
Mamy $\zeta=\sum_{n\ge0}H_n$ z niezależnymi
$H_n\sim\operatorname{Exp}(\lambda_n)$.
Jeżeli $\sum1/\lambda_n<\infty$, to
$\E\zeta=\sum\E H_n<\infty$ na mocy twierdzenia o
zbieżności monotonicznej, skąd $\zeta<\infty$ prawie na pewno.

Załóżmy teraz $\sum1/\lambda_n=\infty$.
Dla $c>0$ niezależność i twierdzenie o zbieżności
zmajoryzowanej (przez stałą $1$) dają
\[
\E e^{-c\zeta}
=\lim_{m\to\infty}\prod_{n=0}^m\E e^{-cH_n}
=\prod_{n=0}^\infty\frac{\lambda_n}{\lambda_n+c}.
\]
Musimy pokazać, że iloczyn jest zerowy. Jego ujemny logarytm
to szereg $\sum_n\log(1+c/\lambda_n)$.
Jeśli nieskończenie wiele $\lambda_n\le c$, to w szeregu
nieskończenie wiele wyrazów jest nie mniejszych od $\log2$.
W przeciwnym razie dla wszystkich dostatecznie dużych $n$
mamy $0<c/\lambda_n<1$, a więc
$\log(1+c/\lambda_n)\ge c/(2\lambda_n)$.
W obu przypadkach szereg logarytmów jest rozbieżny.
Zatem $\E e^{-c\zeta}=0$.
Przy przyjętej konwencji $e^{-c\infty}=0$, a dla każdej
skończonej wartości $\zeta$ wykładnik jest ściśle dodatni.
Stąd $\Prob(\zeta<\infty)=0$, co kończy dowód.
\end{proof}

\section{Macierz przejścia i równania Kołmogorowa}
Teraz ograniczymy się do \emph{skończonej} przestrzeni $E$.
Założenie jest istotne: w nieskończonym zbiorze stanów
potrzebujemy dodatkowej kontroli, aby uzasadniać różniczkowanie
nieskończonych szeregów i wykluczać eksplozję.

\begin{definicja}[Macierz przejścia w czasie ciągłym]
Niech $E$ będzie skończone, a $X$ procesem o intensywnościach
$Q$. Dla $i,j\in E$ i $t\ge0$ definiujemy
$p_{ij}(t)=\Prob(X_t=j\mid X_0=i)$,
przy czym po prawej stronie rozumiemy prawo procesu
\emph{uruchomionego ze stanu $i$}, również wtedy,
gdy w wyjściowym rozkładzie początkowym $\Prob(X_0=i)=0$.
Macierz $P(t)=(p_{ij}(t))_{i,j\in E}$ nazywamy
\emph{macierzą przejścia w czasie $t$}.
\end{definicja}

\section{Zegar Poissona i uniformizacja}
Najwygodniejsza metoda konstrukcji skończonego łańcucha
pozwala zarazem udowodnić równania różniczkowe macierzy
przejścia bez odwoływania się do formalnej ``rachunkowości
prawdopodobieństw $o(h)$''.

\begin{twierdzenie}[Uniformizacja]
\label{thm:uniformization}
Niech $E$ będzie niepustym zbiorem skończonym,
a $Q$ konserwatywną macierzą intensywności.
Wybierzmy $\Lambda>\max_{i\in E}q_i$ i ustalmy
$R:=I+Q/\Lambda$. Wtedy $R$ jest macierzą stochastyczną.
Niech $(Y_n)_{n\ge0}$ będzie łańcuchem dyskretnym o macierzy
przejścia $R$, niezależnym od procesu Poissona $N$ o intensywności
$\Lambda$. Proces $X_t:=Y_{N_t}$ jest łańcuchem w czasie
ciągłym o intensywnościach $Q$, a jego macierz przejścia wynosi
\[
P(t)=e^{-\Lambda t}\sum_{n=0}^\infty
\frac{(\Lambda t)^n}{n!}R^n=e^{tQ},\qquad t\ge0.
\]
Dla każdego $i\ne j$ zachodzi
$\lim_{h\downarrow0}p_{ij}(h)/h=q_{ij}$,
a dla $i=j$ zachodzi
$\lim_{h\downarrow0}(p_{ii}(h)-1)/h=q_{ii}$.
\end{twierdzenie}
\begin{proof}
Dla $i\ne j$ mamy $R_{ij}=q_{ij}/\Lambda\ge0$,
a na przekątnej $R_{ii}=1-q_i/\Lambda>0$.
Suma każdego wiersza jest równa
$1+\Lambda^{-1}\sum_j q_{ij}=1$.
Wprowadzamy niezależny pomocniczy łańcuch dyskretny
i zegar Poissona.
W chwili zegara $Y$ dokonuje kroku macierzy $R$; może
przy tym pozostać w obecnym stanie. Te pozorne skoki
nazywamy \emph{skokami fikcyjnymi}.

Warunkowo względem $N_t=n$ końcowy stan procesu ma
rozkład $n$-krokowy $R^n$. Dlatego
$P(t)=\sum_n\Prob(N_t=n)R^n$, czyli pierwsza suma.
Ponieważ $I$ i $R$ komutują,
\[
e^{-\Lambda t}\sum_{n=0}^\infty\frac{(\Lambda tR)^n}{n!}
=e^{-\Lambda t I}e^{\Lambda tR}
=e^{\Lambda t(R-I)}=e^{tQ}.
\]
W podmultiplikatywnej normie macierzowej z $\norm{I}=1$
(np.\ normie operatorowej) szereg jest bezwzględnie zbieżny, gdyż
$\sum_{n\ge0}(\Lambda t)^n\norm{R^n}/n!$ można oszacować
przez $\exp(\Lambda t\norm{R})$, korzystając z
$\norm{R^n}\le\norm{R}^n$.
To uzasadnia również mnożenie szeregów wykładniczych.
Rozwinięcie wykładnika macierzowego
$e^{hQ}=I+hQ+O(h^2)$ przy $h\downarrow0$
daje oba wskazane limity.

Pozostaje porównać konstrukcję z rzeczywistymi
czasami przebywania. Przy starcie ze stanu $i$ o $q_i>0$
niech $K=\inf\{n\ge1:Y_n\ne i\}$.
Z własności Markowa łańcucha dyskretnego
\[
\Prob_i(K=k,Y_K=j)=R_{ii}^{k-1}R_{ij},
\qquad k\ge1,\ j\ne i.
\]
Stąd $K$ ma rozkład geometryczny z parametrem
$1-R_{ii}=q_i/\Lambda$, a kierunek $Y_K$ jest niezależny
od $K$ i ma rozkład $q_{ij}/q_i$.
Ponieważ cały łańcuch $Y$ jest niezależny od zegara,
czas pierwszego rzeczywistego skoku $H=T_K$, gdzie
$T_k$ są chwilami impulsów Poissona, jest niezależny
od kierunku $Y_K$. Ponadto
\[
\Prob_i(H>t)=\sum_{n\ge0}\Prob(N_t=n)R_{ii}^n
=e^{-\Lambda t(1-R_{ii})}=e^{-q_i t}.
\]
Jeżeli $q_i=0$, wiersz $R$ ma jedynkę na przekątnej
i żadna rzeczywista zmiana stanu nie następuje.
Po rzeczywistym skoku argument można powtórzyć:
łańcuch dyskretny ma mocną własność Markowa w indeksie
$K$ (tom I), a niezależne wykładnicze odstępy zegara
po tym indeksie są niezależne od dotychczasowych.
Ostatnia własność wynika przez rozbicie względem
$\{K=k\}$, ponieważ $K$ zależy wyłącznie od $Y$,
a odstępy zegara są niezależne od $Y$ i między sobą.
Iteracja daje warunkowe prawa kolejnych czasów i kierunków
takie jak w twierdzeniu~\ref{thm:jump-construction}.
Tym samym uniformizacja realizuje ten sam model.
\end{proof}


\begin{uwaga}[Dlaczego fikcyjne impulsy nie zmieniają modelu?]
Jeżeli proces zaczyna w stanie $i$ z $q_i>0$,
liczba impulsów fikcyjnych przed pierwszą rzeczywistą
zmianą stanu jest geometryczna. Można obliczyć
prawdopodobieństwo pozostawania w $i$ \emph{bez
rzeczywistego skoku} do chwili $t$ bez powoływania się
jedynie na intuicję pamięci rozkładu wykładniczego:
\[
\sum_{n=0}^\infty\Prob(N_t=n)R_{ii}^{n}
=e^{-\Lambda t}\exp(\Lambda tR_{ii})
=e^{-q_it}.
\]
Tak otrzymujemy czas oczekiwania $\operatorname{Exp}(q_i)$.
Warunkowo względem tego, że impuls powoduje rzeczywisty skok,
prawdopodobieństwo wybrania stanu $j\ne i$ jest
$R_{ij}/(1-R_{ii})=q_{ij}/q_i$.
Po skoku powtarzamy argument warunkowo względem
nowego stanu, korzystając z uzasadnienia w dowodzie
uniformizacji. Ta obserwacja uzasadnia
równoważność uniformizacji i konstrukcji przez czasy przebywania.
\end{uwaga}

Ważna korzyść uniformizacji jest obliczeniowa: możemy
symulować jeden wspólny zegar Poissona i niezależną sekwencję
kroków macierzy $R$, a przy prawdopodobieństwach przejścia
stosować nieujemny szereg zamiast bezpośrednio obliczać
wykładnik macierzowy.

\begin{twierdzenie}[Równania Chapmana--Kołmogorowa i Kołmogorowa]
\label{thm:kolmogorov-eq}
Dla skończonej przestrzeni stanów i konserwatywnej macierzy
intensywności $Q$ macierz przejścia $P(t)=e^{tQ}$ spełnia
$P(0)=I$ i $P(t+s)=P(t)P(s)$ dla $s,t\ge0$. Ponadto
jest różniczkowalna dla każdego $t\ge0$ (w zerze prawostronnie)
i spełnia \emph{równanie wsteczne} oraz \emph{równanie w przód}:
\[
P'(t)=QP(t),\qquad P'(t)=P(t)Q.
\]
Współrzędnie, dla każdego $i,j\in E$,
\[
\frac{\dd}{\dd t}p_{ij}(t)=\sum_{k\in E}q_{ik}p_{kj}(t),
\qquad
\frac{\dd}{\dd t}p_{ij}(t)=\sum_{k\in E}p_{ik}(t)q_{kj}.
\]
\end{twierdzenie}
\begin{proof}
Wzór $P(t)=e^{tQ}$ otrzymaliśmy z konstrukcji, zatem
$P(t+s)=e^{(t+s)Q}=e^{tQ}e^{sQ}$.
Istnieje też interpretacja probabilistyczna: warunkujemy
stan w chwili $t$ i sumujemy po wszystkich stanach
pośrednich $k$, otrzymując
$p_{ij}(t+s)=\sum_kp_{ik}(t)p_{kj}(s)$.
Na każdym przedziale $[0,T]$ zarówno szereg wykładnika
macierzowego, jak i szereg jego pochodnych są jednostajnie
zbieżne: w normie operatorowej ich wyrazy są ograniczone
odpowiednio przez $(T\norm Q)^n/n!$ i, dla $n\ge1$,
$\norm Q(T\norm Q)^{n-1}/(n-1)!$.
Można więc różniczkować wyraz po wyrazie; dostajemy
$P'(t)=Qe^{tQ}=e^{tQ}Q$.
Przemienność $Q$ z własnym wykładnikiem wynika również
bezpośrednio z definicji szeregu.
Zapisując mnożenie macierzy współrzędnie, uzyskujemy
obydwa wzory skalarne.
\end{proof}

\begin{uwaga}[Dlaczego są dwa równania?]
Równanie wsteczne analizuje \emph{pierwszą decyzję po starcie}:
pierwszy indeks $i$ należy do stanu startowego.
Równanie w przód analizuje \emph{dopływ do stanu końcowego}:
drugi indeks $j$ jest stanem w chwili $t$.
W skończonym jednorodnym modelu obydwa równania są poprawne
i prowadzą do tej samej macierzy. Dla modeli z nieskończenie
wieloma stanami albo z eksplozją przejścia do różniczkowania
w nieskończonych szeregach wymagają osobnego dowodu.
\end{uwaga}


\begin{przyklad}[Powrót do stanu początkowego to nie brak skoku]
W modelu $E=\{0,1\}$ z intensywnościami
$0\to1$ równą $\alpha>0$ i $1\to0$ równą $\beta>0$
czas do pierwszego \emph{rzeczywistego} skoku, przy starcie
z $0$, ma rozkład $\operatorname{Exp}(\alpha)$. Zatem
\[
\Prob_0(\text{brak skoku do }t)=e^{-\alpha t}.
\]
Natomiast $p_{00}(t)=\Prob_0(X_t=0)$ uwzględnia również
ścieżki z dwoma, czterema lub większą parzystą liczbą skoków.
W istocie, z rozwiązania równań w następnym podrozdziale,
\[
 p_{00}(t)=\frac{\beta+\alpha e^{-(\alpha+\beta)t}}
 {\alpha+\beta}>e^{-\alpha t}\qquad(t>0).
\]
Nierówność jest ścisła, bo z dodatnim prawdopodobieństwem
występują dwa skoki $0\to1\to0$ przed chwilą $t$.
Rozróżnienie między ``stan nie zmienił się do chwili $t$''
a ``w chwili $t$ obserwuję stan początkowy'' będzie ważne
przy modelowaniu trajektorii i przy zadaniach o czasach trafienia.
\end{przyklad}

\section{Warunkowanie, własność Markowa i czasy zatrzymania}
\begin{twierdzenie}[Własność Markowa w czasie ciągłym]
Niech $E$ będzie skończone, $X$ procesem o intensywnościach $Q$,
a $\calF_t^X=\sigma(X_s:0\le s\le t)$ filtracją naturalną.
Dla każdego $s,t\ge0$ i każdej funkcji $f:E\to\R$
\[
\E[f(X_{s+t})\mid\calF_s^X]=(P(t)f)(X_s),\qquad
(P(t)f)(i)=\sum_{j\in E}p_{ij}(t)f(j).
\]
W konsekwencji dla $j\in E$
$\Prob(X_{s+t}=j\mid\calF_s^X)=p_{X_sj}(t)$.
\end{twierdzenie}
\begin{proof}
Pracujemy na przestrzeni z konstrukcji uniformizacyjnej:
$X_t=Y_{N_t}$, gdzie łańcuch dyskretny $Y$ o macierzy $R$
jest niezależny od zegara Poissona $N$ o intensywności $\Lambda$.
Niech $\calH_s$ oznacza informację wygenerowaną przez
$(N_u)_{0\le u\le s}$ oraz przez zmienne
$V_j^s$, równe $Y_j$ na $\{j\le N_s\}$ i ustalonemu
dodatkowemu symbolowi poza $E$ w przeciwnym razie.
Dla $u\le s$ zmienna $X_u=Y_{N_u}$ jest
$\calH_s$-mierzalna, gdyż na $\{N_u=j\}$ jest równa
$V_j^s$. Zatem $\calF_s^X\subseteq\calH_s$.
Dla każdego $A\in\calH_s$ i $n\ge0$ zachodzi
$A\cap\{N_s=n\}\in\sigma(\calF_s^N,Y_0,\ldots,Y_n)$:
klasa zdarzeń o tej własności jest $\sigma$-ciałem
i zawiera wszystkie wskazane generatory.

Ustalmy $k\ge0$ i weźmy zdarzenie $\{N_s=n\}$.
Niezależność całego procesu $Y$ od $N$ i własność
Markowa dyskretnego łańcucha dają, dla każdego stanu $i$,
\[
\E[f(Y_{n+k})\mid Y_0,\ldots,Y_n]=(R^kf)(Y_n).
\]
Warunkowo względem $\calH_s$ liczba dodatkowych impulsów
$K=N_{s+t}-N_s$ ma rozkład $\operatorname{Poisson}(\Lambda t)$
i jest niezależna od całego łańcucha $Y$ oraz przeszłości
zegara łącznie. Dokładniej, dla $A\in\calH_s$,
na $\{N_s=n\}$ stosujemy niezależność $K$ od
$\sigma(\calF_s^N,Y_j:j\ge0)$ oraz dyskretną własność
Markowa, rozszerzoną na $\sigma(\calF_s^N,Y_0,\ldots,Y_n)$
dzięki niezależności $Y$ od zegara (argument $\pi$--$\lambda$).
Otrzymujemy
\[
\E[\chi_{A\cap\{N_s=n\}}\chi_{\{K=k\}}f(Y_{n+k})]
=\Prob(K=k)\E[\chi_{A\cap\{N_s=n\}}(R^kf)(Y_n)].
\]
Sumując po $n,k$ i stosując charakterystykę
warunkowej wartości oczekiwanej, dostajemy:
\[
\begin{aligned}
\E[f(X_{s+t})\mid\calH_s]
 &=\sum_{k=0}^\infty
 e^{-\Lambda t}\frac{(\Lambda t)^k}{k!}
 (R^kf)(Y_{N_s})\\
 &= (P(t)f)(X_s).
\end{aligned}
\]
Powyższy rachunek można najpierw przeprowadzić na zdarzeniach
$\{N_s=n\}$, a potem zsumować po $n\ge0$; jest to dozwolone,
ponieważ $f$ jest ograniczona, $E$ jest skończone,
a szereg wag Poissona sumuje się do $1$.
Wreszcie $\calF_s^X\subseteq\calH_s$ i prawa strona
jest $\calF_s^X$-mierzalna. Własność wieży
\emph{warunkowej wartości oczekiwanej} daje
\[
\E[f(X_{s+t})\mid\calF_s^X]
=\E[(P(t)f)(X_s)\mid\calF_s^X]
=(P(t)f)(X_s).
\]
Dla $f=\chi_{\{j\}}$ otrzymujemy zapis
w terminach prawdopodobieństw warunkowych.
\end{proof}


\begin{twierdzenie}[Własność Markowa procesu minimalnego na przeliczalnej przestrzeni stanów]
Niech $X$ będzie procesem minimalnym z
twierdzenia~\ref{thm:jump-construction}, z dodatkowym
stanem pochłaniającym $\partial$. Dla $i\in E\cup\{\partial\}$
oznaczmy przez $P_t f(i)$ wartość oczekiwaną $f(X_t)$
dla procesu uruchomionego ze stanu $i$.
Dla każdej ograniczonej funkcji $f:E\cup\{\partial\}\to\R$
oraz $s,t\ge0$ zachodzi
\[
\E[f(X_{s+t})\mid\calF_s^X]=(P_t f)(X_s)
\quad\text{prawie na pewno}.
\]
\end{twierdzenie}
\begin{proof}
Wybierzmy rosnący ciąg skończonych zbiorów $E_m$
o sumie $E$. Proces $X^{(m)}$ otrzymujemy, wysyłając
$X$ do $\partial$ przy pierwszym wyjściu poza $E_m$
(również w chwili $0$, jeśli $X_0\notin E_m$).
Do wyjścia korzystamy z tych samych stanów i czasów
przebywania co w procesie pierwotnym.
Na $E_m\cup\{\partial\}$ macierz intensywności zachowuje
$q_{ij}$ dla $i,j\in E_m$ i ma
$q_{i\partial}=\sum_{j\notin E_m}q_{ij}$,
a wiersz $\partial$ jest zerowy. Konstrukcja czasów
przebywania i kierunków skoku daje dokładnie ten
skończony model, więc stosujemy udowodnioną już
własność Markowa w przypadku skończonym.

Na wspólnym zbiorze pełnej miary, na którym
$\sum_n Z_n=\infty$, eksplozja wymusza opuszczenie
każdego skończonego zbioru przed chwilą $\zeta$:
gdyby wszystkie odwiedzane stany należały do
pewnego skończonego zbioru $F$, mielibyśmy
$H_n\ge Z_n/\max_{i\in F}q_i$, jeśli maksimum jest
dodatnie, a przy maksimum zero nie byłoby skoków.
To wyklucza skończoną sumę czasów przebywania.
Dla ustalonej chwili $v<\zeta$ proces odwiedza
do $v$ tylko skończenie wiele stanów, toteż
$X_v^{(m)}=X_v$ dla dostatecznie dużego $m$.
Jeśli $\zeta\le v$, wszystkie obcięte procesy
są w $\partial$ w chwili $v$ na wskazanym zbiorze
pełnej miary. W konsekwencji dla każdej ustalonej
chwili $v$ zachodzi $X_v^{(m)}\to X_v$ prawie na pewno,
w topologii dyskretnej. Ten sam argument przy
starcie z ustalonego $i$ i twierdzenie o zbieżności
zmajoryzowanej dają $P_t^{(m)}f(i)\to P_t f(i)$,
gdy $m$ jest dostatecznie duże, by obejmować $i$
(lub $i=\partial$).

Dla $0\le u_1<\cdots<u_l\le s$ oraz dowolnych
zbiorów $B_1,\ldots,B_l\subseteq E\cup\{\partial\}$
oznaczmy $C_m=\bigcap_{j=1}^l\{X_{u_j}^{(m)}\in B_j\}$
i analogicznie $C$ dla $X$.
Skończony model spełnia
\[
\E[\chi_{C_m} f(X_{s+t}^{(m)})]
=\E[\chi_{C_m}(P_t^{(m)}f)(X_s^{(m)})].
\]
Na wspólnym zbiorze pełnej miary dla wymienionych
chwil mamy od pewnego $m$ równość $X_s^{(m)}=X_s$
oraz $\chi_{C_m}=\chi_C$. Obie strony są ograniczone
przez $\norm{f}_\infty$ pod znakiem wartości oczekiwanej,
zatem przejście do granicy daje
$\E[\chi_C f(X_{s+t})]=\E[\chi_C(P_tf)(X_s)]$.
Zdarzenia cylindryczne $C$ tworzą $\pi$-układ
generujący $\calF_s^X$. Argument $\pi$--$\lambda$
(przy funkcji $f$ dowolnego znaku stosujemy równość
całek jako równość miar skończonych po dodaniu stałej)
rozszerza tę tożsamość na całe $\calF_s^X$.
Mierzalność $(P_tf)(X_s)$ kończy dowód.
\end{proof}

\begin{lemat}[Własność Markowa po prawostronnym uzupełnieniu filtracji]
\label{lem:markov-aug}
Niech $X$ będzie minimalnym procesem skokowym
na przeliczalnym $E$ (z dodatkowym stanem $\partial$
w razie eksplozji). Przyjmijmy topologię dyskretną
na przestrzeni stanów $E'=E\cup\{\partial\}$.
Proces ma prawostronnie ciągłe trajektorie:
przed eksplozją wynika to z konstrukcji, a od chwili
eksplozji trajektoria jest stała. Niech $\calF_t^0=\sigma(X_s:0\le s\le t)$
oraz
$\calF_t=\bigcap_{u>t}\sigma(\calF_u^0\cup\calN)$,
gdzie $\calN$ jest rodziną wszystkich podzbiorów zbiorów
zerowych z bazowego $\sigma$-ciała $\mathcal A$
przestrzeni probabilistycznej $(\Omega,\mathcal A,\Prob)$,
którą zakładamy zupełną.
Wówczas dla $t,h\ge0$ i każdej ograniczonej funkcji $f:E'\to\R$
\[
\E[f(X_{t+h})\mid\calF_t]=(P(h)f)(X_t)
\quad\text{prawie na pewno}.
\]
\end{lemat}
\begin{proof}
Ustalmy $t,h$ i $A\in\calF_t$. Dla każdego $n\ge1$
połóżmy $u_n=t+1/n$. Ponieważ $A\in\calF_t\subseteq\sigma(\calF_{u_n}^0\cup\calN)$,
po odrzuceniu nieistotnych zbiorów zerowych możemy
stosować własność Markowa względem $\calF_{u_n}^0$
do ograniczonej zmiennej $f(X_{u_n+h})$:
\[
\E[\chi_A f(X_{u_n+h})]
 =\E[\chi_A(P(h)f)(X_{u_n})].
\]
Uzasadnijmy przejście przez uzupełnienie. Każde zdarzenie
$A\in\sigma(\calF_{u_n}^0\cup\calN)$ różni się
zbiorem zerowym od pewnego $A_n\in\calF_{u_n}^0$;
wynika to z tego, że klasa takich zdarzeń tworzy
$\sigma$-ciało zawierające oba generatory. Zatem
poprzednią równość, początkowo prawdziwą dla $A_n$,
można zastosować także do $A$.
Prawostronna ciągłość daje $X_{u_n+h}\to X_{t+h}$
oraz $X_{u_n}\to X_t$ niemal na pewno.
W topologii dyskretnej zbieżność stanów oznacza ich
równość od pewnego indeksu. Funkcje $f$ i $P(h)f$
są ograniczone (druga przez $\norm{f}_\infty$),
więc twierdzenie o zbieżności zmajoryzowanej pozwala
przejść do granicy.
Otrzymana równość całek dla każdego $A\in\calF_t$
oraz $\calF_t$-mierzalność $(P(h)f)(X_t)$
stanowią charakterystykę warunkowej wartości oczekiwanej.
\end{proof}

\begin{twierdzenie}[Mocna własność Markowa]
\label{thm:ctmc-strong}
Niech $X$ będzie minimalnym procesem skokowym
o macierzy intensywności $Q$ na przeliczalnym $E$,
z dodatkowym stanem $\partial$ w razie eksplozji.
Niech $\mathbb F=(\calF_t)_{t\ge0}$ będzie zwykłym
uzupełnieniem filtracji naturalnej (jak w poprzednim
lemacie), a przestrzeń probabilistyczna niech będzie zupełna.
Jeżeli $\tau$ jest czasem zatrzymania skończonym prawie na
pewno, to dla każdego $t\ge0$ i ograniczonej funkcji
$f:E\cup\{\partial\}\to\R$
\[
\E[f(X_{\tau+t})\mid\calF_\tau]=(P(t)f)(X_\tau)
\quad\text{prawie na pewno}.
\]
\end{twierdzenie}
\begin{proof}
Na zdarzeniu $\{\tau=\infty\}$, które należy do $\calF_0$
dzięki zupełności, możemy zastąpić $\tau$ przez $0$.
Tak otrzymujemy skończony punktowo czas zatrzymania;
$\sigma$-ciała w tych dwóch chwilach są identyczne,
gdyż przecięcia w ich definicjach różnią się wyłącznie
na zdarzeniu zerowym należącym wraz z podzbiorami do $\calF_0$.
Wartości $X_\tau$ i $X_{\tau+t}$ na tym zdarzeniu
ustalamy dowolnie w $E$. Najpierw niech $\tau$ przyjmuje wartości w przeliczalnym
zbiorze $\{s_k:k\ge1\}\subset[0,\infty)$.
Dla $A\in\calF_\tau$ zdarzenie
$A\cap\{\tau=s_k\}$ należy do $\calF_{s_k}$
(ponieważ $\{\tau=s_k\}=\{\tau\le s_k\}\setminus
\{\tau<s_k\}$, gdzie $\{\tau<s_k\}$ jest sumą
$\{\tau\le q\}$ po nieujemnych wymiernych $q<s_k$;
dla $s_k=0$ jest puste).
Stosując własność Markowa w deterministycznej chwili $s_k$
\emph{dla uzupełnionej filtracji} (lemat~\ref{lem:markov-aug})
i sumując po $k$, uzyskujemy
\[
\E[\chi_A f(X_{\tau+t})]
=\E[\chi_A(P(t)f)(X_\tau)].
\]
Dla ogólnego czasu $\tau<\infty$ prawie na pewno niech
$\tau_m=2^{-m}\lceil2^m\tau\rceil$.
Jest to czas zatrzymania przyjmujący wartości w zbiorze
$2^{-m}\N_0$: dla $v\ge0$ zdarzenie $\{\tau_m\le v\}$
jest równe $\{\tau\le2^{-m}\lfloor2^mv\rfloor\}$.
Ponadto $\tau_m\downarrow\tau$
prawie na pewno. Mamy $\calF_\tau\subseteq\calF_{\tau_m}$,
więc dla każdego $A\in\calF_\tau$ poprzednia równość
zastosowana do $\tau_m$ daje
\[
\E[\chi_A f(X_{\tau_m+t})]
=\E[\chi_A(P(t)f)(X_{\tau_m})].
\]
Trajektorie procesu są prawostronnie ciągłe także
w stanie cmentarnym. Dlatego w topologii dyskretnej
$X_{\tau_m+t}\to X_{\tau+t}$ i $X_{\tau_m}\to X_\tau$
prawie na pewno. Ponieważ $f$ i $P(t)f$ są ograniczone,
możemy zastosować twierdzenie o zbieżności
zmajoryzowanej po obu stronach.
W efekcie identyczna równość zachodzi dla $\tau$ i wszystkich
$A\in\calF_\tau$. Proces $X$ jest progresywnie mierzalny
jako adaptowany proces prawostronnie ciągły
(twierdzenie~\ref{thm:prog}), zatem $X_\tau$ jest
$\calF_\tau$-mierzalna na mocy
twierdzenia~\ref{thm:sampled}. Otrzymujemy więc
charakterystykę wskazanej warunkowej wartości oczekiwanej.
\end{proof}

\begin{uwaga}[Prawo całej przyszłej trajektorii]
Mocna własność Markowa w powyższej postaci dotyczy
jednej przyszłej chwili. Dla $0<t_1<\cdots<t_m$
oraz ograniczonych funkcji $f_1,\ldots,f_m$ stosujemy
ją kolejno w czasach zatrzymania $\tau+t_{m-1},
\ldots,\tau+t_1,\tau$ (dodanie deterministycznej
stałej nieujemnej zachowuje własność czasu zatrzymania).
Własność wieży daje
\[
\E\left[\prod_{j=1}^m f_j(X_{\tau+t_j})
\mid\calF_\tau\right]
=
P(t_1)\bigl(f_1 P(t_2-t_1)
(f_2\cdots P(t_m-t_{m-1})f_m)\bigr)(X_\tau),
\]
gdzie zagnieżdżenie rozumiemy rekurencyjnie od
$f_m$ wstecz, a mnożenie funkcji jest punktowe.
Dla $m=1$ prawa strona to $(P(t_1)f_1)(X_\tau)$.
Wybór funkcji charakterystycznych i argument
$\pi$--$\lambda$ dla cylindrów rozszerzają to na
warunkowe prawo całego przesuniętego procesu
względem $\sigma$-ciała cylindrycznego: jest ono
prawem procesu startującego z $X_\tau$.
\end{uwaga}

\begin{przyklad}[Pierwsza naprawa jako chwila losowa]
W modelu sprawny--uszkodzony niech $\tau$ będzie chwilą
pierwszego powrotu do stanu sprawnego po awarii, przy starcie
ze stanu $0$. Jeżeli $\alpha,\beta>0$, to $\tau<\infty$
prawie na pewno. Jest to czas zatrzymania: dla
$t>0$ zdarzenie $\{\tau\le t\}$ jest sumą po
wymiernych $r\in(0,t)$ zdarzeń
$\{X_r=1\}\cap(\{X_t=0\}\cup
\bigcup_{q\in\mathbb Q\cap(r,t)}\{X_q=0\})$.
Równość wynika z prawostronnej ciągłości dyskretnych
trajektorii i dodatnich czasów przebywania; dla $t=0$
zdarzenie jest puste. Mocna własność Markowa mówi, że warunkowo
względem całej historii do $\tau$ dalszy proces jest
nową kopią procesu zaczynającego w stanie $0$.
Nie wolno zastąpić $\tau$ dowolną losową chwilą zależną
od przyszłych awarii: konieczny jest warunek czasu zatrzymania.
\end{przyklad}


\begin{zadanie}[Pierwszy rzeczywisty skok w trzech stanach]
Niech $E=\{0,1,2\}$, $q_{01}=2$, $q_{10}=1$,
$q_{12}=3$, a stan $2$ będzie pochłaniający.
Proces startuje w stanie $1$. Niech $J$ będzie chwilą
pierwszego skoku i $Y=X_J$. Wyznaczyć prawo $J$, prawo
$Y$ i obliczyć $\Prob(J\le1,Y=2)$.

\emph{Rozwiązanie.} Całkowita intensywność wyjścia ze
stanu $1$ wynosi $q_1=q_{10}+q_{12}=4$, wobec czego
$J\sim\operatorname{Exp}(4)$. Kierunek skoku jest
niezależny od $J$ i ma rozkład
$\Prob(Y=0)=1/4$, $\Prob(Y=2)=3/4$.
Miara $B\mapsto\Prob(J\in B,Y=2)$ na $[0,\infty)$
ma względem miary Lebesgue'a gęstość
\[
g(t)=q_{12}e^{-q_1t}=3e^{-4t},\qquad t\ge0.
\]
Zatem
$\Prob(J\le1,Y=2)=\int_0^1 3e^{-4t}\,\mathrm d\lambda_1(t)
=\tfrac34(1-e^{-4})$.
Nie należy tego mylić z $\Prob_1(X_1=2)$:
ten ostatni wynik obejmuje również ścieżki, na których
pierwszy skok prowadzi do $0$, a proces dopiero później
przechodzi przez $1$ do $2$.
\end{zadanie}

\section{Dwa stany: pełne rozwiązanie równań}
\begin{twierdzenie}[Przejścia modelu awarii i napraw]
Niech $E=\{0,1\}$, $\alpha,\beta>0$ oraz
$Q=\left(\begin{smallmatrix}-\alpha&\alpha\\
\beta&-\beta\end{smallmatrix}\right)$.
Wtedy dla $t\ge0$
\[
P(t)=\frac1{\alpha+\beta}
\begin{pmatrix}
\beta+\alpha e^{-(\alpha+\beta)t} &
\alpha(1-e^{-(\alpha+\beta)t})\\
\beta(1-e^{-(\alpha+\beta)t})&
\alpha+\beta e^{-(\alpha+\beta)t}
\end{pmatrix}.
\]
\end{twierdzenie}
\begin{proof}
Oznaczmy $u(t):=p_{01}(t)$, czyli prawdopodobieństwo
uszkodzenia w chwili $t$ przy starcie w stanie sprawnym.
Równanie w przód daje
$u'(t)=\alpha p_{00}(t)-\beta p_{01}(t)$.
Ponieważ suma elementów w pierwszym wierszu $P(t)$ jest $1$,
$p_{00}(t)=1-u(t)$, skąd
$u'(t)=\alpha-(\alpha+\beta)u(t)$ i $u(0)=0$.
Mnożąc obie strony równania liniowego przez czynnik
$e^{(\alpha+\beta)t}$, uzyskujemy
$(e^{(\alpha+\beta)t}u(t))'=\alpha e^{(\alpha+\beta)t}$.
Całkujemy od $0$ do $t$ i dostajemy
$u(t)=\alpha(1-e^{-(\alpha+\beta)t})/(\alpha+\beta)$.
Analogicznie, dla $v(t)=p_{10}(t)$ otrzymujemy
$v'(t)=\beta-(\alpha+\beta)v(t)$ i $v(0)=0$,
a więc $v(t)=\beta(1-e^{-(\alpha+\beta)t})/(\alpha+\beta)$.
Pozostałe wpisy wynikają z sum wierszy równych $1$.
\end{proof}

\begin{zadanie}[Prawdopodobieństwo awarii i informacja o przeszłości]
Urządzenie ma intensywność awarii $\alpha=2$ i naprawy
$\beta=1$ (w jednostkach: godzina$^{-1}$).
W chwili $0$ działa. Obliczyć prawdopodobieństwo,
że będzie uszkodzone po godzinie, a następnie
prawdopodobieństwo, że będzie uszkodzone po dwóch godzinach,
\emph{jeśli wiadomo, że działało po pierwszej godzinie}.

\emph{Rozwiązanie.} Z powyższego twierdzenia
$\Prob(X_1=1\mid X_0=0)=\tfrac23(1-e^{-3})$.
W drugim pytaniu warunkujemy względem stanu $X_1=0$,
a nie stanu początkowego. Własność Markowa daje
\[
\Prob(X_2=1\mid X_1=0,X_0=0)
=\Prob(X_1=1\mid X_0=0)=\frac23(1-e^{-3}).
\]
Gdybyśmy zamiast tego pytali bez warunku $X_1=0$
o stan po dwóch godzinach, odpowiedzią byłoby
$\tfrac23(1-e^{-6})$. To dwa różne pytania probabilistyczne.
\end{zadanie}

\section{Generatory, martyngały i wzór Dynkina}
\begin{definicja}[Generator skończonego łańcucha]
Dla skończonego $E$, macierzy intensywności $Q$ i funkcji
$f:E\to\R$ definiujemy
\[
(Lf)(i):=\sum_{j\in E}q_{ij}f(j)
=\sum_{j\ne i}q_{ij}\bigl(f(j)-f(i)\bigr),\qquad i\in E.
\]
Operator $L$ na przestrzeni funkcji na $E$ nazywamy
generatorem infinitezymalnym procesu. W zapisie wektorowym
$Lf=Qf$. Nazwa jest uzasadniona granicą
\[
Lf=\lim_{h\downarrow0}\frac{P(h)f-f}{h},
\]
w normie supremum: z szeregu wykładnika macierzowego
$P(h)=I+hQ+O(h^2)$ w normie operatorowej, więc
po zastosowaniu do $f$ i podzieleniu przez $h$
otrzymujemy wskazaną granicę.
\end{definicja}

Druga postać wyjaśnia sens operatora: z każdego stanu $i$
ważymy zmiany wartości $f(j)-f(i)$ intensywnościami
bezpośrednich skoków. Dla funkcji stałej generator znika,
co jest analogią do pochodnej funkcji stałej.

\begin{twierdzenie}[Wzór Dynkina i martyngał generatorowy]
\label{thm:dynkin}
Niech $E$ będzie skończone, $X$ łańcuchem w czasie
ciągłym o intensywnościach $Q$ i filtracji naturalnej.
Dla każdej funkcji $f:E\to\R$ proces
\[
M_t^f:=f(X_t)-f(X_0)-\int_0^t (Lf)(X_s)\,\mathrm d\lambda_1(s)
\]
jest całkowalnym martyngałem.
W szczególności dla każdego $t\ge0$
\[
\E f(X_t)=\E f(X_0)+
\int_0^t\E[(Lf)(X_s)]\,\mathrm d\lambda_1(s).
\]
Jeśli $\tau$ jest czasem zatrzymania względem tej filtracji,
ograniczonym punktowo przez pewną stałą,
to również
\[
\E f(X_\tau)=\E f(X_0)+
\E\int_0^\tau (Lf)(X_s)\,\mathrm d\lambda_1(s).
\]
\end{twierdzenie}
\begin{proof}
Wszystkie funkcje na skończonym $E$ są ograniczone.
Z równania w przód, w zapisie operatorowym,
$\frac{\dd}{\dd t}P(t)f=P(t)Lf$.
Całkując od $0$ do $u$ dostajemy
$P(u)f-f=\int_0^u P(v)Lf\,\mathrm d\lambda_1(v)$.
Dla $0\le s\le t$ własność Markowa implikuje
\[
\E[f(X_t)-f(X_s)\mid\calF_s]
=(P(t-s)f-f)(X_s).
\]
Jednocześnie $Lf$ jest funkcją ograniczoną na skończonym $E$, więc
$\int_s^t\E|Lf(X_u)|\,\mathrm d\lambda_1(u)\le(t-s)\norm{Lf}_\infty<\infty$.
Dla każdego $A\in\calF_s$ możemy więc zastosować twierdzenie
Fubiniego do $\E[\chi_A\int_s^t Lf(X_u)\,\mathrm d\lambda_1(u)]$,
a następnie własność Markowa pod całką. Charakterystyka
warunkowej wartości oczekiwanej daje
\[
\E\left[\int_s^t(Lf)(X_u)\,\mathrm d\lambda_1(u)\mid\calF_s\right]
=\int_0^{t-s}(P(v)Lf)(X_s)\,\mathrm d\lambda_1(v).
\]
Obie prawe strony są równe. Po odjęciu otrzymujemy
$\E[M_t^f-M_s^f\mid\calF_s]=0$.
Proces $M^f$ jest adaptowany: $X$ jest progresywnie
mierzalny, więc całka do chwili $t$ z ograniczonej
funkcji $(Lf)(X_s)$ jest $\calF_t^X$-mierzalna
na mocy twierdzenia Fubiniego o mierzalności całki.
Ponadto całka jest ciągła względem górnej granicy,
więc $M^f$ ma prawostronnie ciągłe trajektorie.
Zatem $M^f$ jest martyngałem. Całkowalność jest
natychmiastowa, ponieważ $f$ i $Lf$ są ograniczone,
a rozważamy skończony czas $t$.

Biorąc wartości oczekiwane uzyskujemy pierwszą tożsamość.
Dla $\tau\le T$ twierdzenie o próbkowaniu opcjonalnym
(twierdzenie~\ref{thm:optional-cont}) stosuje się do
prawostronnie ciągłego martyngału $M^f$, ponieważ
$\sup_{0\le t\le T}|M_t^f|\le2\norm{f}_\infty+T\norm{Lf}_\infty$.
Stąd $\E M_\tau^f=\E M_0^f=0$, co daje ostatni wzór.
\end{proof}

\begin{przyklad}[Przewidywana zmiana funkcji]
Dla modelu dwustanowego i $f(0)=0$, $f(1)=1$ mamy
$Lf(0)=\alpha$ i $Lf(1)=-\beta$.
Proces
\[
\chi_{\{X_t=1\}}-\chi_{\{X_0=1\}}
-\int_0^t\bigl(\alpha\chi_{\{X_s=0\}}
-\beta\chi_{\{X_s=1\}}\bigr)\,\mathrm d\lambda_1(s)
\]
jest martyngałem. Człon podcałkowy opisuje
\emph{chwilową przewidywaną zmianę} wskaźnika awarii,
a jego odjęcie usuwa dryf.
\end{przyklad}

\section{Rozkłady stacjonarne, równowaga szczegółowa i ergodyczność}
\begin{definicja}[Rozkład stacjonarny]
Dla nieeksplodującego łańcucha na skończonym lub
przeliczalnym $E$ z macierzą przejścia $P(t)$ rozkład probabilistyczny $\pi=(\pi_i)_{i\in E}$
nazywamy \emph{stacjonarnym}, jeżeli
$\pi P(t)=\pi$ dla każdego $t\ge0$.
Rozkład $\pi$ spełnia \emph{równania bilansu}, jeśli
$\pi Q=0$, przy czym na przeliczalnym $E$ wymagamy
istnienia odpowiednich sum w każdej współrzędnej.
\end{definicja}

\begin{twierdzenie}[Bilans i równowaga szczegółowa]
Niech $E$ będzie skończone i $Q$ macierzą intensywności.
Rozkład probabilistyczny $\pi$ jest stacjonarny wtedy i tylko
wtedy, gdy $\pi Q=0$. Warunkiem wystarczającym jest
\emph{równowaga szczegółowa}
\[
\pi_iq_{ij}=\pi_jq_{ji}\qquad\text{dla wszystkich }i,j\in E.
\]
\end{twierdzenie}
\begin{proof}
Jeśli $\pi P(t)=\pi$ dla każdego $t$, to przez różniczkowanie
w zerze otrzymujemy $\pi Q=0$.
Odwrotnie, gdy $\pi Q=0$, z rozwinięcia szeregu
$e^{tQ}=I+tQ+t^2Q^2/2!+\cdots$ wynika
$\pi e^{tQ}=\pi$, gdyż $\pi Q^k=0$ dla $k\ge1$.
Jeżeli zachodzi równowaga szczegółowa, to dla każdego $j$
\[
(\pi Q)_j=\sum_{i\ne j}\pi_iq_{ij}-\pi_j\sum_{k\ne j}q_{jk}
=\sum_{i\ne j}\pi_jq_{ji}-\pi_j\sum_{k\ne j}q_{jk}=0.
\]
Stąd $\pi Q=0$ i stacjonarność.
\end{proof}

\begin{stwierdzenie}[Równowaga szczegółowa a odwracalność]
Dla skończonego łańcucha o rozkładzie stacjonarnym
$\pi$ równowaga szczegółowa jest równoważna warunkowi
$\pi_i p_{ij}(t)=\pi_j p_{ji}(t)$ dla wszystkich $i,j$
i $t\ge0$. Przy starcie z rozkładu $\pi$ jest to
równoważne odwracalności w czasie: dla każdego $T>0$ i
$0\le t_1<\cdots<t_m\le T$ wektory
$(X_{t_1},\ldots,X_{t_m})$ oraz
$(X_{T-t_1},\ldots,X_{T-t_m})$ mają ten sam rozkład.
\end{stwierdzenie}
\begin{proof}
Niech $D=\operatorname{diag}(\pi_i:i\in E)$.
Równowaga szczegółowa oznacza $DQ=Q^{\mathsf T}D$.
Przez indukcję $DQ^n=(Q^{\mathsf T})^nD$, zatem
z szeregu wykładnika macierzowego
$DP(t)=P(t)^{\mathsf T}D$. Odwrotnie, różniczkowanie
tej tożsamości w zerze daje $DQ=Q^{\mathsf T}D$.
Nie trzeba zakładać odwracalności macierzy $D$.
Dla ustalonego ciągu stanów $i_1,\ldots,i_m$,
ze stacjonarności i własności Markowa mamy
\[
\pi_{i_1}\prod_{k=1}^{m-1}p_{i_k i_{k+1}}(t_{k+1}-t_k)
=\pi_{i_m}\prod_{k=1}^{m-1}p_{i_{k+1} i_k}(t_{k+1}-t_k).
\]
Równość otrzymujemy, kolejno stosując bilans szczegółowy
dla $P(t_{k+1}-t_k)$, bez dzielenia przez $\pi_i$.
Prawa strona jest prawdopodobieństwem odpowiedniego
zdarzenia dla wektora odwróconych obserwacji.
Odwrotnie, odwracalność dla pary chwil $0,t$ przy
$T=t>0$ daje $\pi_i p_{ij}(t)=\pi_j p_{ji}(t)$.
\end{proof}

Nie należy mylić bilansu całkowitego z równowagą szczegółową:
warunek szczegółowy jest silniejszy i nie zachodzi na przykład
dla jednokierunkowego ruchu po cyklu z co najmniej trzema
stanami, mimo istnienia rozkładu stacjonarnego.

\begin{twierdzenie}[Granica w modelu skończonym]
Niech $E$ będzie skończonym zbiorem stanów, $Q$ macierzą
intensywności, a graf przejść z krawędzią $i\to j$ wtedy
i tylko wtedy, gdy $q_{ij}>0$, niech będzie silnie spójny.
Wtedy istnieje dokładnie jeden rozkład stacjonarny $\pi$;
wszystkie jego współrzędne są dodatnie i dla każdego
$i,j\in E$ zachodzi
$\lim_{t\to\infty}p_{ij}(t)=\pi_j$.
\end{twierdzenie}
\begin{proof}
Weźmy $\Lambda>\max_i q_i$ i $R=I+Q/\Lambda$.
Z silnej spójności grafu intensywności wynika
nieprzywiedlność dyskretnego łańcucha o macierzy $R$,
gdyż $R_{ij}>0$ dla $i\ne j$ dokładnie na dozwolonych
krawędziach. Ponieważ $R_{ii}>0$ dla każdego $i$,
łańcuch ten jest nieokresowy.
Twierdzenie o zbieżności skończonych nieprzywiedlnych
i nieokresowych łańcuchów z tomu I zapewnia istnienie
jedynego stacjonarnego $\pi$, dodatniego w każdym stanie,
i $R^n_{ij}\to\pi_j$ dla wszystkich $i,j$.
Warunek $\pi R=\pi$ jest równoważny $\pi Q=0$,
a więc rozkład jest też stacjonarny dla czasu ciągłego.

Z uniformizacji mamy
$p_{ij}(t)=\sum_{n\ge0}\Prob(N_t=n)R^n_{ij}$,
gdzie $N_t\sim\operatorname{Poisson}(\Lambda t)$.
Dla $\varepsilon>0$ wybierzmy $n_0$ tak, by
$|R^n_{ij}-\pi_j|<\varepsilon$ dla $n\ge n_0$.
Wtedy
\[
|p_{ij}(t)-\pi_j|
\le\Prob(N_t<n_0)+\varepsilon.
\]
Ponieważ $N_t\to\infty$ według prawdopodobieństwa
(gdy $t\to\infty$, choćby z nierówności Czebyszewa),
pierwszy składnik dąży do zera. Dajemy $\varepsilon\downarrow0$.
Jedyność rozkładu stacjonarnego wynika także z tego,
że każdy rozkład stacjonarny dla $Q$ jest stacjonarny dla $R$.
\end{proof}

\begin{twierdzenie}[Ergodyczne średnie czasowe w modelu skończonym]
Niech $X$ będzie łańcuchem na skończonej przestrzeni
stanów $E$, o nieprzywiedlnej macierzy intensywności
i jedynym rozkładzie stacjonarnym $\pi$.
Dla każdej funkcji
$f:E\to\R$ i dowolnego rozkładu początkowego
\[
\lim_{T\to\infty}\frac1T\int_0^T f(X_s)\,
\mathrm d\lambda_1(s)=\sum_{i\in E}\pi_i f(i)
\quad\text{prawie na pewno i w }L^2.
\]
\end{twierdzenie}
\begin{proof}
Najpierw uzasadnimy jednostajne tempo zbieżności
macierzy przejścia. Nieprzywiedlność i uniformizacja
dają $p_{ij}(1)>0$ dla wszystkich $i,j$: istnieje
$n$ z $(R^n)_{ij}>0$, a jego waga Poissona jest dodatnia.
Niech $d=|E|$ i
$\varepsilon=\tfrac12 d\min_{i,j}p_{ij}(1)\in(0,1)$.
Możemy zapisać $P(1)=\varepsilon U+(1-\varepsilon)S$,
gdzie każdy wiersz $U$ jest rozkładem jednostajnym
na $E$, a $S$ jest macierzą stochastyczną.
Dla rozkładów wierszowych $\nu,\eta$ mamy
$(\nu-\eta)U=0$ oraz
$\norm{(\nu-\eta)S}_1\le\norm{\nu-\eta}_1$
z nierówności trójkąta i sum wierszy $S$ równych $1$.
Stąd, używając półgrupy przejścia i stacjonarności,
\[
\norm{\delta_i P(t)-\pi}_1\le
2(1-\varepsilon)^{\lfloor t\rfloor},\qquad t\ge0,
\]
gdzie $\delta_i$ jest rozkładem skupionym w $i$;
dla ułamkowej części czasu stosujemy kontraktywność
macierzy stochastycznej w normie $\ell^1$.

Niech $g=f-\sum_i\pi_i f(i)$. Z powyższego wynika
$\norm{P(t)g}_\infty\le
2\norm{g}_\infty(1-\varepsilon)^{\lfloor t\rfloor}$.
Dla $0\le u\le v$ własność Markowa daje
$\E[g(X_u)g(X_v)]=\E[g(X_u)(P(v-u)g)(X_u)]$.
Wobec ograniczoności $g$ i wspólnej mierzalności
procesu (wynikającej z progresywnej mierzalności)
możemy zastosować twierdzenie Fubiniego.
Całka po kwadracie $[0,T]^2$ jest dwukrotnością
całki po trójkącie $0\le u<v\le T$; przekątna ma
miarę Lebesgue'a zero. Po podstawieniu $w=v-u$
otrzymujemy
\[
\E\left|\frac1T\int_0^T g(X_s)\,\mathrm d\lambda_1(s)\right|^2
\le \frac{4\norm{g}_\infty^2}{T}
\int_0^\infty(1-\varepsilon)^{\lfloor w\rfloor}
\,\mathrm d\lambda_1(w)
=\frac{4\norm{g}_\infty^2}{\varepsilon T}.
\]
Otrzymujemy zbieżność w $L^2$.
Dla $T_n=n^2$ nierówność Czebyszewa i pierwszy
lemat Borela--Cantellego dają zbieżność prawie na pewno
wzdłuż tego ciągu (stosujemy je dla progów $1/k$,
$k\ge1$). Dla $n^2\le T<(n+1)^2$ różnica średnich
w chwilach $T$ i $n^2$ ma moduł nie większy niż
$2\norm{g}_\infty(2n+1)/n^2$, co dąży do zera.
To kończy dowód dla wszystkich $T\to\infty$.
\end{proof}

\begin{przyklad}[Łańcuch trzystanowy na cyklu]
Niech $E=\{0,1,2\}$ i jedyne przejścia mają postać
$0\to1$, $1\to2$, $2\to0$, każde z intensywnością
$\lambda>0$. Wtedy
\[
Q=\lambda\begin{pmatrix}-1&1&0\\0&-1&1\\1&0&-1\end{pmatrix}.
\]
Rozkład $\pi=(1/3,1/3,1/3)$ spełnia $\pi Q=0$,
więc jest stacjonarny. Graf jest silnie spójny,
zatem proces zbiega rozkładem do $\pi$.
Jednak równowaga szczegółowa nie zachodzi, bo
$\pi_0q_{01}=\lambda/3>0$ i $\pi_1q_{10}=0$.
Proces ma więc stacjonarny, niezerowy przepływ wokół cyklu.
\end{przyklad}

\section{Łańcuchy urodzin i śmierci}
\begin{definicja}[Proces urodzin i śmierci]
Niech $E=\N_0$. Procesem urodzin i śmierci nazywamy
minimalny łańcuch Markowa o jedynych dozwolonych przejściach
$n\to n+1$ z intensywnością $\lambda_n\ge0$ i
$n\to n-1$ z intensywnością $\mu_n\ge0$ dla $n\ge1$,
przy $\mu_0:=0$. Pozostałe wyrazy pozadiagonalne są zerowe,
a $q_{nn}=-(\lambda_n+\mu_n)$.
\end{definicja}

Nazwa nawiązuje do populacji, w której każda zmiana liczebności
jest zwiększeniem lub zmniejszeniem o jeden. Ta sama
matematyka opisuje jednak liczbę klientów w kolejce, cząstek
w systemie czy liczebność aktywnych urządzeń.

\begin{stwierdzenie}[Rozkład stacjonarny przy równowadze lokalnej]
Załóżmy, że $\lambda_n>0$ dla $n\ge0$,
$\mu_n>0$ dla $n\ge1$, oraz że proces jest nieeksplodujący.
Jeżeli
\[
Z:=1+\sum_{n=1}^\infty
\prod_{k=0}^{n-1}\frac{\lambda_k}{\mu_{k+1}}<\infty,
\]
to liczby
$\pi_0:=1/Z$ i
$\pi_n:=Z^{-1}\prod_{k=0}^{n-1}\lambda_k/\mu_{k+1}$
dla $n\ge1$ tworzą rozkład probabilistyczny spełniający
równowagę szczegółową
$\pi_n\lambda_n=\pi_{n+1}\mu_{n+1}$.
W szczególności $\pi$ jest rozkładem stacjonarnym procesu.
Odwrotnie, istnienie rozkładu stacjonarnego wymaga $Z<\infty$,
a powyższy rozkład jest wówczas jedyny.
\end{stwierdzenie}
\begin{proof}
Wszystkie $\pi_n$ są nieujemne, a ich suma z definicji
jest równa $Z/Z=1$. Przez skrócenie kolejnych iloczynów
otrzymujemy $\pi_n\lambda_n=\pi_{n+1}\mu_{n+1}$.
Wykazaliśmy zatem równowagę szczegółową.
Pozostaje udowodnić stacjonarność \emph{dla procesu
o przeliczalnej przestrzeni stanów}; nie wystarczy tylko zapisać
formalnej równości $\pi Q=0$.

Dla $K\in\N$ rozważmy łańcuch na $\{0,1,\ldots,K\}$
z tymi samymi intensywnościami, z wyjątkiem tego,
że tłumimy przejście $K\to K+1$.
Jest to skończony, nieeksplodujący proces,
a jego rozkład
\[
\pi_n^{(K)}:=\frac{\pi_n}{\sum_{k=0}^K\pi_k},
\qquad 0\le n\le K,
\]
spełnia równowagę szczegółową dla wszystkich jego krawędzi,
a więc jest stacjonarny w myśl twierdzenia skończonego.

Dla ustalonego stanu początkowego $i\in\N_0$ skonstruujmy
sprzężenie procesu pierwotnego i obciętego przy użyciu
wspólnych zegarów potencjalnych przejść w górę i w dół.
W stanie $n$ zegary mają niezależne czasy wykładnicze
o parametrach $\lambda_n,\mu_n$ (parametr zero oznacza
czas nieskończony). Minimum ma parametr sumy, a kierunek
jest niezależny od minimum, z prawdopodobieństwami
proporcjonalnymi do parametrów, co wynika z iloczynu
funkcji przeżycia i gęstości. W stanie $K$ proces
obcięty ignoruje zegar przejścia w górę; zachowuje
zegar przejścia w dół. Do pierwszego wyjścia procesu pierwotnego poza
$\{0,\ldots,K\}$ oba procesy są jednakowe (bez
samej chwili wyjścia). Po rozdzieleniu proces
obcięty może dokończyć oczekiwanie według już
wylosowanego zegara w dół; jego czas przebywania
od wejścia do $K$ jest więc dokładnie wykładniczy
z parametrem $\mu_K$. Po zakończeniu tego pobytu
używamy nowych niezależnych zegarów dla procesu
obciętego. W ten sposób dla każdego $K$ mamy sprzężenie
o właściwych rozkładach brzegowych. Ponieważ oryginalny
proces jest nieeksplodujący,
na każdym $[0,t]$ wykonuje tylko skończenie wiele skoków
prawie na pewno; jego trajektoria ma zatem na $[0,t]$
skończone maksimum. W opisanym sprzężeniu różnica
prawdopodobieństw przejścia jest oszacowana przez
$\Prob_i(\sup_{0\le v\le t}X_v>K)$, które dąży do zera.
Stąd dla każdego $i,j$ i $t$
\[
P_t^{(K)}(i,j)\longrightarrow P_t(i,j)
\quad (K\to\infty,\ K\ge i),
\]
gdzie $P_t^{(K)}$ oznacza macierz przejścia łańcucha
obciętego.

Teraz ustalmy $j$ i $t$. Stacjonarność $\pi^{(K)}$ daje
\[
\pi_j^{(K)}=\sum_{i=0}^K\pi_i^{(K)}P_t^{(K)}(i,j).
\]
Dla ustalonego $i$ mamy $\pi_i^{(K)}\to\pi_i$
i zbieżność odpowiednich prawdopodobieństw przejścia.
Ponieważ $\sum_{i\ge0}\pi_i=1$, ogon
$\sum_{i>M}\pi_i$ jest dowolnie mały, a dla dużych $K$
$\pi_i^{(K)}\le2\pi_i$. Stąd możemy najpierw przepuścić
$K\to\infty$ w skończonej sumie po $i\le M$,
a następnie $M\to\infty$ kontrolując ogon przez
$2\sum_{i>M}\pi_i$. Otrzymujemy
$\pi_j=\sum_{i\ge0}\pi_iP_t(i,j)$.
Ponieważ $j,t$ były dowolne, $\pi$ jest stacjonarna.

Dla konieczności niech $\nu$ będzie dowolnym rozkładem
stacjonarnym. Ustalmy $n\ge0$ i funkcję
$f=\chi_{\{0,\ldots,n\}}$ na $\N_0$.
Jej generator, liczony za pomocą macierzy intensywności
procesu urodzin i śmierci, jest ograniczony:
\[
g(i)=
\begin{cases}
-\lambda_n,&i=n,\\
\mu_{n+1},&i=n+1,\\
0,&\text{w pozostałych stanach}.
\end{cases}
\]
Dla każdego $K\ge n+1$ generator skończonego procesu
obciętego, zastosowany do ograniczenia $f$, jest
ograniczeniem tej samej funkcji $g$. Ustalmy stan
początkowy $i$ i bierzmy $K\ge\max(i,n+1)$.
Wzór Dynkina (twierdzenie~\ref{thm:dynkin}) daje
\[
\E_i f(X_t^{(K)})-f(i)
=\int_0^t\E_i g(X_s^{(K)})\,\mathrm d\lambda_1(s).
\]
Z powyższego sprzężenia wynika zbieżność wartości
oczekiwanych każdej ograniczonej funkcji przy ustalonej
chwili $s$: moduł różnicy jest ograniczony przez
dwukrotność jej normy supremum pomnożoną przez
$\Prob_i(\sup_{0\le v\le s}X_v>K)$.
Stosując twierdzenie o zbieżności zmajoryzowanej
pod całką, otrzymujemy wzór Dynkina dla tej konkretnej
funkcji $f$ i procesu nieobciętego.
Uśredniamy teraz względem $\nu$; zamiana sumy z całką
jest uzasadniona ograniczonością $g$ i $\sum_i\nu_i=1$.
Stacjonarność daje dla $t>0$
\[
0=t(-\nu_n\lambda_n+\nu_{n+1}\mu_{n+1}).
\]
Zatem $\nu$ spełnia równowagę szczegółową na każdej
krawędzi, a indukcja prowadzi do
$\nu_n=\nu_0\prod_{k=0}^{n-1}\lambda_k/\mu_{k+1}$.
Gdyby $\nu_0=0$, wszystkie współrzędne byłyby zerowe.
Normalizacja wymaga więc $Z<\infty$ oraz $\nu_0=1/Z$,
co dowodzi konieczności i jedyności. Nie zakładaliśmy
ograniczoności wszystkich intensywności.
\end{proof}

\begin{uwaga}[Wersja nieskończona wymaga ostrożności]
W poprzednim stwierdzeniu założyliśmy nieeksplodowanie.
W ogólnych modelach na przestrzeni przeliczalnej
formalna równość $\pi Q=0$ nie jest automatycznie
równoważna stacjonarności minimalnego procesu,
gdy masa może uciekać przez eksplozję. Z tego powodu
nie przenosimy bez komentarza wyników z przypadku
skończonego na wszystkie macierze $Q$.
\end{uwaga}


\begin{uwaga}[Dlaczego obcięcie nieskończonego łańcucha wymaga sprzężenia?]
W dowodzie stacjonarności procesu urodzin i śmierci
przybliżamy łańcuch na $\N_0$ przez łańcuchy na
$\{0,\ldots,K\}$. Nie wystarczy tylko stwierdzić, że
macierze intensywności są podobne: w stanie brzegowym $K$
mają inne czasy przebywania. Można jednak \emph{sprzęgnąć}
procesy, korzystając z tych samych niezależnych zegarów
potencjalnych przejść $n\to n+1$ i $n\to n-1$.
W łańcuchu obciętym ignorujemy wyłącznie impulsy
$K\to K+1$; pozostałe są wspólne dla obu procesów.
Do pierwszego wejścia oryginalnego procesu w $K+1$
trajektorie obu procesów są więc identyczne.
Przy założonym braku eksplozji liczba skoków procesu
oryginalnego na $[0,t]$ jest prawie na pewno skończona,
co oznacza, że na tym przedziale osiąga tylko skończenie
wiele stanów. Prawdopodobieństwo rozbieżności trajektorii w tym
sprzężeniu na $[0,t]$ nie przekracza
$\Prob_i(\sup_{0\le v\le t}X_v>K)$ i dąży do zera.
Nie potrzebujemy zakładać wspólnej realizacji wszystkich
przybliżeń dla różnych $K$.
To uzasadnia zbieżność ich praw przejścia, a dopiero
następnie można przejść do granicy w równaniu stacjonarności.
\end{uwaga}

\begin{uwaga}[Zanik pamięci a własność Markowa]
Własność Markowa mówi, że po znajomości $X_t$ cała
wcześniejsza historia nie jest potrzebna do prognozowania
przyszłości. Nie oznacza, że rozkład $X_t$ przestaje
zależeć od stanu początkowego w krótkim czasie.
Asymptotyczny zanik tej zależności wynika w modelu
skończonym z nieprzywiedlności i twierdzenia o granicy
$P(t)$; dla niektórych procesów na przestrzeniach
nieskończonych taki wynik może być fałszywy.
\end{uwaga}

\section{Model kolejki \texorpdfstring{$M/M/1$}{M/M/1}}
\begin{przyklad}[Pojedynczy serwer]
Klienci napływają zgodnie z procesem Poissona o intensywności
$\lambda>0$. Czasy obsługi są niezależne, mają rozkład
$\operatorname{Exp}(\mu)$ z $\mu>0$ i są niezależne
od napływu. Jeden serwer obsługuje co najwyżej jednego klienta
naraz, a klienci czekają bez ograniczenia pojemności.
Stanem $X_t$ jest liczba osób w systemie (obsługiwanych
i oczekujących). Przy $X_t=n$ nowy klient zwiększa stan do
$n+1$ z intensywnością $\lambda$, a przy $n\ge1$
zakończenie obsługi zmniejsza stan do $n-1$ z intensywnością
$\mu$. Jest to proces urodzin i śmierci z
$\lambda_n=\lambda$ i $\mu_n=\mu$ dla $n\ge1$.

Na mocy twierdzenia~\ref{thm:bounded-rates} proces
jest nieeksplodujący, ponieważ $q_n\le\lambda+\mu$.
Z poprzedniego stwierdzenia otrzymujemy dla
$\rho:=\lambda/\mu<1$:
\[
\pi_n=(1-\rho)\rho^n,\qquad n\in\N_0.
\]
Konieczność warunku $\rho<1$ oraz jedyność rozkładu
wynikają już z ogólnego kryterium: tutaj $Z=\sum_{n\ge0}\rho^n$.
Dla ilustracji można też wyprowadzić równania bilansu
bezpośrednio, korzystając z ograniczoności intensywności.
Każdy rozkład stacjonarny spełnia $\pi Q=0$: dla
$h\downarrow0$ prawdopodobieństwa przejścia spełniają
$p_{ij}(h)=\chi_{\{i=j\}}+hq_{ij}+O(h^2)$,
z błędem jednostajnym względem $i,j$.
Istotnie, uniformizacja z zegarem o intensywności
$\Lambda>\lambda+\mu$ działa także w tym modelu
przeliczalnym: $R=I+Q/\Lambda$ jest stochastyczna,
a prawdopodobieństwo co najmniej dwóch impulsów jest
$O(h^2)$, niezależnie od stanu początkowego.
Rozkład stacjonarny spełnia $\pi P(h)=\pi$,
więc po odjęciu $\pi$, podzieleniu przez $h$ i przejściu
do granicy można sumować po $i$ dzięki jednostajnemu
oszacowaniu błędu i $\sum_i\pi_i=1$.
Równanie bilansu w stanie $0$ brzmi:
$\lambda\pi_0=\mu\pi_1$.
Bilans w stanie $n\ge1$ brzmi
$\lambda\pi_{n-1}+\mu\pi_{n+1}=(\lambda+\mu)\pi_n$.
Indukcyjnie wynika stąd $\pi_{n+1}=\rho\pi_n$, gdyż
równość dla $n=0$ jest początkiem indukcji, a podstawienie
$\lambda\pi_{n-1}=\mu\pi_n$ do równania dla $n$
prowadzi do $\mu\pi_{n+1}=\lambda\pi_n$.
Stąd $\pi_n=\rho^n\pi_0$. Musi być $\pi_0>0$,
gdyż w przeciwnym razie wszystkie $\pi_n$ byłyby zerowe.
Możliwość normalizacji wymaga zatem $\rho<1$. Wtedy $\pi_0=1-\rho$.
W tym stanie równowagi
$\Prob(X_t=0)=1-\rho$ i
$\E X_t=\sum_{n\ge0}n(1-\rho)\rho^n=\rho/(1-\rho)$.
Ostatni wzór otrzymujemy różniczkując szereg geometryczny
$\sum_{n\ge0}\rho^n=(1-\rho)^{-1}$ dla $0<\rho<1$.

Przy $\rho\ge1$
proces nadal istnieje, lecz nie ma rozkładu stacjonarnego
normalizowalnego na $\N_0$.
\end{przyklad}

\begin{zadanie}[Kolejka i liczba klientów]
Do stanowiska przychodzi średnio $3$ klientów na godzinę,
a jeden serwer obsługuje średnio $5$ klientów na godzinę;
napływ i czasy obsługi spełniają założenia modelu $M/M/1$.
Obliczyć w stanie stacjonarnym: (a) prawdopodobieństwo
pustego systemu; (b) prawdopodobieństwo obecności co
najmniej trzech klientów; (c) średnią liczbę klientów.

\emph{Rozwiązanie.} Mamy $\rho=3/5$. Zatem
$\pi_0=1-\rho=2/5$.
Dalej
$\Prob(X\ge3)=\sum_{n=3}^\infty(1-\rho)\rho^n
=\rho^3=27/125$.
Średnia wynosi $\rho/(1-\rho)=\frac{3/5}{2/5}=3/2$.
Wszystkie te wyniki dotyczą \emph{stanu stacjonarnego};
nie są poprawnymi odpowiedziami na pytania o układ,
który właśnie został uruchomiony z zerową liczbą klientów.
\end{zadanie}

\section{Model Ehrenfesta w czasie ciągłym}
\begin{przyklad}[Dwie komory]
Niech $K\ge1$ będzie liczbą cząstek niezależnie
przemieszczających się
między komorami $A$ i $B$. Każda cząstka w $B$ przechodzi
do $A$ z intensywnością $\alpha>0$, a każda cząstka w $A$
przechodzi do $B$ z intensywnością $\beta>0$.
Niech $X_t$ oznacza liczbę cząstek w $A$ w chwili $t$.
Gdy $X_t=i$, w $B$ znajduje się $K-i$ cząstek, więc
$\lambda_i=(K-i)\alpha$, a $\mu_i=i\beta$.
Istotnie, minimum niezależnych wykładniczych czasów
do zmian komór poszczególnych cząstek ma intensywność
$(K-i)\alpha+i\beta$, a prawdopodobieństwa kierunków
są proporcjonalne do odpowiednich sum intensywności.
Po zmianie stanu pozostałe zegary zachowują
własność braku pamięci. Prawo następnego czasu
i kierunku zależy więc jedynie od liczby $i$,
nie od tożsamości cząstek w komorze.
Łańcuch na $\{0,\ldots,K\}$ jest nieprzywiedlny.

Dla $p:=\alpha/(\alpha+\beta)$ rozkład
\[
\pi_i=\binom Ki p^i(1-p)^{K-i},\qquad 0\le i\le K,
\]
jest stacjonarny. Rzeczywiście, dla $0\le i<K$
\begin{align*}
\frac{\pi_{i+1}}{\pi_i}
&=\frac{K-i}{i+1}\frac{p}{1-p}
=\frac{(K-i)\alpha}{(i+1)\beta}
=\frac{\lambda_i}{\mu_{i+1}},
\end{align*}
co dokładnie oznacza równowagę szczegółową
$\pi_i\lambda_i=\pi_{i+1}\mu_{i+1}$.
Rozkład dwumianowy pojawia się zatem nie tylko
przy niezależnych próbach, lecz także jako prawo
stacjonarne łańcucha w czasie ciągłym.
\end{przyklad}

\section{Czas pochłonięcia i macierz fundamentalna}
\begin{twierdzenie}[Oczekiwany czas pochłonięcia]
\label{thm:ctmc-absorption}
Niech $E=T\mathbin{\dot\cup}A$ będzie skończonym zbiorem stanów,
$A\ne\varnothing$ zbiorem stanów pochłaniających,
a $T\ne\varnothing$ zbiorem stanów niepochłaniających. Przy uporządkowaniu
najpierw stanów $T$ macierz intensywności ma blokową postać
\[
Q=\begin{pmatrix}B&C\\0&0\end{pmatrix}.
\]
Załóżmy, że startując z każdego $i\in T$, proces
osiąga $A$ z prawdopodobieństwem jeden.
Wtedy $B$ jest odwracalna, macierz
$G:=(-B)^{-1}$ ma nieujemne wyrazy i
\[
G_{ij}=\E_i\left[\int_0^\tau\chi_{\{X_t=j\}}\,\mathrm d\lambda_1(t)\right],
\qquad
\E_i\tau=\sum_{j\in T}G_{ij},
\]
gdzie $\tau=\inf\{t\ge0:X_t\in A\}$, a
$\E_i$ oznacza wartość oczekiwaną dla procesu
startującego ze stanu $i$.
\end{twierdzenie}
\begin{proof}
Stany z $A$ są pochłaniające, dlatego dla $i,j\in T$
$\Prob_i(X_t=j,\tau>t)=\Prob_i(X_t=j)$.
Macierz przejścia ograniczona do $T$ jest równa $e^{tB}$:
wynika to z blokowo-trójkątnej postaci $e^{tQ}$,
którą widać bezpośrednio z potęg macierzy $Q$.
Stąd całkując funkcję wskaźnikową i stosując Tonellego,
\[
\E_i\int_0^\tau\chi_{\{X_t=j\}}\,\mathrm d\lambda_1(t)
=\int_0^\infty\Prob_i(X_t=j,\tau>t)\,\mathrm d\lambda_1(t)
=\int_0^\infty(e^{tB})_{ij}\,\mathrm d\lambda_1(t).
\]

Musimy wykazać zbieżność tej całki oraz odwracalność $B$.
Ponieważ proces z każdego stanu $i\in T$ zostaje pochłonięty
prawie na pewno, istnieje $t_i>0$ takie, że
$\Prob_i(\tau>t_i)<1/2$. Biorąc $t_0=\max_i t_i$,
z monotoniczności funkcji przeżycia dostajemy
$\sup_{i\in T}\Prob_i(\tau>t_0)\le1/2$.
Ponieważ $A$ jest pochłaniający, $\{\tau>t\}$
jest równoważne $\{X_t\in T\}$. Własność Markowa
w chwilach $nt_0$ daje dla $i\in T$
\[
\Prob_i(\tau>(n+1)t_0)
=\E_i[\chi_{\{\tau>nt_0\}}\Prob_{X_{nt_0}}(\tau>t_0)]
\le\tfrac12\Prob_i(\tau>nt_0).
\]
Indukcyjnie $\sup_{i\in T}\Prob_i(\tau>nt_0)\le2^{-n}$.
Zatem wierszowe sumy macierzy $e^{tB}$ maleją
wykładniczo do zera, a nieujemne wyrazy tej macierzy
są całkowalne na $[0,\infty)$. Niech
$G=\int_0^\infty e^{tB}\,\mathrm d\lambda_1(t)$.
Z równania $\frac{\dd}{\dd t}e^{tB}=Be^{tB}$ wynika
\[
BG=\lim_{t\to\infty}e^{tB}-I=-I.
\]
W skończonym wymiarze istnienie prawej odwrotności
macierzy kwadratowej oznacza odwracalność, a
$G=(-B)^{-1}$. Nieujemność wyrazów jest widoczna
z całkowej reprezentacji. Korzystając z tożsamości
$\int_0^\tau\sum_{j\in T}\chi_{\{X_t=j\}}\,\mathrm d\lambda_1(t)=\tau$
(na $[0,\tau)$ proces pozostaje w $T$)
i liniowości wartości oczekiwanej dla skończonej sumy,
otrzymujemy drugi wzór.
\end{proof}

\begin{zadanie}[Pełne obliczenie czasu pochłonięcia]
Rozważmy łańcuch z trzema stanami i intensywnościami
$0\to1$ z intensywnością $2$, $1\to0$ z intensywnością $1$
oraz $1\to2$ z intensywnością $3$.
Stan $2$ jest pochłaniający. Obliczyć oczekiwany czas
osiągnięcia stanu $2$ przy starcie ze stanu $0$ i $1$,
a także oczekiwany łączny czas spędzony w stanach $0$ i $1$
przed pochłonięciem przy starcie z $0$.

\emph{Rozwiązanie.} Macierz intensywności jest równa
\[
Q=\begin{pmatrix}-2&2&0\\1&-4&3\\0&0&0\end{pmatrix},
\quad
B=\begin{pmatrix}-2&2\\1&-4\end{pmatrix}.
\]
Pochłonięcie zachodzi prawie na pewno: przy każdym wejściu
do stanu $1$ następny skok prowadzi do $2$ z
prawdopodobieństwem $3/4$, a powroty do stanu $1$ po skoku
do $0$ następują prawie na pewno. Prawdopodobieństwo
uniknięcia stanu $2$ po $m$ wizytach w $1$ wynosi $(1/4)^m$,
co dąży do zera.
Z twierdzenia~\ref{thm:ctmc-absorption}
\[
G=(-B)^{-1}
=\begin{pmatrix}2&-2\\-1&4\end{pmatrix}^{-1}
=\frac16\begin{pmatrix}4&2\\1&2\end{pmatrix}.
\]
Przy starcie z $0$ oczekiwane czasy zajmowania kolejno
stanu $0$ i stanu $1$ wynoszą $4/6=2/3$ oraz $2/6=1/3$.
Ich suma to $\E_0\tau=1$. Przy starcie z $1$ suma
elementów drugiego wiersza daje $\E_1\tau=1/2$.

\emph{Sprawdzenie bez macierzy.} Oznaczmy
$u_i=\E_i\tau$, $u_2=0$.
Ze stanu $0$ oczekujemy średnio $1/2$ jednostki czasu
i przechodzimy do $1$, więc $u_0=1/2+u_1$.
Ze stanu $1$ czekamy średnio $1/4$, po czym z
prawdopodobieństwem $1/4$ przechodzimy do $0$,
a z $3/4$ do $2$. Stąd
$u_1=1/4+(1/4)u_0$. Rozwiązanie daje
$u_0=1$, $u_1=1/2$ zgodnie z rachunkiem macierzowym.
To jest praktyczna metoda \emph{analizy pierwszego skoku}.
\end{zadanie}

\section{Jak czytać prawdopodobieństwa trajektorii?}
Przy procesie dyskretnym używaliśmy wzoru
$\Prob(X_0=i_0,\ldots,X_n=i_n)$.
W czasie ciągłym należy rozróżnić obserwacje
w \emph{ustalonych chwilach} od zarejestrowania
\emph{dokładnych czasów skoków}.

\begin{stwierdzenie}[Wspólny rozkład obserwacji w ustalonych chwilach]
Niech $X$ będzie skończonym łańcuchem w czasie ciągłym,
$0=t_0<t_1<\cdots<t_n$, a $\mu$ rozkładem początkowym.
Wtedy dla każdego $i_0,\ldots,i_n\in E$
\[
\Prob(X_{t_0}=i_0,\ldots,X_{t_n}=i_n)
=\mu_{i_0}\prod_{k=1}^n
 p_{i_{k-1},i_k}(t_k-t_{k-1}).
\]
\end{stwierdzenie}
\begin{proof}
Jeżeli prawdopodobieństwo wcześniejszej części zdarzenia
jest dodatnie, stosujemy łańcuchowy wzór na prawdopodobieństwo
warunkowe. Własność Markowa daje
\[
\Prob(X_{t_k}=i_k\mid X_{t_0}=i_0,\ldots,X_{t_{k-1}}=i_{k-1})
=p_{i_{k-1},i_k}(t_k-t_{k-1}).
\]
Mnożymy kolejne czynniki, rozpoczynając od
$\Prob(X_{t_0}=i_0)=\mu_{i_0}$.
Jeżeli prawdopodobieństwo wcześniejszego zdarzenia jest zero,
to prawdopodobieństwo całej ścieżki obserwacji również
jest zero, a iloczyn po prawej zawiera zerowy czynnik
na pierwszym etapie, na którym otrzymujemy zerowy iloczyn.
W ten sposób wzór obejmuje także zdarzenia zerowe.
\end{proof}

\begin{uwaga}[Prawdopodobieństwo dokładnej chwili skoku]
Jeżeli $J_1$ jest pierwszą chwilą skoku ze stanu $i$ o
$q_i>0$, to $J_1\sim\operatorname{Exp}(q_i)$ i
$\Prob(J_1=t)=0$ dla każdego ustalonego $t\ge0$.
Nie pytamy więc o \emph{prawdopodobieństwo} dokładnego czasu,
lecz o jego \emph{gęstość}. Dla ściśle rosnących
$0<s_1<\cdots<s_m<t$ oraz ciągu rzeczywistych zmian
$i_0\to i_1\to\cdots\to i_m$ gęstość zdarzenia:
``skoki nastąpiły w chwilach $s_1,\ldots,s_m$,
a do $t$ nie nastąpił następny skok'' jest równa
\[
\mu_{i_0}\left[\prod_{k=1}^m
q_{i_{k-1},i_k}e^{-q_{i_{k-1}}(s_k-s_{k-1})}\right]
e^{-q_{i_m}(t-s_m)},\qquad s_0:=0.
\]
Wynika to z mnożenia gęstości wykładniczych czasów
przebywania i prawdopodobieństw wyboru kolejnych stanów:
$[q_i e^{-q_i u}](q_{ij}/q_i)=q_{ij}e^{-q_i u}$.
To zapis, od którego zaczyna się matematyczne formułowanie
wiarygodności ciągłych trajektorii procesów skokowych;
sama statystyczna estymacja należy jednak do tomu IV.
\end{uwaga}

\section{Podsumowanie i połączenie z procesem Wienera}
Proces Poissona miał \emph{stałą} intensywność skoków,
a jego stan zawsze wzrastał o jeden. Łańcuch w czasie
ciągłym dopuszcza zależną od stanu intensywność $q_i$,
różne kierunki zmian oraz pochłanianie. Procesy te nadal
mają trajektorie schodkowe.

W następnych rozdziałach zajmiemy się procesami gaussowskimi
i procesem Wienera, których ciągłe trajektorie różnią się
zasadniczo od trajektorii Poissona. Pomimo tej różnicy
oba światy połączą wspólne narzędzia: filtracje, martyngały,
generatory i równania Kołmogorowa. Jeszcze później
zbudujemy całkę stochastyczną i miarę Wienera na przestrzeni
trajektorii, a następnie zbadamy jej związek z całkami
funkcyjnymi w mechanice kwantowej.


% ==============================================================
% TOM II — ROZDZIAŁ 6: PROCESY GAUSSOWSKIE
% ==============================================================

\chapter{Procesy gaussowskie: od kowariancji do prawa procesu}
\label{ch:gauss}

\section{Motywacja: czy średnia i kowariancja opisują cały proces?}
W rozdziale pierwszym zobaczyliśmy, że znajomość rozkładu $X_t$ dla
każdego $t$ nie wystarcza do określenia prawa procesu: potrzebujemy
rozkładów \emph{łącznych} dla wszystkich skończonych zbiorów chwil.
W ogólności obliczanie tych rozkładów jest bardzo trudne. Procesy
gaussowskie są niezwykłe właśnie dlatego, że wszystkie ich rozkłady
skończeniewymiarowe można odtworzyć z dwóch funkcji:
$m(t)=\E X_t$ oraz $K(s,t)=\Cov(X_s,X_t)$.

Warto od razu oddzielić \emph{gaussowskość pojedynczych zmiennych}
od \emph{gaussowskości całego procesu}. Można mieć $X_s\sim\Normal(0,1)$
i $X_t\sim\Normal(0,1)$, choć para $(X_s,X_t)$ nie ma rozkładu
normalnego dwuwymiarowego. Dopiero odpowiedni warunek na każdą
skończoną rodzinę rozstrzyga o gaussowskości procesu.

\begin{definicja}[Wektor gaussowski]
Niech $d\in\N$ i $Z=(Z_1,\ldots,Z_d)$ będzie wektorem losowym
o wartościach w $\R^d$. Nazywamy go \emph{wektorem gaussowskim},
jeżeli dla każdego $a\in\R^d$ zmienna $\sum_{j=1}^d a_jZ_j$
ma rozkład normalny (dopuszczamy rozkład zdegenerowany, czyli
skoncentrowany w jednym punkcie). Jeżeli $m=\E Z$ oraz
$\Sigma=(\Cov(Z_i,Z_j))_{i,j=1}^d$, zapisujemy
$Z\sim\Normal_d(m,\Sigma)$.
\end{definicja}

Dopuszczenie degeneracji nie jest tylko wygodą techniczną. Dla
$Z=(U,U)$, gdzie $U\sim\Normal(0,1)$, wektor $Z$ jest gaussowski,
ale jego rozkład jest skoncentrowany na przekątnej $\{(x,x):x\in\R\}$,
więc nie ma gęstości względem dwuwymiarowej miary Lebesgue'a.

\begin{stwierdzenie}[Charakterystyka wektora gaussowskiego]
Dla każdego $d\ge1$, $m\in\R^d$ i symetrycznej nieujemnie
określonej macierzy $\Sigma\in\R^{d\times d}$ istnieje wektor
$Z\sim\Normal_d(m,\Sigma)$, którego funkcja charakterystyczna ma postać
\[
\E e^{i\langle u,Z\rangle}
=\exp\left(i\langle u,m\rangle-\frac12 u^\top\Sigma u\right),
\qquad u\in\R^d.
\]
Prawo tego wektora jest jednoznacznie wyznaczone przez $m$ i $\Sigma$.
\end{stwierdzenie}
\begin{proof}
Z twierdzenia spektralnego mamy $\Sigma=ADA^\top$, gdzie $A$
jest ortogonalna, $D=\operatorname{diag}(\lambda_1,\ldots,\lambda_d)$
i $\lambda_j\ge0$. Niech $G_1,\ldots,G_d$ będą niezależnymi
zmiennymi $\Normal(0,1)$ i połóżmy
$Z=m+AD^{1/2}G$. Wówczas
$\langle u,Z-m\rangle$ jest liniową kombinacją niezależnych
zmiennych normalnych, a zatem jest normalna o wariancji
$u^\top AD A^\top u=u^\top\Sigma u$. Jego funkcję charakterystyczną
otrzymujemy, mnożąc funkcje charakterystyczne składników.
Z jednoznaczności funkcji charakterystycznej wynika jednoznaczność
rozkładu wektora $Z$. Istotnie, dla dowolnego wektora
gaussowskiego o średniej $m$ i kowariancji $\Sigma$ zmienna
$\langle u,Z\rangle$ ma z definicji rozkład normalny,
o średniej $\langle u,m\rangle$ i wariancji $u^\top\Sigma u$.
Jej jednowymiarowa funkcja charakterystyczna w punkcie $1$
daje ten sam wzór, także gdy wariancja jest zerowa.
\end{proof}

\begin{definicja}[Proces gaussowski]
Niech $T\ne\varnothing$ będzie zbiorem indeksów oraz
$(\Omega,\calF,\Prob)$ przestrzenią probabilistyczną. Proces
rzeczywisty $(X_t)_{t\in T}$ na tej przestrzeni nazywamy
\emph{gaussowskim}, jeżeli dla dowolnego $n\ge1$ oraz
$t_1,\ldots,t_n\in T$ wektor $(X_{t_1},\ldots,X_{t_n})$
jest gaussowski. Definiujemy
$m(t):=\E X_t$ i $K(s,t):=\E[(X_s-m(s))(X_t-m(t))]$.
\end{definicja}

\begin{twierdzenie}[Średnia i kowariancja wyznaczają prawo]
Niech $X=(X_t)_{t\in T}$ i $Y=(Y_t)_{t\in T}$ będą rzeczywistymi
procesami gaussowskimi (niekoniecznie na tej samej przestrzeni).
Jeżeli dla wszystkich $s,t\in T$ zachodzą równości
$\E X_t=\E Y_t$ i $\Cov(X_s,X_t)=\Cov(Y_s,Y_t)$, to $X$ i $Y$
mają te same rozkłady skończeniewymiarowe, a więc to samo
prawo na $\R^T$ z $\sigma$-ciałem cylindrycznym.
\end{twierdzenie}
\begin{proof}
Ustalmy $t_1,\ldots,t_n$. Oba odpowiadające im wektory są
gaussowskie, mają wektor średnich
$(m(t_1),\ldots,m(t_n))$ i tę samą macierz kowariancji
$(K(t_i,t_j))_{i,j}$. Z poprzedniego stwierdzenia ich funkcje
charakterystyczne są identyczne. Jednoznaczność funkcji
charakterystycznej daje równość rozkładów. Ponieważ wybór chwil był
dowolny, otrzymujemy równość rozkładów skończeniewymiarowych.
Zdarzenia cylindryczne tworzą $\pi$-układ generujący
$\sigma$-ciało cylindryczne, więc jednoznaczność miar
na tym układzie daje równość praw procesów.
\end{proof}

\begin{uwaga}[Dlaczego warunek wspólnej normalności jest istotny?]
Niech $U\sim\Normal(0,1)$, a $\varepsilon$ będzie od niego
niezależnym symetrycznym znakiem o wartościach $\pm1$.
Wtedy $X=U$, $Y=\varepsilon U$ mają standardowe rozkłady normalne,
a $\Cov(X,Y)=0$, ponieważ $\E\varepsilon=0$.
Jednak $|X|=|Y|$ prawie na pewno. Dla $p=\Prob(|U|\le1)\in(0,1)$
mamy $\Prob(|X|\le1,|Y|\le1)=p\ne p^2$,
zatem zmienne są zależne,
a para $(X,Y)$ nie jest wektorem gaussowskim.
Wektor gaussowski o takiej macierzy kowariancji miałby niezależne
współrzędne, co jest sprzeczne z $|X|=|Y|$.
\end{uwaga}

\section{Jakie funkcje mogą być kowariancjami?}
Zachodzi $K(s,t)=K(t,s)$, ale sama symetria nie wystarcza.
Dla dowolnych $t_1,\ldots,t_n\in T$ i $a_1,\ldots,a_n\in\R$
mamy
\[
\sum_{i,j=1}^n a_i a_jK(t_i,t_j)
=\Var\left(\sum_{i=1}^n a_iX_{t_i}\right)\ge0.
\]
Warunek ten nazywamy \emph{nieujemną określonością jądra} $K$.

\begin{twierdzenie}[Istnienie procesu o zadanej kowariancji]
Niech $T\ne\varnothing$ będzie dowolnym zbiorem,
$m:T\to\R$ dowolną funkcją, zaś $K:T\times T\to\R$
niech spełnia warunki: $K(s,t)=K(t,s)$ oraz
$\sum_{i,j=1}^n a_i a_jK(t_i,t_j)\ge0$ dla wszystkich
$n\in\N$, $t_1,\ldots,t_n\in T$ i $a_1,\ldots,a_n\in\R$.
Wtedy istnieje przestrzeń probabilistyczna i rzeczywisty proces
gaussowski $X=(X_t)_{t\in T}$, dla którego
$\E X_t=m(t)$ oraz $\Cov(X_s,X_t)=K(s,t)$.
\end{twierdzenie}
\begin{proof}
Dla każdej skończonej uporządkowanej rodziny chwil
$\mathbf t=(t_1,\ldots,t_n)$ wybierzmy rozkład gaussowski
$\mu_{\mathbf t}=\Normal_n(m_{\mathbf t},\Sigma_{\mathbf t})$,
gdzie $m_{\mathbf t}=(m(t_1),\ldots,m(t_n))$,
$\Sigma_{\mathbf t}=(K(t_i,t_j))_{i,j}$. Jest to dobrze określony
rozkład, bo macierz kowariancji jest nieujemnie określona.

Trzeba sprawdzić zgodność tej rodziny. Usunięcie części współrzędnych
z wektora gaussowskiego daje wektor gaussowski, którego średnia i
kowariancja są odpowiednimi podwektorami i podmacierzami.
Zmiana kolejności współrzędnych jedynie permutuje średnią i macierz.
Przy powtórzeniu chwili $t_i=t_j$ różnica odpowiednich współrzędnych
ma wariancję $K(t_i,t_i)+K(t_j,t_j)-2K(t_i,t_j)=0$ i średnią zero;
zatem współrzędne są równe prawie na pewno.
Warunki zgodności obowiązują więc także dla powtórzonych chwil.

Możemy teraz zastosować udowodnione w rozdziale~\ref{ch:konstrukcja}
twierdzenie Kołmogorowa o istnieniu procesu, które daje proces
z zadanymi rozkładami skończeniewymiarowymi. Wszystkie te rozkłady
są gaussowskie, a ich średnie i kowariancje są równe $m,K$.
\end{proof}

\begin{stwierdzenie}[Reprezentacja jądra kowariancji w przestrzeni Hilberta]
Dla symetrycznego jądra $K:T\times T\to\R$ nieujemna
określoność jest równoważna istnieniu rzeczywistej
przestrzeni Hilberta $H$ i odwzorowania $\Phi:T\to H$
takiego, że
\[
K(s,t)=\langle\Phi(s),\Phi(t)\rangle_H.
\]
\end{stwierdzenie}
\begin{proof}
Dla jądra nieujemnie określonego poprzednie twierdzenie
daje wycentrowany proces $X$ o kowariancji $K$.
Wystarczy przyjąć $H=L^2(\Omega,\calF,\Prob)$
i $\Phi(t)=X_t$ (klasa równości prawie na pewno).
Wówczas $\langle X_s,X_t\rangle_H=K(s,t)$.
Odwrotnie, dla takiej reprezentacji
\[
\sum_{i,j=1}^n a_ia_jK(t_i,t_j)
=\left\|\sum_{i=1}^n a_i\Phi(t_i)\right\|_H^2\ge0.
\]
Możemy ograniczyć $H$ do domkniętej przestrzeni
liniowej rozpiętej przez $\Phi(T)$.
\end{proof}

\begin{definicja}[Kanoniczna pseudometryka kowariancji]
Dla procesu $X$ o skończonych drugich momentach
definiujemy
\[
d_K(s,t)=\sqrt{K(s,s)+K(t,t)-2K(s,t)}.
\]
Jest to pseudometryka: jest symetryczna, nieujemna,
znika na przekątnej i spełnia nierówność trójkąta,
lecz może znikać dla różnych indeksów.
\end{definicja}
Istotnie, po wycentrowaniu procesu mamy
$d_K(s,t)=\|(X_s-m(s))-(X_t-m(t))\|_{L^2}$,
więc wszystkie własności wynikają z własności normy.
Równość $d_K(s,t)=0$ oznacza równość wycentrowanych
zmiennych prawie na pewno, niekoniecznie równość
ich średnich. Na zbiorze klas równoważności relacji
$s\sim t$ wtedy i tylko wtedy, gdy $d_K(s,t)=0$,
pseudometryka indukuje metrykę: jej niezależność
od reprezentantów wynika z nierówności trójkąta.

\begin{przyklad}[Losowa prosta i losowa fala]
Niech $A,B$ będą niezależne o rozkładzie $\Normal(0,1)$.
Proces $X_t=A+tB$ ma średnią zero i kowariancję
$K_X(s,t)=1+st$. Jego trajektorie są prostymi; proces jest
gaussowski, ponieważ każda skończona rodzina jego wartości
jest liniowym obrazem wektora $(A,B)$.

Dla ustalonej $\omega_0>0$ proces
$Y_t=A\cos(\omega_0t)+B\sin(\omega_0t)$ również jest gaussowski,
a z tożsamości $\cos a\cos b+\sin a\sin b=\cos(a-b)$
wynika $K_Y(s,t)=\cos(\omega_0(t-s))$. Średnia jest zerowa,
a kowariancja nie zmienia się przy wspólnym przesunięciu
argumentów, więc z jednoznaczności praw gaussowskich
wynika stacjonarność w sensie rozkładów skończeniewymiarowych, choć każda trajektoria
jest gładką, okresową funkcją. Stacjonarność nie wymaga więc
ani niezależnych wartości, ani nieregularnych trajektorii.
\end{przyklad}

\begin{przyklad}[Kowariancja wykładnicza: proces typu Ornsteina--Uhlenbecka]
Ustalmy $\sigma>0$ i $\alpha>0$. Zdefiniujmy
$K(s,t)=\sigma^2 e^{-\alpha|t-s|}$ dla $s,t\ge0$.
Rzeczywiście jest to jądro nieujemnie określone, ponieważ
\[
K(s,t)=2\alpha\sigma^2\int_{-\infty}^{\min(s,t)}
 e^{-\alpha(s-u)}e^{-\alpha(t-u)}\,\mathrm d\lambda_1(u).
\]
Dla dowolnych $a_i,t_i$ odpowiadająca mu forma kwadratowa jest
całką z nieujemnego kwadratu:
\[
\sum_{i,j}a_ia_jK(t_i,t_j)=2\alpha\sigma^2\int_{-\infty}^\infty
\left(\sum_i a_i\chi_{\{u\le t_i\}}e^{-\alpha(t_i-u)}\right)^2\mathrm d\lambda_1(u)\ge0.
\]
Poprzednie twierdzenie, zastosowane z $m\equiv0$, daje
wycentrowany proces gaussowski $Y$ o tej kowariancji.
Niezmienniczość średniej i kowariancji względem przesunięć
czasu dowodzi jego stacjonarności. Co więcej,
$\E|Y_t-Y_s|^4=3\,[2\sigma^2(1-e^{-\alpha|t-s|})]^2
\le12\sigma^4\alpha^2|t-s|^2$.
Kryterium ciągłości Kołmogorowa
(twierdzenie~\ref{thm:kol-cont}) zapewnia istnienie ciągłej
modyfikacji na każdym $[0,n]$. Na częściach wspólnych
modyfikacje są nierozróżnialne, co wynika z ich ciągłości
i równości prawie na pewno w każdej wymiernej chwili.
Po przecięciu przeliczalnej rodziny zbiorów pełnej miary
sklejamy je w jedną ciągłą modyfikację na $[0,\infty)$,
a poza tym zbiorem przyjmujemy trajektorię zerową. Nie użyliśmy przy tym całki Itô, którą zdefiniujemy
dopiero w dalszych rozdziałach.
\end{przyklad}

\section{Niezależność i warunkowanie dla procesów gaussowskich}
\begin{twierdzenie}[Nieskorelowane podwektory gaussowskie są niezależne]
Niech $(U,V)$ będzie wektorem gaussowskim o wartościach
w $\R^k\times\R^m$. Jeżeli $\Cov(U_i,V_j)=0$
dla wszystkich $i,j$, to wektory $U$ i $V$ są niezależne.
\end{twierdzenie}
\begin{proof}
Macierz kowariancji wektora $(U,V)$ jest blokowo diagonalna:
$\Sigma=\operatorname{diag}(\Sigma_U,\Sigma_V)$.
Jego funkcja charakterystyczna w punkcie $(u,v)$ jest zatem
iloczynem funkcji charakterystycznych $U$ i $V$,
ponieważ wyrażenie kwadratowe rozdziela się na sumę
$u^\top\Sigma_Uu+v^\top\Sigma_Vv$.
Z jednoznaczności funkcji charakterystycznych wnioskujemy,
że rozkład $(U,V)$ jest iloczynem praw brzegowych.
\end{proof}

\begin{twierdzenie}[Warunkowy rozkład gaussowski]
Niech $(X,Y)$ będzie dwuwymiarowym wektorem gaussowskim,
$\E X=m_X$, $\E Y=m_Y$, $\Var X=\sigma_X^2$,
$\Var Y=\sigma_Y^2>0$ i $\Cov(X,Y)=c$.
Wówczas
\[
X\mid(Y=y)\ \sim\ \Normal\left(
 m_X+\frac{c}{\sigma_Y^2}(y-m_Y),\,
 \sigma_X^2-\frac{c^2}{\sigma_Y^2}\right)
\]
stanowi wersję regularnego rozkładu warunkowego $X$ przy zadanej
wartości $Y=y$. W szczególności
$\E[X\mid\sigma(Y)]=m_X+\frac{c}{\sigma_Y^2}(Y-m_Y)$
i $\Var(X\mid\sigma(Y))=\sigma_X^2-c^2/\sigma_Y^2$ p.n.
\end{twierdzenie}
\begin{proof}
Połóżmy $a=c/\sigma_Y^2$ oraz
$Z=X-m_X-a(Y-m_Y)$. Para $(Z,Y)$ jest gaussowska,
a $\Cov(Z,Y)=c-a\sigma_Y^2=0$.
Z poprzedniego twierdzenia $Z$ i $Y$ są niezależne.
Bezpośredni rachunek daje $\E Z=0$ i
$\Var Z=\sigma_X^2-c^2/\sigma_Y^2\ge0$.
Otrzymaliśmy rozkład $X=m_X+a(Y-m_Y)+Z$,
w którym $Z$ jest niezależna od $Y$.
Dla dowolnej ograniczonej mierzalnej funkcji $f$ i dowolnego
zbioru borelowskiego $A$ niezależność oraz twierdzenie Fubiniego
dają
\[
\E[f(X)\chi_{\{Y\in A\}}]
=\int_A\left(\int_\R f(m_X+a(y-m_Y)+z)\,\mathrm d\Law(Z)(z)\right)
\mathrm d\Law(Y)(y).
\]
Dla $f=\chi_B$ z borelowskim $B$ wewnętrzna całka
jest borelowską funkcją $y$ (z mierzalności funkcji
podcałkowej i twierdzenia Tonellego). Dla ustalonego $y$
jest to prawdopodobieństwo zbioru $B$ względem rozkładu
normalnego. Otrzymujemy zatem jądro probabilistyczne.
Dla ogólnego $f$ wewnętrzna całka jest całką względem
zadeklarowanego rozkładu normalnego (także w przypadku wariancji zero), co dowodzi
własności rozkładu warunkowego. Dla $f(x)=x$ i $x^2$
stosujemy obcięcia i przejście do granicy, wykorzystując
całkowalność drugich momentów, uzyskując obie formuły.
\end{proof}

\begin{uwaga}[Wiele obserwacji i problem najlepszego przewidywania]
Jeżeli $Y\in\R^m$ jest wektorem gaussowskim z odwracalną macierzą
kowariancji $\Sigma_Y$, a $(X,Y)$ jest wspólnie gaussowski,
to ten sam argument z resztą
$Z=X-m_X-\Cov(X,Y)\Sigma_Y^{-1}(Y-m_Y)$ dowodzi wzoru
\[
\E[X\mid\sigma(Y)]
=m_X+\Cov(X,Y)\Sigma_Y^{-1}(Y-m_Y).
\]
Reszta jest nieskorelowana z każdą współrzędną $Y$,
więc niezależna od całego $Y$. To objaśnia, dlaczego
w modelu gaussowskim \emph{najlepsza predykcja średniokwadratowa}
jest funkcją afiniczną zaobserwowanych wartości. Istotnie,
dla $g(Y)\in L^2$ różnica $X-\E[X\mid\sigma(Y)]$
jest ortogonalna do każdej $\sigma(Y)$-mierzalnej
zmiennej z $L^2$, zatem
\[
\E|X-g(Y)|^2=\E|X-\E[X\mid\sigma(Y)]|^2
+\E|\E[X\mid\sigma(Y)]-g(Y)|^2.
\]
Minimum jest osiągane dokładnie przez wersje warunkowej
wartości oczekiwanej.
Gdy $\Sigma_Y$ jest osobliwa, należy uwzględnić liniowe
zależności między obserwacjami. Zamiast odwrotności
zastosujemy poniżej pseudoodwrotność Moore'a--Penrose'a.
\end{uwaga}

\begin{stwierdzenie}[Kryterium Markowa dla procesu gaussowskiego]
Niech $X=(X_t)_{t\in I}$ będzie rzeczywistym, wycentrowanym
procesem gaussowskim na przedziale $I$.
Wówczas $X$ jest procesem Markowa względem swojej
filtracji naturalnej wtedy i tylko wtedy, gdy
\[
K(s,u)=K(s,t)K(t,t)^+K(t,u)
\quad\text{dla wszystkich }s<t<u\text{ w }I,
\]
gdzie $a^+=1/a$ dla $a>0$, a $0^+=0$.
W szczególności przy $K(t,t)=0$ warunek ma postać
$K(s,u)=0$ dla wszystkich $s<t<u$.
\end{stwierdzenie}
\begin{proof}
Ustalmy $t$ i oznaczmy $a_v=K(v,t)K(t,t)^+$.
Jeżeli $K(t,t)=0$, to $X_t=0$ prawie na pewno
i $K(v,t)=0$ dla wszystkich $v$ z nierówności
Cauchy'ego--Schwarza. W obu przypadkach
$a_v K(t,t)=K(v,t)$.
Reszty $R_s=X_s-a_sX_t$ oraz $R_u=X_u-a_uX_t$
tworzą wektor wspólnie gaussowski z $X_t$ i jako para
są niezależne od $X_t$ (obie kowariancje z $X_t$ są zerowe). Obliczenie daje
$\Cov(R_s,R_u)=K(s,u)-K(s,t)K(t,t)^+K(t,u)$.

Jeżeli zachodzi wskazana tożsamość dla każdych $s<t<u$,
to dla dowolnych skończonych zbiorów chwil
$s_1,\ldots,s_k<t<u_1,\ldots,u_m$ trzy bloki
$(R_{s_j})_j$, $(R_{u_l})_l$ oraz $X_t$ mają wspólnie
gaussowski rozkład o blokowo diagonalnej macierzy
kowariancji. Funkcja charakterystyczna rozkłada się
na iloczyn trzech czynników, więc bloki są łącznie
niezależne. W szczególności wektor przyszłych reszt
jest niezależny od $\sigma(X_t,R_s:s<t)$: dla
zdarzeń cylindrycznych wynika to z powyższej
niezależności, a argument $\pi$--$\lambda$ rozszerza
ją na całe $\sigma$-ciało. Ponieważ
$\sigma(X_t,R_s:s<t)=\sigma(X_s:s\le t)$,
dla każdej ograniczonej borelowskiej funkcji $h$ mamy
\[
\E[h(X_u)\mid\sigma(X_s:s\le t)]
=H(X_t),\quad H(x)=\E h\left(
a_ux+R_u\right).
\]
Mierzalność $H$ wynika z twierdzenia Fubiniego,
a tożsamość z niezależności $R_u$ od przeszłości.
Warunkując jeszcze względem $\sigma(X_t)$,
otrzymujemy własność Markowa.

Odwrotnie, własność Markowa daje warunkową niezależność
$X_s$ i $X_u$ względem $X_t$ dla $s<t<u$: dla
ograniczonych funkcji przeszłości i przyszłości stosujemy
własność wieży i warunkowanie przyszłości względem
$\sigma(X_r:r\le t)$; funkcje nieograniczone dopuszczamy
następnie przez obcięcie w $L^2$. W szczególności $\E[X_u\mid\calF_t^X]
=\E[X_u\mid\sigma(X_t)]=a_uX_t$,
gdzie $\calF_t^X=\sigma(X_s:s\le t)$, a przejście
od funkcji ograniczonych do $X_u$ jest dozwolone przez
zbieżność w $L^2$ i kontraktywność warunkowej wartości
oczekiwanej. Ostatnia równość wynika z niezależności
wycentrowanej reszty $R_u$ od $X_t$, także gdy $X_t$
jest stała prawie na pewno. Zatem $\E[R_u\mid\calF_t^X]=0$.
Ponieważ $R_s$ jest $\calF_t^X$-mierzalna i należy do $L^2$,
przez obcięcie $R_s$ i nierówność Cauchy'ego--Schwarza
otrzymujemy $\E[R_sR_u]=0$. Reszty są wycentrowane,
więc $\Cov(R_s,R_u)=0$, czyli żądana tożsamość.
\end{proof}


\begin{twierdzenie}[Warunkowanie wielowymiarowe: pełny wzór]
Niech $(X,Y)$ będzie wspólnie gaussowskim wektorem,
$X\in\R^k$, $Y\in\R^m$, o średnich $m_X,m_Y$ i blokowej
macierzy kowariancji
$\left(\begin{smallmatrix}\Sigma_X&C\\ C^\top&\Sigma_Y\end{smallmatrix}\right)$.
Niech $\Sigma_Y^+$ oznacza pseudoodwrotność
Moore'a--Penrose'a macierzy kowariancji $\Sigma_Y$:
jeżeli $\Sigma_Y=U\operatorname{diag}(d_1,\ldots,d_m)U^\top$
jest rozkładem spektralnym, to
\[
\Sigma_Y^+=U\operatorname{diag}(d_1^+,\ldots,d_m^+)U^\top,
\qquad
d_j^+=\begin{cases}1/d_j,&d_j>0,\\0,&d_j=0.\end{cases}
\]
Definicja nie zależy od wyboru bazy własnej, ponieważ
operator działa skalarnie na każdej przestrzeni własnej.
Wtedy regularny rozkład warunkowy $X$ przy $Y=y$ ma
wersję normalną z parametrami
\[
\mu(y)=m_X+C\Sigma_Y^+(y-m_Y),\qquad
V=\Sigma_X-C\Sigma_Y^+C^\top.
\]
Macierz $V$ jest nieujemnie określona, także wtedy,
gdy rozkład warunkowy jest zdegenerowany.
Jeżeli $\Sigma_Y$ jest odwracalna, to $\Sigma_Y^+=\Sigma_Y^{-1}$.
\end{twierdzenie}
\begin{proof}
Dla $v\in\ker\Sigma_Y$ zmienna $v^\top(Y-m_Y)$
ma średnią i wariancję zero, więc jest zerowa prawie
na pewno. W konsekwencji
$Cv=\E[(X-m_X)v^\top(Y-m_Y)]=0$.
Zatem $\ker\Sigma_Y\subseteq\ker C$.
Ze wzoru spektralnego wynika, że
$\Sigma_Y^+\Sigma_Y=\Sigma_Y\Sigma_Y^+$ jest rzutem
ortogonalnym na obraz $\Sigma_Y$, a
$\Sigma_Y^+\Sigma_Y\Sigma_Y^+=\Sigma_Y^+$.
Ponieważ $C$ znika na jądrze $\Sigma_Y$,
mamy $C\Sigma_Y^+\Sigma_Y=C$.

Połóżmy $A=C\Sigma_Y^+$ i
$R=X-m_X-A(Y-m_Y)$. Wektor $(R,Y)$ jest gaussowski,
bo jest afinicznym obrazem wektora $(X,Y)$.
Ponieważ
$\Cov(R,Y)=C-A\Sigma_Y=0$,
wektory $R,Y$ są niezależne. Ponadto $\E R=0$.
Rozwijając macierz kowariancji, otrzymujemy
\[
\Cov(R,R)=\Sigma_X-AC^\top-CA^\top+A\Sigma_YA^\top
=\Sigma_X-C\Sigma_Y^+C^\top=V.
\]
Jako macierz kowariancji $V$ musi być nieujemnie określona.
Ponadto $X=m_X+A(Y-m_Y)+R$. Dla ograniczonej funkcji
borelowskiej $h$ na $\R^k$ i zbioru borelowskiego
$D\subset\R^m$, z niezależności oraz twierdzenia Fubiniego,
\[
\E[h(X)\chi_{\{Y\in D\}}]
=\int_D\int_{\R^k}h(m_X+A(y-m_Y)+r)
\,\mathrm d\Law(R)(r)\,\mathrm d\Law(Y)(y).
\]
Dla funkcji charakterystycznych zbiorów borelowskich
mierzalność wewnętrznej całki względem $y$ wynika
z twierdzenia Tonellego; dla ustalonego $y$ jest ona
prawdopodobieństwem względem prawa przesuniętego $R$.
To dokładnie własność regularnego rozkładu warunkowego
o podanych parametrach. Rozkład $Y$ jest skoncentrowany
na $m_Y+\operatorname{im}\Sigma_Y$: wystarczy wybrać
skończoną bazę jądra i przeciąć zbiory pełnej miary,
na których odpowiednie iloczyny skalarne z $Y-m_Y$
są zerowe. Wartości jądra poza tą podprzestrzenią
nie są wyznaczone przez warunek rozkładu warunkowego;
podany wzór wybiera jedną wersję na całym $\R^m$.
\end{proof}

\begin{przyklad}[Powtórzona obserwacja]
Niech $U,V$ będą niezależnymi zmiennymi standardowo
normalnymi, $X=U+V$, a $Y=(U,U)^\top$.
Wtedy
\[
\Sigma_Y=\begin{pmatrix}1&1\\1&1\end{pmatrix},\qquad
\Sigma_Y^+=\tfrac14\begin{pmatrix}1&1\\1&1\end{pmatrix},
\qquad C=\begin{pmatrix}1&1\end{pmatrix}.
\]
Wzór daje $X\mid(Y=y)\sim
\Normal((y_1+y_2)/2,1)$. Na prostej $y_1=y_2$
jest to rozkład $\Normal(y_1,1)$, zgodnie z niezależnością
$V$ od $U$. Powtórzenie tej samej obserwacji nie
zmniejsza wariancji błędu predykcji.
\end{przyklad}

\begin{twierdzenie}[Predykcja względem dowolnej rodziny obserwacji gaussowskich]
\label{thm:gauss-prediction}
Niech $X$ oraz $(Y_j)_{j\in J}$ będą wspólnie gaussowską
rodziną rzeczywistych zmiennych losowych, to znaczy
każda skończona podrodzina jest wektorem gaussowskim.
Niech $\mathcal G=\sigma(Y_j:j\in J)$ i niech
\[
H=\overline{\operatorname{span}\{Y_j-\E Y_j:j\in J\}}^{L^2}
\]
będzie domkniętą podprzestrzenią rzeczywistej przestrzeni
Hilberta $L^2(\Omega,\calF,\Prob)$. Oznaczmy przez $P_H$
rzut ortogonalny na $H$. Wtedy
\[
\E[X\mid\mathcal G]=\E X+P_H(X-\E X).
\]
Reszta $R=X-\E X-P_H(X-\E X)$ jest wycentrowana,
gaussowska i niezależna od $\mathcal G$. W szczególności
wariancja warunkowa jest stała i równa $\E R^2$.
\end{twierdzenie}
\begin{proof}
Niech $Z=P_H(X-\E X)$. Z definicji domknięcia
istnieje ciąg skończonych kombinacji liniowych
wycentrowanych obserwacji $Z_n\to Z$ w $L^2$.
Z podciągu zbieżnego prawie na pewno otrzymujemy
$\mathcal G$-mierzalną wersję $Z$ (definiujemy
granicę na mierzalnym zbiorze jej istnienia, a poza nim
przyjmujemy zero). Ponieważ $\E Z_n=0$, ciągłość
wartości oczekiwanej w $L^2$ daje $\E Z=0$.

Ustalmy skończenie wiele indeksów $j_1,\ldots,j_m$.
Każdy wektor $(X-\E X-Z_n,Y_{j_1},\ldots,Y_{j_m})$
jest gaussowski. Dowolna jego kombinacja liniowa
$S_n$ zbiega w $L^2$ do odpowiadającej kombinacji $S$
wektora $(R,Y_{j_1},\ldots,Y_{j_m})$.
Wtedy $\E S_n\to\E S$ oraz $\Var S_n\to\Var S$,
a dla $u\in\R$ nierówność
$|e^{iuS_n}-e^{iuS}|\le |u||S_n-S|$
daje zbieżność funkcji charakterystycznych.
W granicy otrzymujemy funkcję charakterystyczną
$\Normal(\E S,\Var S)$, także przy wariancji zero.
Z jednoznaczności funkcji charakterystycznej $S$
jest gaussowska. Tym samym $(R,Y_{j_1},\ldots,Y_{j_m})$
jest wspólnie gaussowski.

Z własności rzutu ortogonalnego mamy
$\E[R(Y_j-\E Y_j)]=0$ dla wszystkich $j\in J$.
Nieskorelowanie i wspólna gaussowskość dają niezależność
$R$ od każdej skończonej rodziny obserwacji. Zdarzenia
cylindryczne tworzą $\pi$-układ generujący $\mathcal G$.
Dla ustalonego borelowskiego zbioru $B$ stosujemy
argument $\pi$--$\lambda$ do miar
$A\mapsto\Prob(A\cap\{R\in B\})$ oraz
$A\mapsto\Prob(A)\Prob(R\in B)$ na $\mathcal G$.
Otrzymujemy niezależność $R$ od $\mathcal G$,
a więc $\E[R\mid\mathcal G]=0$ oraz
$\E[R^2\mid\mathcal G]=\E R^2$.
Z rozkładu $X=\E X+Z+R$ wynika teza.
\end{proof}

\begin{uwaga}[Interpretacja geometryczna]
Predykcja gaussowska jest rzutem ortogonalnym nawet wtedy,
gdy rodzina obserwacji jest nieprzeliczalna. Konkretny rzut
można przybliżać kombinacjami skończenie wielu obserwacji.
Nie jest potrzebne założenie o przeliczalności całej
rodziny. W przypadku skończonym współczynniki rzutu
spełniają równania normalne $A\Sigma_Y=C$;
rozwiązanie $A=C\Sigma_Y^+$ prowadzi do poprzedniego
wzoru. Niezależność reszty od wszystkich obserwacji
jest mocniejsza niż sama ortogonalność w $L^2$.
\end{uwaga}

\begin{przyklad}[Prawo przejścia procesu Ornsteina--Uhlenbecka]
Dla wycentrowanego procesu $Y$ o kowariancji
$K(s,t)=\sigma^2e^{-\alpha|t-s|}$, przy $s<t<u$,
\[
K(s,t)K(t,t)^{-1}K(t,u)
=\sigma^2e^{-\alpha(t-s)}e^{-\alpha(u-t)}
=K(s,u).
\]
Kryterium Markowa daje więc własność Markowa względem
filtracji naturalnej. Dla $h\ge0$ rozkład warunkowy
$Y_{t+h}$ względem całej przeszłości do $t$ ma postać
\[
\Normal\bigl(e^{-\alpha h}Y_t,\,
\sigma^2(1-e^{-2\alpha h})\bigr).
\]
Wynika to z warunkowego rozkładu pary $(Y_{t+h},Y_t)$.
Przy $h=0$ otrzymujemy rozkład skupiony w $Y_t$.
Dla $h\to\infty$ średnia warunkowa dąży do zera,
a wariancja do $\sigma^2$. Warunkowy rozkład przy
$Y_t=x$ dąży słabo do $\Normal(0,\sigma^2)$
dla każdego $x$, co można zobaczyć, przedstawiając go
jako prawo $e^{-\alpha h}x+\sigma\sqrt{1-e^{-2\alpha h}}G$,
gdzie $G\sim\Normal(0,1)$, i stosując twierdzenie
o zbieżności zmajoryzowanej do ograniczonych funkcji
ciągłych. Jest to przykład zaniku pamięci o stanie
początkowym w modelu stacjonarnym.
\end{przyklad}

\section{Od kowariancji do ciągłości trajektorii}

\begin{stwierdzenie}[Ciągłość średniokwadratowa a kowariancja]
Niech $I$ będzie przedziałem i niech $X=(X_t)_{t\in I}$
będzie procesem o skończonych drugich momentach.
Następujące warunki są równoważne:
\begin{itemize}
\item odwzorowanie $t\mapsto X_t$ jest ciągłe w $L^2$;
\item funkcja średniej $m$ i jądro kowariancji $K$
są ciągłe (odpowiednio na $I$ i $I\times I$).
\end{itemize}
\end{stwierdzenie}
\begin{proof}
Rozkład przyrostu na część wycentrowaną i średnią daje
\[
\E|X_t-X_s|^2=|m(t)-m(s)|^2+d_K(s,t)^2.
\]
Stąd ciągłość $m,K$ implikuje ciągłość w $L^2$.
Odwrotnie, z nierówności Cauchy'ego--Schwarza
$|\E X_t-\E X_s|\le\|X_t-X_s\|_{L^2}$.
Ciągły jest więc również proces wycentrowany
$Z_t=X_t-m(t)$ jako odwzorowanie do $L^2$.
Jeżeli $(s_n,t_n)\to(s,t)$, to
\[
|K(s_n,t_n)-K(s,t)|
\le\|Z_{s_n}-Z_s\|_{L^2}\|Z_{t_n}\|_{L^2}
+\|Z_s\|_{L^2}\|Z_{t_n}-Z_t\|_{L^2}.
\]
Normy $\|Z_{t_n}\|_{L^2}$ są ograniczone, gdyż
$Z_{t_n}\to Z_t$ w $L^2$. Prawa strona dąży do zera.
\end{proof}

\begin{wniosek}[Ciągłość stochastyczna procesu gaussowskiego]
Dla rzeczywistego procesu gaussowskiego na przedziale
ciągłość stochastyczna jest równoważna ciągłości
średniokwadratowej.
\end{wniosek}
\begin{proof}
Ciągłość w $L^2$ implikuje ciągłość stochastyczną
na mocy nierówności Markowa. Odwrotnie, niech
$t_n\to t$ i $X_{t_n}-X_t\to0$ według prawdopodobieństwa.
Oznaczmy średnią i wariancję przyrostu przez $a_n,b_n$.
Zbieżność według prawdopodobieństwa daje zbieżność
jego funkcji charakterystycznej do $1$ w każdym
punkcie $u\in\R$ (można zastosować ograniczoność
funkcji wykładniczej i kryterium podciągów).
Dla $u=1$ moduł tej funkcji wynosi $e^{-b_n/2}$,
więc $b_n\to0$. Wycentrowane przyrosty zbiegają
zatem do zera w $L^2$, a więc również według
prawdopodobieństwa. Odejmując je od przyrostów,
otrzymujemy $a_n\to0$ według prawdopodobieństwa;
ponieważ $a_n$ są deterministyczne, jest to
zwykła zbieżność liczb. W konsekwencji
$\E|X_{t_n}-X_t|^2=a_n^2+b_n\to0$.
\end{proof}

\begin{twierdzenie}[Różniczkowanie średniokwadratowe]
\label{thm:gauss-ms-derivative}
Niech $X$ będzie procesem o skończonych drugich
momentach na otwartym przedziale $I$. Załóżmy, że
$m\in C^1(I)$ oraz $K\in C^2(I\times I)$.
Wówczas dla każdego $t\in I$ istnieje pochodna
średniokwadratowa
\[
D_t=\lim_{h\to0}\frac{X_{t+h}-X_t}{h}
\quad\text{w }L^2.
\]
Zachodzą wzory
\[
\E D_t=m'(t),\qquad
\Cov(D_s,D_t)=\partial_s\partial_t K(s,t),
\qquad
\Cov(X_s,D_t)=\partial_t K(s,t).
\]
Proces $D$ jest ciągły średniokwadratowo.
Jeżeli $X$ jest gaussowski, to rodzina jego wartości
i pochodnych średniokwadratowych jest wspólnie
gaussowska.
\end{twierdzenie}
\begin{proof}
Niech $Z_t=X_t-m(t)$ i
$A_h(t)=(Z_{t+h}-Z_t)/h$ dla $h\ne0$, $t+h\in I$.
Dwukrotne zastosowanie podstawowego twierdzenia
rachunku całkowego daje dla dopuszczalnych $h,k$
\[
\E[A_h(s)A_k(t)]
=\int_0^1\int_0^1
\partial_s\partial_t K(s+uh,t+vk)
\,\mathrm d\lambda_1(u)\,\mathrm d\lambda_1(v).
\]
Wzór obowiązuje także dla ujemnych $h,k$, gdyż
stosujemy parametryzację odcinków przez $[0,1]$.
Z ciągłości pochodnej mieszanej prawa strona dąży
do $\partial_s\partial_t K(s,t)$ przy $h,k\to0$.
W szczególności dla $s=t$, rozwijając
$\E|A_h(t)-A_k(t)|^2$, otrzymujemy granicę zero.
Z zupełności $L^2$ wynika istnienie granicy $A(t)$
w $L^2$ (najpierw dla dowolnego ciągu $h_n\to0$,
a warunek Cauchy'ego zapewnia tę samą granicę dla
wszystkich $h\to0$). Mamy $\E A(t)=0$.
Dodając iloraz różnicowy $m$, otrzymujemy
$D_t=A(t)+m'(t)$. Granica iloczynów skalarnych
w $L^2$ daje wzór na kowariancję pochodnych.
Ponadto z ciągłości iloczynu skalarnego
\[
\Cov(X_s,D_t)=\lim_{h\to0}\frac{K(s,t+h)-K(s,t)}{h}
=\partial_t K(s,t).
\]
Funkcja średniej procesu $D$ oraz jego jądro kowariancji
są ciągłe na mocy założeń, więc poprzednie stwierdzenie
daje ciągłość średniokwadratową $D$.
Jeżeli $X$ jest gaussowski, każda skończona kombinacja
liniowa wartości i ilorazów różnicowych jest normalna.
Przechodząc do granicy w $L^2$, stosujemy argument
z funkcjami charakterystycznymi użyty w dowodzie
twierdzenia~\ref{thm:gauss-prediction}. Otrzymujemy wspólną gaussowskość.
\end{proof}

\begin{przyklad}[Gładkość średniokwadratowa losowej fali i procesu Ornsteina--Uhlenbecka]
Dla losowej fali o kowariancji
$K(s,t)=\cos(\omega_0(t-s))$ mamy
$\partial_s\partial_t K(s,t)=\omega_0^2\cos(\omega_0(t-s))$.
Pochodna średniokwadratowa jest tu także pochodną
trajektorii, gdyż bezpośrednio różniczkujemy
$A\cos(\omega_0t)+B\sin(\omega_0t)$, a zbieżność
współczynników daje zbieżność w $L^2$.
Dla procesu Ornsteina--Uhlenbecka natomiast
\[
\E\left|\frac{Y_{t+h}-Y_t}{h}\right|^2
=\frac{2\sigma^2(1-e^{-\alpha|h|})}{h^2}
\longrightarrow\infty\qquad(h\to0).
\]
Wynika to z granicy $(1-e^{-x})/x\to1$ przy
$x\downarrow0$. Ilorazy różnicowe nie mogą więc
zbiegać w $L^2$, choć proces ma ciągłą modyfikację.
Brak pochodnej w $L^2$ sam w sobie, dla dowolnego
procesu, nie jest dowodem braku różniczkowalności
trajektorii. Tutaj można jednak powiedzieć więcej:
ilorazy różnicowe procesu Ornsteina--Uhlenbecka
są wycentrowanymi zmiennymi normalnymi o wariancjach
dążących do nieskończoności. Dla każdego $M>0$
prawdopodobieństwo, że ich moduł nie przekracza $M$,
dąży do zera, ponieważ po standaryzacji otrzymujemy
przedziały o długościach dążących do zera, a gęstość
standardowego rozkładu normalnego jest ograniczona.
Nie mogą więc zbiegać według prawdopodobieństwa do
skończonej zmiennej losowej. Co więcej, dla dowolnego
deterministycznego ciągu $h_n\to0$, $h_n\ne0$,
zdarzenie zbieżności ilorazów do skończonej granicy
ma prawdopodobieństwo zero: zawiera się w sumie po
$M\in\N$ zdarzeń
$\bigcup_N\bigcap_{n\ge N}\{|(Y_{t+h_n}-Y_t)/h_n|\le M\}$,
z których każde ma miarę zero.
Zatem dla każdego ustalonego $t>0$ prawdopodobieństwo
istnienia pochodnej trajektorii w $t$ wynosi zero.
Nie jest to jeszcze twierdzenie o jednoczesnym braku
pochodnej we wszystkich chwilach.
\end{przyklad}

\begin{lemat}[Moment przyrostu gaussowskiego]
Jeżeli $X$ jest rzeczywistym procesem gaussowskim,
a $X_t-X_s$ ma średnią zero oraz wariancję $v(s,t)$,
to dla każdego $p>0$
\[
\E|X_t-X_s|^p=c_p\,v(s,t)^{p/2},\qquad
c_p=\frac{2^{p/2}\Gamma((p+1)/2)}{\sqrt\pi}.
\]
\end{lemat}
\begin{proof}
Gdy $v(s,t)=0$, przyrost wynosi zero prawie na pewno.
W przeciwnym razie $(X_t-X_s)/\sqrt{v(s,t)}$ ma standardowy
rozkład normalny. Podstawienie w całce względem jego gęstości
i zastosowanie definicji funkcji gamma daje
$\E|G|^p=(2/\sqrt{2\pi})\int_0^\infty
 x^p e^{-x^2/2}\mathrm d\lambda_1(x)=c_p$.
\end{proof}

\begin{wniosek}[Praktyczne kryterium ciągłości procesu gaussowskiego]
Niech $X=(X_t)_{t\in[0,T]}$, $T>0$, będzie procesem gaussowskim.
Załóżmy, że istnieją $C>0$, $\eta>0$ takie, że
$\E|X_t-X_s|^2\le C|t-s|^{\eta}$ dla wszystkich $s,t\in[0,T]$.
Nie wymagamy stałej funkcji średniej.
Wtedy $X$ ma ciągłą modyfikację. Co więcej, ma modyfikację,
której trajektorie są h\"olderowskie na $[0,T]$
każdego rzędu $0<\gamma<\min(\eta/2,1)$
na jednym zbiorze pełnej miary.
\end{wniosek}
\begin{proof}
Niech $m(t)=\E X_t$. Nierówność Cauchy'ego--Schwarza
daje $|m(t)-m(s)|\le C^{1/2}|t-s|^{\eta/2}$.
Wariancja przyrostu jest nie większa od jego drugiego
momentu. Stosując poprzedni lemat do wycentrowanego
przyrostu oraz nierówność
$|a+b|^p\le2^{p-1}(|a|^p+|b|^p)$, dla $p\ge2$ dostajemy
\[
\E|X_t-X_s|^p\le
2^{p-1}(1+c_p)C^{p/2}|t-s|^{\eta p/2}.
\]
Wybierzmy $p$ tak duże, by $\eta p/2>1$.
Twierdzenie~\ref{thm:kol-cont} daje ciągłą modyfikację,
a także wykładniki H\"oldera mniejsze niż
$(\eta p/2-1)/p=\eta/2-1/p$ (oraz mniejsze niż $1$).
Dążąc z $p$ do nieskończoności i korzystając z
nierozróżnialności ciągłych modyfikacji, otrzymujemy
jedną modyfikację o wszystkich wskazanych wykładnikach
(po ograniczeniu do przeliczalnego ciągu wykładników
zbiegającego do $\min(\eta/2,1)$).
\end{proof}

\begin{uwaga}[Ciągłość kowariancji to nie całe rozwiązanie]
W szczególności, jeżeli proces gaussowski jest wycentrowany,
to $\E|X_t-X_s|^2=K(t,t)+K(s,s)-2K(s,t)$.
Oszacowanie tej różnicy daje konkretne kryterium regularności.
Nie należy jednak wnioskować automatycznie o ciągłości
trajektorii wyłącznie z samej ciągłości funkcji $K$:
kryterium momentowe wymaga kontroli \emph{szybkości}
zanikania wariancji przyrostów. Bardziej ogólne kryteria
wykorzystują między innymi metryczną entropię.
\end{uwaga}

\begin{przyklad}[Ciągła kowariancja bez ciągłej modyfikacji]
Niech $(G_n)_{n\ge1}$ będą niezależnymi zmiennymi
$\Normal(0,1)$. Na każdym przedziale
$I_n=[2^{-(n+1)},2^{-n}]$ niech $\varphi_n$ będzie
ciągłą funkcją łamaną, równą zero na końcach,
równą $1$ w środku $c_n$ i zerową poza $I_n$.
Na $[0,1]$ definiujemy
\[
X_t=\sum_{n\ge1}\frac{G_n\varphi_n(t)}{\sqrt{\log(n+1)}}.
\]
Dla każdego $t$ co najwyżej jeden składnik jest
niezerowy, więc każda skończona rodzina wartości
jest liniowym obrazem skończonego wektora niezależnych
zmiennych normalnych. Proces jest wycentrowany, a
\[
K(s,t)=\sum_{n\ge1}\frac{\varphi_n(s)\varphi_n(t)}{\log(n+1)}.
\]
Szereg jest jednostajnie zbieżny: dla każdego $(s,t)$
co najwyżej jeden składnik jest niezerowy; ponieważ
$0\le\varphi_n\le1$, supremum ogona po $n>N$
nie przekracza $1/\log(N+2)$.
Zatem $K$ jest ciągłe na $[0,1]^2$ i proces jest
ciągły średniokwadratowo.

Pokażemy, że nie istnieje ciągła modyfikacja.
Niech $a_n=\sqrt{\log(n+1)}$. Dla standardowej
zmiennej normalnej $G$ i $a\ge1$ mamy
\[
\Prob(G>a)\ge
\frac1{\sqrt{2\pi}}\int_a^{a+1/a}e^{-x^2/2}\,
\mathrm d\lambda_1(x)
\ge \frac{e^{-3/2}}{\sqrt{2\pi}a}e^{-a^2/2}.
\]
Ostatnia nierówność wynika z
$(a+1/a)^2/2\le a^2/2+3/2$.
Stąd szereg $\sum_n\Prob(G_n>a_n)$ jest rozbieżny,
gdyż jego wyrazy dla dużych $n$ są ograniczone od dołu
przez dodatnią stałą razy
$1/\sqrt{(n+1)\log(n+1)}$, a $\log(n+1)\le n+1$.
Z niezależności zdarzeń i drugiego lematu
Borela--Cantellego otrzymujemy $X_{c_n}>1$
dla nieskończenie wielu $n$ prawie na pewno.
Ponieważ $c_n\to0$ oraz $X_0=0$, ciągłość w zerze
nie zachodzi. Każda modyfikacja zgadza się z $X$
w chwilach $0,c_1,c_2,\ldots$ na jednym zbiorze
pełnej miary, więc również nie może być ciągła
w zerze prawie na pewno.
\end{przyklad}

\section{Most gaussowski: obserwacja końca zmienia cały proces}
W rozdziale~\ref{ch:konstrukcja} skonstruowaliśmy proces
Wienera $W$; w następnym rozdziale zbadamy jego własności.
Przypomnijmy jego charakterystyczną kowariancję
$\Cov(W_s,W_t)=\min(s,t)$ dla $s,t\ge0$.
Ustalmy $T>0$ i zdefiniujmy na $[0,T]$
\[
\beta_t:=W_t-\frac{t}{T}W_T.
\]
Wówczas $\beta_0=\beta_T=0$, a zatem każda trajektoria
\emph{mostu Browna} zaczyna i kończy się w zerze.

\begin{twierdzenie}[Most Browna i warunkowanie na punkt końcowy]
Niech $W$ będzie standardowym procesem Wienera i $T>0$.
Proces $\beta=(\beta_t)_{0\le t\le T}$ jest wycentrowanym
procesem gaussowskim, niezależnym od zmiennej $W_T$,
z funkcją kowariancji
\[
K_\beta(s,t)=\min(s,t)-\frac{st}{T},
\qquad 0\le s,t\le T.
\]
Ponadto prawo procesu $t\mapsto (t/T)y+\beta_t$
na $C([0,T],\R)$ z borelowskim $\sigma$-ciałem
normy supremum, jako funkcja $y$, stanowi wersję
regularnego rozkładu warunkowego całej trajektorii $W$
przy warunku $W_T=y$.
\end{twierdzenie}
\begin{proof}
Każdy skończony wektor wartości $\beta$ jest liniowym obrazem
wektora wartości $W$, więc jest gaussowski.
Ponieważ $\Cov(W_t,W_T)=t$ oraz $\Var W_T=T$,
otrzymujemy $\Cov(\beta_t,W_T)=t-t=0$.
Wspólna gaussowskość wszystkich skończonych wektorów daje
niezależność $(\beta_{t_1},\ldots,\beta_{t_n})$ od $W_T$.
Przechodząc ze zdarzeń cylindrycznych do $\sigma$-ciała generowanego
przez proces, otrzymujemy niezależność całego $\beta$ od $W_T$.
Rozwijając kowariancję, otrzymujemy
\[
\Cov(\beta_s,\beta_t)=\min(s,t)
-\frac{t}{T}s-\frac{s}{T}t+\frac{st}{T^2}T
=\min(s,t)-\frac{st}{T}.
\]
Na koniec równość $W_t=(t/T)W_T+\beta_t$ i niezależność
pozwalają rozwarunkować względem $W_T$ dokładnie tak,
jak w dowodzie wzoru na warunkowy rozkład gaussowski.
Dla każdej ograniczonej borelowskiej funkcji $F$ skończonej
liczby współrzędnych oraz zbioru $A$ zachodzi
\[
\E\bigl[F(W_{t_1},\ldots,W_{t_n})\chi_{\{W_T\in A\}}\bigr]
=\int_A \E\bigl[F((t_1/T)y+\beta_{t_1},\ldots,
(t_n/T)y+\beta_{t_n})\bigr]\,\mathrm d\Law(W_T)(y),
\]
co dowodzi tezy dla cylindrów. Przyjmujemy ciągłe
wersje procesów, z $W_0=0$ punktowo.
Borelowskie $\sigma$-ciało $C([0,T],\R)$ jest generowane
przez ewaluacje w przeliczalnym gęstym zbiorze chwil:
normę supremum różnicy funkcji ciągłych można obliczać
po tym zbiorze, więc każda otwarta kula jest mierzalna
względem ewaluacji. Przestrzeń $C([0,T],\R)$ jest
ośrodkowa (na przykład łamane na siatkach
$kT/n$ o wymiernych wartościach w węzłach tworzą
przeliczalny zbiór gęsty), zatem każdy zbiór otwarty
jest przeliczalną sumą takich kul. Odwrotna inkluzja
wynika z ciągłości ewaluacji. W szczególności $W$
i $\beta$ są mierzalnymi zmiennymi losowymi
o wartościach w tej przestrzeni.

Oznaczmy przez $\nu$ prawo $\beta$ na tej przestrzeni.
Odwzorowanie
$(y,f)\mapsto f+(\cdot/T)y$ jest ciągłe względem
normy supremum. Dla borelowskiego $B\subset C([0,T],\R)$
definiujemy
\[
Q(y,B)=\int\chi_B(f+(\cdot/T)y)\,\mathrm d\nu(f).
\]
Jest to jądro probabilistyczne: mierzalność względem $y$
wynika z twierdzenia Tonellego, a dla ustalonego $y$
jest to miara będąca obrazem $\nu$ przez przesunięcie.
Dla borelowskiego $A\subseteq\R$ równość
\[
\Prob(W\in B,\ W_T\in A)
=\int_A Q(y,B)\,\mathrm d\Law(W_T)(y)
\]
zachodzi dla cylindrów na mocy wcześniejszego rachunku.
Dla ustalonego $A$ obie strony są miarami skończonymi
względem $B$, więc argument $\pi$--$\lambda$ rozszerza
równość na całe borelowskie $\sigma$-ciało.
To dowodzi tezy dla całych trajektorii. Ponieważ
$\Prob(W_T=y)=0$ dla każdego $y$, chodzi o wersję
jądra, a nie o iloraz prawdopodobieństw zdarzeń.
\end{proof}

\begin{center}
\begin{tikzpicture}[x=0.9cm,y=0.8cm,>=Latex]
\draw[->] (-0.3,0)--(10.3,0) node[right] {$t$};
\draw[->] (0,-1.4)--(0,2.0) node[above] {$\beta_t$};
\draw[dashed,gray] (0,0)--(10,0);
\draw[very thick,blue!70!black]
 plot coordinates {(0,0)(0.6,0.5)(1.2,0.2)(1.8,1.1)(2.5,0.35)(3,0.8)(3.7,0.2)(4.5,-0.8)(5.1,-0.35)(5.9,0.9)(6.6,0.2)(7.3,0.85)(8.1,0.25)(8.8,-0.35)(9.4,0.2)(10,0)};
\node[below] at (0,0) {$0$};
\node[below] at (10,0) {$T$};
\node[above right,blue!70!black] at (2,1.15) {przykładowy most};
\end{tikzpicture}
\par\smallskip\small Schematyczna łamana ilustrująca most Browna: oba końce są
przypięte do zera. Łamana jest jedynie ilustracją, nie
rzeczywistą trajektorią procesu.
\end{center}

\begin{zadanie}[Predykcja wartości procesu Wienera na podstawie dwóch obserwacji]
Niech $W$ będzie procesem Wienera, $0<s<t$ oraz $s\le u\le t$.
Obliczyć $\E[W_u\mid\sigma(W_s,W_t)]$ oraz wariancję warunkową.

\emph{Rozwiązanie.} Z niezależności przyrostów $W_s$, $W_u-W_s$ i $W_t-W_u$
są łącznie niezależne, z uwzględnieniem zerowych
przyrostów na końcach przedziału. Zatem pierwsza
współrzędna wektora $(W_s,W_u-W_s,W_t-W_s)$
jest niezależna od pary pozostałych. W szczególności
$\E[W_u\mid W_s,W_t]=W_s+\E[W_u-W_s\mid W_t-W_s]$.
Obserwowane dwa przyrosty są wspólnie gaussowskie,
a $\Cov(W_u-W_s,W_t-W_s)=u-s$ i
$\Var(W_t-W_s)=t-s$.
Wzór na warunkowy rozkład gaussowski daje
\[
\E[W_u\mid W_s,W_t]
=W_s+\frac{u-s}{t-s}(W_t-W_s),\qquad
\Var(W_u\mid W_s,W_t)=\frac{(u-s)(t-u)}{t-s}.
\]
Zatem predykcją jest interpolacja liniowa między pomiarami,
a wariancja błędu jest zerowa na końcach przedziału
i osiąga maksimum $(t-s)/4$ w jego środku.
\end{zadanie}

\section{Podsumowanie: dlaczego procesy gaussowskie przygotowują ruch Browna?}
Mamy już komplet narzędzi: jądro kowariancji pozwala skonstruować
prawo procesu, formuły warunkowe opisują predykcję, a
oszacowania momentów dają regularność trajektorii.
Następny rozdział pokaże, że prosty wzór $K(s,t)=\min(s,t)$
prowadzi do procesu o bardzo bogatych, a nawet zaskakujących
własnościach geometrycznych.

% ==============================================================
% TOM II — ROZDZIAŁ 7: PROCES WIENERA
% ==============================================================

\chapter{Proces Wienera: własności probabilistyczne i geometria trajektorii}
\label{ch:wiener}
\paragraph{Odsyłacze międzytomowe.} Rachunek całkowy względem Wienera zaczyna się w tomie III, rozdziale~\ref{III-ch:wiener-integral}; zastosowania symulacyjne znajdują się w tomie IV, rozdziale~\ref{IV-ch:monte-carlo}.


\section{Od sum niezależnych przyrostów do ruchu Browna}
Rozważmy błąd pomiaru, który narasta jako suma bardzo wielu drobnych,
w przybliżeniu niezależnych zakłóceń. W granicy czasowej spodziewamy
się procesu o przyrostach gaussowskich, którego wariancja rośnie
proporcjonalnie do długości przedziału. Tak powstaje matematyczna
motywacja standardowego procesu Wienera, czyli ruchu Browna.
Dla każdego $t$ mamy rozkład normalny, ale znacznie silniejszy
jest warunek \emph{niezależności przyrostów na rozłącznych
przedziałach}. To właśnie on odróżnia Wienera od dowolnego
procesu gaussowskiego.

\begin{definicja}[Standardowy proces Wienera]
Niech $(\Omega,\calF,\Prob)$ będzie przestrzenią probabilistyczną.
Proces rzeczywisty $W=(W_t)_{t\ge0}$ nazywamy
\emph{standardowym procesem Wienera} (\emph{standardowym ruchem Browna}),
jeżeli:
\begin{enumerate}
\item $W_0=0$ prawie na pewno;
\item jego trajektorie $t\mapsto W_t(\omega)$ są ciągłe dla prawie
wszystkich $\omega$;
\item dla dowolnych $0\le t_0<t_1<\cdots<t_n$ przyrosty
$W_{t_1}-W_{t_0},\ldots,W_{t_n}-W_{t_{n-1}}$
są wzajemnie niezależne;
\item $W_t-W_s\sim\Normal(0,t-s)$ dla wszystkich $0\le s<t$.
\end{enumerate}
Jeżeli ponadto zadano filtrację $(\calF_t)_{t\ge0}$, to $W$
nazywamy procesem Wienera \emph{względem tej filtracji},
gdy jest do niej adaptowany oraz $W_t-W_s$ jest niezależny
od $\calF_s$ dla wszystkich $s<t$.
\end{definicja}

Rozdział~\ref{ch:konstrukcja} dowiódł już istnienia takiego
procesu na pewnej przestrzeni probabilistycznej.
W dalszym ciągu możemy bez straty ogólności przyjąć,
że wszystkie jego trajektorie są ciągłe: modyfikujemy proces
na jednym zbiorze zerowym, przypisując na nim trajektorię zerową.
Takie ustalenie nie zmienia rozkładów skończeniewymiarowych.

\begin{twierdzenie}[Równoważna charakteryzacja gaussowska]
Niech $X=(X_t)_{t\ge0}$ będzie rzeczywistym procesem
o prawie na pewno ciągłych trajektoriach. Wówczas $X$
jest standardowym procesem Wienera wtedy i tylko wtedy,
gdy jest procesem gaussowskim i dla wszystkich $s,t\ge0$
zachodzi
\[
\E X_t=0,\qquad \Cov(X_s,X_t)=\min(s,t).
\]
\end{twierdzenie}
\begin{proof}
Jeżeli $X=W$ jest procesem Wienera, to dla
$0=t_0<t_1<\cdots<t_n$ wektor $(W_{t_1},\ldots,W_{t_n})$
jest liniowym obrazem wektora niezależnych przyrostów normalnych,
więc jest gaussowski. Przy $s\le t$ korzystamy z rozkładu
$W_t=W_s+(W_t-W_s)$ oraz niezależności przyrostu od $W_s$:
$\Cov(W_s,W_t)=\Var(W_s)=s$. Średnie wynoszą zero.

Odwrotnie, $\Var(X_0)=K(0,0)=0$ i $\E X_0=0$,
a więc $X_0=0$ p.n. Dla $s<t$
$\Var(X_t-X_s)=t+s-2s=t-s$ i średnia przyrostu
jest zero. Jeżeli $[a,b]$ i $[c,d]$ są rozłączne oraz $b\le c$,
to $\Cov(X_b-X_a,X_d-X_c)=0$, co wynika z podstawienia
czterech wartości funkcji $\min(\cdot,\cdot)$.
Dowolne skończone rodziny takich przyrostów są wspólnie
gaussowskie, zatem ich zerowe kowariancje oznaczają niezależność.
Dostajemy wszystkie punkty definicji.
\end{proof}

\begin{przyklad}[Rozkład kilku obserwacji]
Dla $0<s<t$ wektor $(W_s,W_t)$ jest normalny dwuwymiarowo
ze średnią zero i macierzą kowariancji
$\left(\begin{smallmatrix}s&s\\s&t\end{smallmatrix}\right)$.
Jego współrzędne są zależne: $\Cov(W_s,W_t)=s$.
Jednak $W_s$ i $W_t-W_s$ są niezależne.
Dla $0<t_1<t_2<t_3$ mamy na przykład
\[
\Prob(W_{t_3}-W_{t_2}>a\mid W_{t_1}=x)
=1-\Phi\left(\frac{a}{\sqrt{t_3-t_2}}\right),
\]
gdzie równość rozumiana jest jako odpowiednia wersja
regularnego prawdopodobieństwa warunkowego.
Warunek na $W_{t_1}$ nie wpływa na późniejszy rozłączny przyrost.
Nie należy tego mylić z pytaniem o rozkład $W_{t_3}$,
który oczywiście zależy od stanu w chwili $t_1$.
\end{przyklad}

\section{Skalowanie, symetria i niezależność przyszłości}
\begin{twierdzenie}[Symetrie i samopodobieństwo]
Niech $W$ będzie procesem Wienera, a $c>0$.
Procesy $(c^{-1/2}W_{ct})_{t\ge0}$ i
$(-W_t)_{t\ge0}$ mają to samo prawo na przestrzeni
ciągłych trajektorii co $W$. Dla każdego $T>0$ proces
$(W_T-W_{T-t})_{0\le t\le T}$ ma prawo procesu Wienera
ograniczonego do $[0,T]$.
\end{twierdzenie}
\begin{proof}
Każdy z wymienionych procesów ma ciągłe trajektorie,
wartość początkową zero i jest gaussowski.
Dla pierwszego procesu kowariancja wynosi
$c^{-1}\min(cs,ct)=\min(s,t)$; dla drugiego
jest taka sama. Dla trzeciego, gdy $0\le s\le t\le T$,
przyrosty $W_T-W_{T-s}$ oraz $W_T-W_{T-t}$ mają
kowariancję $s$: wspólnym składnikiem są przyrosty
na przedziale $[T-s,T]$. Z poprzedniego twierdzenia
i jednoznaczności prawa procesu otrzymujemy tezy.
\end{proof}

\begin{uwaga}[Jak rozumieć skalowanie?]
Jeśli skala czasu rośnie $c$ razy, typowe wychylenie
rośnie $\sqrt c$ razy, a nie $c$ razy.
To odpowiada zależności $\Var(W_t)=t$.
Samopodobieństwo stanowi podstawę późniejszej zasady
niezmienniczości Donskera, wiążącej spacery losowe z
ruchem Browna.
\end{uwaga}

\begin{twierdzenie}[Własność Markowa i półgrupa przejścia]
Niech $W$ będzie procesem Wienera względem filtracji
$(\calF_t)_{t\ge0}$. Dla dowolnych $s,t\ge0$
i ograniczonej mierzalnej funkcji $f:\R\to\R$ zachodzi
\[
\E[f(W_{s+t})\mid\calF_s]
=(P_tf)(W_s),\qquad
(P_tf)(x):=\int_\R f(x+y)p_t(y)\,\dd\lambda_1(y),
\]
gdzie $p_t(y)=(2\pi t)^{-1/2}\exp(-y^2/(2t))$ dla $t>0$,
a $P_0f=f$. Ponadto $P_{s+t}=P_sP_t$.
\end{twierdzenie}
\begin{proof}
Zapiszmy $W_{s+t}=W_s+Z$, gdzie $Z=W_{s+t}-W_s$
jest niezależna od $\calF_s$ i ma rozkład $\Normal(0,t)$.
Dla funkcji $f(x)=\ind{\{x\in A\}}$ własność
warunkowania niezależnej zmiennej daje żądany wzór.
Przez liniowość uzyskujemy go dla funkcji prostych, a
następnie przez zbieżność monotoniczną i rozkład na
części dodatnią i ujemną dla wszystkich ograniczonych
funkcji mierzalnych.
Aby udowodnić własność półgrupy, bierzemy niezależne
$Z_s\sim\Normal(0,s)$ i $Z_t\sim\Normal(0,t)$.
Ich suma ma rozkład $\Normal(0,s+t)$. Stąd
$P_s(P_tf)(x)=\E f(x+Z_s+Z_t)=P_{s+t}f(x)$;
przejścia $s=0$ lub $t=0$ są oczywiste.
\end{proof}

\section{Mocna własność Markowa: można zacząć od losowej chwili}
Dla deterministycznej chwili $s$ przyszły przyrost
$W_{s+t}-W_s$ jest niezależny od $\calF_s$.
Prawdziwie nowa treść pojawia się przy \emph{losowej}
chwili $\tau$, np. pierwszym trafieniu w poziom $a$.
Nie można po prostu podstawić $s=\tau(\omega)$ w definicji
niezależności: musimy najpierw przybliżyć czas losowy
czasami przyjmującymi przeliczalnie wiele wartości.

\begin{twierdzenie}[Mocna własność Markowa procesu Wienera]
Niech $W$ będzie procesem Wienera względem filtracji
$(\calF_t)_{t\ge0}$, a $\tau$ niech będzie względem niej
czasem zatrzymania takim, że $\Prob(\tau<\infty)=1$.
Wtedy proces $(W_{\tau+t}-W_\tau)_{t\ge0}$ jest
standardowym procesem Wienera i jest niezależny od
sigma-ciała $\calF_\tau$.
\end{twierdzenie}
\begin{proof}
Dla $n\ge1$ niech $h_n=2^{-n}$ i
$\tau_n=h_n\lceil\tau/h_n\rceil$.
Ponieważ $\{\tau_n\le t\}$ jest zdarzeniem określonym
przez warunek $\{\tau\le h_n\lfloor t/h_n\rfloor\}$,
czas $\tau_n$ jest czasem zatrzymania, przyjmuje wartości
w siatce $h_n\N_0$, spełnia $\tau\le\tau_n<\tau+h_n$
i zbiega do $\tau$ prawie na pewno.

Niech $A\in\calF_\tau$. Z nierówności $\tau\le\tau_n$
wynika $A\in\calF_{\tau_n}$. Ustalmy
$0\le t_1<\cdots<t_m$ i $u=(u_1,\ldots,u_m)\in\R^m$.
Dla zdarzenia $A\cap\{\tau_n=kh_n\}\in\calF_{kh_n}$
niezależność przyrostów po chwili $kh_n$ od tej
sigma-algebry daje
\begin{align*}
&\E\left[\ind{A\cap\{\tau_n=kh_n\}}
\exp\left(i\sum_{j=1}^m u_j(W_{\tau_n+t_j}-W_{\tau_n})\right)\right]\\
&\qquad=\Prob(A\cap\{\tau_n=kh_n\})
\exp\left(-\frac12\sum_{j,l=1}^m u_ju_l\min(t_j,t_l)\right).
\end{align*}
Sumujemy po $k\ge0$; jest to dozwolone przez
zbieżność zdominowaną, gdyż moduł każdego czynnika
wykładniczego nie przekracza $1$. Otrzymujemy
\[
\E\left[\ind A\exp\left(i\sum_j u_j
(W_{\tau_n+t_j}-W_{\tau_n})\right)\right]
=\Prob(A)\exp\left(-\frac12\sum_{j,l}
 u_ju_l\min(t_j,t_l)\right).
\]
Ciągłość trajektorii i $\tau_n\to\tau$ zapewniają
$W_{\tau_n+t_j}-W_{\tau_n}\to W_{\tau+t_j}-W_\tau$
prawie na pewno. Ponownie stosujemy zbieżność zdominowaną.
Otrzymana równość dla wszystkich $u$ oznacza, z jednoznaczności
funkcji charakterystycznych, że wektor przyszłych przyrostów
ma standardowy rozkład gaussowski i jest niezależny
od każdego $A\in\calF_\tau$. Niezależność
całej ciągłej trajektorii od $\calF_\tau$ wynika z faktu,
że borelowskie sigma-ciało przestrzeni ciągłych trajektorii
jest generowane przez odczyty w chwilach wymiernych.
Ciągłość przyszłych trajektorii wynika z ciągłości $W$.
\end{proof}

\begin{przyklad}[Co znaczy ``zapominanie'' w losowej chwili?]
Niech $a>0$ i $\tau_a=\inf\{t\ge0:W_t=a\}$.
Na zdarzeniu $\{\tau_a<\infty\}$ obserwator zna
całą przeszłą drogę oraz fakt, że cząstka właśnie dotarła
do $a$. Mocna własność Markowa mówi, że przyszły \emph{przyrost}
$W_{\tau_a+t}-a$ jest niezależnym, nowym ruchem Browna.
Nie oznacza to niezależności samego położenia $W_{\tau_a+t}$
od $W_{\tau_a}$, bo położenie zawiera znany punkt startowy.
\end{przyklad}

\section{Regularność trajektorii: ciągłość bez gładkości}
Dla ustalonych $s<t$ przyrost ma skalę $\sqrt{t-s}$,
a nie $t-s$. Ta pozornie drobna różnica ma zasadniczy
skutek: trajektorie są ciągłe, ale prawie na pewno nie
mają pochodnej w żadnym punkcie.

\begin{twierdzenie}[Regularność H\"oldera poniżej $1/2$]
Dla każdego $T>0$ i standardowego procesu Wienera $W$
zachodzi z prawdopodobieństwem $1$: dla każdego
$\gamma\in(0,1/2)$ istnieje losowa stała
$C_{\gamma,T}(\omega)<\infty$ taka, że
\[
|W_t(\omega)-W_s(\omega)|
\le C_{\gamma,T}(\omega)|t-s|^\gamma,
\qquad s,t\in[0,T].
\]
\end{twierdzenie}
\begin{proof}
Dla każdego $p>2$ i $s,t\in[0,T]$ mamy
$\E|W_t-W_s|^p=c_p|t-s|^{p/2}$.
Kryterium ciągłości Kołmogorowa daje modyfikację
o trajektoriach H\"oldera każdego rzędu
$\gamma<1/2-1/p$. Ponieważ zarówno $W$, jak i ta
modyfikacja mają ciągłe trajektorie, są nierozróżnialne;
możemy więc przypisać tę własność $W$.
Wybierając przeliczalny ciąg $p_k\to\infty$,
przecinamy zbiory o prawdopodobieństwie $1$ i
otrzymujemy jednocześnie wszystkie $\gamma<1/2$.
\end{proof}

\section{Wariacja kwadratowa: niezwykła wielkość trajektorii}
Dla gładkiej funkcji $f$ i podziałów o coraz mniejszej
średnicy suma kwadratów przyrostów $\sum(\Delta f)^2$
dąży do zera. Dla Wienera dzieje się coś innego:
suma ta dąży do długości przedziału. To najważniejszy
mechanizm odróżniający rachunek Itô od klasycznego.

\begin{twierdzenie}[Wariacja kwadratowa procesu Wienera]
Niech $W$ będzie standardowym procesem Wienera, $T>0$, zaś
$\pi_n=\{0=t_{0,n}<\cdots<t_{m_n,n}=T\}$
ciągiem deterministycznych podziałów $[0,T]$ o średnicy
$|\pi_n|:=\max_k(t_{k+1,n}-t_{k,n})\to0$.
Wówczas
\[
Q(\pi_n):=\sum_{k=0}^{m_n-1}
(W_{t_{k+1,n}}-W_{t_{k,n}})^2
\longrightarrow T\quad\text{w }L^2.
\]
Jeżeli $\pi_n$ jest podziałem dyadycznym, to zbieżność
zachodzi również prawie na pewno.
\end{twierdzenie}
\begin{proof}
Oznaczmy $\Delta_{k,n}=W_{t_{k+1,n}}-W_{t_{k,n}}$
oraz $h_{k,n}=t_{k+1,n}-t_{k,n}$.
Przyrosty są niezależne, $\Delta_{k,n}\sim\Normal(0,h_{k,n})$,
a $\E\Delta_{k,n}^2=h_{k,n}$ i
$\E\Delta_{k,n}^4=3h_{k,n}^2$. Zatem
$\Var(\Delta_{k,n}^2)=2h_{k,n}^2$.
Dzięki niezależności
\[
\E Q(\pi_n)=\sum_k h_{k,n}=T,
\qquad
\E\bigl|Q(\pi_n)-T\bigr|^2
=2\sum_k h_{k,n}^2\le 2T|\pi_n|\to0.
\]
Dla siatek dyadycznych $|\pi_n|=T2^{-n}$.
Nierówność Czebyszewa daje, dla każdego $\varepsilon>0$,
$\Prob(|Q(\pi_n)-T|>\varepsilon)
\le 2T^2 2^{-n}/\varepsilon^2$.
Suma po $n$ jest skończona, więc z pierwszego lematu
Borela--Cantellego odchylenie większe niż $\varepsilon$
występuje tylko skończenie wiele razy prawie na pewno.
Stosując to do $\varepsilon=1,1/2,1/3,\ldots$
i przecinając przeliczalnie wiele zdarzeń pełnej miary,
uzyskujemy $Q(\pi_n)\to T$ p.n.
\end{proof}

\begin{wniosek}[Jednostajna zbieżność wariacji kwadratowej]
Dla podziałów dyadycznych $t_{k,n}=kT2^{-n}$
zdefiniujmy dla $t\in[0,T]$
\[
Q_n(t)=\sum_{k=0}^{2^n-1}
\bigl(W_{t\wedge t_{k+1,n}}-W_{t\wedge t_{k,n}}\bigr)^2.
\]
Wówczas $\sup_{0\le t\le T}|Q_n(t)-t|\to0$
prawie na pewno.
\end{wniosek}
\begin{proof}
Niech $h=T2^{-n}$ i
$Y_{k,n}=(W_{(k+1)h}-W_{kh})^2-h$.
Zmienne $Y_{k,n}$ są niezależne, wycentrowane i
$\sum_k\Var Y_{k,n}=2Th$.
Nierówność maksymalna Kołmogorowa daje
\[
\Prob\left(\max_{0\le j\le 2^n}
\left|\sum_{k<j}Y_{k,n}\right|>\varepsilon\right)
\le \frac{2Th}{\varepsilon^2}.
\]
Prawa strona jest sumowalna po $n$, więc maksimum
sum cząstkowych dąży do zera p.n. Gdy
$t\in[jh,(j+1)h]$, różnica między $Q_n(t)-t$
a sumą $\sum_{k<j}Y_{k,n}$ jest równa
$(W_t-W_{jh})^2-(t-jh)$. Jej moduł nie przekracza
$\max_j\sup_{u\in[jh,(j+1)h]}|W_u-W_{jh}|^2+h$.
Ta wielkość dąży do zera p.n. z jednostajnej
ciągłości trajektorii $W$ na $[0,T]$.
Połączenie obu oszacowań dowodzi tezy.
\end{proof}

\begin{wniosek}[Trajektorie nie mają skończonej wariacji całkowitej]
Dla każdego $T>0$ trajektorie $t\mapsto W_t(\omega)$
mają z prawdopodobieństwem $1$ nieskończoną wariację
całkowitą na $[0,T]$. Prawie na pewno jest tak nawet
na każdym przedziale o dodatniej długości.
\end{wniosek}
\begin{proof}
Gdyby pewna ciągła trajektoria $w$ miała skończoną
wariację całkowitą $V$, to dla dowolnego podziału
$\sum_k(w(t_{k+1})-w(t_k))^2
\le\max_k|w(t_{k+1})-w(t_k)|\,V$.
Jednostajna ciągłość $w$ zapewnia zbieżność maksimum
przyrostów do zera wraz ze średnicą podziału.
Stąd suma kwadratów dążyłaby do zera, co przeczy
poprzedniemu twierdzeniu, dającemu granicę $T>0$.
Aby uzyskać twierdzenie dla wszystkich nietrywialnych
przedziałów jednocześnie, stosujemy ten sam dowód do
przedziałów o wymiernych końcach i przecinamy
przeliczalnie wiele zbiorów pełnej miary.
Każdy przedział o dodatniej długości zawiera taki
wymierny podprzedział; wariacja na większym przedziale
nie może być mniejsza od wariacji na mniejszym.
\end{proof}

\begin{przyklad}[Dlaczego lewy i środkowy punkt przedziału dają inne całki?]
Niech $t_k=kT/n$. Tożsamość algebraiczna
$W_{t_{k+1}}^2-W_{t_k}^2
=2W_{t_k}\Delta_kW+(\Delta_kW)^2$
daje
\[
\sum_{k=0}^{n-1}W_{t_k}(W_{t_{k+1}}-W_{t_k})
=\frac12W_T^2-\frac12\sum_{k=0}^{n-1}(\Delta_kW)^2
\longrightarrow \frac12(W_T^2-T)
\quad\text{w }L^2.
\]
Z kolei suma z uśrednieniem wartości na końcach
$\frac12(W_{t_k}+W_{t_{k+1}})$ spełnia
\[
\sum_{k=0}^{n-1}\frac{W_{t_k}+W_{t_{k+1}}}{2}
 (W_{t_{k+1}}-W_{t_k})=\frac12W_T^2
\]
\emph{dokładnie} dla każdego $n$.
To algebraiczny początek rozróżnienia całki Itô i
Stratonovicha. Pełna konstrukcja obu całek pojawi
się w części o rachunku stochastycznym.
\end{przyklad}

\section{Dlaczego ruch Browna nie ma pochodnej?}
Brak skończonej wariacji całkowitej na każdym przedziale nie
wyklucza jeszcze pochodnej w pojedynczym punkcie. Potrzebny
jest silniejszy argument. Pokażemy, że na drobnej siatce
prawdopodobieństwo wystąpienia choćby trzech \emph{jednocześnie
bardzo małych} kolejnych przyrostów maleje wystarczająco szybko.
Gdyby trajektoria miała pochodną, takie trzy małe przyrosty
musiałyby się regularnie pojawiać blisko punktu różniczkowalności.

\begin{twierdzenie}[Nieróżniczkowalność trajektorii Wienera]
Niech $W$ będzie standardowym procesem Wienera i $T>0$.
Wówczas z prawdopodobieństwem $1$ trajektoria
$t\mapsto W_t(\omega)$ nie jest różniczkowalna
w żadnym punkcie $t\in(0,T)$ i nie ma skończonej
pochodnej jednostronnej w punktach $0,T$.
\end{twierdzenie}
\begin{proof}
Ustalmy liczbę całkowitą $M\ge1$ i siatkę dyadyczną
o kroku $h_n=T2^{-n}$. Niech $A_n(M)$ oznacza zdarzenie,
że istnieje $k\in\{0,\ldots,2^n-3\}$, dla którego
trzy kolejne przyrosty
$\Delta_{k,n}W,\Delta_{k+1,n}W,\Delta_{k+2,n}W$
mają moduł nie większy niż $8Mh_n$.
Każdy z tych przyrostów ma rozkład $\Normal(0,h_n)$,
a są one niezależne. Z oszacowania gęstości standardowej
normalnej przez $1/\sqrt{2\pi}$ dostajemy
\[
\Prob(|\Delta_{k,n}W|\le8Mh_n)
=\Prob(|G|\le8M\sqrt{h_n})
\le \frac{16M}{\sqrt{2\pi}}\sqrt{h_n}.
\]
Stąd, z niezależności trzech przyrostów i sumowania
prawdopodobieństw po najwyżej $T/h_n$ pozycjach,
\[
\Prob(A_n(M))\le C\,M^3T\sqrt{h_n},
\]
gdzie $C$ jest stałą liczbową. Szereg po $n$ jest zbieżny,
więc lemat Borela--Cantellego daje
$\Prob(A_n(M)\text{ nieskończenie często})=0$.

Przypuśćmy teraz, że ciągła trajektoria $w$ ma skończoną
pochodną w pewnym $t\in(0,T)$. Z definicji pochodnej
istnieją $M\in\N$ i $\delta>0$ takie, że
$|w(s)-w(t)|\le M|s-t|$ dla $|s-t|<\delta$.
Dla dostatecznie dużego $n$ istnieje $k$ z
$kh_n\le t<(k+1)h_n$, a punkty
$kh_n,(k+1)h_n,(k+2)h_n,(k+3)h_n$
leżą w $(t-\delta,t+\delta)\cap[0,T]$.
Dla $j=0,1,2$ nierówność trójkąta daje
\[
|w((k+j+1)h_n)-w((k+j)h_n)|
\le M(|(k+j+1)h_n-t|+|(k+j)h_n-t|)
\le 8Mh_n.
\]
Zatem $A_n(M)$ zachodzi dla wszystkich dostatecznie
dużych $n$, w sprzeczności z poprzednim wnioskiem.
Przecinając zbiory pełnej miary po $M\in\N$,
wykluczamy wszystkie punkty różniczkowalności wewnątrz
$[0,T]$ naraz.

Dla końca $t=0$, istnienie skończonej pochodnej
prawostronnej oznaczałoby $|W_{h_n}|\le Mh_n$
dla pewnego całkowitego $M$ i wszystkich dużych $n$.
Jednak $\Prob(|W_{h_n}|\le Mh_n)\le C M\sqrt{h_n}$,
co jest sumowalne. Borel--Cantelli wyklucza taką
sytuację. W punkcie $T$ stosujemy analogiczny argument
do $W_T-W_{T-h_n}$. To kończy dowód.
\end{proof}

\begin{wniosek}[Próg wykładnika H\"oldera]
Z prawdopodobieństwem $1$ na żadnym przedziale
$[a,b]\subset[0,\infty)$ o $a<b$ trajektoria Wienera
nie spełnia oszacowania
$|W_t-W_s|\le C|t-s|^\gamma$ dla wszystkich $s,t\in[a,b]$
z jakimkolwiek $C<\infty$ i $\gamma>1/2$.
\end{wniosek}
\begin{proof}
Gdyby takie oszacowanie zachodziło, to dla równomiernego
podziału na $n$ przedziałów suma kwadratów przyrostów
nie przekraczałaby $C^2(b-a)^{2\gamma}n^{1-2\gamma}$,
a więc dążyłaby do zera. To sprzeczne z dodatnią
wariacją kwadratową na każdym przedziale o wymiernych
końcach (przy niemal pewnym wspólnym zdarzeniu).
Każdy przedział $[a,b]$ z $a<b$ zawiera przedział
$[r,s]$ o wymiernych końcach i dodatniej długości;
oszacowanie H\"oldera na $[a,b]$ przeniosłoby się na
$[r,s]$. To daje sprzeczność.
\end{proof}

\begin{center}
\begin{tikzpicture}[x=0.9cm,y=0.8cm,>=Latex]
\draw[->] (-0.2,0)--(10.5,0) node[right] {$t$};
\draw[->] (0,-1.3)--(0,2.6) node[above] {$W_t$};
\draw[blue!75!black,thick] (0,0)--(0.45,0.55)--(0.8,0.18)--(1.18,1.05)--(1.55,0.82)--(1.83,1.5)--(2.14,0.9)--(2.5,1.1)--(2.86,0.22)--(3.2,0.67)--(3.45,0.08)--(3.75,-0.45)--(4.1,0.36)--(4.4,-0.05)--(4.7,0.78)--(5.04,0.42)--(5.37,1.2)--(5.65,0.66)--(5.95,1.3)--(6.25,0.55)--(6.56,0.1)--(6.86,0.88)--(7.21,0.47)--(7.5,1.22)--(7.83,1.0)--(8.15,1.72)--(8.53,1.03)--(8.85,1.4)--(9.2,0.82)--(9.6,1.65)--(10,1.15);
\draw[dashed,red!70!black] (0,1.42)--(10,1.42);
\node[red!70!black,above left] at (10,1.42) {$a$};
\draw[densely dotted,red!70!black] (1.80,0)--(1.80,1.42);
\node[below] at (1.80,0) {$\tau_a$};
\end{tikzpicture}
\par\smallskip\small Schemat trajektorii i pierwszego trafienia w poziom $a$.
Rysunek ma charakter poglądowy: prawdziwa trajektoria Wienera
jest ciągła, ale nie jest łamaną o skończonej liczbie odcinków.
\end{center}

\section{Pierwsze trafienie i zasada odbicia}
\begin{definicja}[Czas pierwszego trafienia]
Niech $W$ będzie ciągłym procesem Wienera i $a>0$.
Definiujemy $\tau_a:=\inf\{t\ge0:W_t=a\}$,
przyjmując $\inf\varnothing=\infty$.
Dla naturalnej filtracji (lub dowolnej filtracji, do której
$W$ jest adaptowany) jest to czas zatrzymania, gdyż
$\{\tau_a\le t\}=\{\sup_{0\le s\le t}W_s\ge a\}$,
a supremum ciągłej trajektorii można obliczać po
chwilach wymiernych z $[0,t]$ wraz z końcem przedziału.
\end{definicja}

\begin{twierdzenie}[Zasada odbicia]
Dla standardowego procesu Wienera $W$, $t>0$ i $a>0$
zachodzi
\[
\Prob\left(\sup_{0\le s\le t}W_s\ge a\right)
=2\Prob(W_t\ge a)
=2\left(1-\Phi\left(\frac{a}{\sqrt t}\right)\right),
\]
gdzie $\Phi$ jest dystrybuantą standardowego rozkładu normalnego.
\end{twierdzenie}
\begin{proof}
Niech $\tau=\tau_a$. Z ciągłości trajektorii
$\{\sup_{s\le t}W_s\ge a\}=\{\tau\le t\}$.
Jeśli $W_t>a$, to z ciągłości i $W_0=0<a$ musieliśmy
wcześniej osiągnąć $a$, zatem $\{W_t>a\}\subset\{\tau\le t\}$.

Na zdarzeniu $\{\tau\le t\}$ mamy $W_\tau=a$.
Z mocnej własności Markowa przyszły przyrost
$W_t-W_\tau$, przy zadanym $\calF_\tau$ i czasie
pozostałym $t-\tau$, ma symetryczny rozkład
$\Normal(0,t-\tau)$. Możemy to sformalizować,
warunkując najpierw względem $\calF_\tau$ i stosując
mocną własność Markowa do przyszłego procesu
$B_u=W_{\tau+u}-W_\tau$ oraz mierzalnej funkcji
$u=t-\tau$; fakt dla prostych $u$ przenosimy do
ogólnego $u$ przez aproksymację i ciągłość trajektorii.
Stąd
\[
\Prob(\tau\le t,\ W_t<a)
=\Prob(\tau\le t,\ W_t>a).
\]
Równość jest prawdziwa także na granicy $\tau=t$,
gdyż oba zdarzenia tam nie zachodzą.
Ponieważ $W_t$ ma rozkład ciągły,
$\Prob(W_t=a)=0$ i
$\Prob(\tau\le t,\ W_t>a)=\Prob(W_t>a)$.
Dzieląc $\{\tau\le t\}$ według zdarzeń
$\{W_t<a\},\{W_t=a\},\{W_t>a\}$,
dostajemy żądany wzór.
\end{proof}

\begin{wniosek}[Rozkład czasu trafienia]
Dla $a>0$ czas $\tau_a$ jest skończony prawie na pewno,
a jego dystrybuanta i gęstość na $(0,\infty)$ mają postać
\[
\Prob(\tau_a\le t)=2\left(1-\Phi(a/\sqrt t)\right),
\qquad
f_{\tau_a}(t)=\frac{a}{\sqrt{2\pi}\,t^{3/2}}
\exp\left(-\frac{a^2}{2t}\right).
\]
Jednak $\E\tau_a=\infty$.
\end{wniosek}
\begin{proof}
Pierwszy wzór jest zasadą odbicia. Gdy $t\to\infty$,
argument $a/\sqrt t$ dąży do zera, więc dystrybuanta
dąży do $2(1-\Phi(0))=1$: czas jest skończony p.n.
Różniczkując dystrybuantę po $t>0$, dostajemy gęstość.
Ponadto, z twierdzenia o wartości średniej dla $\Phi$,
\[
\Prob(\tau_a>t)=2(\Phi(a/\sqrt t)-1/2)
\sim \sqrt{\frac2\pi}\frac{a}{\sqrt t}
\quad\text{gdy }t\to\infty.
\]
Dla nieujemnej zmiennej losowej zachodzi
$\E\tau_a=\int_0^\infty\Prob(\tau_a>t)\,\dd t$
(z twierdzenia Tonellego). Całka jest rozbieżna,
gdyż $\int_1^\infty t^{-1/2}\,\dd t=\infty$.
\end{proof}

\begin{uwaga}[Prawie pewne zakończenie nie daje skończonego średniego czasu]
Ruch Browna dociera do ustalonego dodatniego poziomu
z prawdopodobieństwem $1$, lecz wartość oczekiwana
czasu oczekiwania jest nieskończona.
To ważny kontrprzykład: nie wolno bez dodatkowych
warunków stosować próbkowania opcjonalnego bezpośrednio
w chwili $\tau_a$. Istnienie granicy $W_{\tau_a}=a$
prawie na pewno nie oznacza, że
$\E W_{\tau_a}=\E W_0$.
\end{uwaga}

\begin{zadanie}[Transformata Laplace'a czasu pierwszego trafienia]
Niech $W$ będzie procesem Wienera, $a>0$,
$\tau_a=\inf\{t\ge0:W_t=a\}$.
Obliczyć $\E e^{-\lambda\tau_a}$ dla $\lambda>0$.

\emph{Rozwiązanie.} Dla $\theta=\sqrt{2\lambda}>0$
proces $M_t=\exp(\theta W_t-\theta^2t/2)$
jest martyngałem: wynika to z niezależności przyrostów i
$\E e^{\theta(W_t-W_s)}=e^{\theta^2(t-s)/2}$.
Dla ograniczonego czasu $\tau_a\wedge n$ twierdzenie o
próbkowaniu opcjonalnym daje
$1=\E M_{\tau_a\wedge n}$.
Przed trafieniem w $a$ mamy $W_t<a$, a w chwili
trafienia $W_{\tau_a}=a$, więc
$0\le M_{\tau_a\wedge n}\le e^{\theta a}$.
Wiemy już, że $\tau_a<\infty$ p.n.; zatem
$M_{\tau_a\wedge n}\to\exp(\theta a-\lambda\tau_a)$
p.n. Zbieżność zdominowana pozwala przejść do granicy:
$1=e^{\theta a}\E e^{-\lambda\tau_a}$, skąd
\[
\E e^{-\lambda\tau_a}=e^{-a\sqrt{2\lambda}}.
\]
Ta transformata będzie w przyszłości użyteczna przy
porównywaniu ruchu Browna z procesami Lévy'ego.
\end{zadanie}

\section{Wyjście z przedziału: prawdopodobieństwa i średni czas}
\begin{twierdzenie}[Wyjście procesu Wienera z $(-a,b)$]
Niech $a,b>0$ oraz
$\tau:=\inf\{t\ge0:W_t\notin(-a,b)\}$.
Wtedy $\tau<\infty$ prawie na pewno,
$\E\tau=ab$ oraz
\[
\Prob(W_\tau=b)=\frac{a}{a+b},\qquad
\Prob(W_\tau=-a)=\frac{b}{a+b}.
\]
\end{twierdzenie}
\begin{proof}
Z ciągłości $W_{t\wedge\tau}\in[-a,b]$ i
$W_\tau\in\{-a,b\}$ na zdarzeniu $\{\tau<\infty\}$.
Z przyrostów gaussowskich wynika, że zarówno $W_t$,
jak i $W_t^2-t$ są martyngałami. Dla drugiego procesu
sprawdzamy bowiem, że przy $s<t$:
\[
\E[W_t^2-t\mid\calF_s]
=W_s^2+\E[(W_t-W_s)^2]-t=W_s^2-s,
\]
gdyż mieszany wyraz ma zerową warunkową wartość oczekiwaną.
Z próbkowania opcjonalnego w ograniczonej chwili
$t\wedge\tau$ dostajemy
$\E W_{t\wedge\tau}=0$ oraz
$\E W_{t\wedge\tau}^2=\E(t\wedge\tau)$.
Wprowadzając $M=\max(a,b)$, otrzymujemy
$\E(t\wedge\tau)\le M^2$. Zbieżność monotoniczna
przy $t\to\infty$ daje $\E\tau\le M^2<\infty$,
a więc $\tau<\infty$ p.n.

Ze zbieżności zdominowanej,
$W_{t\wedge\tau}\to W_\tau$ zarówno w $L^1$,
jak i w $L^2$; zatem
$\E W_\tau=0$ i
$\E W_\tau^2=\E\tau$.
Niech $p=\Prob(W_\tau=b)$. Ponieważ
$W_\tau$ przyjmuje tylko dwie wartości,
$0=\E W_\tau=bp-a(1-p)$, skąd $p=a/(a+b)$.
Następnie
\[
\E\tau=\E W_\tau^2
=b^2\frac{a}{a+b}+a^2\frac{b}{a+b}=ab.
\]
Wszystkie przejścia do granicy zostały więc uzasadnione
odpowiednio zbieżnością monotoniczną i zdominowaną.
\end{proof}

\begin{zadanie}[Start z dowolnego punktu przedziału]
Niech $x\in(-a,b)$, $a,b>0$, a proces
$X_t=x+W_t$. Obliczyć prawdopodobieństwo wyjścia
przez prawą krawędź i oczekiwany czas wyjścia.

\emph{Rozwiązanie.} Wyjście $X$ z $(-a,b)$ jest
równoważne wyjściu $W$ z przedziału
$(-a-x,b-x)$. W poprzednim twierdzeniu podstawiamy
$a'=a+x$, $b'=b-x$, obie liczby są dodatnie. Dostajemy
\[
\Prob(X_\tau=b)=\frac{a+x}{a+b},\qquad
\E\tau=(a+x)(b-x).
\]
Zwróćmy uwagę, że funkcja $u(x)=(a+x)(b-x)$
spełnia równanie $\frac12u''(x)=-1$ na $(-a,b)$
i warunki brzegowe $u(-a)=u(b)=0$.
Jest to pierwsza zapowiedź ścisłego związku procesów
Markowa z równaniami różniczkowymi.
\end{zadanie}

\section{Miara Wienera: całkowanie po wszystkich trajektoriach}
Proces $W$ można opisywać dwojako. Dla ustalonej chwili $t$
zmienna $W_t$ jest normalna. Ale pojedynczy wynik
$\omega\in\Omega$ wyznacza także całą ciągłą funkcję
$t\mapsto W_t(\omega)$. Możemy zatem zapytać nie tylko
o prawdopodobieństwo $\{W_t\in A\}$, lecz także o
prawdopodobieństwo zdarzenia zależącego od całej drogi,
na przykład: ``trajektoria przez cały czas pozostaje
poniżej bariery'' lub ``pole pod trajektorią jest dodatnie''.

\begin{definicja}[Przestrzeń trajektorii i miara Wienera]
Dla $T>0$ definiujemy
\[
\calC_0([0,T]):=\{w\in C([0,T];\R):w(0)=0\},
\qquad \norm{w}_\infty=\sup_{0\le t\le T}|w(t)|.
\]
Niech $W$ będzie standardowym procesem Wienera o ciągłych
trajektoriach, traktowanym jako odwzorowanie
$W:\Omega\to\calC_0([0,T])$,
$W(\omega)(t)=W_t(\omega)$.
Jego \emph{prawem na przestrzeni trajektorii}, nazywanym
\emph{miarą Wienera}, jest miara probabilistyczna
$\mu_T:=\Prob\circ W^{-1}$ na
$(\calC_0([0,T]),\calB(\calC_0([0,T])))$.
\end{definicja}

\begin{stwierdzenie}[Mierzalność odwzorowania trajektorii]
Jeżeli $W$ ma ciągłe trajektorie i każda zmienna $W_t$
jest mierzalna, to odwzorowanie
$\omega\mapsto W(\omega)\in\calC_0([0,T])$ jest mierzalne.
\end{stwierdzenie}
\begin{proof}
Dla każdej wymiernej chwili $r\in[0,T]$ odwzorowanie
ewaluacji $e_r(w)=w(r)$ jest ciągłe w normie supremum,
a $e_r\circ W=W_r$ jest mierzalne.
Przestrzeń $\calC_0([0,T])$ jest ośrodkową przestrzenią
metryczną (można przybliżać funkcje ciągłe łamanymi
o wymiernych węzłach i wymiernych wartościach).
Borelowskie sigma-ciało tej przestrzeni jest generowane
przez ewaluacje w przeliczalnym gęstym zbiorze chwil:
rzeczywiście $\norm{w-v}_\infty=\sup_{r\in D}|w(r)-v(r)|$
dla gęstego przeliczalnego $D$ zawierającego końce,
więc każda kula otwarta jest mierzalna względem
$\sigma(e_r:r\in D)$; dzięki ośrodkowości każdy zbiór
otwarty jest przeliczalną sumą takich kul.
Zatem przeciwobrazy zbiorów borelowskich przez $W$
są mierzalne.
\end{proof}

\begin{twierdzenie}[Całki cylindryczne względem miary Wienera]
Niech $T>0$, $0<t_1<\cdots<t_m\le T$ oraz
$F:\R^m\to\R$ będzie funkcją borelowską, nieujemną
lub całkowalną względem rozkładu wektora
$(W_{t_1},\ldots,W_{t_m})$. Wówczas
\begin{align*}
&\int_{\calC_0([0,T])}F(w(t_1),\ldots,w(t_m))\,\dd\mu_T(w)\\
&\quad=\int_{\R^m}F(x_1,\ldots,x_m)
\prod_{k=1}^m p_{t_k-t_{k-1}}(x_k-x_{k-1})
\,\dd\lambda_m(x_1,\ldots,x_m),
\end{align*}
gdzie $t_0=0$, $x_0=0$ i $p_h(z)$ jest gęstością
$\Normal(0,h)$.
\end{twierdzenie}
\begin{proof}
Z definicji miary obrazowej całka po lewej jest równa
$\E F(W_{t_1},\ldots,W_{t_m})$.
Przeprowadzamy liniowe przekształcenie od przyrostów
$Z_k=W_{t_k}-W_{t_{k-1}}$ do kolejnych położeń
$x_k=Z_1+\cdots+Z_k$.
Przyrosty są niezależne, $Z_k$ mają gęstości
$p_{t_k-t_{k-1}}$, więc wektor $Z$ ma gęstość
będącą ich iloczynem. Macierz przekształcenia z $Z$
do $x$ jest trójkątna z jedynkami na przekątnej,
a zatem jej wyznacznik ma moduł $1$.
Wzór na zmianę zmiennych w całce Lebesgue'a
daje tezę. Dla $F\ge0$ dopuszczalność całkowania
wynika z twierdzenia Tonellego; w przypadku funkcji
całkowalnej stosujemy rozkład na części dodatnią i ujemną.
\end{proof}

\begin{zadanie}[Całka z całej trajektorii jako zmienna gaussowska]
Niech $A_T:=\int_0^T W_s\,\dd s$, gdzie całka jest zwykłą
całką Riemanna (równoważnie Lebesgue'a) z ciągłej trajektorii.
Obliczyć rozkład $A_T$ i jego funkcję charakterystyczną.

\emph{Rozwiązanie.} Aproksymujmy całkę sumami Riemanna
$A_T^{(n)}=\sum_{k=0}^{n-1}W_{kT/n}\,(T/n)$.
Każda suma jest gaussowska, ponieważ jest liniową
kombinacją wspólnie gaussowskich zmiennych.
Z ciągłości trajektorii $A_T^{(n)}\to A_T$ p.n.
Co więcej, jeśli $\varphi_n(s)$ jest lewym końcem
przedziału siatki zawierającego $s$, to nierówność
Cauchy'ego--Schwarza i twierdzenie Fubiniego dają
\begin{align*}
\E|A_T^{(n)}-A_T|^2
&\le T\int_0^T\E|W_{\varphi_n(s)}-W_s|^2\,\dd s\\
&\le T^2\frac{T}{n}\longrightarrow0.
\end{align*}
Zatem $A_T^{(n)}$ zbiega w $L^2$ do $A_T$,
a jego średnie i wariancje zbiegają do odpowiednich
wielkości granicy. Granica rozkładów normalnych o
zbieżnych parametrach jest normalna (wynika to wprost
ze zbieżności funkcji charakterystycznych).
Ponieważ $\E A_T=0$ i
\[
\Var(A_T)=\int_0^T\!\int_0^T\Cov(W_s,W_t)
\,\dd s\,\dd t
=\int_0^T\!\int_0^T\min(s,t)\,\dd s\,\dd t
=\frac{T^3}{3},
\]
dostajemy $A_T\sim\Normal(0,T^3/3)$ oraz
\[
\E e^{iuA_T}=\exp\left(-\frac{u^2T^3}{6}\right).
\]
W języku miary Wienera jest to dokładnie obliczenie
całki po trajektoriach
\[
\int_{\calC_0([0,T])}
\exp\left(iu\int_0^T w(s)\,\dd s\right)\,\dd\mu_T(w).
\]
\end{zadanie}

\begin{uwaga}[Trzy całki, których nie wolno utożsamiać]
Wyrażenie $\int_0^T W_s\,\dd s$ jest zwykłą całką
po czasie z jednej losowej, ciągłej funkcji.
Wyrażenie $\int_0^T H_s\,\dd W_s$ będzie całką
stochastyczną Itô, skonstruowaną później z przyrostów
procesu Wienera. Wreszcie
$\int_{\calC_0([0,T])}F(w)\,\dd\mu_T(w)$
całkuje funkcjonał $F$ \emph{po wszystkich możliwych
trajektoriach}. Jest to ścisła matematyczna całka Wienera
po drogach, stanowiąca punkt wyjścia do wzoru
Feynmana--Kaca. Oscylacyjna całka Feynmana w czasie
rzeczywistym wymaga odrębnej analizy i nie jest
zwykłą całką względem miary probabilistycznej.
\end{uwaga}


\begin{stwierdzenie}[Most Browna i miara Wienera mają różne prawa]
Dla $T>0$ niech $\mu_T$ będzie miarą Wienera na
$\calC_0([0,T])$, a $\nu_T$ prawem mostu Browna
$\beta_t=W_t-(t/T)W_T$. Wtedy $\mu_T$ i $\nu_T$
są wzajemnie singularne.
\end{stwierdzenie}
\begin{proof}
Zbiór $A=\{w\in\calC_0([0,T]):w(T)=0\}$ jest domknięty,
więc borelowski. Ponieważ $\beta_T=0$, mamy $\nu_T(A)=1$.
Natomiast $W_T\sim\Normal(0,T)$ ma rozkład bezatomowy,
więc $\mu_T(A)=\Prob(W_T=0)=0$. Jest to definicja
wzajemnej singularności miar.
\end{proof}

\begin{uwaga}[Jaki typ zbieżności wariacji kwadratowej udowodniono?]
Dla \emph{dowolnych deterministycznych} podziałów
$[0,T]$ o średnicy zmierzającej do zera wykazaliśmy
zbieżność sum kwadratów przyrostów w $L^2$, a zatem
według prawdopodobieństwa. Zbieżność prawie na pewno
uzyskaliśmy dla \emph{konkretnego ciągu siatek dyadycznych}.
Z tych dwóch zdań nie wolno wyciągać wniosku, że
zbieżność prawie na pewno zachodzi dla każdej dowolnie
wybranej sekwencji podziałów. Ważne jest również,
czy podziały wolno wybierać zależnie od trajektorii.
\end{uwaga}

\section{Równanie ciepła i generator ruchu Browna}
\begin{twierdzenie}[Generator i równanie ciepła]
Niech $(P_t)_{t\ge0}$ będzie półgrupą zdefiniowaną wzorem
$(P_tf)(x)=\E f(x+W_t)$ dla ograniczonych borelowskich
$f:\R\to\R$. Wówczas:
\begin{enumerate}
\item dla każdej $f\in C_b^3(\R)$, czyli funkcji o
ciągłych i ograniczonych pochodnych do rzędu trzeciego,
zachodzi dla każdego $x\in\R$
\[
\lim_{t\downarrow0}\frac{(P_tf)(x)-f(x)}{t}
=\frac12 f''(x);
\]
\item jeżeli $f\in C_b(\R)$, to funkcja
$u(t,x)=(P_tf)(x)$ jest klasy $C^{1,2}$ na
$(0,\infty)\times\R$, spełnia
$\partial_tu=\frac12\partial_{xx}u$
oraz $u(t,x)\to f(x)$ dla każdego $x$ przy $t\downarrow0$.
\end{enumerate}
\end{twierdzenie}
\begin{proof}
Dla $f\in C_b^3(\R)$ wzór Taylora z resztą daje
\[
f(x+h)=f(x)+hf'(x)+\frac12h^2f''(x)+R(x,h),\qquad
|R(x,h)|\le\frac{\norm{f'''}_\infty}{6}|h|^3.
\]
Podstawmy $h=W_t\sim\Normal(0,t)$ i uśrednijmy.
Mamy $\E W_t=0$, $\E W_t^2=t$ i
$\E|W_t|^3=c_3t^{3/2}$. Zatem
$|(P_tf)(x)-f(x)-tf''(x)/2|
\le \norm{f'''}_\infty c_3t^{3/2}/6$.
Po podzieleniu przez $t$ reszta dąży do zera,
co dowodzi pierwszej części.

Dla $t>0$ możemy zapisać
$u(t,x)=\int_\R f(y)p_t(y-x)\,\dd\lambda_1(y)$.
Na każdym przedziale czasowym $t\in[\delta,R]$
($0<\delta<R$) pochodne jądra gaussowskiego
$\partial_t p_t(z)$, $\partial_zp_t(z)$ i
$\partial_{zz}p_t(z)$ są ograniczone przez całkowalne
funkcje zmiennej $z$, niezależne od $t$ w tym przedziale.
Wynika to z jawnej postaci pochodnych, będących
wielomianami w $z$ pomnożonymi przez
$\exp(-z^2/(2R))$, z odpowiednimi stałymi.
Twierdzenie o różniczkowaniu pod znakiem całki daje
istnienie i ciągłość wskazanych pochodnych $u$.
Bezpośredni rachunek gęstości wykazuje
$\partial_t p_t(z)=\frac12\partial_{zz}p_t(z)$,
a więc po przeniesieniu pochodnych pod całkę
$\partial_tu=\frac12\partial_{xx}u$.

Na koniec $W_t\to0$ według prawdopodobieństwa,
bo $\E W_t^2=t\to0$. Dla ustalonego $x$ i
$\varepsilon>0$ ciągłość $f$ w $x$ daje $\delta>0$
takie, że $|f(x+z)-f(x)|<\varepsilon$ dla $|z|<\delta$.
Dzieląc wartość oczekiwaną według $|W_t|<\delta$
i jego dopełnienia, otrzymujemy
$|P_tf(x)-f(x)|\le\varepsilon
+2\norm f_\infty\Prob(|W_t|\ge\delta)\to\varepsilon$.
Dowolność $\varepsilon$ daje wymaganą zbieżność.
\end{proof}

\begin{przyklad}[Od losowej trajektorii do równania różniczkowego]
Dla $f(x)=\cos(kx)$, $k\in\R$, mamy
$P_tf(x)=\E\cos(kx+kW_t)
=e^{-k^2t/2}\cos(kx)$.
Rzeczywiście, część urojona funkcji charakterystycznej
rozkładu $\Normal(0,t)$ znika, a część rzeczywista
wynosi $e^{-k^2t/2}$. Jednocześnie
$\partial_tu=-(k^2/2)u$ i
$\frac12\partial_{xx}u=-(k^2/2)u$.
Jeden rachunek probabilistyczny dostarcza więc
jawnego rozwiązania równania ciepła.
\end{przyklad}

\section{Podsumowanie i kierunek dalszej teorii}
Proces Wienera łączy zaskakujące fakty: ma ciągłe
trajektorie, ale żadna z nich nie jest różniczkowalna;
ma nieskończoną wariację całkowitą, lecz jego wariacja
kwadratowa jest dokładnie równa upływającemu czasowi.
Jednocześnie można go analizować za pomocą elementarnych
narzędzi gaussowskich, mocnej własności Markowa i
równania ciepła.

Przygotowaliśmy również przestrzeń trajektorii i miarę Wienera.
Dalsze rozdziały rozwiną całki po drogach, ruch Browna
w wielu wymiarach i ścisły rachunek stochastyczny.
W szczególności wzór Feynmana--Kaca połączy
ważone miary trajektorii z półgrupami operatorów
Schrödingera, a całka Itô zbuduje matematyczne
narzędzia do równań różniczkowych stochastycznych.


% ==============================================================
% ROZDZIAŁ 8. ROZSZERZENIE TEORII RUCHU BROWNA
% ==============================================================

\chapter{Ruch Browna w wielu wymiarach i dalsze własności trajektorii}
\label{ch:brown-multi}

\section{Dlaczego potrzebujemy procesu wielowymiarowego?}
W ruchu jednowymiarowym możemy badać przekroczenie bariery przez
cząstkę. Ruch w przestrzeni fizycznej wymaga jednak opisu wielu
współrzędnych naraz. Pytania zmieniają się wtedy jakościowo:
jak długo cząstka pozostaje w kuli? Czy dociera do wyróżnionego
punktu? Jak zmieniają się wyniki po obróceniu układu współrzędnych?
Zacznijmy od konstrukcji, w której odpowiedź na ostatnie pytanie
wynika ze szczególnej postaci macierzy kowariancji.

\begin{definicja}[Standardowy wielowymiarowy ruch Browna]
Niech $d\ge1$ będzie liczbą całkowitą. Proces $B=(B_t)_{t\ge0}$ o wartościach
w $\R^d$, z prawie na pewno ciągłymi trajektoriami, nazywamy
\emph{standardowym $d$-wymiarowym ruchem Browna}, jeżeli
$B_0=0$ prawie na pewno, ma niezależne przyrosty oraz
$B_t-B_s\sim\Normal_d(0,(t-s)I_d)$ dla wszystkich $0\le s<t$.
Jeżeli $x\in\R^d$, proces $X_t=x+B_t$ nazywamy ruchem Browna
rozpoczynającym się w $x$; jego prawo i wartość oczekiwaną oznaczamy
odpowiednio przez $\Prob_x$ oraz $\E_x$.
Gdy zadano filtrację $(\calF_t)_{t\ge0}$, wymagamy
dodatkowo adaptacji oraz niezależności przyrostu
$B_t-B_s$ od $\calF_s$ dla $s<t$.
W dalszym ciągu, o ile nie zaznaczono inaczej, wybieramy
wersję o wszystkich trajektoriach ciągłych i $B_0=0$
punktowo oraz jej filtrację naturalną.
\end{definicja}

Ruch Browna względem swojej filtracji naturalnej
spełnia warunek niezależności przyszłego przyrostu:
skończony wektor obserwacji do chwili $s$ jest
prawie na pewno funkcją przyrostów do tej chwili.
Niezależność przyrostu $B_t-B_s$ od zdarzeń
cylindrycznych rozszerzamy argumentem $\pi$--$\lambda$
na $\calF_s=\sigma(B_v:0\le v\le s)$.

\begin{twierdzenie}[Konstrukcja i izotropia]
Dla każdego $d\ge1$ istnieje standardowy $d$-wymiarowy ruch
Browna $B$. Jego współrzędne są niezależnymi standardowymi
procesami Wienera, a dla $1\le i,j\le d$ i $s,t\ge0$
\[
\E B_t^{(i)}=0,\qquad
\Cov(B_s^{(i)},B_t^{(j)})=\delta_{ij}\min(s,t).
\]
Jeżeli $U\in\R^{d\times d}$ jest macierzą ortogonalną,
to proces $(UB_t)_{t\ge0}$ ma to samo prawo na
$C_0([0,T];\R^d)$ co $B$, dla każdego $T>0$.
\end{twierdzenie}
\begin{proof}
Na przestrzeni iloczynowej $d$ kopii przestrzeni probabilistycznej
z ruchem Browna weźmy niezależne procesy $W^{(1)},\ldots,W^{(d)}$
i połóżmy $B_t=(W_t^{(1)},\ldots,W_t^{(d)})$.
Wektor każdego przyrostu ma niezależne współrzędne o rozkładzie
$\Normal(0,t-s)$, czyli ma rozkład $\Normal_d(0,(t-s)I_d)$.
Współrzędne na różnych rozłącznych przedziałach także są niezależne;
można to zobaczyć, grupując skończone wektory przyrostów po
współrzędnych. Ciągłość wynika z przecięcia $d$ zdarzeń pełnej miary.

Dla dowolnego procesu spełniającego definicję wektor
wartości w skończonym zbiorze chwil jest prawie na pewno
liniowym obrazem wektora niezależnych gaussowskich przyrostów
(zaczynamy od chwili zero). Jest zatem gaussowski.
Dla $s\le t$ niezależność $B_t-B_s$ od $B_s$ daje
$\Cov(B_s^{(i)},B_t^{(j)})=\Cov(B_s^{(i)},B_s^{(j)})
=s\delta_{ij}$; średnie są zerowe. Dla dowolnych skończonych
rodzin obserwacji poszczególnych współrzędnych macierz
kowariancji jest blokowo diagonalna względem współrzędnych,
więc odpowiednie bloki są łącznie niezależne. Iterując
argument $\pi$--$\lambda$ po $d$ współrzędnych, otrzymujemy
niezależność $\sigma$-ciał generowanych przez całe procesy
współrzędnych, nie tylko przez ich wartości w jednej chwili.
Każdy proces współrzędnej jest standardowym procesem Wienera.

Ponieważ $U((t-s)I_d)U^\top=(t-s)I_d$, przyrosty procesu $UB$
mają te same rozkłady co przyrosty $B$ i pozostają niezależne.
Stąd równość rozkładów skończeniewymiarowych, a przez generowanie borelowskiego $\sigma$-ciała
przestrzeni $C_0([0,T];\R^d)$ przez ewaluacje w przeliczalnym
gęstym zbiorze chwil także równość praw trajektorii.
Wynika to z jednowymiarowego argumentu z rozdziału~\ref{ch:wiener}
i identyfikacji przestrzeni trajektorii wektorowych z iloczynem
$d$ przestrzeni trajektorii rzeczywistych.
\end{proof}

\begin{stwierdzenie}[Równoważna charakteryzacja gaussowska]
Rzeczywisty proces wektorowy $B$ o prawie na pewno ciągłych
trajektoriach jest standardowym ruchem Browna w $\R^d$
wtedy i tylko wtedy, gdy jest wycentrowanym procesem
gaussowskim o macierzach kowariancji
$\Cov(B_s,B_t)=\min(s,t)I_d$.
\end{stwierdzenie}
\begin{proof}
Konieczność wykazano w poprzednim twierdzeniu.
Odwrotnie, z zerowej średniej i kowariancji w chwili zero
wynika $B_0=0$ p.n. Przyrost $B_t-B_s$ jest gaussowski
o średniej zero i kowariancji $(t-s)I_d$.
Dla $0\le a<b\le c<e$ kowariancja przyrostów
$B_b-B_a$ i $B_e-B_c$ jest równa
$(b-b-a+a)I_d=0$. Dowolny skończony wektor takich
przyrostów jest gaussowski o blokowo diagonalnej
macierzy kowariancji, co daje ich łączną niezależność.
\end{proof}

\begin{uwaga}[Niezależność współrzędnych a niezależność położenia]
Współrzędne $B_t^{(1)},\ldots,B_t^{(d)}$ są niezależne w ustalonej
chwili $t$. Nie wynika stąd, że położenia $B_s$ i $B_t$ są niezależne
dla $0<s<t$: przeciwnie, mają wspólną przeszłość, a ich kowariancja
wynosi $sI_d$. Przy $s=0$ położenie początkowe jest stałe p.n. Niezależne są \emph{przyrosty} na rozłącznych
przedziałach czasu.
\end{uwaga}

\begin{twierdzenie}[Mocna własność Markowa w wielu wymiarach]
\label{thm:brown-multi-strong}
Niech $B$ będzie standardowym ruchem Browna względem
filtracji $(\calF_t)_{t\ge0}$, o wszystkich trajektoriach
ciągłych, a $\tau$ czasem zatrzymania skończonym prawie
na pewno. Wówczas proces
$(B_{\tau+v}-B_\tau)_{v\ge0}$ (z trajektorią zerową
przy $\tau=\infty$) jest standardowym $d$-wymiarowym
ruchem Browna niezależnym od $\calF_\tau$.
\end{twierdzenie}
\begin{proof}
Dowód jest wektorową wersją dowodu
twierdzenia~\ref{thm:wiener-strong-markov}.
Dla $h_n=2^{-n}$ zaokrąglamy $\tau$ w górę do
$\tau_n=h_n\lceil\tau/h_n\rceil$ (z wartością
$\infty$ przy $\tau=\infty$). Są to czasy zatrzymania,
$\tau_n\downarrow\tau$, a $\calF_\tau\subseteq\calF_{\tau_n}$.
Obserwacje w losowych chwilach są mierzalne na mocy
twierdzeń~\ref{thm:prog} i~\ref{thm:sampled}, ponieważ
$\tau+v$ i $\tau_n+v$ są czasami zatrzymania dla $v\ge0$.
Przy nieskończonej chwili przyjmujemy wartość zero
dla wszystkich takich obserwacji.
Dla $A\in\calF_\tau$, chwil $0\le v_1<\cdots<v_m$
i wektorów $u_1,\ldots,u_m\in\R^d$ warunkujemy na
zdarzeniach $A\cap\{\tau_n=kh_n\}\in\calF_{kh_n}$.
Kolejne przyrosty po $kh_n$ są łącznie niezależne
od $\calF_{kh_n}$ przez iterowane warunkowanie
względem chwil ich rozpoczęcia. Z gaussowskości
i wzoru na kowariancję dostajemy po zsumowaniu
po $k\ge0$
\[
\E\left[\chi_A\exp\left(
i\sum_{j=1}^m\langle u_j,B_{\tau_n+v_j}-B_{\tau_n}\rangle
\right)\right]
=\Prob(A)\exp\left(
-\frac12\sum_{j,l=1}^m\min(v_j,v_l)\langle u_j,u_l\rangle
\right).
\]
Sumowanie jest dozwolone przez ograniczenie modułu
funkcji wykładniczej przez $1$ i fakt, że zdarzenie
$\{\tau_n=\infty\}$ ma prawdopodobieństwo zero.
Z ciągłości trajektorii i zbieżności zmajoryzowanej
ta sama równość zachodzi z $\tau$ w miejsce $\tau_n$.
Dla $A=\Omega$ identyfikuje ona rozkłady
skończeniewymiarowe ruchu Browna, a dla dowolnego
$A$ dowodzi ich niezależności od $\calF_\tau$
przez jednoznaczność funkcji charakterystycznych.
Ciągłość przyszłych trajektorii i mierzalność
poszczególnych obserwacji zapewniają mierzalność
odwzorowania trajektorii, ponieważ jego borelowskie
$\sigma$-ciało jest generowane przez ewaluacje
w chwilach wymiernych. Argument $\pi$--$\lambda$
dla cylindrów rozszerza niezależność na całe to
$\sigma$-ciało i kończy dowód.
\end{proof}

\begin{twierdzenie}[Generator i wielowymiarowe równanie ciepła]
\label{thm:brown-multi-generator}
Niech $B$ będzie standardowym ruchem Browna w $\R^d$.
Dla $t\ge0$ oraz ograniczonej funkcji borelowskiej
$f:\R^d\to\R$ połóżmy $P_tf(x)=\E f(x+B_t)$.
Operatory te tworzą półgrupę przejścia: $P_0f=f$,
$P_sP_t=P_{s+t}$, a względem filtracji naturalnej
dla $X=x+B$ zachodzi
$\E[f(X_{s+t})\mid\calF_s]=(P_tf)(X_s)$ p.n.
Jeżeli $f\in C_b^3(\R^d)$ (wszystkie pochodne do rzędu trzeciego
są ciągłe i ograniczone, włącznie z samą funkcją), to
\[
\lim_{t\downarrow0}\left\|\frac{P_tf-f}{t}-\frac12\Delta f
\right\|_\infty=0.
\]
Jeżeli $f\in C_b(\R^d)$, to $u(t,x)=P_tf(x)$ jest klasy
$C^{1,2}$ na $(0,\infty)\times\R^d$, spełnia
$\partial_tu=\frac12\Delta u$ oraz $u(t,x)\to f(x)$ dla $t\downarrow0$.
\end{twierdzenie}
\begin{proof}
Mierzalność $P_tf$ oraz wzór na warunkowanie
wynikają z twierdzenia Fubiniego i niezależności
$B_{s+t}-B_s$ od $\calF_s$, tak jak w rozdziale
\ref{ch:wiener}. Własność półgrupy wynika z faktu,
że suma niezależnych wektorów o rozkładach
$\Normal_d(0,sI_d)$ i $\Normal_d(0,tI_d)$
ma rozkład $\Normal_d(0,(s+t)I_d)$.

Z wielowymiarowego wzoru Taylora z resztą, ograniczeniem
pochodnych trzeciego rzędu i normą euklidesową dostajemy
$f(x+h)=f(x)+\langle\nabla f(x),h\rangle+
\frac12 h^\top D^2f(x)h+R_x(h)$, gdzie
$|R_x(h)|\le C_d\max_{|\alpha|=3}\|\partial^\alpha f\|_\infty|h|^3$.
Ponieważ $\E B_t=0$,
$\E[B_t^{(i)}B_t^{(j)}]=t\delta_{ij}$ i
$\E|B_t|^3=t^{3/2}\E|B_1|^3$, otrzymujemy
$P_tf(x)-f(x)=\frac{t}{2}\sum_i\partial_{ii}f(x)+O(t^{3/2})$.
Stała w reszcie jest niezależna od $x$, a trzeci moment
normy $B_1$ jest skończony, ponieważ
$|B_1|^3\le d^{1/2}\sum_i|B_1^{(i)}|^3$.
To dowodzi granicy w normie supremum po podzieleniu przez $t$.
Dla $t>0$ rozkład $B_t$ ma gęstość
$p_t^{(d)}(z)=(2\pi t)^{-d/2}\exp(-|z|^2/(2t))$.
Na każdym pasie czasowym $0<\delta\le t\le R$ mamy
\[
\partial_t p_t^{(d)}(z)=
\left(-\frac d{2t}+\frac{|z|^2}{2t^2}\right)p_t^{(d)}(z),
\qquad
\partial_{ij}p_t^{(d)}(z)=
\left(\frac{z_i z_j}{t^2}-\frac{\delta_{ij}}{t}\right)p_t^{(d)}(z).
\]
Ponadto $\partial_i p_t^{(d)}(z)=-(z_i/t)p_t^{(d)}(z)$.
Piszemy $u(t,x)=\int_{\R^d} f(y)p_t^{(d)}(y-x)\,
\mathrm d\lambda_d(y)$, aby różniczkować jądro, a nie $f$.
Dla $|x-x_0|\le1$ mamy $|y-x|^2\ge |y-x_0|^2/2-1$.
Zatem wszystkie wskazane pochodne pod znakiem całki
mają wspólną całkowalną majorantę, będącą stałą razy
$(1+|y-x_0|^2)\exp(-|y-x_0|^2/(4R))$.
Różniczkowanie pod znakiem całki i twierdzenie o zbieżności
zmajoryzowanej dają ciągłość pochodnych oraz
$\partial_tu=\frac12\Delta u$, gdyż
$\partial_t p_t^{(d)}=\frac12\Delta p_t^{(d)}$.
Zbieżność do warunku początkowego wynika z ciągłości $f$ w $x$,
ograniczenia $f$ i faktu $\E|B_t|^2=dt\to0$, dokładnie jak
w dowodzie jednowymiarowym.
\end{proof}

\begin{przyklad}[Ruch Browna z anizotropową kowariancją]
Niech $A\in\R^{m\times d}$, $b\in\R^m$ i $X_t=b+AB_t$.
Wówczas $X_t-X_s\sim\Normal_m(0,(t-s)AA^\top)$ i
$\E X_t=b$. Jeżeli $AA^\top$ jest dodatnio określona, ale nie jest wielokrotnością
macierzy jednostkowej, kierunki nie są równoważne: powierzchnie
jednakowej gęstości są elipsoidami. Dla
$f\in C_b^3(\R^m)$ generator ma postać
$Lf=\frac12\sum_{i,j}(AA^\top)_{ij}\partial_{ij}f$.
Wzór wynika z tego samego rachunku Taylora, po zastąpieniu
macierzy kowariancji $tI_m$ przez $tAA^\top$, dla półgrupy
$P_t^Af(x)=\E f(x+AB_t)$. Przy osobliwej macierzy
$AA^\top$ mamy $\operatorname{rank} A<m$, ponieważ
$\ker(AA^\top)=\ker A^\top$. Proces startujący w $b$
porusza się w $b+\operatorname{im}A$, podprzestrzeni
afinicznej o wymiarze mniejszym od $m$. Jej miara
Lebesgue'a jest zerowa, więc rozkład $X_t$ jest singularny
względem $\lambda_m$.
\end{przyklad}


\begin{zadanie}[Anizotropowa dyfuzja w płaszczyźnie]
Niech $B=(B^{(1)},B^{(2)})$ będzie standardowym ruchem Browna
w $\R^2$ i niech $X_t=(2B_t^{(1)},B_t^{(2)})$.
Wyznaczyć rozkład $X_t$, generator procesu i prawdopodobieństwo
$\Prob(|X_t^{(1)}|>2a)$ dla $a,t>0$.

\emph{Rozwiązanie.} Macierz $A=\operatorname{diag}(2,1)$ daje
$AA^\top=\operatorname{diag}(4,1)$. Zatem
$X_t\sim\Normal_2(0,t\operatorname{diag}(4,1))$,
a jego gęstość wynosi
$ (4\pi t)^{-1}\exp(-x_1^2/(8t)-x_2^2/(2t))$.
Generator ma postać $Lf=2\partial_{11}f+\tfrac12\partial_{22}f$.
Wreszcie
$\Prob(|X_t^{(1)}|>2a)=\Prob(|B_t^{(1)}|>a)
=2(1-\Phi(a/\sqrt t))$.
Współczynnik dyfuzji w pierwszym kierunku jest cztery razy
większy niż w drugim; nie można go odczytać z samej
wartości oczekiwanej $\E X_t=0$.
\end{zadanie}

\section{Czas wyjścia z kuli i martyngał kwadratu promienia}
\begin{twierdzenie}[Średni czas wyjścia z kuli]
Niech $d\ge1$ będzie liczbą całkowitą, $R>0$,
$x\in\R^d$ i $|x|<R$.
Dla standardowego $d$-wymiarowego ruchu Browna $B$ połóżmy
$X_t=x+B_t$ oraz $\tau=\inf\{t\ge0:|X_t|\ge R\}$.
Wtedy $\tau<\infty$ prawie na pewno i
\[
\E\tau=\frac{R^2-|x|^2}{d}.
\]
\end{twierdzenie}
\begin{proof}
Pracujemy z ciągłą wersją $B$ zaczynającą się w zerze
i jej filtracją naturalną. Niezależność przyrostu od
tej filtracji wynika ze zdarzeń cylindrycznych i argumentu
$\pi$--$\lambda$, tak jak w przypadku jednowymiarowym.
Czas $\tau$ jest czasem zatrzymania na mocy
twierdzenia~\ref{thm:hitting-closed}, zastosowanego do
domkniętego zbioru $\{y:|y|\ge R\}$.
Dla $s<t$ i $Z=B_t-B_s$, niezależnego od przeszłości,
$\E Z=0$ i $\E|Z|^2=d(t-s)$. Rozwinięcie kwadratu daje
$\E[|X_t|^2\mid\calF_s]=|X_s|^2+d(t-s)$.
Zatem $M_t=|X_t|^2-dt$ jest całkowalnym martyngałem
o ciągłych trajektoriach. Dla każdego całkowitego $n\ge1$
twierdzenie o próbkowaniu opcjonalnym
(twierdzenie~\ref{thm:optional-cont}) daje
\[
\E|X_{\tau\wedge n}|^2-|x|^2=d\E(\tau\wedge n).
\]
Z ciągłości trajektorii zawsze $|X_{\tau\wedge n}|\le R$:
przed $\tau$ proces jest w kuli, a w chwili $\tau<\infty$
osiąga jej sferę. Stąd $d\E(\tau\wedge n)\le R^2-|x|^2$.
Zbieżność monotoniczna daje $\E\tau\le (R^2-|x|^2)/d<\infty$,
co wymusza $\tau<\infty$ p.n. Wtedy
$|X_{\tau\wedge n}|^2\to|X_\tau|^2=R^2$ p.n., a
zmienne są ograniczone przez $R^2$. Zbieżność zmajoryzowana
i monotoniczna po obu stronach poprzedniej tożsamości dają
$R^2-|x|^2=d\E\tau$.
\end{proof}

\begin{zadanie}[Wyjście z dysku]
Rozważmy standardowy dwuwymiarowy ruch Browna, którego
położenie początkowe ma odległość $r<R$ od środka dysku
promienia $R$. Ile średnio trzeba czekać na pierwsze osiągnięcie
okręgu?

\emph{Rozwiązanie.} Z twierdzenia dla $d=2$ dostajemy
$\E\tau=(R^2-r^2)/2$. Gdy $r=0$, średni czas to $R^2/2$.
Zależność od $r$ jest kwadratowa, mimo że typowe wychylenie
po czasie $t$ ma rząd $\sqrt t$.
\end{zadanie}

\section{Trafianie w kule i punkty: zależność od wymiaru}

Aby zbadać trafianie w małe zbiory, potrzebujemy
martyngałów związanych z funkcjami harmonicznymi.
Wyprowadzimy je bez stosowania wzoru Itô, który
zostanie wprowadzony w tomie III.

\begin{lemat}[Martyngał związany z generatorem]
\label{lem:brown-generator-martingale}
Dla $f\in C_c^2(\R^d)$ i $X_t=x+B_t$ proces
\[
M_t^f=f(X_t)-f(x)-\frac12\int_0^t
\Delta f(X_v)\,\mathrm d\lambda_1(v)
\]
jest ciągłym martyngałem względem filtracji naturalnej.
\end{lemat}
\begin{proof}
Najpierw przyjmijmy $f\in C_c^\infty(\R^d)$.
Oznaczmy $g=\frac12\Delta f$. Dla $0\le s<t$
weźmy podział $v_j=s+jh$, $h=(t-s)/n$.
Niezależność przyrostów od przeszłości daje
\[
\E[f(X_{v_{j+1}})-f(X_{v_j})\mid\calF_{v_j}]
=(P_h f-f)(X_{v_j}).
\]
Z twierdzenia~\ref{thm:brown-multi-generator}
mamy $P_h f-f=hg+h r_h$, gdzie
$\norm{r_h}_\infty\to0$. Stosując własność wieży
i sumując, otrzymujemy
\[
\E[f(X_t)-f(X_s)\mid\calF_s]
=\E\left[\sum_{j=0}^{n-1}h g(X_{v_j})
+\sum_{j=0}^{n-1}h r_h(X_{v_j})\mid\calF_s\right].
\]
Druga suma jest ograniczona w module przez
$(t-s)\norm{r_h}_\infty\to0$. Pierwsza zbiega
punktowo do całki po $[s,t]$, a ponieważ $g$
jest ograniczona, również w $L^1$.
Kontraktywność warunkowej wartości oczekiwanej
w $L^1$ pozwala przejść do granicy. Otrzymujemy
$\E[M_t^f-M_s^f\mid\calF_s]=0$.
Adaptacja całki wynika z aproksymacji sumami
Riemanna, a jej ciągłość i całkowalność
z ciągłości trajektorii i ograniczoności $g$.

Dla ogólnego $f\in C_c^2(\R^d)$ weźmy sploty
$f_\varepsilon=f*\varphi_\varepsilon$ z nieujemną
gładką funkcją aproksymującą jedynkę, o całce $1$
i nośniku w kuli promienia $\varepsilon$.
Funkcje $f_\varepsilon$ są gładkie o zwartym nośniku,
a $f_\varepsilon\to f$ oraz
$\Delta f_\varepsilon\to\Delta f$ jednostajnie:
wynika to z jednostajnej ciągłości $f$ i jej drugich
pochodnych oraz różniczkowania splotu.
Dla ustalonego $T$ różnica procesów $M^{f_\varepsilon}$
i $M^f$ jest na $[0,T]$ ograniczona w module przez
$2\norm{f_\varepsilon-f}_\infty+
(T/2)\norm{\Delta f_\varepsilon-\Delta f}_\infty$.
Przechodząc do granicy w tożsamościach
martyngałowych w $L^1$, dostajemy tezę.
\end{proof}

\begin{twierdzenie}[Rozkład miejsca wyjścia z kuli ze środka]
\label{thm:brown-exit-sphere}
Niech $B$ będzie standardowym ruchem Browna w $\R^d$,
$R>0$ i $\sigma_R=\inf\{t\ge0:|B_t|\ge R\}$.
Zmienna $B_{\sigma_R}$ ma jednostajny rozkład na
sferze $S_R=\{y\in\R^d:|y|=R\}$, to znaczy
\[
\Prob(B_{\sigma_R}\in D)=
\frac{\mathcal H^{d-1}(D\cap S_R)}
{\mathcal H^{d-1}(S_R)}
\]
dla każdego zbioru borelowskiego $D\subseteq\R^d$.
W wymiarze $d=1$ miara $\mathcal H^0$ jest miarą
liczącą, a obydwa końce przedziału mają masę $1/2$.
Ponadto czas wyjścia $\sigma_R$ i miejsce wyjścia
$B_{\sigma_R}$ są niezależne.
\end{twierdzenie}
\begin{proof}
Czas $\sigma_R$ jest skończony prawie na pewno
na mocy twierdzenia o średnim czasie wyjścia.
Mierzalność $B_{\sigma_R}$ wynika z twierdzenia
\ref{thm:sampled}; na zbiorze zerowym nieskończonego
czasu można przyjąć dowolną stałą wartość.
Ciągłość trajektorii daje $|B_{\sigma_R}|=R$ p.n.
Dla dowolnej macierzy ortogonalnej $U$ czas wyjścia
procesu $UB$ jest taki sam jak czas wyjścia $B$.
Izotropia całego procesu daje więc
$\Law(UB_{\sigma_R})=\Law(B_{\sigma_R})$.
Wyjście jako funkcjonał trajektorii jest mierzalne:
czas wyjścia wyrażamy przez maksimum normy na
zwartych przedziałach, a miejsce wyjścia przez
granice ewaluacji w chwilach zaokrąglonych w górę.

Wykażemy jednoznaczność miary probabilistycznej
na $S_R$ niezmienniczej względem wszystkich
przekształceń ortogonalnych. Niech $Y$ ma taką miarę,
a $V=G/|G|$, gdzie $G\sim\Normal_d(0,I_d)$ jest
niezależny od $Y$ (w razie potrzeby na przestrzeni
iloczynowej) (na zdarzeniu $G=0$ przyjmujemy
dowolny wektor jednostkowy). Prawo $V$ jest
niezmiennicze względem przekształceń ortogonalnych.
Dla $u\in\R^d$ funkcja charakterystyczna $Y$
zależy tylko od $|u|$, więc
\[
\E e^{i\langle u,Y\rangle}
=\E e^{i|u|\langle V,Y\rangle}
=\E e^{i|u|R V_1}.
\]
W drugiej równości warunkujemy względem $Y$;
dla każdego $y\in S_R$ izotropia $V$ daje
$\langle V,y\rangle\sim R V_1$.
Prawa dowolnych dwóch miar niezmienniczych mają
zatem tę samą funkcję charakterystyczną.
Znormalizowana miara Hausdorffa na sferze
jest taką miarą, ponieważ izometrie zachowują
miarę Hausdorffa, a jej całkowita masa na sferze
jest skończona i dodatnia (dla $d\ge2$ wynika to
z lokalnych gładkich map sfery i wzoru na pole,
dla $d=1$ z liczenia punktów).
Jednoznaczność funkcji charakterystycznych dowodzi
pierwszej tezy.

Aby wykazać niezależność, ustalmy borelowski zbiór
$A\subseteq[0,\infty)$ o $\Prob(\sigma_R\in A)>0$.
Miara na sferze
\[
\nu_A(D)=
\frac{\Prob(\sigma_R\in A,\ B_{\sigma_R}\in D)}
{\Prob(\sigma_R\in A)}
\]
jest również niezmiennicza względem macierzy
ortogonalnych, gdyż obrót lub odbicie zachowuje
czas wyjścia, a przekształca położenie zgodnie z $U$.
Udowodniona jednoznaczność miary niezmienniczej
daje $\nu_A=\Law(B_{\sigma_R})$.
Dla $A$ o prawdopodobieństwie zero faktoryzacja
jest automatyczna. Otrzymujemy niezależność.
\end{proof}

\begin{wniosek}[Własność wartości średniej funkcji harmonicznej]
Niech $u$ będzie funkcją klasy $C^2$ na otwartym
otoczeniu domkniętej kuli $\{y:|y-x|\le R\}$,
harmoniczną w jej wnętrzu, czyli $\Delta u=0$.
Wtedy
\[
u(x)=
\frac{1}{\mathcal H^{d-1}(S_R)}
\int_{S_R}u(x+y)\,\mathrm d\mathcal H^{d-1}(y).
\]
\end{wniosek}
\begin{proof}
Wybieramy gładką funkcję odcinającą równą $1$
w otoczeniu domkniętej kuli, o zwartym nośniku
w obszarze definicji $u$. Po pomnożeniu przez $u$
i przedłużeniu przez zero otrzymujemy funkcję
$f\in C_c^2(\R^d)$ równą $u$ na kuli.
Niech $X_t=x+B_t$ i $\tau=\sigma_R$.
Z lematu~\ref{lem:brown-generator-martingale}
oraz twierdzenia~\ref{thm:optional-cont} wynika
$\E f(X_{\tau\wedge n})=f(x)$, ponieważ
$\Delta f(X_s)=0$ dla $s<\tau$.
Czas $\tau$ jest skończony p.n., a $f$ jest ograniczona.
Przechodząc do granicy przez twierdzenie
o zbieżności zmajoryzowanej, dostajemy
$u(x)=\E u(X_\tau)$. Twierdzenie
\ref{thm:brown-exit-sphere} daje wskazaną całkę.
\end{proof}

\begin{twierdzenie}[Trafienie w sferę wewnętrzną przed zewnętrzną]
\label{thm:brown-annulus}
Niech $0<r<|x|<R$ i $X_t=x+B_t$. Oznaczmy
$\tau_r=\inf\{t\ge0:|X_t|\le r\}$ oraz
$\sigma_R=\inf\{t\ge0:|X_t|\ge R\}$.
Wtedy
\[
\Prob_x(\tau_r<\sigma_R)=
\begin{cases}
\dfrac{R-|x|}{R-r},&d=1,\\[6pt]
\dfrac{\log(R/|x|)}{\log(R/r)},&d=2,\\[6pt]
\dfrac{|x|^{2-d}-R^{2-d}}{r^{2-d}-R^{2-d}},&d\ge3.
\end{cases}
\]
\end{twierdzenie}
\begin{proof}
Dla radialnej funkcji $u(y)=h(|y|)$ poza zerem
bezpośrednie różniczkowanie daje
$\Delta u(y)=h''(|y|)+(d-1)h'(|y|)/|y|$.
W istocie $\partial_i u=h'(\rho)y_i/\rho$,
gdzie $\rho=|y|$, oraz
\[
\partial_{ii}u(y)=h''(\rho)\frac{y_i^2}{\rho^2}
+h'(\rho)\left(\frac1\rho-\frac{y_i^2}{\rho^3}\right).
\]
Zsumowanie po $i$ daje wskazany wzór. Wybieramy $h(\rho)=\rho$
dla $d=1$, $h(\rho)=\log\rho$ dla $d=2$ i
$h(\rho)=\rho^{2-d}$ dla $d\ge3$.
W każdym przypadku $\Delta u=0$ poza zerem.
Mnożąc $u$ przez gładką funkcję odcinającą,
równą $1$ w otoczeniu domkniętego pierścienia
$\{r\le |y|\le R\}$, o nośniku zwartym poza zerem,
otrzymujemy $f\in C_c^\infty(\R^d)$ zgodną z $u$
na tym otoczeniu.

Czas $\eta=\tau_r\wedge\sigma_R$ jest czasem
zatrzymania na mocy twierdzenia~\ref{thm:hitting-closed}
i jest skończony p.n., gdyż $\eta\le\sigma_R$,
a skończoność $\sigma_R$ wynika z twierdzenia
o średnim czasie wyjścia z kuli.
Przed $\eta$ mamy $\Delta f(X_v)=0$.
Stosując lemat~\ref{lem:brown-generator-martingale}
oraz twierdzenie~\ref{thm:optional-cont} do
$\eta\wedge n$, otrzymujemy $\E_x f(X_{\eta\wedge n})=f(x)$.
Przez zbieżność zmajoryzowaną
$\E_x f(X_\eta)=f(x)$. Ciągłość trajektorii daje
$|X_\eta|\in\{r,R\}$; oba czasy nie mogą być
jednocześnie równe $\eta<\infty$.
Dla $p=\Prob_x(\tau_r<\sigma_R)$ otrzymujemy
$h(|x|)=p h(r)+(1-p)h(R)$.
Rozwiązując równanie względem $p$, dostajemy wszystkie
trzy wzory.
\end{proof}

\begin{wniosek}[Prawdopodobieństwo trafienia w domkniętą kulę]
Dla $|x|>r>0$ zachodzi
\[
\Prob_x(\tau_r<\infty)=
\begin{cases}
1,&d=1,2,\\
(r/|x|)^{d-2},&d\ge3.
\end{cases}
\]
\end{wniosek}
\begin{proof}
Zdarzenia $\{\tau_r<\sigma_R\}$ rosną przy
$R\to\infty$ wzdłuż rosnącego ciągu promieni.
Ich suma jest równa $\{\tau_r<\infty\}$:
jeśli $\tau_r<\infty$, trajektoria na $[0,\tau_r]$
jest ograniczona i mieści się w dostatecznie dużej
kuli. Wystarczy zatem przejść do granicy we
wzorach twierdzenia~\ref{thm:brown-annulus}.
\end{proof}

\begin{twierdzenie}[Trafianie w pojedyncze punkty]
Niech $X_t=x+B_t$ i niech $y$ będzie ustalonym
punktem różnym od $x$.
Jeżeli $d=1$, to $X$ trafia w $y$ prawie na pewno.
Jeżeli $d\ge2$, to z prawdopodobieństwem $1$
nie trafia w ten punkt $y$ w żadnej chwili.
Ponadto przy $d\ge2$, startując w $y$, proces
prawie na pewno nie wraca do $y$ w żadnej chwili $t>0$.
\end{twierdzenie}
\begin{proof}
Przez przesunięcie możemy przyjąć $y=0$.
Dla $d=1$ wniosek wynika z rozkładu czasu
trafienia i symetrii w rozdziale~\ref{ch:wiener}.
Dla $d\ge2$ i ustalonego $R>|x|>0$
trafienie w zero przed wyjściem z kuli promienia $R$
wymusza trafienie w każdą kulę o promieniu
$0<r<|x|$ przed tym wyjściem. Z twierdzenia
\ref{thm:brown-annulus} prawdopodobieństwo
takiego zdarzenia jest nie większe niż
$\Prob_x(\tau_r<\sigma_R)\to0$ przy $r\downarrow0$.
Każde trafienie w zero w skończonej chwili
musiałoby zajść przed wyjściem z dostatecznie
dużej kuli, więc suma po całkowitych $R>|x|$
ma prawdopodobieństwo zero.

W przypadku powrotu do punktu startowego $x=y$,
dla każdej ustalonej chwili $q>0$
mamy $\Prob_y(X_q=y)=0$, ponieważ rozkład $X_q$
ma gęstość. Warunkując względem $\calF_q$ i
korzystając z niezależności przyszłego ruchu Browna
od przeszłości (twierdzenie~\ref{thm:brown-multi-strong}
dla $\tau=q$), wnioskujemy, że prawdopodobieństwo
trafienia w $y$ w chwili $t\ge q$ jest zero.
Zastosowane warunkowanie jest uzasadnione twierdzeniem
Fubiniego: dla $z\in\R^d$ i ciągłej trajektorii $w$
funkcja $\inf_{0\le v\le k}|z+w(v)-y|$ jest ciągła
w parze $(z,w)$ z normą supremum na $[0,k]$.
Zdarzenie trafienia w chwili $t\ge q$ jest zatem
borelowskie jako przeliczalna suma zdarzeń
zerowania takiego minimum dla $k=1,2,\ldots$.
Dla $z\ne y$ jego prawdopodobieństwo wynosi zero
na mocy pierwszej części twierdzenia, a stan $z=y$
ma w chwili $q$ prawdopodobieństwo zero. Suma po
$q=1/n$ wyklucza wszystkie powroty w chwilach $t>0$.
\end{proof}

\begin{twierdzenie}[Powracalność i przejściowość ruchu Browna]
Dla dowolnego ustalonego punktu startowego $x$
i $d=1,2$ z prawdopodobieństwem $1$ ruch Browna
odwiedza każdy niepusty zbiór otwarty w dowolnie
późnych chwilach. Dla $d\ge3$ zachodzi prawie na pewno
\[
\lim_{t\to\infty}|x+B_t|=\infty.
\]
\end{twierdzenie}
\begin{proof}
Dla $d=1,2$, dowolnej otwartej kuli $U$
i całkowitej chwili $n\ge1$ prawdopodobieństwo
wejścia do $U$ w chwili $t\ge n$ wynosi $1$:
warunkujemy względem $\calF_n$ i stosujemy wniosek
o trafieniu w domkniętą kulę zawartą w $U$
(z dowolnego punktu startowego). Jeśli stan należy
już do tej mniejszej kuli, trafienie jest natychmiastowe.
Mierzalność zdarzeń wejścia do otwartej kuli wynika
z ciągłości i wystarczy sprawdzać chwile wymierne.
Przecinamy zbiory pełnej miary po $n$ oraz
po przeliczalnej bazie kul o wymiernych środkach
i dodatnich wymiernych promieniach.

Dla $d\ge3$ ustalmy $r>0$ i wybierzmy rosnące
promienie $R_n>\max(r,|x|)$, $R_n\to\infty$.
Czasy $\sigma_{R_n}$ są skończone p.n. i
$|X_{\sigma_{R_n}}|=R_n$.
Mocna własność Markowa
(twierdzenie~\ref{thm:brown-multi-strong}) oraz
wzór na trafienie w kulę dają (na mocy niezależności
przyszłej trajektorii od $\calF_{\sigma_{R_n}}$,
warunkujemy po położeniu $X_{\sigma_{R_n}}$;
zdarzenie trafienia jest borelowskie, gdyż można
je opisać minimum odległości na przedziałach
$[0,k]$, $k\ge1$)
\[
\Prob_x(\exists t\ge\sigma_{R_n}:|X_t|\le r)
=(r/R_n)^{d-2}\longrightarrow0.
\]
Zdarzenia po lewej są malejące (czasy wyjścia
rosną), zatem ich przecięcie ma miarę zero.
Prawie na pewno dla pewnego $n$ nie ma już
wizyt w kuli promienia $r$ po $\sigma_{R_n}$.
Przecinając zbiory pełnej miary dla $r=1,2,\ldots$,
otrzymujemy ucieczkę do nieskończoności.
\end{proof}

\begin{uwaga}[Powracalność w płaszczyźnie a trafianie w punkty]
Dwuwymiarowy ruch Browna prawie na pewno wchodzi
dowolnie późno do każdego otoczenia każdego punktu,
lecz ustalony punkt różny od punktu startowego
prawie na pewno nie zostaje trafiony.
Powracalność dotyczy otoczeń, nie pojedynczych punktów.
\end{uwaga}

\section{Maksimum ruchu Browna i mostu Browna}
Pełną zasadę odbicia i łączny rozkład maksimum
$M_T=\sup_{0\le s\le T}W_s$ oraz położenia $W_T$
udowodniono w twierdzeniu~\ref{thm:wiener-reflection}
i wniosku~\ref{cor:wiener-max-joint}.
Przypomnijmy ich konsekwencję: dla $a>0$, $T>0$ i
borelowskiego $A\subset(-\infty,a)$
\[
\Prob(M_T<a,\ W_T\in A)
=\int_A(p_T(y)-p_T(2a-y))\,\mathrm d\lambda_1(y),
\]
gdzie $p_T$ jest gęstością $\Normal(0,T)$.
Ponieważ $M_T$ ma prawo $|W_T|$, z momentu bezwzględnego
standardowej zmiennej normalnej otrzymujemy
$\E M_T=\sqrt{2T/\pi}$.
Wykorzystamy teraz te wyniki do analizy mostu Browna.

\begin{wniosek}[Maksimum mostu Browna]
Niech $W$ będzie standardowym procesem Wienera,
$T>0$ i $\beta_t=W_t-(t/T)W_T$ dla $0\le t\le T$.
Dla każdego $a>0$ zachodzi
\[
\Prob\left(\sup_{0\le t\le T}\beta_t\le a\right)
=1-\exp\left(-\frac{2a^2}{T}\right).
\]
\end{wniosek}
\begin{proof}
Z rozdziału~\ref{ch:gauss} wiemy, że $\beta$ jest
niezależny od $W_T$. Dla każdego $y\in\R$ proces
$\beta_t+(t/T)y$ ma rozkład warunkowy całej trajektorii
Wienera przy $W_T=y$; jest to \emph{regularna wersja}
rozkładu warunkowego, a nie warunkowanie na zdarzeniu o
prawdopodobieństwie dodatnim.

Niech $V=\sup_{0\le t\le T}\beta_t$ i dla $b>0$, $y<b$ połóżmy
$H(y,b)=\Prob(\sup_{0\le t\le T}(\beta_t+ty/T)<b)$.
Z twierdzenia~\ref{thm:wiener-reflection} i powyższego wzoru,
przez podzielenie gęstości
wspólnej zdarzenia braku przekroczenia i położenia końcowego
przez $p_T(y)>0$, otrzymujemy dla prawie każdego $y<b$
\[
H(y,b)=1-\frac{p_T(2b-y)}{p_T(y)}.
\]
Wyrażenie po prawej stronie ma dla $y\to0$ granicę
$g(b):=1-\exp(-2b^2/T)$, o ile $b>0$.

Ustalmy $a>0$ i $0<\varepsilon<a$. Jeśli $|y|<\min(\varepsilon,a-\varepsilon)$,
to oba progi $a-\varepsilon,a+\varepsilon$ są większe
od $y$, a dla każdej trajektorii $\sup_t|(t/T)y|\le\varepsilon$.
Dlatego mamy zawieranie zdarzeń prowadzące do nierówności
\[
H(y,a-\varepsilon)\le\Prob(V<a)
\le H(y,a+\varepsilon).
\]
Możemy wybrać ciąg $y_k\to0$ spoza sumy dwóch zbiorów
wyjątkowych miary Lebesgue'a zero, odpowiadających progom
$a-\varepsilon$ i $a+\varepsilon$. Podstawienie wzoru
na $H$ i przejście do granicy daje
$g(a-\varepsilon)\le\Prob(V<a)\le g(a+\varepsilon)$.
Ponieważ $g$ jest ciągła, po $\varepsilon\downarrow0$
otrzymujemy $\Prob(V<a)=g(a)$. Wreszcie
$\{V\le a\}=\bigcap_{n\ge1}\{V<a+1/n\}$;
ciągłość miary z góry i ciągłość $g$ dają
$\Prob(V\le a)=g(a)$. Dowód przy okazji wykazał,
że $\Prob(V=a)=0$ dla każdego $a>0$.
\end{proof}


\begin{uwaga}[Dlaczego przejście $y\to0$ nie jest tautologią?]
Dla ciągłej zmiennej $W_T$ zdarzenie $\{W_T=0\}$ ma
prawdopodobieństwo zero. Reguła Bayesa dla zdarzeń o dodatniej
mierze nie daje zatem rozkładu mostu przy zerowym punkcie
końcowym. Poprzedni dowód najpierw konstruuje \emph{konkretną}
wersję warunkowego prawa przy każdym $y$ przez rozkład
gaussowski, a następnie wykorzystuje ciągłość przesunięcia
trajektorii o funkcję $ty/T$. Wartość w $y=0$ jest tym
samym uzasadniona przez tę wersję, nie przez arbitralne
podstawienie do tożsamości obowiązującej tylko prawie wszędzie.
\end{uwaga}

\begin{zadanie}[Maksimum ruchu Browna w jednostce czasu]
Niech $W$ będzie standardowym procesem Wienera.
Obliczyć $\Prob(\sup_{0\le s\le1}W_s\ge1)$ oraz
$\E\sup_{0\le s\le1}W_s$.

\emph{Rozwiązanie.} Z zasady odbicia pierwsze
prawdopodobieństwo wynosi $2(1-\Phi(1))\approx0{,}3173$.
Wartość oczekiwana maksimum to $\sqrt{2/\pi}$.
Obserwujemy, że $\Prob(W_1\ge1)=1-\Phi(1)$ jest dwa razy
mniejsze: część trajektorii zdążyła osiągnąć próg,
a później wróciła poniżej niego.
\end{zadanie}


\begin{zadanie}[Warunkowy rozkład położenia mostu w dwóch wymiarach]
Niech $B$ będzie standardowym ruchem Browna w $\R^2$.
Wyznaczyć regularny rozkład warunkowy $B_1$ przy $B_4=y$,
gdzie $y=(2,-2)$, a następnie obliczyć prawdopodobieństwo,
że pierwsza współrzędna $B_1$ jest dodatnia.

\emph{Rozwiązanie.} Nie wolno dzielić przez $\Prob(B_4=y)=0$.
Korzystamy z rozkładu warunkowego wektorów gaussowskich.
Mamy $\Cov(B_1,B_4)=I_2$ oraz $\Var(B_4)=4I_2$.
Średnia warunkowa jest więc równa $I_2(4I_2)^{-1}y=y/4$,
a kowariancja wynosi
$I_2-I_2(4I_2)^{-1}I_2=\tfrac34 I_2$.
Zatem
\[
\Law(B_1\mid B_4=y)=\Normal_2\!\left(\frac{y}{4},\frac34 I_2\right).
\]
Dla $y=(2,-2)$ pierwsza współrzędna ma rozkład
$\Normal(1/2,3/4)$, skąd
$\Prob(B_1^{(1)}>0\mid B_4=y)=\Phi(1/\sqrt3)$.
Wynik opisuje \emph{wersję} warunkowania określoną
przez ciągłą funkcję parametru $y$.
\end{zadanie}

\begin{zadanie}[Warunkowanie położenia procesu Wienera w środku przedziału]
Niech $W$ będzie standardowym procesem Wienera. Wyznaczyć
regularny rozkład warunkowy $W_1$ przy $W_2=y$ oraz
obliczyć $\Prob(W_1>0\mid W_2=2)$ w wybranej wersji
rozkładu warunkowego.

\emph{Rozwiązanie.} Wektor $(W_1,W_2)$ jest gaussowski,
$\Var(W_1)=1$, $\Var(W_2)=2$ i
$\Cov(W_1,W_2)=1$. Wzór na warunkowanie wektorów
gaussowskich daje
\[
\Law(W_1\mid W_2=y)=\Normal(y/2,1/2).
\]
Nie dzielimy przez $\Prob(W_2=y)$, bo jest ono zerowe:
wykorzystujemy regularną wersję warunkowego rozkładu.
Przy $y=2$ otrzymujemy
\[
\Prob(W_1>0\mid W_2=2)=\Phi(\sqrt 2).
\]
Średnia warunkowa $y/2$ jest liniową interpolacją
między punktem początkowym $0$ i zaobserwowanym punktem
$y$; nie znaczy to jednak, że cała trajektoria mostu
jest deterministyczną prostą. Warunkowa wariancja jest
dodatnia i wynosi $1/2$.
\end{zadanie}

\section{Oscylacje drogi na dowolnie małych przedziałach}
\begin{twierdzenie}[Natychmiastowe oscylacje przy zadanej chwili]
Niech $W$ będzie standardowym procesem Wienera i niech
$t_0\ge0$ będzie ustaloną deterministyczną chwilą.
Wtedy z prawdopodobieństwem $1$ dla każdego $\varepsilon>0$
istnieją $s_+,s_-\in(t_0,t_0+\varepsilon)$ takie, że
$W_{s_+}>W_{t_0}$ oraz $W_{s_-}<W_{t_0}$.
\end{twierdzenie}
\begin{proof}
Proces $B_s=W_{t_0+s}-W_{t_0}$ jest standardowym
ruchem Browna. Dla każdego $r>0$ i $a>0$ zasada odbicia
daje
$\Prob(\sup_{0\le s\le r}B_s\ge a)=2(1-\Phi(a/\sqrt r))$.
Gdy $a\downarrow0$, prawdopodobieństwo to dąży do $1$.
Z ciągłości miary dla rosnących zdarzeń dostajemy
$\Prob(\sup_{0\le s\le r}B_s>0)=1$.
Ta sama własność dla $-B$ daje
$\Prob(\inf_{0\le s\le r}B_s<0)=1$.
Przecinamy zdarzenia pełnej miary dla $r=1/n$,
$n\in\N$. Z tego wynika teza dla wszystkich
$\varepsilon>0$, przez wybór $1/n<\varepsilon$.
\end{proof}

\begin{wniosek}[Natychmiastowe oscylacje po czasie zatrzymania]
Niech $W$ będzie procesem Wienera względem filtracji
$(\calF_t)_{t\ge0}$ o ciągłych trajektoriach i niech
$\tau$ będzie czasem zatrzymania skończonym p.n.
Wtedy prawie na pewno dla każdego $\varepsilon>0$
istnieją $u_+,u_-\in(0,\varepsilon)$ takie, że
\[
W_{\tau+u_+}>W_\tau,\qquad
W_{\tau+u_-}<W_\tau.
\]
\end{wniosek}
\begin{proof}
Twierdzenie~\ref{thm:wiener-strong-markov} daje
nowy proces Wienera $V_u=W_{\tau+u}-W_\tau$
na zdarzeniu skończoności $\tau$, z trajektorią
zerową na jego dopełnieniu. Stosujemy twierdzenie
o natychmiastowych oscylacjach do $V$ przy $t_0=0$.
Wspólny zbiór pełnej miary dla wszystkich
$\varepsilon>0$ otrzymujemy z przeliczalnej
rodziny przedziałów o długości $1/n$.
\end{proof}

\begin{uwaga}[Nie wolno zamienić ``każda ustalona chwila'' na ``wszystkie chwile'']
Udowodnione zdarzenie pełnej miary zależy od $t_0$.
Na ogół nie wolno przecinać nieprzeliczalnie wielu
zbiorów pełnej miary. Własność stwierdzona w twierdzeniu
nie zachodzi jednocześnie dla wszystkich chwil trajektorii:
na przykład maksimum na $[0,T]$ jest osiągane
w pewnej chwili $t_*<T$ prawie na pewno: proces odwrócony
$(W_T-W_{T-s})_{0\le s\le T}$ ma prawo Wienera
i przyjmuje ujemne wartości prawie na pewno, więc $W_T$
nie jest maksimum. Dla takiej chwili $t_*$ nie może
wystąpić większa wartość przed końcem przedziału. To ważna lekcja dotycząca kwantyfikatorów
w twierdzeniach o trajektoriach.
W szczególności chwila pierwszego osiągnięcia
maksimum na całym $[0,T]$ nie jest czasem zatrzymania
względem naturalnej filtracji: gdyby była, poprzedni
wniosek wymuszałby przekroczenie tego maksimum
przed $T$ na zdarzeniu $\{t_*<T\}$.
Mierzalność pierwszej chwili maksimum jako funkcjonału
całej trajektorii wynika z
$\{t_*\le s\}=\{\max_{[0,s]}W=\max_{[0,T]}W\}$
dla $0\le s\le T$.
\end{uwaga}

\section{Wielowymiarowa wariacja kwadratowa}
\begin{twierdzenie}[Macierz wariacji kwadratowej ruchu Browna]
\label{thm:brown-multi-bracket}
Niech $T>0$, $B$ będzie standardowym ruchem Browna
w $\R^d$ i niech $\pi=\{0=t_0<\cdots<t_m=T\}$.
Dla $t\in[0,T]$ oznaczmy
$D_k(t)=B_{t\wedge t_{k+1}}-B_{t\wedge t_k}$ oraz
\[
Q^\pi_t=\sum_{k=0}^{m-1}D_k(t)D_k(t)^\top.
\]
Dla dowolnego deterministycznego ciągu podziałów
$\pi_n$ o średnicach dążących do zera zachodzi
\[
\sup_{0\le t\le T}\norm{Q^{\pi_n}_t-tI_d}
\longrightarrow0
\quad\text{według prawdopodobieństwa},
\]
gdzie można użyć dowolnej normy macierzowej.
Jeżeli ponadto $\sum_n|\pi_n|<\infty$, zbieżność
jest prawie pewna.
\end{twierdzenie}
\begin{proof}
Dla ustalonej współrzędnej stosujemy jednowymiarowy
rachunek wariacji kwadratowej z rozdziału~\ref{ch:wiener}.
Dokładniej, dla dowolnego podziału, $h_k=t_{k+1}-t_k$,
zmienne $Y_k=(B^{(i)}_{t_{k+1}}-B^{(i)}_{t_k})^2-h_k$
są niezależne, wycentrowane i mają wariancje $2h_k^2$.
Nierówność maksymalna Kołmogorowa daje
\[
\Prob\left(\max_{j\le m}\left|\sum_{k<j}Y_k\right|
>\varepsilon\right)
\le \frac{2T|\pi|}{\varepsilon^2}.
\]
Między węzłami błąd względem lewego węzła
nie przekracza kwadratu modułu ciągłości trajektorii
na skali $|\pi|$ powiększonego o $|\pi|$.
Zatem mamy jednostajną zbieżność według
prawdopodobieństwa, a przy sumowalnych średnicach
także prawie pewną (Borel--Cantelli dla
$\varepsilon=1/l$, $l\ge1$).
Dla $i\ne j$ proces
$W=(B^{(i)}+B^{(j)})/\sqrt2$ jest standardowym
procesem Wienera na mocy charakterystyki gaussowskiej.
Z tożsamości polaryzacyjnej
\[
Q^\pi_{t,ij}
=V_t^\pi(W)-\tfrac12V_t^\pi(B^{(i)})
-\tfrac12V_t^\pi(B^{(j)}),
\]
gdzie $V^\pi$ oznacza jednowymiarową sumę kwadratów,
wynika zbieżność jednostajna do $t-t/2-t/2=0$.
Na przekątnej granica wynosi $t$.
Ponieważ współrzędnych jest skończenie wiele,
otrzymujemy tezę w dowolnej normie macierzowej
przez równoważność norm. W rozdziale
\ref{ch:quadratic-cont} powrócimy do tego pojęcia
dla ogólnych ciągłych martyngałów.
\end{proof}

\section{Rozkład promienia i oszacowania ryzyka wyjścia}
\begin{twierdzenie}[Rozkład odległości od początku]
Niech $B$ będzie standardowym ruchem Browna w $\R^d$
i $t>0$. Wtedy $|B_t|^2/t$ ma rozkład $\chi^2_d$,
a w szczególności
$\E|B_t|^2=dt$ i $\Var(|B_t|^2)=2dt^2$.
\end{twierdzenie}
\begin{proof}
Z definicji współrzędne $B_t^{(i)}/\sqrt t$ są niezależnymi
zmiennymi o rozkładzie $\Normal(0,1)$. Zatem
$|B_t|^2/t=\sum_{i=1}^d(B_t^{(i)}/\sqrt t)^2$,
co daje rozkład $\chi^2_d$ na mocy twierdzenia
o sumie kwadratów niezależnych zmiennych normalnych
z tomu I, rozdziału~\ref{I-ch:distributions}.
Każdy składnik kwadratowy ma średnią $1$ i wariancję $2$,
a niezależność pozwala dodać wartości oczekiwane
i wariancje.
\end{proof}

\begin{wniosek}[Gęstość promienia i jego momenty]
Dla $t>0$ zmienna $R_t=|B_t|$ ma na $(0,\infty)$
gęstość
\[
f_{R_t}(r)=
\frac{2^{1-d/2}}{\Gamma(d/2)t^{d/2}}
r^{d-1}e^{-r^2/(2t)}.
\]
Dla rzeczywistego $q$ moment $\E R_t^q$ jest
skończony wtedy i tylko wtedy, gdy $q>-d$;
w tym przypadku
\[
\E R_t^q=(2t)^{q/2}
\frac{\Gamma((d+q)/2)}{\Gamma(d/2)}.
\]
\end{wniosek}
\begin{proof}
Zmienna $Z=R_t^2/t$ ma rozkład $\chi_d^2$
o gęstości
$z^{d/2-1}e^{-z/2}/(2^{d/2}\Gamma(d/2))$ dla $z>0$
(por.~\eqref{I-eq:chisq}).
Podstawienie $z=r^2/t$ daje czynnik Jacobiego
$2r/t$ i wskazaną gęstość. Nie ma atomu w zerze,
ponieważ $B_t$ ma gęstość w $\R^d$.
W całce na moment podstawiamy $u=r^2/(2t)$:
\[
\E R_t^q=\frac{(2t)^{q/2}}{\Gamma(d/2)}
\int_0^\infty u^{(d+q)/2-1}e^{-u}\,\mathrm d\lambda_1(u).
\]
Dla każdego $q$ całka jest skończona przy
nieskończoności dzięki czynnikowi wykładniczemu.
Przy zerze jest skończona dokładnie wtedy,
gdy $(d+q)/2>0$. Otrzymujemy wzór z definicji
funkcji gamma.
\end{proof}

\begin{zadanie}[Oszacowanie prawdopodobieństwa wyjścia z sześcianu]
Niech $B$ będzie standardowym ruchem Browna w $\R^d$, $a,T>0$.
Oszacować $\Prob(\exists t\in[0,T]:\max_i|B_t^{(i)}|\ge a)$.

\emph{Rozwiązanie.} Jeśli proces opuszcza sześcian $(-a,a)^d$,
to dla pewnej współrzędnej $i$ zachodzi
$\sup_{t\le T} B_t^{(i)}\ge a$ lub
$\inf_{t\le T} B_t^{(i)}\le-a$.
Z nierówności sumy i zasady odbicia,
\[
\Prob\bigl(\sup_{t\le T}\|B_t\|_\infty\ge a\bigr)
\le4d\bigl(1-\Phi(a/\sqrt T)\bigr).
\]
Możemy zastąpić prawą stronę jej minimum z $1$.
Oszacowanie jest proste, choć nie musi być ostre,
ponieważ wielokrotnie zlicza trajektorie przekraczające
kilka ścian sześcianu. Z niezależności całych procesów
współrzędnych otrzymujemy również dokładną redukcję
do problemu jednowymiarowego:
\[
\Prob\left(\sup_{t\le T}\norm{B_t}_\infty\ge a\right)
=1-(1-p)^d,\qquad
p=\Prob\left(\sup_{t\le T}|W_t|\ge a\right).
\]
Zdarzenia pozostawania poszczególnych współrzędnych
w $(-a,a)$ przez cały $[0,T]$ są niezależne,
a każde z nich ma prawdopodobieństwo $1-p$.
\end{zadanie}

\section{Schematyczny rysunek procesu dwuwymiarowego}
\begin{center}
\begin{tikzpicture}[x=0.85cm,y=0.85cm,>=Latex]
 \draw[->,gray] (-3.2,0)--(3.7,0) node[right] {$x_1$};
 \draw[->,gray] (0,-2.5)--(0,2.7) node[above] {$x_2$};
 \draw[blue!70!black,very thick,rounded corners=1pt]
 (0,0)--(.15,.3)--(.45,.1)--(.1,.6)--(.3,.9)--(-.1,1.15)--(-.7,.9)
 --(-.55,.4)--(-.9,.2)--(-.4,-.2)--(-.8,-.7)--(-1.2,-.5)
 --(-1.45,-.85)--(-1.1,-1.1)--(-1.7,-1.4)--(-1.3,-1.7)
 --(-2.1,-1.5)--(-1.85,-1.1)--(-2.25,-.65)--(-1.9,-.25)
 --(-2.6,.2)--(-2.15,.6)--(-2.55,1.05)--(-2.1,1.3);
 \fill[red!75!black] (0,0) circle (2pt) node[below right] {$B_0$};
 \fill[blue!70!black] (-2.1,1.3) circle (2pt) node[above] {$B_T$};
 \draw[->,blue!70!black,thick] (-2.55,1.05)--(-2.1,1.3);
\end{tikzpicture}
\par\smallskip\small Schemat dwóch współrzędnych ruchu Browna w płaszczyźnie.
Łamana jest jedynie ilustracją: niemal każda rzeczywista trajektoria
jest ciągła i nieróżniczkowalna.
\end{center}


% ==============================================================
% ROZDZIAŁ 9. PROCESY ODNOWY
% ==============================================================

\chapter{Procesy odnowy: zliczanie zdarzeń bez własności braku pamięci}
\label{ch:renewal}
\paragraph{Odsyłacze międzytomowe.} Modele intensywności i eksponowania czasu obserwacji znajdziesz w tomie IV, rozdziale~\ref{IV-ch:glm}.


\section{Dlaczego proces Poissona nie opisuje każdego strumienia zdarzeń?}
W procesie Poissona po dowolnej deterministycznej chwili obserwacji oczekiwanie na
kolejne zdarzenie ma nadal rozkład wykładniczy. To wynika z własności
braku pamięci. Awaria mechanizmu, którego prawdopodobieństwo
uszkodzenia rośnie z wiekiem, nie ma jednak na ogół tej własności.
Potrzebujemy więc modelu z niezależnymi, ale niekoniecznie
wykładniczymi czasami między kolejnymi zdarzeniami.
\emph{Proces odnowy} opisuje liczbę ukończonych cykli, w których
po każdym odnowieniu system zaczyna działać jak nowy.

\begin{definicja}[Proces odnowy]
Niech $(T_n)_{n\ge1}$ będzie ciągiem niezależnych zmiennych
losowych o wspólnym prawie $\nu$ na $(0,\infty)$.
Niech $F(t)=\Prob(T_1\le t)$ będzie jego dystrybuantą,
a $\overline F(t)=1-F(t)$ funkcją przeżycia.
Na wspólnym mierzalnym zbiorze zerowym możemy
zastąpić wszystkie odstępy wartością $1$, aby
$0<T_n<\infty$ zachodziło punktowo dla każdego $n$.
Nie zmienia to rozkładu całego ciągu.
Definiujemy $S_0=0$,
$S_n=\sum_{k=1}^nT_k$ oraz, dla $t\ge0$,
\[
N_t:=\sup\{n\ge0:S_n\le t\}\in\N_0\cup\{\infty\}.
\]
Proces $N=(N_t)_{t\ge0}$ nazywamy \emph{procesem odnowy}.
Liczby $S_n$ to chwile odnowień, a $T_n=S_n-S_{n-1}$
to czasy pomiędzy odnowieniami. Gdy $T_n\sim\mathrm{Exp}(\lambda)$,
odzyskujemy proces Poissona o intensywności $\lambda$
na mocy konstrukcji z rozdziału~\ref{ch:poisson-process}.
\end{definicja}

\begin{twierdzenie}[Brak eksplozji i podstawowe prawdopodobieństwa]
Dla procesu odnowy z poprzedniej definicji istnieje
mierzalny zbiór pełnej miary, na którym $N_t<\infty$
jednocześnie dla wszystkich $t\ge0$. Na tym zbiorze
trajektorie są niemalejące, prawostronnie ciągłe,
mają granice lewostronne i na każdym ograniczonym
przedziale wykonują skończenie wiele skoków wielkości $1$.
Dla $n\in\N_0$ zachodzi
\[
\Prob(N_t\ge n)=\Prob(S_n\le t),\qquad
\Prob(N_t=n)=\Prob(S_n\le t)-\Prob(S_{n+1}\le t).
\]
Dla $n\ge1$ prawem $S_n$ jest $n$-krotny splot miar
$\nu^{*n}$. Jego dystrybuantę oznaczamy przez
$F^{*n}(t)=\nu^{*n}([0,t])$; przyjmujemy też
$\nu^{*0}=\delta_0$.
\end{twierdzenie}
\begin{proof}
Ponieważ $\Prob(T_1>0)=1$, istnieje $\varepsilon>0$
takie, że $p:=\Prob(T_1\ge\varepsilon)>0$.
Zmienne $I_k=\chi_{\{T_k\ge\varepsilon\}}$ są niezależne
i mają wspólny rozkład Bernoulliego z parametrem $p$.
Z mocnego prawa wielkich liczb
$n^{-1}\sum_{k=1}^n I_k\to p>0$ prawie na pewno.
Ponieważ $S_n\ge\varepsilon\sum_{k=1}^nI_k$, wynika stąd
$S_n\to\infty$ prawie na pewno; dlatego na jednym
zbiorze pełnej miary $N_t$ jest skończone dla każdego
skończonego $t$. Na każdym ograniczonym przedziale
trajektoria jest funkcją schodkową z izolowanymi
punktami skoku $S_n$. Tożsamość
$N_t=\sum_{n\ge1}\chi_{\{S_n\le t\}}$ dowodzi
mierzalności $N_t$, a włączenie punktu $S_n$
do kolejnego przedziału stałości daje prawostronną
ciągłość. Ścisła dodatniość odstępów zapewnia
skoki wielkości $1$ i $N_0=0$.
Równość zdarzeń $\{N_t\ge n\}=\{S_n\le t\}$ jest bezpośrednia;
odejmując prawdopodobieństwa kolejnych zagnieżdżonych zdarzeń,
uzyskujemy drugi wzór. Rozkład sumy niezależnych zmiennych
jest splotem ich miar probabilistycznych.
\end{proof}

W dalszej części przyjmujemy wersję bez eksplozji
dla wszystkich trajektorii: na mierzalnym zbiorze
$\{S_n\not\to\infty\}$ miary zero zastępujemy
cały ciąg odstępów jedynkami. Zachowujemy wszystkie
jego prawa skończeniewymiarowe oraz niezależność.

\begin{uwaga}[Warunek $T_n>0$ jest istotny]
Gdyby dopuszczano $\Prob(T_1=0)>0$, lecz nadal
$\Prob(T_1>0)>0$, nie doszłoby do eksplozji, ale mogłyby
wystąpić wielokrotne skoki w tej samej chwili.
Założenie dodatniości usuwa ten problem i upraszcza zapis.
Jeżeli zaś $T_1=0$ prawie na pewno, wszystkie odnowienia
nastąpiłyby w chwili zerowej, a $N_0=\infty$.
\end{uwaga}

\begin{stwierdzenie}[Skończona średnia liczby odnowień]
Dla każdego procesu odnowy z $T_1>0$ prawie na pewno i każdego
$t<\infty$ zachodzi $\E N_t<\infty$, nawet gdy $\E T_1=\infty$.
\end{stwierdzenie}
\begin{proof}
Ustalmy $\theta>0$. Liczba
$r=\E e^{-\theta T_1}$ spełnia $0<r<1$.
Z nierówności Markowa i niezależności
\[
\Prob(S_n\le t)=\Prob(e^{-\theta S_n}\ge e^{-\theta t})
\le e^{\theta t}\E e^{-\theta S_n}=e^{\theta t}r^n.
\]
Z tożsamości $N_t=\sum_{n\ge1}\chi_{\{S_n\le t\}}$
i twierdzenia Tonellego otrzymujemy
$\E N_t=\sum_{n\ge1}\Prob(S_n\le t)
\le e^{\theta t}r/(1-r)<\infty$.
\end{proof}

\begin{center}
\begin{tikzpicture}[x=1cm,y=0.55cm,>=Latex]
\draw[->] (0,0)--(11,0) node[right] {$t$};
\draw[->] (0,0)--(0,5.8) node[above] {$N_t$};
\draw[very thick,blue!70!black]
(0,0)--(1.6,0) (1.6,1)--(3.8,1) (3.8,2)--(4.8,2)
(4.8,3)--(8,3) (8,4)--(10.2,4);
\foreach \x/\y in {1.6/0,3.8/1,4.8/2,8/3}{
 \draw[blue!70!black,very thick] (\x,\y)--(\x,{\y+1});
 \fill[blue!70!black] (\x,{\y+1}) circle (2pt);
 \draw[dashed,gray] (\x,0)--(\x,\y);
}
\foreach \x/\t in {1.6/1,3.8/2,4.8/3,8/4}{
 \node[below] at (\x,0) {$S_{\t}$};
}
\foreach \y in {1,2,3,4}{\node[left] at (0,\y) {$\y$};}
\end{tikzpicture}
\par\smallskip\small Schemat procesu odnowy. Odstępy
$T_1=S_1$, $T_2=S_2-S_1$, $T_3=S_3-S_2$ nie muszą mieć rozkładu
wykładniczego; dlatego liczba odnowień na rozłącznych przedziałach
zwykle nie jest niezależna.
\end{center}

\begin{przyklad}[Deterministyczny cykl]
Jeżeli $T_n\equiv c>0$, to $S_n=nc$ i
$N_t=\lfloor t/c\rfloor$. W modelu Poissona
$\Var(N_t)=\E N_t$, tutaj zaś $\Var(N_t)=0$.
Proces odnowy obejmuje więc zarówno modele losowe, jak i
całkowicie regularne cykle. Przykład pokazuje, że sam parametr
$1/\E T_1$ nie opisuje charakteru fluktuacji.
Ten deterministyczny proces jest jednak procesem Markowa
(w sensie niejednorodnym w czasie): znając aktualną chwilę,
znamy bez losowości wszystkie jego przyszłe wartości.
Nie może więc służyć jako dowód, że ogólny proces odnowy
nie ma własności Markowa. Rzeczywisty kontrprzykład podajemy dalej.
\end{przyklad}


\begin{twierdzenie}[Proces odnowy nie musi być procesem Markowa]
Niech $(T_n)_{n\ge1}$ będą niezależnymi zmiennymi, dla których
$\Prob(T_n=1)=\Prob(T_n=2)=1/2$, i niech
$N_t=\sup\{n\ge0:\sum_{k=1}^nT_k\le t\}$.
Wtedy proces $(N_t)_{t\ge0}$ \emph{nie jest} procesem Markowa
względem swojej filtracji naturalnej
$\mathcal F_t^N=\sigma(N_s:0\le s\le t)$.
\end{twierdzenie}
\begin{proof}
Spójrzmy na chwile $1$, $5/2$ i $13/4$.
Rozważmy dwa zdarzenia należące do $\mathcal F_{5/2}^N$:
\[
A=\{N_1=1,N_{5/2}=1\},\qquad
D=\{N_1=0,N_{5/2}=1\}.
\]
Oba mają dodatnie prawdopodobieństwo.
Na $A$ koniecznie $T_1=1$ i $T_2=2$: gdyby
$T_2=1$, druga odnowa nastąpiłaby już w chwili $2$.
Zatem $S_2=3<13/4$ i
$\Prob(N_{13/4}\ge2\mid A)=1$.
Na $D$ mamy $T_1=2$, a $T_2$ nadal przyjmuje wartości
$1$ i $2$ z jednakowymi prawdopodobieństwami,
niezależnie od $D$. Stąd
$\Prob(N_{13/4}\ge2\mid D)=1/2$.

Na obu zdarzeniach obecny stan $N_{5/2}$ jest taki sam,
mianowicie równy $1$. Gdyby $N$ miał własność Markowa,
warunkowe prawdopodobieństwo zdarzenia
$\{N_{13/4}\ge2\}$ przy całej informacji do chwili $5/2$
byłoby funkcją wyłącznie $N_{5/2}$. Jego średnie
warunkowe na $A$ i $D$ musiałyby więc być równe.
Obliczone wartości $1$ i $1/2$ przeczą temu warunkowi.
\end{proof}

\begin{uwaga}[Stan trzeba czasem poszerzyć]
Sama liczba odnowień $N_t$ nie mówi, ile czasu upłynęło
od ostatniej odnowy. Wprowadzamy wiek
$A_t=t-S_{N_t}$ i czas resztkowy $R_t=S_{N_t+1}-t$. To właśnie brakująca informacja w
powyższym dowodzie. Para $(N_t,A_t)$, gdzie
$A_t=t-S_{N_t}$, jest naturalnym kandydatem na stan
procesu Markowa: po znajomości wieku rozkład czasu
pozostałego do nowego cyklu można obliczyć z ogona
rozkładu odstępów. Własność Markowa takiej pary
wymaga osobnego sformułowania i dowodu; jej nie
utożsamiamy z własnością Markowa samego $N_t$.
\end{uwaga}


\begin{twierdzenie}[Wiek odnowy przywraca własność Markowa]
\label{thm:wiek-markow}
Niech $T_1,T_2,\ldots$ będą niezależnymi, jednakowo rozłożonymi
zmiennymi o wartościach w $(0,\infty)$, o dystrybuancie $F$,
a $N_t$ i $A_t=t-S_{N_t}$ oznaczają liczbę odnowień i wiek
bieżącego cyklu. Wówczas $Y_t=(N_t,A_t)$ jest jednorodnym
w czasie procesem Markowa względem filtracji naturalnej
$\calF_t^N=\sigma(N_s:0\le s\le t)$.
Dokładniej, warunkowe prawo czasu resztkowego
względem całej $\calF_t^N$ zależy tylko od $A_t$:
jego funkcja przeżycia dla $u\ge0$ wynosi
$\overline F(A_t+u)/\overline F(A_t)$ prawie na pewno,
a po następnym odnowieniu proces jest zbudowany z nowych,
niezależnych odstępów o prawie $F$. Opis ten nie zależy od $t$.
\end{twierdzenie}
\begin{proof}
Ustalmy deterministyczne $t$. Z trajektorii $(N_s)_{0\le s\le t}$
można odzyskać wszystkie chwile skoków $S_k\le t$, a więc
$N_t$ i $S_{N_t}$ są $\calF_t^N$-mierzalne. Na zdarzeniu
$\{N_t=n\}$ znamy $T_1,\ldots,T_n$, natomiast o $T_{n+1}$
wiemy tylko, że $T_{n+1}>t-S_n=A_t$. Istotnie,
$\{N_t=n\}=\{S_n\le t,\ T_{n+1}>t-S_n\}$.
Ściślej, na zdarzeniu $\{N_t=n\}$ ślad
$\calF_t^N$ jest równy śladowi
$\sigma(T_1,\ldots,T_n)$: każda obserwacja $N_s$,
$s\le t$, jest tam funkcją $S_1,\ldots,S_n$.
Odwrotnie, chwile $S_k$ dla $k\le n$ odzyskujemy
jako kresy dolne chwil osiągnięcia poziomu $k$,
sprawdzając chwile wymierne do $t$ oraz samą chwilę $t$.
Dla $n=0$ informacja na tym zdarzeniu jest trywialna.

Dla ograniczonej borelowskiej $\psi:(0,\infty)\to\R$
definiujemy borelowską funkcję
\[
g_\psi(a)=\frac{\int_{(a,\infty)}\psi(v-a)\,\mathrm d\nu(v)}
{\overline F(a)},\qquad \overline F(a)>0.
\]
Mierzalność licznika wynika z twierdzenia Fubiniego
dla parametru $a$. Na zdarzeniu $\{N_t=n\}$
warunkowe prawo $R_t$ spełnia
$\E[\psi(R_t)\mid\calF_t^N]=g_\psi(A_t)$
prawie na pewno. Aby zweryfikować tę równość,
dla ograniczonej borelowskiej funkcji $H$ wektora
$(T_1,\ldots,T_n)$ całkujemy po jego prawie:
\[
\E[H\chi_{\{N_t=n\}}\psi(R_t)]
=\E\left[H\chi_{\{S_n\le t\}}
\int_{(t-S_n,\infty)}\psi(v-(t-S_n))\,\mathrm d\nu(v)\right].
\]
Zastąpienie całki przez
$\overline F(t-S_n)g_\psi(t-S_n)$ daje
$\E[H\chi_{\{N_t=n\}}g_\psi(A_t)]$.
Równość śladów $\sigma$-ciał i sumowanie po $n$
uzasadniają warunkowanie względem całej historii.
Nie dzielimy przez prawdopodobieństwo zdarzenia
$\{A_t=a\}$.
Na zdarzeniu $\{N_t=n\}$ przyszłe odstępy
$T_{n+2},T_{n+3},\ldots$ zachowują prawo
produktowe $\nu^{\otimes\N}$ i są warunkowo niezależne
od $R_t$: powyższy rachunek powtarzamy z dodatkowym
czynnikiem będącym ograniczoną funkcją skończonej
liczby przyszłych odstępów. Niezależność pozwala
wyłączyć jej wartość oczekiwaną przed całkę;
argument $\pi$--$\lambda$ rozszerza wynik na cały ciąg. Po upływie resztkowego
czasu kolejna liczba odnowień oraz wiek procesu są zatem
jednoznacznymi funkcjami tych niezależnych odstępów i
aktualnej pary $(n,a)$. Dla dowolnej ograniczonej funkcji
borelowskiej $h$ na $\N_0\times[0,\infty)$ istnieje więc
funkcja $P_uh(n,a)$, która nie zależy od chwili $t$ i spełnia
$\E[h(Y_{t+u})\mid\calF_t^N]=P_uh(Y_t)$.
To jest żądana jednorodna własność Markowa.
Co istotne, nie wystarczy znać tylko rozkładu $R_t$:
po następnej odnowie całą przyszłość wyznaczają nowe,
niezależne odstępy. Właśnie ta niezależność pozwala
przejść od opisu pierwszego kolejnego skoku do dowolnej
liczby dalszych odnowień.
W szczególności, jeżeli do chwili $t+u$ nie ma nowej odnowy,
to $Y_{t+u}=(n,a+u)$; prawdopodobieństwo tego zdarzenia
jest równe $\overline F(a+u)/\overline F(a)$ przy
$\overline F(a)>0$. Zdarzenie $\{\overline F(A_t)=0\}$
ma jednak prawdopodobieństwo zero: dla każdego $n$ z niezależności
$T_{n+1}$ od $(T_1,\ldots,T_n)$ wynika
\[
\begin{aligned}
&\Prob\bigl(N_t=n,\,\overline F(A_t)=0\bigr)\\
&=\E\bigl[\chi_{\{S_n\le t,\,\overline F(t-S_n)=0\}}
     \Prob(T_{n+1}>t-S_n\mid T_1,\ldots,T_n)\bigr]=0.
\end{aligned}
\]
Sumując po $n$, uzasadniamy, że iloraz wystarcza na całym
zdarzeniu pełnej miary. Naturalną przestrzenią stanów jest
$E=\N_0\times\{a\ge0:\overline F(a)>0\}$.
Dla każdego jej punktu jądro pierwszego odstępu
jest dane powyższą całką, a dalsze odstępy mają
prawo produktowe $\nu^{\otimes\N}$.
Liczbę nowych odnowień i wiek w chwili $u$
otrzymujemy jako mierzalne funkcje tych odstępów
(przez zdarzenia wyznaczone sumami częściowymi).
Całkowanie względem tak skonstruowanego jądra daje
borelowską funkcję $P_uh$.
Po pierwszym skoku wiek zeruje się; jeśli skoku nie
ma, warunkowanie po przeżyciu dodatkowego czasu
$v$ daje ogon $\overline F(a+v+u)/\overline F(a+v)$.
Ten sam argument odnowy zastosowany do procesu
startującego w dowolnym $(n,a)\in E$ daje
$P_vP_u=P_{v+u}$.
\end{proof}

\begin{zadanie}[Dlaczego sam licznik nie wystarcza?]
Niech odstępy między odnowieniami mają rozkład jednostajny
na $[1,2]$. Ustalmy $t=2$ i rozważmy dwa możliwe przebiegi
z $N_t=1$: w pierwszym ostatnia odnowa nastąpiła w chwili
$t-1/4$, w drugim w chwili $t-3/4$. Porównać prawdopodobieństwa
braku odnowy w następnych $1/2$ jednostki czasu.

\emph{Rozwiązanie.} Dla wieku $a$ ogon pozostałego czasu
wynosi $\overline F(a+1/2)/\overline F(a)$.
Przy $a=1/4$ mamy $\overline F(3/4)=\overline F(1/4)=1$,
więc szansa braku odnowy jest $1$. Przy $a=3/4$ mamy
$\overline F(5/4)=3/4$ i $\overline F(3/4)=1$, zatem
szansa wynosi $3/4$. Zdarzenia warunkujące dokładnym
wiekiem mogą mieć prawdopodobieństwo zero; obliczenie
dotyczy regularnej wersji prawa warunkowego. Właśnie z tego
powodu para $(N_t,A_t)$ przenosi więcej informacji niż $N_t$.
\end{zadanie}

\section{Równanie odnowy i jego jednoznaczność}

Dla lokalnie ograniczonej funkcji borelowskiej $g$
na $[0,\infty)$ i lokalnie skończonej miary dodatniej
$V$ na tej półprostej piszemy
$(g*V)(t)=\int_{[0,t]}g(t-s)\,\mathrm dV(s)$.
Miara odnowy
\[
U=\sum_{n=0}^\infty \nu^{*n}
\]
obejmuje atom $\delta_0$ odpowiadający chwili początkowej.
Jest lokalnie skończona, gdyż $U([0,t])=1+\E N_t$.
Przeliczalna suma miar jest miarą z twierdzenia Tonellego.

\begin{twierdzenie}[Ogólne równanie odnowy]
\label{thm:renewal-equation}
Dla każdej lokalnie ograniczonej funkcji borelowskiej $g$
równanie
\[
u=g+u*\nu
\]
ma dokładnie jedno lokalnie ograniczone rozwiązanie
borelowskie, dane przez $u=g*U$.
\end{twierdzenie}
\begin{proof}
Dla $t\le T$ szereg
$\sum_{n\ge0}(g*\nu^{*n})(t)$ jest bezwzględnie
i jednostajnie zbieżny, gdyż suma oszacowań modułów
nie przekracza
$\sup_{[0,T]}|g|\,U([0,T])<\infty$.
Każdy składnik jest borelowski na mocy twierdzenia
Fubiniego. Z łączności splotu (Fubini na ograniczonym
trójkącie sumowania) i bezwzględnej zbieżności
otrzymujemy $(g*U)*\nu=g*(U-\delta_0)$,
co dowodzi istnienia.
Różnica $w$ dwóch rozwiązań spełnia
$w=w*\nu=w*\nu^{*n}$ dla każdego $n$.
Stąd $|w(t)|\le\sup_{[0,t]}|w|\,\Prob(S_n\le t)\to0$.
Dowolność $t$ daje jednoznaczność.
\end{proof}

\begin{twierdzenie}[Równanie dla średniej liczby odnowień]
Niech $N$ będzie procesem odnowy o dystrybuancie
odstępów $F$. Funkcja $m(t)=\E N_t$ jest skończona
dla $t\ge0$ i spełnia równanie
\[
m(t)=F(t)+\int_{(0,t]}m(t-s)\,\mathrm d\nu(s),\qquad t\ge0.
\]
Wśród lokalnie ograniczonych funkcji borelowskich
$u:[0,\infty)\to\R$ jest to jedyne rozwiązanie tego równania.
\end{twierdzenie}
\begin{proof}
Pierwsza odnowa następuje w chwili $T_1$. Jeżeli $T_1>t$,
liczba odnowień jest zero. Jeśli $T_1=s\le t$,
to pierwszą odnowę już zaliczamy, a kolejna liczba
odnowień ma rozkład $N'_{t-s}$, gdzie proces $N'$ jest
zbudowany z niezależnych odstępów $T_2,T_3,\ldots$
i ma to samo prawo co $N$. Stąd, warunkując względem
$T_1$, otrzymujemy
$m(t)=\E[\chi_{\{T_1\le t\}}(1+m((t-T_1)_+))]
=F(t)+\int_{(0,t]}m(t-s)\,\mathrm d\nu(s)$.
Dla ścisłości w drugim składniku używamy niezależności
$T_1$ od całego ciągu przyszłych odstępów, a nie tylko
równości rozkładów poszczególnych $T_n$.

Załóżmy, że $u_1,u_2$ są lokalnie ograniczonymi rozwiązaniami.
Dla $d=u_1-u_2$ mamy $d=d*\nu$,
a przez iterację $d=d*\nu^{*n}$.
Funkcja $m$ jest lokalnie ograniczona i borelowska,
ponieważ jest niemalejąca i skończona.
Można więc również zastosować
twierdzenie~\ref{thm:renewal-equation}.
Dla ustalonego $t$ niech $M_t=\sup_{0\le s\le t}|d(s)|<\infty$.
Wówczas
$|d(t)|\le M_t F^{*n}(t)=M_t\Prob(S_n\le t)\to0$,
bo $S_n\to\infty$ p.n. Stąd $d(t)=0$ i dowolność $t$
kończy dowód.
\end{proof}

\begin{przyklad}[Proces Poissona jako rozwiązanie równania odnowy]
Jeżeli $T_1\sim\mathrm{Exp}(\lambda)$, to $m(t)=\lambda t$.
Możemy sprawdzić ten wynik bez odwoływania się do rozkładu Poissona:
$F(t)=1-e^{-\lambda t}$, zatem prawa strona równania wynosi
\[
1-e^{-\lambda t}
+\int_0^t\lambda(t-s)\lambda e^{-\lambda s}\,\mathrm d\lambda_1(s)
=1-e^{-\lambda t}+\lambda t-(1-e^{-\lambda t})
=\lambda t.
\]
Z jednoznaczności rozwiązania równania odnowy otrzymujemy
$\E N_t=\lambda t$.
\end{przyklad}

\section{Mocne i elementarne twierdzenie odnowy}
\begin{twierdzenie}[Mocne prawo dla liczby odnowień]
Niech $N$ będzie procesem odnowy z $T_n>0$ p.n.
i $\mu=\E T_1\in(0,\infty]$.
Wtedy
\[
\frac{N_t}{t}\longrightarrow
\begin{cases}1/\mu,&\mu<\infty,\\0,&\mu=\infty,\end{cases}
\qquad\text{prawie na pewno, gdy }t\to\infty.
\]
\end{twierdzenie}
\begin{proof}
Najpierw przyjmijmy $0<\mu<\infty$.
Z MPWL $S_n/n\to\mu$ prawie na pewno.
Na zdarzeniu tej zbieżności $N_t\to\infty$ i
$S_{N_t}\le t<S_{N_t+1}$. Po podzieleniu nierówności
przez $N_t$ otrzymujemy
\[
\frac{S_{N_t}}{N_t}\le \frac{t}{N_t}
<\frac{N_t+1}{N_t}\frac{S_{N_t+1}}{N_t+1}.
\]
Obie skrajne wielkości dążą do $\mu$, więc
$t/N_t\to\mu$, a zatem $N_t/t\to1/\mu$.
Jeśli $\mu=\infty$, dla dowolnego $K>0$ zmienne
$T_n^{(K)}=\min(T_n,K)$ są niezależne, jednakowo rozłożone
i całkowalne. MPWL daje
$\liminf_{n\to\infty}S_n/n\ge\E T_1^{(K)}$.
Względem przeliczalnego ciągu $K\in\N$ i po przejściu
$K\to\infty$ twierdzenie o zbieżności monotonicznej
daje $S_n/n\to\infty$ prawie na pewno.
Ponownie $S_{N_t}\le t$ i $N_t\to\infty$, więc
$t/N_t\ge S_{N_t}/N_t\to\infty$; stąd $N_t/t\to0$.
\end{proof}

\begin{lemat}[Tożsamość Walda dla chwili przekroczenia]
Niech $T_1,T_2,\ldots$ będą niezależne i jednakowo rozłożone,
$T_n>0$ prawie na pewno, z $\mu=\E T_1<\infty$.
Dla $t\ge0$ połóżmy $K_t=N_t+1=\inf\{n\ge1:S_n>t\}$.
Wówczas $\E K_t<\infty$ i
$\E S_{K_t}=\mu\E K_t$.
\end{lemat}
\begin{proof}
Skończoność $\E K_t$ wynika z wcześniej udowodnionego
oszacowania $\E N_t<\infty$. Zauważmy, że
$\{K_t\ge i\}=\{S_{i-1}\le t\}$ zależy tylko od
$T_1,\ldots,T_{i-1}$, zatem jest niezależne od $T_i$.
Ponieważ wszystkie składniki są nieujemne, twierdzenie Tonellego
pozwala zapisać
\[
\E S_{K_t}=\E\sum_{i=1}^\infty T_i\chi_{\{K_t\ge i\}}
=\sum_{i=1}^\infty\mu\Prob(K_t\ge i)
=\mu\E K_t.
\]
Nie używamy tutaj błędnej zamiany $K_t$ na $N_t$:
liczba $N_t$ na ogół nie jest czasem zatrzymania
względem $\mathcal G_i=\sigma(T_1,\ldots,T_i)$,
gdyż $\{N_t\le i\}=\{S_{i+1}>t\}$ może zależeć
od niezaobserwowanego jeszcze $T_{i+1}$.
Na przykład dla odstępów $1$ lub $2$ z dodatnimi
prawdopodobieństwami i $t=3/2$, zdarzenie
$\{N_t=0\}=\{T_1=2\}$ nie należy do trywialnej
$\mathcal G_0$.
\end{proof}

\begin{twierdzenie}[Elementarne twierdzenie odnowy]
Przy założeniach definicji procesu odnowy niech
$\mu=\E T_1\in(0,\infty]$. Wtedy
\[
\frac{\E N_t}{t}\longrightarrow
\begin{cases}1/\mu,&\mu<\infty,\\0,&\mu=\infty,\end{cases}
\qquad\text{gdy }t\to\infty.
\]
\end{twierdzenie}
\begin{proof}
Dla $\mu<\infty$ niech $K_t=N_t+1$.
Ponieważ $S_{K_t}>t$, z lematu Walda dostajemy
$\mu(\E N_t+1)=\E S_{K_t}\ge t$,
czyli $\E N_t/t\ge1/\mu-1/t$.

Dla oszacowania od góry ustalmy $L>0$ i połóżmy
$T_i^{(L)}=\min(T_i,L)$, $\mu_L=\E T_1^{(L)}>0$.
Odnowienia w modelu z obciętymi odstępami występują
nie później, więc $N_t\le N_t^{(L)}$.
W tym modelu chwila pierwszego przekroczenia
$K_t^{(L)}=N_t^{(L)}+1$ spełnia
$t<S_{K_t^{(L)}}^{(L)}\le t+L$.
Tożsamość Walda (stosuje się do obciętych odstępów) daje
$\mu_L(\E N_t^{(L)}+1)\le t+L$. Wobec tego
\[
\limsup_{t\to\infty}\frac{\E N_t}{t}\le\frac{1}{\mu_L}.
\]
Przejście $L\to\infty$ i zbieżność monotoniczna
$\mu_L\uparrow\mu$ dają oszacowanie od góry $1/\mu$,
które wraz z oszacowaniem dolnym kończy dowód.
Jeśli $\mu=\infty$, stosujemy wyłącznie ostatnie
oszacowanie i otrzymujemy
$0\le\limsup_{t\to\infty}\E N_t/t\le1/\mu_L\to0$.
\end{proof}

\begin{uwaga}[Dlaczego silne prawo nie wystarcza?]
Zbieżność $N_t/t\to1/\mu$ prawie na pewno nie uprawnia
sama przez się do zamiany granicy i wartości oczekiwanej.
W ogólności potrzebna jest jednostajna całkowalność
albo odrębne oszacowania. W powyższym dowodzie rolę takich
oszacowań pełnią lemat Walda i obcięcie odstępów.
\end{uwaga}

\section{Twierdzenie Blackwella i kluczowe twierdzenie odnowy}

Elementarne twierdzenie odnowy opisuje średnią na
przedziale $[0,t]$, lecz nie mówi, co dzieje się
na przesuwanym przedziale o ustalonej długości.
Wprowadzamy rozróżnienie zależne od nośnika prawa
odstępu, a nie od tego, czy prawo ma gęstość.

\begin{definicja}[Rozkład arytmetyczny]
Prawo $\nu$ na $(0,\infty)$ nazywamy arytmetycznym,
jeżeli istnieje $h>0$ takie, że
$\nu(h\N)=1$. Największe takie $h$ nazywamy krokiem
arytmetycznym. W przeciwnym razie prawo jest
niearytmetyczne. Równoważnie, domknięta podgrupa
addytywna $\R$ generowana przez nośnik $\nu$
jest odpowiednio $h\Z$ lub $\R$.
\end{definicja}
Istotnie, każda domknięta niezerowa podgrupa
addytywna $\R$ ma tę postać: jeśli kres dolny
jej dodatnich elementów jest zerem, wielokrotności
dowolnie małych dodatnich elementów tworzą dowolnie
gęste siatki, więc domknięcie jest całą prostą.
Jeśli kres dolny jest dodatni, należy do podgrupy
przez domkniętość, a dzielenie z resztą dowodzi,
że wszystkie elementy są jego wielokrotnościami.

\begin{lemat}[Jednostajne ograniczenie miary odnowy]
\label{lem:renewal-local-bound}
Dla każdego $a>0$ istnieje skończona stała $C_a$
taka, że $U([x,x+a])\le C_a$ dla wszystkich $x\in\R$,
gdzie $U$ przedłużamy przez zero na $(-\infty,0)$.
\end{lemat}
\begin{proof}
Dla $x>0$ niech $J=\inf\{n\ge0:S_n\ge x\}$.
Jest to skończony czas zatrzymania względem
$\sigma(T_1,\ldots,T_n)$.
Jeśli $S_J\le x+a$, odnowień w $[x,x+a]$ jest
nie więcej niż $1+N'_a$, gdzie $N'$ jest procesem
zbudowanym z odstępów po indeksie $J$.
Warunkując na zdarzeniach $\{J=j\}$, które zależą
tylko od pierwszych $j$ odstępów, widzimy,
że przyszłe odstępy mają to samo prawo produktowe
co pierwotny ciąg. Zatem wartość oczekiwana
liczby odnowień jest nie większa niż $1+\E N_a$.
Dla $x\le0$ to samo ograniczenie wynika bezpośrednio
z $[x,x+a]\cap[0,\infty)\subseteq[0,a]$.
\end{proof}

\begin{twierdzenie}[Twierdzenie Blackwella]
\label{thm:blackwell-renewal}
Niech $\mu=\E T_1\in(0,\infty]$ i przyjmijmy
$1/\infty=0$.
Jeżeli $\nu$ jest niearytmetyczne, to dla każdego $a>0$
\[
\lim_{t\to\infty}U((t,t+a])=\frac a\mu.
\]
Jeżeli $\nu$ jest arytmetyczne o kroku $h$, to
\[
\lim_{n\to\infty}U(\{nh\})=\frac h\mu.
\]
\end{twierdzenie}
\begin{proof}
Podamy dowód metodą przesuwanych miar.
Z lematu~\ref{lem:renewal-local-bound} miary
$U_t(D)=U(t+D)$ mają jednostajnie ograniczoną
masę na każdym ograniczonym przedziale.
Z każdego ciągu $t_j\to\infty$ można wybrać
podciąg zbieżny mgliście (tj. w sensie zbieżności
całek wszystkich funkcji ciągłych o zwartym nośniku) do pewnej
lokalnie skończonej miary $V$.
Dla ścisłości wybieramy przeliczalny zbiór
funkcji gęsty w normie supremum w każdej
przestrzeni funkcji ciągłych o nośniku w $[-m,m]$,
$m\in\N$. Argument diagonalny daje zbieżność
całek na tym zbiorze, a lokalne ograniczenia
mas rozszerzają ją na całe $C_c(\R)$.
Dodatni funkcjonał graniczny reprezentuje
miara Radona na mocy twierdzenia Riesza--Markowa.
Granica $V$ również ma jednostajnie ograniczone
masy przedziałów jednostkowych: ograniczamy
ich funkcje charakterystyczne od góry przez
ciągłe funkcje o wartościach w $[0,1]$ i nośnikach
w przedziałach o długości $3$.

Z równania miar $U=\delta_0+U*\nu$ otrzymujemy
$V=V*\nu$. Rzeczywiście, dla $f\in C_c(\R)$
\[
\int f\,\mathrm dU_{t_j}
=f(-t_j)+\int_{(0,\infty)}
\left(\int f(z+s)\,\mathrm dU_{t_j}(z)\right)\mathrm d\nu(s).
\]
Wewnętrzne całki są jednostajnie ograniczone
niezależnie od $s,j$ dzięki lematowi
\ref{lem:renewal-local-bound}. Zbieżność mglista
dla ustalonego $s$ i zbieżność zmajoryzowana
dają równanie dla $V$.

Ustalmy $f\in C_c(\R)$ i połóżmy
$H(x)=\int f(x-y)\,\mathrm dV(y)$.
Jest to ograniczona ciągła funkcja
(na mocy jednostajnej ciągłości $f$ i
jednostajnych ograniczeń mas).
Z $V=V*\nu$ wynika
$H(x)=\int H(x-s)\,\mathrm d\nu(s)$.
Dla ustalonego $x$ ciąg $H(x-S_n)$ jest
ograniczonym martyngałem względem odstępów.
Ma więc granicę $L_x$ prawie na pewno i w $L^1$
(z twierdzenia o zbieżności martyngałów ograniczonych).
Granica jest niezmiennicza względem skończonych
permutacji odstępów, ponieważ od pewnego indeksu
takie permutacje nie zmieniają sum częściowych.

Przypomnijmy krótki argument zero-jedynkowy.
Jeżeli zdarzenie $A$ zależne od niezależnych,
jednakowo rozłożonych zmiennych jest niezmiennicze
względem skończonych permutacji, to dla każdego
$\varepsilon>0$ można przybliżyć je zdarzeniem $D$
zależnym od pierwszych $m$ zmiennych tak, że
$\Prob(A\mathbin{\triangle}D)<\varepsilon$.
Istnienie takiego przybliżenia wynika z faktu,
że zdarzenia przybliżalne w mierze tworzą
$\sigma$-ciało zawierające cylindry (dla sumy
przeliczalnej najpierw bierzemy sumę skończoną).
Zamieniając pierwsze $m$ zmiennych z następnymi
$m$, dostajemy niezależną kopię $D'$ zdarzenia $D$
i nadal $\Prob(A\mathbin{\triangle}D')<\varepsilon$.
Stąd $|\Prob(A)-\Prob(D\cap D')|<2\varepsilon$,
a $\Prob(D\cap D')=\Prob(D)^2$.
Przejście $\varepsilon\downarrow0$ daje
$\Prob(A)=\Prob(A)^2$.
Stosując ten argument do zdarzeń
$\{L_x\le q\}$ dla wymiernych $q$, widzimy,
że $L_x$ jest stała p.n.
Możemy pracować na kanonicznej przestrzeni
produktowej odstępów i zdefiniować granicę
przez $\limsup$, dzięki czemu niezmienniczość
jest punktowa. Ponieważ $\E L_x=H(x)$,
z własności martyngału i zbieżności w $L^1$
wynika $H(x-T_1)=\E[L_x\mid T_1]=H(x)$ p.n.
Ciągłość $H$ daje $H(x-s)=H(x)$ dla każdego
$s$ z nośnika $\nu$, dla każdego ustalonego $x$.
W konsekwencji każdy element domkniętej
podgrupy generowanej przez nośnik jest okresem $H$.
Ponieważ dotyczy to każdej $f\in C_c(\R)$,
miara $V$ jest niezmiennicza względem przesunięć
o elementy tej podgrupy.

W przypadku niearytmetycznym oznacza to
$V=c\lambda_1$ dla $c\ge0$. Istotnie,
niezmienniczość wymusza brak atomów
(lokalna skończoność), a masa przedziału
jest addytywną niemalejącą funkcją jego długości,
zatem jest proporcjonalna do długości.
W przypadku arytmetycznym rozważamy wyłącznie
$t_j=n_jh$. Miary $U_{t_j}$ i ich granica
są skupione na domkniętym zbiorze $h\Z$,
a niezmienniczość względem $h$ daje
$V=c\sum_{k\in\Z}\delta_{kh}$.

Pozostaje ustalić $c$. Dokładnie jeden cykl
zawiera chwilę $t$, więc niezależność
$T_{n+1}$ od $S_n$ daje tożsamość
\[
1=\sum_{n\ge0}\Prob(S_n\le t<S_{n+1})
=\int_{[0,t]}\overline F(t-s)\,\mathrm dU(s).
\]
W zmiennej przesuniętej $z=s-t$ jest to
całka $\overline F(-z)$ po $[-t,0]$
względem $U_t$. Dla $\mu<\infty$
ogon tej całki po $z<-M$ jest jednostajnie
mały, ponieważ z lokalnego ograniczenia
jest nie większy niż
$C_1\sum_{k\ge\lfloor M\rfloor}\overline F(k)$,
a szereg jest zbieżny z
$\int_0^\infty\overline F(v)\,\mathrm d\lambda_1(v)=\mu$
(tożsamość wynika z twierdzenia Tonellego zastosowanego
do $T_1=\int_0^\infty\chi_{\{v<T_1\}}\,\mathrm d\lambda_1(v)$).
W przypadku niearytmetycznym na każdym
$[-M,0]$ możemy przejść do granicy:
funkcja monotoniczna $\overline F(-z)$ ma
co najwyżej przeliczalnie wiele nieciągłości,
które mają miarę $V$ zero. Uzasadnienie
dają górne i dolne sumy schodkowe na podziałach
o końcach bez atomów $V$ oraz zbieżność
mas tych przedziałów wynikająca z mglistej
zbieżności przez aproksymację ciągłą.
Otrzymujemy $1=c\mu$.
Gdy $\mu=\infty$, używamy wyłącznie obciętych
całek: $1\ge c\int_0^M\overline F(v)\,\mathrm d\lambda_1(v)$,
więc $c=0$.

W przypadku arytmetycznym na każdym skończonym
zbiorze punktów siatki masy poszczególnych punktów
zbiegają do $c$:
każdy punkt można izolować funkcją ciągłą
o nośniku w przedziale długości mniejszej niż $h$.
Analogicznie dla $\mu<\infty$ dostajemy
$1=c\sum_{k\ge0}\overline F(kh)=c\mu/h$,
gdzie ostatnia równość wynika z wzoru na średnią
dodatniej całkowitoliczbowej zmiennej $T_1/h$.
Przy nieskończonej średniej skończone sumy
wymuszają $c=0$.

Każda granica podciągowa jest więc taka sama.
W przypadku niearytmetycznym końce $(0,a]$
nie mają masy granicznej, więc zbieżność mglista
daje $U_{t_j}((0,a])\to a/\mu$.
W przypadku arytmetycznym izolujemy punkt zero
i dostajemy $U_{n_jh}(\{0\})\to h/\mu$.
Argument podciągowy kończy dowód.
\end{proof}

\begin{definicja}[Bezpośrednia całkowalność Riemanna]
Nieujemną funkcję $g:[0,\infty)\to\R$ nazywamy
bezpośrednio całkowalną w sensie Riemanna, jeżeli
dla przedziałów $I_{k,\delta}=[k\delta,(k+1)\delta)$
górne sumy
$\delta\sum_{k\ge0}\sup_{I_{k,\delta}}g$
są skończone dla wszystkich dostatecznie małych
$\delta>0$, a różnica górnej i dolnej sumy
dąży do zera przy $\delta\downarrow0$.
Dla funkcji ze znakiem wymagamy tej własności
od części dodatniej i ujemnej.
\end{definicja}

\begin{twierdzenie}[Kluczowe twierdzenie odnowy]
\label{thm:key-renewal}
Jeżeli $\nu$ jest niearytmetyczne, a $g$ jest
borelowska i bezpośrednio całkowalna w sensie
Riemanna, to
\[
\lim_{t\to\infty}(g*U)(t)
=\frac1\mu\int_0^\infty g(v)\,\mathrm d\lambda_1(v).
\]
W przypadku arytmetycznym o kroku $h$,
jeśli $\sum_{k\ge0}|g(kh)|<\infty$, to
\[
\lim_{n\to\infty}(g*U)(nh)
=\frac h\mu\sum_{k\ge0}g(kh).
\]
\end{twierdzenie}
\begin{proof}
Dla $g\ge0$ w przypadku niearytmetycznym
ustalmy $\delta>0$ z definicji.
Rozbijamy całkę względem $U_t$ na przedziały
$(- (k+1)\delta,-k\delta]$ i ograniczamy ją
górnymi oraz dolnymi wartościami $g$ na
$I_{k,\delta}$. Miara każdego takiego przedziału
dąży do $\delta/\mu$ na mocy twierdzenia
\ref{thm:blackwell-renewal}, a jest jednostajnie
ograniczona przez $C_\delta$ z lematu
\ref{lem:renewal-local-bound}. Przedłużenie
$U_t$ przez zero poza jego nośnikiem pozwala
sumować po wszystkich $k\ge0$.
Sumowalność górnych wartości daje zbieżność
sum przez zmajoryzowanie szeregu.
Dolna i górna granica splotu leżą więc między
dolną i górną sumą Riemanna podzielonymi przez
$\mu$. Przejście $\delta\downarrow0$ daje
wskazaną całkę; gdy $\mu=\infty$, już dla
jednego $\delta$ górna granica wynosi zero.
Dla funkcji ze znakiem stosujemy dowód do obu części.
W przypadku arytmetycznym splot w chwili $nh$
jest sumą $\sum_{k=0}^n g(kh)U(\{(n-k)h\})$.
Każdy czynnik miary zbiega do $h/\mu$,
a lokalne ograniczenie i bezwzględna sumowalność
pozwalają przejść do granicy w szeregu.
\end{proof}

\begin{przyklad}[Dlaczego arytmetyczność ma znaczenie?]
Dla $T_n=h$ punktowo miara odnowy jest
$\sum_{k\ge0}\delta_{kh}$.
Jeśli $0<a<h$, to $U((t,t+a])$ przyjmuje
wartości $0$ i $1$ w dowolnie późnych chwilach,
więc nie może zbiegać do $a/h$.
Natomiast $U(\{nh\})=1=h/\mu$.
Prawo o dwóch atomach $1,\sqrt2$ jest
niearytmetyczne, mimo że nie ma gęstości:
żadna dodatnia liczba nie jest wspólnym krokiem
dla obu tych liczb, ponieważ ich stosunek
jest niewymierny.
\end{przyklad}

\section{Wiek odnowy i czas pozostały do następnego zdarzenia}
\begin{definicja}[Wiek i czas resztkowy]
Dla $t\ge0$ oraz procesu odnowy o chwilach $S_n$ definiujemy
$A_t:=t-S_{N_t}$ oraz $R_t:=S_{N_t+1}-t$.
Nazywamy je odpowiednio \emph{wiekiem bieżącego cyklu}
i \emph{pozostałym czasem do następnego odnowienia}.
\end{definicja}

\begin{stwierdzenie}[Warunkowy rozkład czasu resztkowego]
Niech $\nu$ będzie prawem ściśle dodatnich odstępów,
a $F$ jego dystrybuantą.
Ustalmy $t\ge0$, $n\ge0$ oraz $s\in[0,t]$ takie, że
$\overline F(t-s)>0$. Wtedy możemy wybrać regularną wersję rozkładu
warunkowego względem $(N_t,S_n)$, która dla
warunków $N_t=n$, $S_n=s$ spełnia dla $u\ge0$
\[
\Prob(R_t>u\mid N_t=n,\ S_n=s)
=\frac{\overline F(t-s+u)}{\overline F(t-s)}.
\]
\end{stwierdzenie}
\begin{proof}
Na zdarzeniu $\{S_n\le t\}$ warunek $N_t=n$
jest równoważny $T_{n+1}>t-S_n$.
Zmienna $T_{n+1}$ jest niezależna od $S_n$, gdyż
$S_n$ zależy tylko od poprzednich odstępów.
Warunkujemy zatem zmienną $T_{n+1}$ na przekroczenie
progu $t-s$ i dostajemy iloraz
$\Prob(T_{n+1}>t-s+u)/\Prob(T_{n+1}>t-s)$.
Formalnie dla dowolnego zbioru borelowskiego $D$
\[
\Prob(N_t=n,S_n\in D,R_t>u)
=\int_{D\cap[0,t]}\overline F(t-s+u)\,\mathrm d\nu^{*n}(s).
\]
Z kolei miara rozkładu $S_n$ ograniczona
do zdarzenia $\{N_t=n\}$ wynosi
$\chi_{[0,t]}(s)\overline F(t-s)\,\mathrm d\nu^{*n}(s)$.
Całkowanie podanego ilorazu względem tej miary
daje dokładnie poprzednią równość, także dla $n=0$.
Jest to uzasadnienie regularnej wersji warunkowej
z twierdzenia~\ref{thm:wiek-markow}.
\end{proof}

\begin{przyklad}[Dwie zupełnie różne pamięci]
Dla odstępów wykładniczych $\overline F(v)=e^{-\lambda v}$,
więc iloraz wynosi $e^{-\lambda u}$, niezależnie od wieku
$t-s$. Stąd proces Poissona zapomina czas ostatniego zgłoszenia.
Dla deterministycznych odstępów $T_i\equiv c$ pozostały czas
jest natomiast równy $c-A_t$ we wszystkich chwilach,
w tym w chwilach odnowień, gdy $A_t=0$ i $R_t=c$. Znając wiek, dokładnie znamy czas kolejnej
odnowy. Nie należy więc oczekiwać własności Markowa dla
samego procesu $N_t$ w dowolnym procesie odnowy; często
naturalnym stanem jest para $(N_t,A_t)$.
\end{przyklad}

\section{Graniczny wiek, czas resztkowy i paradoks inspekcji}

\begin{twierdzenie}[Graniczny rozkład wieku i czasu resztkowego]
\label{thm:renewal-age-limit}
Załóżmy, że $\mu=\E T_1<\infty$ i prawo $\nu$
jest niearytmetyczne. Wówczas dla $a,b\ge0$
\[
\lim_{t\to\infty}\Prob(A_t\ge a,\ R_t>b)
=\frac1\mu\int_{a+b}^\infty\overline F(v)\,
\mathrm d\lambda_1(v).
\]
Para $(A_t,R_t)$ zbiega według rozkładu do pary
$(A_*,R_*)$, której obie współrzędne mają
rozkład równowagowy o gęstości
$\overline F(v)/\mu$ na $[0,\infty)$.
\end{twierdzenie}
\begin{proof}
Dla ustalonego $t\ge0$ dokładnie jeden indeks $n$
spełnia $S_n\le t<S_{n+1}$.
Na tym zdarzeniu warunki $A_t\ge a$, $R_t>b$
są równoważne $S_n\le t-a$ i $T_{n+1}>t+b-S_n$.
Z niezależności odstępów, sumowania po $n$ i
definicji $U$ otrzymujemy
\[
\Prob(A_t\ge a,R_t>b)
=\int_{[0,t]}\chi_{[a,\infty)}(t-s)
\overline F(t-s+b)\,\mathrm dU(s).
\]
Funkcja $g(v)=\chi_{[a,\infty)}(v)\overline F(v+b)$
jest bezpośrednio całkowalna w sensie Riemanna:
poza jednym punktem skoku jest niemalejąca
na początkowym przedziale (gdzie jest równa zeru)
i nierosnąca po $a$. Dla każdej siatki górna
suma jest skończona z całkowalności ogona,
a różnica górnej i dolnej sumy jest ograniczona
przez $2\delta$, z oszacowania całkowitej wariacji.
Twierdzenie~\ref{thm:key-renewal} daje żądaną granicę.
Aby zidentyfikować prawo graniczne, wybierzmy
$L$ o prawie
\[
\nu_{\rm len}(D)=\frac1\mu\int_D v\,\mathrm d\nu(v),
\]
oraz niezależną zmienną $Q$ jednostajną na $(0,1)$.
Połóżmy $A_*=QL$ i $R_*=(1-Q)L$.
Warunkując względem $L$ otrzymujemy
\[
\Prob(A_*\ge a,R_*>b)
=\frac1\mu\int_{(a+b,\infty)}(v-a-b)\,\mathrm d\nu(v)
=\frac1\mu\int_{a+b}^\infty\overline F(v)\,\mathrm d\lambda_1(v).
\]
Ostatnia równość wynika z twierdzenia Tonellego.
Dla $a=0$ lub $b=0$ dostajemy rozkłady brzegowe
o ciągłych dystrybuantach. Granice wspólnych
ogonów i ogonów brzegowych wyznaczają granice
dystrybuant łącznych (dla ścisłych nierówności
dotyczących $A_t$ stosujemy oszacowanie przez
progi $a$ i $a+\varepsilon$ oraz ciągłość
rozkładu granicznego). Stąd zbieżność według
rozkładu całej pary.
\end{proof}

\begin{wniosek}[Długość cyklu widzianego w późnej chwili]
Przy założeniach twierdzenia~\ref{thm:renewal-age-limit}
długość bieżącego cyklu $L_t=A_t+R_t=T_{N_t+1}$
zbiega według rozkładu do $L$ o prawie
$\nu_{\rm len}(\mathrm dv)=v\,\nu(\mathrm dv)/\mu$.
Ponadto
\[
\E L=\frac{\E T_1^2}{\mu},\qquad
\E A_*=\E R_*=\frac{\E T_1^2}{2\mu},
\]
z dopuszczeniem nieskończonych wartości.
\end{wniosek}
\begin{proof}
Zbieżność wynika z twierdzenia o odwzorowaniu
ciągłym zastosowanego do sumy współrzędnych.
Wzór na $\E L$ wynika z definicji jego prawa,
a niezależność $Q$ i $L$ oraz $\E Q=1/2$
dają pozostałe wzory przez Tonellego.
Nie twierdzimy na tej podstawie, że wartości
oczekiwane $L_t,A_t,R_t$ zbiegają: do tego
potrzebne jest dodatkowe uzasadnienie.
\end{proof}

\begin{uwaga}[Paradoks inspekcji]
Losowy cykl obserwowany w późnej deterministycznej
chwili nie ma na ogół prawa zwykłego odstępu $\nu$.
Cykl długości $v$ ma wagę proporcjonalną do $v$.
Jeżeli drugi moment odstępu jest skończony, to
\[
\E L=\mu+\frac{\Var(T_1)}{\mu}\ge\mu.
\]
Równość zachodzi jedynie dla odstępów stałych.
Interpretacja przez rozkład długościowo obciążony
obowiązuje również przy nieskończonym drugim
momencie, kiedy średnia obserwowanej długości
jest nieskończona.
\end{uwaga}

\begin{przyklad}[Odstępy jednostajne]
Jeżeli $T_1$ ma rozkład jednostajny na $[1,2]$,
to $\mu=3/2$ i $\E T_1^2=7/3$.
Graniczna długość obserwowanego cyklu ma gęstość
$2v/3$ na $[1,2]$, a jej średnia wynosi $14/9$.
Średni graniczny czas resztkowy wynosi $7/9$,
a nie $\mu/2=3/4$. Jego gęstość jest równa
$2/3$ na $[0,1)$ i $2(2-v)/3$ na $[1,2]$,
a poza tymi przedziałami zero.
\end{przyklad}

\begin{stwierdzenie}[Graniczny rozkład w przypadku arytmetycznym]
Niech $\nu$ będzie arytmetyczne o kroku $h$,
$\mu<\infty$. Ustalmy fazę $\theta\in[0,h)$.
Dla $a,b\ge0$ zachodzi
\[
\lim_{n\to\infty}\Prob(A_{nh+\theta}\ge a,R_{nh+\theta}>b)
=\frac h\mu\sum_{k\ge0}
\chi_{[a,\infty)}(kh+\theta)\overline F(kh+\theta+b).
\]
\end{stwierdzenie}
\begin{proof}
Wykorzystujemy dokładną tożsamość splotową
z dowodu twierdzenia~\ref{thm:renewal-age-limit}.
Wszystkie odnowienia są na siatce $h\N_0$,
więc dla $t=nh+\theta$ otrzymujemy sumę
\[
\sum_{k=0}^n
\chi_{[a,\infty)}(kh+\theta)\overline F(kh+\theta+b)
U(\{(n-k)h\}).
\]
Wagi ogonowe są sumowalne, gdyż
$\sum_{k\ge0}\overline F(kh)=\mu/h$.
Z twierdzenia~\ref{thm:blackwell-renewal}
i jednostajnego ograniczenia mas punktów siatki
wynika teza przez zmajoryzowanie szeregu.
W szczególności nie wolno w ogólności oczekiwać
jednej granicy dla wszystkich rzeczywistych $t$:
dla $T_1=h$ mamy dokładnie $A_{nh+\theta}=\theta$
i $R_{nh+\theta}=h-\theta$.
\end{proof}

\section{Proces odnowy w równowadze}

\begin{definicja}[Opóźniony proces odnowy]
W opóźnionym procesie odnowy pierwsze oczekiwanie
$D>0$ może mieć inne prawo niż następne odstępy.
Zakładamy niezależność $D$ od niezależnych odstępów
$T_1,T_2,\ldots$ o wspólnym prawie $\nu$.
Chwile odnowień wynoszą
$D+S_n$, $n\ge0$, gdzie $S_0=0$.
Gdy $\mu<\infty$, proces nazywamy równowagowym,
jeśli $D$ ma gęstość
$\overline F(v)/\mu$ dla $v\ge0$.
\end{definicja}

\begin{twierdzenie}[Stacjonarność przyrostów procesu równowagowego]
\label{thm:renewal-equilibrium}
Niech $\widetilde N_t$ zlicza odnowienia procesu
równowagowego na $(0,t]$.
Wówczas dla każdego $t,u\ge0$
\[
\E\widetilde N_t=\frac t\mu,\qquad
\Prob(\widetilde R_t>u)=
\frac1\mu\int_u^\infty\overline F(v)\,\mathrm d\lambda_1(v),
\]
gdzie $\widetilde R_t$ jest czasem do najbliższego
odnowienia ściśle po $t$. Ponadto proces
$(\widetilde N_{t+u}-\widetilde N_t)_{u\ge0}$
ma to samo prawo co $(\widetilde N_u)_{u\ge0}$.
Nie zakładamy niearytmetyczności $\nu$.
\end{twierdzenie}
\begin{proof}
Niech $\eta$ będzie prawem $D$. Miara zliczająca
oczekiwane odnowienia opóźnione wynosi
$\widetilde U=\eta*U$, przez niezależność i
sumowanie praw chwil $D+S_n$.
Ma ona gęstość względem miary Lebesgue'a:
z twierdzenia Tonellego
\[
\frac{\mathrm d\widetilde U}{\mathrm d\lambda_1}(t)
=\frac1\mu\int_{[0,t]}\overline F(t-s)\,\mathrm dU(s)
=\frac1\mu.
\]
Ostatnia równość to tożsamość wykorzystana
w dowodzie twierdzenia~\ref{thm:blackwell-renewal};
zachodzi dla każdego $t\ge0$.
Zatem $\E\widetilde N_t=\widetilde U((0,t])=t/\mu$.

Zdarzenie $\{\widetilde R_t>u\}$ zachodzi albo
gdy $D>t+u$, albo gdy jedno z odnowień
$D+S_n\le t$ jest ostatnim przed $t$ i następujący
po nim odstęp przekracza $t+u-D-S_n$.
Zdarzenia te są rozłączne. Z niezależności i
otrzymanej postaci $\widetilde U$ wynika
\[
\begin{aligned}
\Prob(\widetilde R_t>u)
&=\Prob(D>t+u)+
\int_{[0,t]}\overline F(t+u-s)\,\mathrm d\widetilde U(s)\\
&=\frac1\mu\left(
\int_{t+u}^\infty\overline F(v)\,\mathrm d\lambda_1(v)
+\int_u^{t+u}\overline F(v)\,\mathrm d\lambda_1(v)\right).
\end{aligned}
\]
To jest ogon prawa $D$.

Odstępy po pierwszym odnowieniu ściśle po $t$
są niezależne od $\widetilde R_t$ i mają prawo
produktowe $\nu^{\otimes\N}$.
Weryfikujemy to, rozbijając przestrzeń na
$\{D>t\}$ oraz $\{D+S_n\le t<D+S_{n+1}\}$,
$n\ge0$. Na każdym z tych zdarzeń czas resztkowy
i samo zdarzenie zależą tylko od odstępów do
najbliższego odnowienia włącznie; dalszy ciąg
jest niezależny od tych zmiennych i ma ustalone
prawo produktowe. Faktoryzacja dla ograniczonych
funkcji czasu resztkowego i cylindrów dalszych
odstępów, a następnie sumowanie oraz argument
$\pi$--$\lambda$, dowodzą niezależności.
Zatem proces liczby odnowień po chwili $t$
ma tę samą konstrukcję co pierwotny proces
równowagowy. Otrzymujemy równość praw procesów,
w szczególności wszystkich rozkładów
skończeniewymiarowych.
\end{proof}

\begin{uwaga}[Stacjonarne przyrosty a stacjonarność licznika]
Licznik $\widetilde N_t$ nie jest procesem
stacjonarnym: zaczyna się w zerze, a jego średnia
rośnie jak $t/\mu$. Stacjonarne są jego przyrosty,
czyli liczby odnowień na przesuwanych przedziałach
oraz ich wspólne rozkłady. Nie oznacza to
niezależności przyrostów. Dla stałych odstępów $h$
pierwsze oczekiwanie w równowadze jest jednostajne
na $(0,h)$, co losuje fazę siatki odnowień.
\end{uwaga}

\section{Nagrody w cyklach: twierdzenie odnowy z nagrodą}
\begin{twierdzenie}[Długookresowa średnia nagrody]
Niech $(T_n,R_n)_{n\ge1}$ będą niezależnymi parami o jednakowym
rozkładzie w $(0,\infty)\times\R$, gdzie
$0<\mu=\E T_1<\infty$ oraz $\E|R_1|<\infty$.
Nie wymagamy niezależności $T_n$ i $R_n$ wewnątrz pary.
Niech $S_n=\sum_{k=1}^nT_k$, $N_t=\max\{n:S_n\le t\}$ i
$C_t=\sum_{k=1}^{N_t}R_k$. Wówczas
\[
\frac{C_t}{t}\longrightarrow\frac{\E R_1}{\E T_1}
\qquad\text{prawie na pewno, gdy }t\to\infty.
\]
\end{twierdzenie}
\begin{proof}
Na wspólnym zdarzeniu pełnej miary mocne prawo wielkich liczb
dla obu współrzędnych par daje
$S_n/n\to\mu$ oraz $n^{-1}\sum_{k=1}^nR_k\to\rho:=\E R_1$.
Z twierdzenia odnowy $N_t\to\infty$ oraz
$N_t/t\to1/\mu$ prawie na pewno. Wobec tego
\[
\frac{C_t}{t}=
\left(\frac{1}{N_t}\sum_{k=1}^{N_t}R_k\right)
\frac{N_t}{t}\longrightarrow\rho/\mu.
\]
Jeżeli przez pewien początkowy czas $N_t=0$, przyjmujemy
$C_t=0$; nie ma to wpływu na granicę.
\end{proof}

\begin{zadanie}[Udział czasu naprawy maszyny]
Maszyna pracuje przez czas $U_n\sim\mathrm{Exp}(\alpha)$,
a następnie jest naprawiana przez czas
$V_n\sim\mathrm{Exp}(\beta)$. Załóżmy, że wszystkie
zmienne $(U_n,V_n)$ są niezależne pomiędzy cyklami,
a w każdym cyklu czasy pracy i naprawy są niezależne.
Obliczyć długookresowy udział czasu napraw.

\emph{Rozwiązanie.} Czas jednego cyklu wynosi
$T_n=U_n+V_n$, a nagrodą jest długość naprawy $R_n=V_n$.
Zatem $\E T_1=1/\alpha+1/\beta$ i $\E R_1=1/\beta$.
Twierdzenie daje udział napraw w ukończonych cyklach
\[
\frac{\E V_1}{\E(U_1+V_1)}
=\frac{1/\beta}{1/\alpha+1/\beta}
=\frac{\alpha}{\alpha+\beta}.
\]
Aby uzyskać udział rzeczywistego czasu napraw na przedziale
$[0,t]$, należy uwzględnić nieukończony cykl. Różnica
między nagrodą z ukończonych cykli a rzeczywistym czasem
napraw nie przekracza $T_{N_t+1}$. Ponieważ
$S_n/n\to\mu$, mamy
$T_n/n=(S_n/n)-((n-1)/n)(S_{n-1}/(n-1))\to0$ p.n.,
a stąd $T_{N_t+1}/t\to0$ po połączeniu z $N_t/t\to1/\mu$.
Dostajemy ten sam udział dla całego czasu $[0,t]$.
\end{zadanie}


\section{Transformaty procesu odnowy i przykład odstępów gamma}
\begin{twierdzenie}[Transformata miary odnowy]
Niech $S_0=0$ i $S_n=\sum_{k=1}^nT_k$ dla niezależnych,
ściśle dodatnich odstępów o jednakowym rozkładzie.
Zdefiniujmy lokalnie skończoną miarę na $[0,\infty)$
\[
U(A)=\sum_{n=0}^\infty\Prob(S_n\in A),\qquad
A\in\calB([0,\infty)).
\]
Jej funkcją dystrybucyjną jest $U([0,t])=1+\E N_t$.
Jeżeli $\varphi(\theta)=\E e^{-\theta T_1}$, to dla $\theta>0$
\[
\int_{[0,\infty)}e^{-\theta t}\,\mathrm dU(t)
=\frac{1}{1-\varphi(\theta)}.
\]
\end{twierdzenie}
\begin{proof}
Dla $t<\infty$ równość
$U([0,t])=1+\sum_{n\ge1}\Prob(S_n\le t)=1+\E N_t$
i poprzednie oszacowanie $\E N_t<\infty$ dowodzą
lokalnej skończoności miary. Z twierdzenia Tonellego,
nieujemności funkcji podcałkowej i niezależności odstępów
\[
\int e^{-\theta t}\,\mathrm dU(t)
=\sum_{n=0}^\infty\E e^{-\theta S_n}
=\sum_{n=0}^\infty \varphi(\theta)^n.
\]
Ponieważ $T_1>0$ p.n., mamy $0\le\varphi(\theta)<1$.
Suma szeregu geometrycznego daje tezę.
\end{proof}

\begin{zadanie}[Każde drugie zdarzenie Poissona]
Niech $P_t$ będzie procesem Poissona o intensywności $\lambda>0$.
Interesuje nas tylko co drugie zdarzenie: czasem odnowy jest
suma dwóch kolejnych odstępów wykładniczych. Obliczyć
rozkład $N_t$ i $\E N_t$.

\emph{Rozwiązanie.} Kolejne pary niezależnych
odstępów wykładniczych dają niezależne czasy
$T_n\sim\mathrm{Gamma}(2,\lambda)$ (parametry: kształt $2$,
intensywność $\lambda$). Liczba zakończonych par to
$N_t=\lfloor P_t/2\rfloor$. Jeśli $u=\lambda t$, mamy
\[
\Prob(N_t=n)=\Prob(P_t=2n)+\Prob(P_t=2n+1)
=e^{-u}\left(\frac{u^{2n}}{(2n)!}+
\frac{u^{2n+1}}{(2n+1)!}\right).
\]
Ponieważ $\lfloor k/2\rfloor=(k-\chi_{\{k\text{ nieparzyste}\}})/2$,
\[
\E N_t=\frac{\E P_t}{2}-\frac12\Prob(P_t\text{ nieparzyste}).
\]
Z funkcji tworzącej rozkładu Poissona
$\E[(-1)^{P_t}]=e^{-2u}$, skąd
$\Prob(P_t\text{ nieparzyste})=(1-e^{-2u})/2$.
Ostatecznie
\[
\E N_t=\frac{\lambda t}{2}-\frac{1-e^{-2\lambda t}}{4}.
\]
W szczególności $\E N_t/t\to\lambda/2=1/\E T_1$,
zgodnie z elementarnym twierdzeniem odnowy.
Ale $N_t$ nie jest procesem Poissona: odstępy są gamma,
a nie wykładnicze.
\end{zadanie}

\section{Podsumowanie i dalsze pytania}
W procesie odnowy niezależne są całe kolejne cykle, a nie
zwykle liczby zdarzeń w rozłącznych przedziałach. Wzór
odnowy opisuje średnią liczbę zakończonych cykli, a oba
prawa graniczne podają długookresową częstość odnowień.
Zagadnienia o wieku i czasie resztkowym wyjaśniają, dlaczego
sama liczba zdarzeń nie musi wystarczać do uzyskania własności Markowa.
Twierdzenie Blackwella oraz kluczowe twierdzenie
odnowy opisują lokalne granice miary odnowy
i jej splotów, z rozróżnieniem przypadku
arytmetycznego i niearytmetycznego.
Graniczne rozkłady wieku i czasu resztkowego
wyjaśniają paradoks inspekcji, a odpowiedni
dobór pierwszego oczekiwania daje proces
odnowy o stacjonarnych przyrostach.


% ==============================================================
% ROZDZIAŁ 10. PROCESY ROZGAŁĘZIAJĄCE
% ==============================================================

\chapter{Procesy rozgałęziające: losowy wzrost i wymieranie populacji}
\label{ch:branching}

\section{Od niezależnych osobników do zależnych pokoleń}
W procesie odnowy po każdej wymianie mechanizmu zaczyna się nowy
cykl. W procesie rozgałęziającym każdy osobnik produkuje losową
liczbę potomków, którzy w następnym pokoleniu rozmnażają się niezależnie.
Liczba osobników w kolejnym pokoleniu zależy od obecnej liczby,
ale cała genealogia powstaje z niezależnych lokalnych decyzji.
Zasadnicze pytanie brzmi: kiedy populacja wymiera, a kiedy ma
dodatnie prawdopodobieństwo przetrwania przez wszystkie pokolenia?

\begin{definicja}[Proces Galtona--Watsona]
Niech $\xi$ będzie zmienną losową o wartościach w $\N_0$,
z rozkładem $p_k=\Prob(\xi=k)$, $k\ge0$.
Niech $(\xi_{n,i})_{n\ge0,i\ge1}$ będą niezależnymi kopiami
$\xi$ oraz niech $Z_0=z\in\N_0$ będzie deterministyczne.
Definiujemy rekurencyjnie
\[
Z_{n+1}=\sum_{i=1}^{Z_n}\xi_{n,i},\qquad n\ge0,
\]
przyjmując sumę pustą za zero. Proces $(Z_n)_{n\ge0}$
nazywamy \emph{procesem rozgałęziającym Galtona--Watsona}.
Przyjmujemy filtrację
$\calF_n=\sigma(\xi_{k,i}:0\le k<n,\ i\ge1)$,
z trywialnym $\calF_0$.
\end{definicja}

Istnienie tablicy niezależnych zmiennych zapewnia
przeliczalny iloczyn przestrzeni z prawem $(p_k)_{k\ge0}$.
Indukcyjnie $Z_n$ jest skończoną zmienną całkowitoliczbową
i jest $\calF_n$-mierzalna: na zdarzeniu $\{Z_n=k\}$
suma definiująca $Z_{n+1}$ ma dokładnie $k$ składników,
a zdarzenia te tworzą przeliczalny podział przestrzeni.
Wiersz $(\xi_{n,i})_{i\ge1}$ jest niezależny od $\calF_n$.

\begin{uwaga}[Dlaczego występują dwa indeksy?]
Liczba $n$ identyfikuje pokolenie, natomiast $i$ rozróżnia
osobników w tym pokoleniu. Zmienna $\xi_{n,i}$ określa liczbę
dzieci $i$-tego osobnika w pokoleniu $n$. Niezależność jest
założeniem modelu: jeżeli wszystkie osobniki w danym pokoleniu
poddane są temu samemu losowemu wpływowi środowiska, zwykły
proces Galtona--Watsona może być nieadekwatny.
\end{uwaga}

\begin{center}
\begin{tikzpicture}[>=Latex,x=1cm,y=0.7cm]
\node[circle,draw,fill=blue!10,inner sep=2.5pt] (a) at (4,0) {};
\node[circle,draw,fill=blue!10,inner sep=2.5pt] (b) at (2,-1.8) {};
\node[circle,draw,fill=blue!10,inner sep=2.5pt] (c) at (4,-1.8) {};
\node[circle,draw,fill=blue!10,inner sep=2.5pt] (d) at (6,-1.8) {};
\node[circle,draw,fill=blue!10,inner sep=2.5pt] (e) at (1,-3.6) {};
\node[circle,draw,fill=blue!10,inner sep=2.5pt] (f) at (2,-3.6) {};
\node[circle,draw,fill=blue!10,inner sep=2.5pt] (g) at (5.3,-3.6) {};
\node[circle,draw,fill=blue!10,inner sep=2.5pt] (h) at (6.6,-3.6) {};
\foreach \x/\y in {a/b,a/c,a/d,b/e,b/f,d/g,d/h}{\draw[thick,blue!60!black] (\x)--(\y);}
\node[left] at (-.2,0) {$n=0$};
\node[left] at (-.2,-1.8) {$n=1$};
\node[left] at (-.2,-3.6) {$n=2$};
\node[right] at (7.3,0) {$Z_0=1$};
\node[right] at (7.3,-1.8) {$Z_1=3$};
\node[right] at (7.3,-3.6) {$Z_2=4$};
\end{tikzpicture}
\par\smallskip\small Przykładowe drzewo genealogiczne.
Środkowy osobnik pierwszego pokolenia nie ma potomstwa.
Proces zlicza liczby osobników, nie zapamiętuje pełnego drzewa.
\end{center}

\section{Własność Markowa i funkcje tworzące}
\begin{definicja}[Funkcja tworząca potomstwa]
Dla zmiennej $\xi$ o wartościach w $\N_0$ jej funkcją tworzącą
prawdopodobieństwa na $[0,1]$ jest
$f(s)=\E[s^\xi]=\sum_{k=0}^\infty p_ks^k$.
Przyjmujemy $s^0=1$, także dla $s=0$.
Średnią liczbę potomków oznaczamy
$m=\E\xi\in[0,\infty]$, a gdy istnieje drugi moment,
wariancję przez $\sigma^2=\Var(\xi)$.
Szereg definiujący $f$ jest jednostajnie zbieżny na
$[0,1]$ przez oszacowanie wyrazów przez $p_k$;
w szczególności $f$ jest ciągła i niemalejąca.
\end{definicja}

\begin{twierdzenie}[Rozkład przejścia i składanie funkcji tworzących]
Proces Galtona--Watsona z dowolnym deterministycznym
$Z_0=z\in\N_0$ jest jednorodnym łańcuchem Markowa
na $\N_0$. Jeżeli $Z_0=1$, to dla
$0\le s\le1$ jego funkcja tworząca $g_n(s)=\E[s^{Z_n}]$
spełnia
\[
g_0(s)=s,\qquad g_{n+1}(s)=g_n(f(s)),
\qquad g_n=f^{\circ n}.
\]
Przy starcie od $Z_0=z$ zachodzi $\E[s^{Z_n}]=(f^{\circ n}(s))^z$.
\end{twierdzenie}
\begin{proof}
Dla $k,j\in\N_0$ niech $P(k,j)$ będzie masą w $j$
$k$-krotnego splotu prawa $\xi$, przy czym splot zerowy
jest masą jednostkową w zerze.
Na zdarzeniu $\{Z_n=k\}$ niezależność $n$-tego wiersza
od $\calF_n$ daje
$\Prob(Z_{n+1}=j\mid\calF_n)=P(k,j)$.
Formalnie mnożymy przez funkcję charakterystyczną
dowolnego zdarzenia $A\in\calF_n$ i sumujemy po $k$.
Otrzymujemy $\Prob(Z_{n+1}=j\mid\calF_n)=P(Z_n,j)$,
czyli jednorodną własność Markowa, również względem
filtracji naturalnej (przez własność wieży).
Analogicznie dla $s\in[0,1]$ niezależność składników
daje $\E[s^{Z_{n+1}}\mid\calF_n]=f(s)^{Z_n}$.
Biorąc wartość oczekiwaną, otrzymujemy
$g_{n+1}(s)=\E[f(s)^{Z_n}]=g_n(f(s))$.
Indukcja i $g_0(s)=s$ dają $g_n=f^{\circ n}$.
Dla dowolnego $z$ ta sama rekurencja zaczyna się
od $g_0(s)=s^z$, więc daje $(f^{\circ n}(s))^z$.
\end{proof}

\begin{stwierdzenie}[Własność rozgałęziania]
\label{prop:gw-branching}
Prawo całego procesu z $z$ przodkami jest prawem sumy
$z$ niezależnych procesów Galtona--Watsona z jednym
przodkiem. Ponadto warunkowo względem $\calF_n$,
na $\{Z_n=k\}$ przyszła populacja ma prawo sumy
$k$ takich niezależnych procesów, z czasem liczonym od $n$.
\end{stwierdzenie}
\begin{proof}
Na przestrzeni iloczynowej budujemy niezależne kopie
procesu z jednym przodkiem. Jeśli ich obecne liczebności
wynoszą $k_1,\ldots,k_z$, rozkład sumy następnych
liczebności jest splotem $k_1+\cdots+k_z$ kopii
prawa $\xi$, niezależnie od podziału tej sumy.
Warunkowanie względem wspólnej historii i następnie
względem historii samej sumy daje macierz przejścia
$P$ z poprzedniego dowodu.
Stan początkowy sumy wynosi $z$, więc indukcja
po liczbie chwil identyfikuje jej rozkłady
skończeniewymiarowe z rozkładami pierwotnego procesu.
Dla drugiej tezy niezależność wszystkich przyszłych
wierszy tablicy od $\calF_n$ pozwala powtórzyć
ten rachunek przy ustalonym $Z_n=k$. Tożsamości dla
cylindrów rozszerzamy argumentem $\pi$--$\lambda$
na $\sigma$-ciało produktowe przyszłych trajektorii.
\end{proof}

\begin{twierdzenie}[Średnia liczba osobników]
Niech $m=\E\xi<\infty$ i $Z_0=z\in\N_0$.
Wtedy dla każdego $n\in\N_0$
\[
\E[Z_{n+1}\mid\calF_n]=mZ_n,\qquad
\E Z_n=z\,m^n,
\]
z konwencją $m^0=1$ także przy $m=0$.
Jeżeli $m>0$, proces $(Z_n/m^n)_{n\ge0}$ jest
nieujemnym martyngałem (dla $z\ge0$).
\end{twierdzenie}
\begin{proof}
Warunkowo względem $\calF_n$, po ustaleniu $Z_n=k$,
wartość oczekiwana sumy $k$ niezależnych kopii $\xi$
wynosi $km$. Ponieważ $\E Z_n<\infty$ wynika
indukcyjnie z $\E Z_0=z$, otrzymujemy
$\E Z_{n+1}=m\E Z_n$ i formułę $zm^n$.
Dla $m>0$, dzieląc równość warunkową przez $m^{n+1}$,
dostajemy $\E[Z_{n+1}/m^{n+1}\mid\calF_n]=Z_n/m^n$.
Całkowalność każdego wyrazu wykazaliśmy wcześniej.
\end{proof}

\section{Prawdopodobieństwo wymarcia jako punkt stały}
\begin{definicja}[Wymarcie]
Dla procesu Galtona--Watsona z $Z_0=1$ definiujemy
$\mathrm{Ext}:=\{\exists n\ge0:Z_n=0\}$ oraz
$q:=\Prob(\mathrm{Ext})$.
Ponieważ stan $0$ jest pochłaniający,
zdarzenie wymarcia jest sumą rosnącą zdarzeń $\{Z_n=0\}$.
\end{definicja}

\begin{twierdzenie}[Najmniejszy punkt stały]
Niech $f(s)=\E[s^\xi]$ będzie funkcją tworzącą potomstwa
procesu Galtona--Watsona z jednym przodkiem.
Wtedy prawdopodobieństwo wymarcia $q$ jest
\emph{najmniejszym} rozwiązaniem równania $f(s)=s$
w przedziale $[0,1]$. Można je obliczać rekurencją
$q_0=0$, $q_{n+1}=f(q_n)$, przy czym
$q_n\uparrow q$.
\end{twierdzenie}
\begin{proof}
Z poprzedniego twierdzenia
$q_n=\Prob(Z_n=0)=g_n(0)=f^{\circ n}(0)$,
co daje rekurencję. Zdarzenia $\{Z_n=0\}$ rosną,
więc $q_n$ rośnie i ma granicę $q\in[0,1]$.
Równość $q=\Prob(\mathrm{Ext})$ wynika z ciągłości miary
prawdopodobieństwa z dołu.
Funkcja tworząca $f$ jest ciągła na $[0,1]$:
ciągłość z lewej w $1$ otrzymujemy ze zbieżności monotonicznej
$s^\xi\uparrow1$, a w pozostałych punktach ze zbieżności
zmajoryzowanej. Dlatego przechodząc do granicy w
$q_{n+1}=f(q_n)$, dostajemy $q=f(q)$.
Jeżeli $r\in[0,1]$ spełnia $f(r)=r$, to $q_0=0\le r$.
Monotoniczność $f$ i indukcja dają $q_n\le r$
dla każdego $n$, a więc $q\le r$. To dowodzi minimalności.
\end{proof}

\begin{twierdzenie}[Kryterium wymierania]
Niech $\xi$ będzie zmienną o wartościach w $\N_0$, z
$m=\E\xi\in[0,\infty]$ oraz $p_0=\Prob(\xi=0)>0$.
Dla procesu Galtona--Watsona z jednym przodkiem:
\begin{enumerate}
\item jeśli $m\le1$, to $q=1$;
\item jeśli $m>1$, to $0<q<1$.
\end{enumerate}
Jeżeli $p_0=0$, wymarcie jest niemożliwe i $q=0$;
przypadek $\xi\equiv1$ ma $m=1$, lecz również $q=0$.
\end{twierdzenie}
\begin{proof}
Funkcja $f$ jest wypukła na $[0,1]$, bo każda
funkcja $s\mapsto s^k$ jest wypukła dla $k\ge0$.
Jej pochodna lewostronna w $1$ wynosi $m$,
również gdy $m=\infty$. Istotnie,
\[
\frac{1-f(s)}{1-s}
=\sum_{k\ge1}p_k\sum_{j=0}^{k-1}s^j
\uparrow\sum_{k\ge1}kp_k=m
\quad(s\uparrow1)
\]
na mocy twierdzenia o zbieżności monotonicznej.
Jeżeli $m\le1$, wypukłość daje dla $0\le s<1$
\[
f(s)\ge f(1)+m(s-1)=1-m(1-s)\ge s.
\]
Różnica $f(s)-s$ jest ściśle dodatnia dla $m<1$.
Jeżeli $m=1$ i $p_0>0$, warunek średniej wymusza
$\Prob(\xi\ge2)>0$, a więc $f$ jest ściśle wypukła
na $(0,1)$ i $f(s)>s$ również w tym przypadku.
Zatem jedynym punktem stałym jest $1$, stąd $q=1$.
Jeżeli $m>1$, to
$f(s)-s<0$ dla $s<1$ dostatecznie bliskich $1$:
istotnie iloraz $(f(s)-1)/(s-1)$ dąży do $m>1$.
Jednocześnie $f(0)-0=p_0>0$.
Ciągłość zapewnia istnienie punktu stałego $q\in(0,1)$.
Najmniejszy punkt stały, będący prawdopodobieństwem wymarcia,
leży w $(0,1)$: mniejszy od $1$ na mocy istnienia punktu stałego,
a większy od zera, gdyż $\Prob(Z_1=0)=p_0>0$.
Ponieważ $m>1$ wymusza $\Prob(\xi\ge2)>0$,
funkcja $f(s)-s$ jest ściśle wypukła na $[0,1]$.
Nie może mieć trzech różnych zer (wartość w środkowym
byłaby ściśle mniejsza od zera). Zatem $q$ jest
jedynym punktem stałym w $[0,1)$.
Gdy $p_0=0$, z prawdopodobieństwem $1$ wszystkie
zmienne z przeliczalnej tablicy są dodatnie,
więc $Z_n\ge1$ dla wszystkich $n$ i $q=0$.
\end{proof}

\begin{wniosek}[Wymarcie przy wielu przodkach i po zadanym pokoleniu]
Przy starcie z $z$ przodkami prawdopodobieństwo
wymarcia wynosi $q^z$ (także dla $z=0$, z $q^0=1$).
Dla dowolnego $n$ warunkowe prawdopodobieństwo
ostatecznego wymarcia względem $\calF_n$ wynosi $q^{Z_n}$.
\end{wniosek}
\begin{proof}
Dla startu z $z$ mamy
$\Prob(Z_n=0)=(f^{\circ n}(0))^z$.
Przejście do granicy daje pierwszą tezę.
Druga wynika z własności rozgałęziania
(stwierdzenie~\ref{prop:gw-branching}) albo z własności
Markowa dla zdarzeń $\{Z_{n+k}=0\}$ i warunkowej
zbieżności monotonicznej przy $k\to\infty$.
\end{proof}

\begin{zadanie}[Dwoje dzieci albo brak]
Każdy osobnik ma dwoje dzieci z prawdopodobieństwem $2/3$
albo nie ma dzieci z prawdopodobieństwem $1/3$.
Obliczyć prawdopodobieństwo wymarcia i średnią liczbę
osobników w pokoleniu $n$.

\emph{Rozwiązanie.} Funkcja tworząca ma postać
$f(s)=\frac13+\frac23s^2$, a średnia potomstwa
$m=4/3>1$. Równanie $f(q)=q$ daje
$2q^2-3q+1=0$, czyli pierwiastki $q=1/2$ i $q=1$.
Najmniejszy to $q=1/2$. Pomimo tego, że średnia
liczba osobników rośnie jak $(4/3)^n$, prawdopodobieństwo
całkowitego wymarcia wynosi aż $1/2$.
\end{zadanie}


\begin{zadanie}[Prawdopodobieństwo wymarcia dla trzech możliwości]
Niech liczba potomków ma rozkład
$\Prob(\xi=0)=1/4$, $\Prob(\xi=1)=1/4$,
$\Prob(\xi=2)=1/2$. Obliczyć średnią liczbę potomków
oraz prawdopodobieństwo ostatecznego wymarcia populacji
rozpoczętej od dwóch osobników.

\emph{Rozwiązanie.} Średnia wynosi
$m=0\cdot\tfrac14+1\cdot\tfrac14+2\cdot\tfrac12=5/4>1$.
Funkcja tworząca to
$f(s)=\tfrac14+\tfrac14s+\tfrac12s^2$. Równanie
$f(s)=s$ jest równoważne $2s^2-3s+1=0$,
a więc ma pierwiastki $1/2$ i $1$. Dla jednego przodka
$q=1/2$. Dwie niezależne linie genealogiczne wymierają
jednocześnie z prawdopodobieństwem $q^2=1/4$.
Przykład pokazuje, dlaczego na ogół nie wolno zastąpić
liczby przodków przez zmianę średniej $m$.
\end{zadanie}


\begin{zadanie}[Wymarcie w skończonej liczbie pokoleń]
Niech $\Prob(\xi=0)=\Prob(\xi=2)=1/2$ i $Z_0=1$.
Obliczyć prawdopodobieństwa wymarcia do drugiego i trzeciego
pokolenia; wyjaśnić różnicę między nimi a ostatecznym wymarciem.

\emph{Rozwiązanie.} Funkcja tworząca wynosi
$f(s)=(1+s^2)/2$. W pierwszym pokoleniu
$q_1=f(0)=1/2$. W drugim $q_2=f(q_1)=5/8$,
a w trzecim $q_3=f(q_2)=\tfrac12(1+25/64)=89/128$.
To konkretne kolejne przybliżenia rosną do $q=1$,
gdyż $m=1$ i $p_0>0$. Zdarzenie wymarcia do pokolenia
$n$ jest zdarzeniem $\{Z_n=0\}$, a ciąg tych zdarzeń
rośnie: raz wymarła populacja już się nie odradza.
\end{zadanie}

\section{Wariancja, normalizacja i granica martyngału}
\begin{twierdzenie}[Wariancja pokolenia]
Załóżmy, że $Z_0=1$, $0<m=\E\xi<\infty$ oraz
$\sigma^2=\Var(\xi)<\infty$. Wtedy
\[
\Var(Z_n)=
\begin{cases}
\displaystyle \sigma^2m^{n-1}\frac{m^n-1}{m-1},&m\ne1,\ n\ge1,\\[6pt]
 n\sigma^2,&m=1,\ n\ge0.
\end{cases}
\]
Jeżeli $m>1$, martyngał $W_n=Z_n/m^n$ jest ograniczony
w $L^2$, zatem zbiega w $L^2$ i prawie na pewno do
pewnej zmiennej $W_\infty\ge0$ o $\E W_\infty=1$ oraz
\[
\Var(W_\infty)=\frac{\sigma^2}{m(m-1)}.
\]
\end{twierdzenie}
\begin{proof}
Z niezależności potomstwa, na zdarzeniu $Z_n=k$,
$\E[Z_{n+1}^2\mid\calF_n]=k\sigma^2+k^2m^2$.
Indukcyjnie $\E Z_n^2<\infty$, zaczynając od $Z_0=1$.
W szczególności
$\Var(Z_{n+1}\mid\calF_n)=\sigma^2Z_n$.
Wzór na wariancję całkowitą daje dla $v_n=\Var(Z_n)$
\[
v_{n+1}=\E[\sigma^2Z_n]
+\Var(mZ_n)=\sigma^2m^n+m^2v_n,
\quad v_0=0.
\]
Indukcja daje
$v_n=\sigma^2\sum_{j=0}^{n-1}m^{2(n-1-j)+j}
=\sigma^2m^{n-1}\sum_{j=0}^{n-1}m^j$,
co po zsumowaniu szeregu geometrycznego daje wzór.
Dla $n=0$ wariancja wynosi zero. Dla $m>1$ otrzymujemy
$\Var(W_n)=m^{-2n}v_n
=\frac{\sigma^2}{m(m-1)}(1-m^{-n})$.
Ponieważ $\E W_n=1$, drugie momenty są jednostajnie
ograniczone. Z twierdzenia o zbieżności martyngałów
ograniczonych w $L^2$ dostajemy granicę w $L^2$ i p.n.;
przejście do granicy w pierwszych i drugich momentach daje
deklarowane wzory.
\end{proof}

\begin{uwaga}[Wartość oczekiwana granicy nie zawsze zachowuje się sama]
Dla $m>1$ nieujemny martyngał $W_n$ ma granicę prawie na pewno
nawet bez założenia skończonej wariancji, ale z tego
\emph{nie wynika} automatycznie $\E W_\infty=1$.
Skończoność wariancji jest tutaj wygodnym dostatecznym
warunkiem jednostajnej całkowalności. Bardziej ogólny
warunek zachowania średniej w procesie nadkrytycznym
wiąże się z twierdzeniem Kestena--Stiguma i momentem
$\E[\xi\log^+\xi]$; ten głębszy wynik można omówić
przy bardziej zaawansowanych procesach rozgałęziających.
Poniższe twierdzenie wyjaśnia, co można już
wywnioskować z niezerowej granicy martyngału.
\end{uwaga}

\begin{twierdzenie}[Granica martyngału a wymarcie]
\label{thm:gw-limit-extinction}
Niech $Z_0=1$ oraz $1<m<\infty$.
Martyngał $W_n=Z_n/m^n$ ma skończoną granicę
$W$ prawie na pewno. Jeżeli $\Prob(W>0)>0$, to
\[
\Prob(W=0)=q,\qquad
\{W=0\}=\mathrm{Ext}\quad\text{z dokładnością do zbioru zerowego}.
\]
W szczególności zachodzi to przy skończonej wariancji
potomstwa lub, ogólniej, przy jednostajnej całkowalności
$(W_n)_{n\ge0}$. Na zdarzeniu przetrwania w tych
przypadkach $Z_n/m^n\to W>0$ prawie na pewno.
\end{twierdzenie}
\begin{proof}
Twierdzenie o zbieżności nieujemnych martyngałów
daje granicę $W\ge0$ skończoną prawie na pewno,
z $\E W\le1$ na mocy lematu Fatou.
Na zdarzeniu wymarcia $Z_n=0$ od pewnego indeksu,
więc $W=0$.

Z własności rozgałęziania
(stwierdzenie~\ref{prop:gw-branching}) warunkowo
względem $Z_1=k$ proces od pierwszego pokolenia
ma prawo sumy $k$ niezależnych procesów z jednym
przodkiem. Ponieważ $k$ jest skończone, przechodząc
do granicy prawie pewnych granic na przestrzeni
produktowej, otrzymujemy równość rozkładów
\[
W\overset{d}=\frac1m\sum_{j=1}^{\xi} W^{(j)},
\]
gdzie $W^{(j)}$ są niezależnymi kopiami $W$,
niezależnymi od $\xi$. Można ją także uzyskać
z rozkładów warunkowych dla każdego skończonego $k$
i zsumować po $k$.
Nieujemność składników daje dla $r=\Prob(W=0)$
równanie $r=f(r)$. Jeśli $\Prob(W>0)>0$, to $r<1$.
Dla $p_0>0$ jedynym punktem stałym poniżej $1$
jest $q$. Dla $p_0=0$ i $m>1$ mamy dla $0<s<1$
$f(s)<s$, ponieważ $s^k\le s$ dla $k\ge1$
i $\Prob(\xi\ge2)>0$; jedynym punktem stałym
poniżej $1$ jest wtedy $q=0$.
Stąd $r=q$. Z zawierania $\mathrm{Ext}\subseteq\{W=0\}$
z dokładnością do zbioru zerowego i równości miar
wynika równość zdarzeń prawie na pewno.
Jeśli martyngał jest jednostajnie całkowalny,
zbiega w $L^1$, a $\E W=1$, więc $r<1$.
\end{proof}

\begin{wniosek}[Liczebność na zdarzeniu przetrwania]
Jeżeli $m>1$ (także $m=\infty$), to na zdarzeniu
przetrwania proces Galtona--Watsona z jednym
przodkiem spełnia $Z_n\to\infty$ prawie na pewno.
Nie trzeba zakładać skończonego drugiego momentu.
\end{wniosek}
\begin{proof}
Gdy $p_0>0$, mamy $0<q<1$. Warunkowe
prawdopodobieństwa wymarcia są równe $q^{Z_n}$.
Z twierdzenia o zbieżności warunkowych wartości
oczekiwanych względem rosnących $\sigma$-ciał
zbiegają prawie na pewno do
$\chi_{\mathrm{Ext}}$, ponieważ zdarzenie
wymarcia jest mierzalne względem
$\sigma(\bigcup_n\calF_n)$.
Na zdarzeniu przetrwania otrzymujemy więc
$q^{Z_n}\to0$, czyli $Z_n\to\infty$.
Jeżeli $p_0=0$, wszystkie osobniki mają co najmniej
jedno dziecko prawie na pewno, więc $Z_n$ jest
niemalejący. Przy $m>1$ mamy $p_1<1$.
Dla ustalonych $k\ge1$ i $n\ge0$ prawdopodobieństwo
pozostawania na poziomie $k$ przez następne $l$
pokoleń, pod warunkiem $Z_n=k$, wynosi $p_1^{kl}$.
Przejście $l\to\infty$ i suma po $k,n$ wykluczają
ostateczne pozostawanie na skończonym poziomie.
\end{proof}

\section{Całkowita liczba urodzonych osobników}
\begin{twierdzenie}[Średnia całkowita liczba osobników]
Niech $Z_0=1$ i $m=\E\xi<1$.
Zdefiniujmy całkowitą liczebność drzewa genealogicznego
$H=\sum_{n=0}^\infty Z_n$, gdzie suma może przyjmować $\infty$.
Wtedy $H<\infty$ prawie na pewno i
$\E H=1/(1-m)$.
\end{twierdzenie}
\begin{proof}
Zmienne $Z_n$ są nieujemne. Twierdzenie Tonellego
i wzór $\E Z_n=m^n$ dają
\[
\E H=\sum_{n=0}^\infty\E Z_n=\sum_{n=0}^\infty m^n
=\frac{1}{1-m}<\infty.
\]
Gdyby $\Prob(H=\infty)>0$, wartość oczekiwana $H$
byłaby nieskończona. Stąd $H<\infty$ prawie na pewno.
\end{proof}

\begin{stwierdzenie}[Funkcja tworząca całkowitej liczebności]
Niech $Z_0=1$ i $\Prob(H<\infty)=1$.
Dla $s\in[0,1]$ funkcja $h(s)=\E[s^H]$ spełnia
\[
h(s)=s f(h(s)),
\]
gdzie $f$ jest funkcją tworzącą liczby potomków.
\end{stwierdzenie}
\begin{proof}
Warunkowo względem $\xi=Z_1=k$, z własności
rozgałęziania (stwierdzenie~\ref{prop:gw-branching})
proces od pokolenia pierwszego ma prawo sumy
$k$ niezależnych drzew z jednym przodkiem.
Sumując liczebności po wszystkich pokoleniach,
otrzymujemy równość rozkładów
$H\overset{d}=1+\sum_{j=1}^\xi H_j$, gdzie po prawej
stronie możemy użyć niezależnych kopii $H$,
niezależnych od $\xi$, na przestrzeni iloczynowej.
Warunkowanie względem $\xi$ daje
$\E[s^H\mid\xi]=s\,h(s)^\xi$ prawie na pewno
(z warunkowej własności rozgałęziania), skąd po
uśrednieniu $h(s)=sf(h(s))$.
\end{proof}

\begin{zadanie}[Potomstwo Poissona]
Niech liczba potomków każdego osobnika ma rozkład
$\mathrm{Poisson}(\lambda)$, $0<\lambda<1$.
Obliczyć średnią całkowitą liczebność i równanie
dla jej funkcji tworzącej.

\emph{Rozwiązanie.} Mamy
$f(s)=\exp(\lambda(s-1))$ i $m=\lambda$.
Stąd $\E H=1/(1-\lambda)$ oraz
$h(s)=s\exp(\lambda(h(s)-1))$.
To równanie pozwala numerycznie obliczać
$\Prob(H=k)$, bez budowania całego losowego drzewa.
\end{zadanie}

\section{Model rozgałęziania w czasie ciągłym: proces Yule'a}
\begin{definicja}[Proces czystych urodzin Yule'a]
Niech $\lambda>0$. Procesem Yule'a rozpoczynającym się
od jednego osobnika nazywamy minimalny łańcuch Markowa
w czasie ciągłym na $\{1,2,3,\ldots\}$ z intensywnościami
$q_{n,n+1}=\lambda n$, $q_{nn}=-\lambda n$ oraz
$q_{n,j}=0$ dla pozostałych $j$.
Interpretacja: każdy z $n$ obecnych osobników rodzi
kolejnego niezależnie z intensywnością $\lambda$.
\end{definicja}

\begin{twierdzenie}[Brak eksplozji i rozkład geometryczny]
Proces Yule'a rozpoczynający się od $Z_0=1$ nie wybucha
w skończonym czasie. Dla wszystkich $t\ge0$ i $n\ge1$
\[
\Prob(Z_t=n)=e^{-\lambda t}(1-e^{-\lambda t})^{n-1},
\qquad \E Z_t=e^{\lambda t}.
\]
\end{twierdzenie}
\begin{proof}
Po przejściu ze stanu $n$ proces musi przejść do $n+1$.
Czas oczekiwania w stanie $n$ ma rozkład
$\mathrm{Exp}(\lambda n)$; kolejne czasy są niezależne.
Suma $n$ pierwszych czasów oczekiwania jest czasem przejścia
od $1$ do $n+1$. Jej transformata Laplace'a dla $c>0$ to
$\prod_{k=1}^n (\lambda k)/(\lambda k+c)$.
Gdyby czas eksplozji $\zeta=\sum_{k\ge1}H_k$ był
skończony z dodatnim prawdopodobieństwem, to
$\E e^{-c\zeta}>0$. Tymczasem
$\sum_{k\ge1}\log(1+c/(\lambda k))=\infty$,
ponieważ dla dużych $k$ wyrazy są nie mniejsze od
$c/(2\lambda k)$, a szereg harmoniczny jest rozbieżny.
Zbieżność zmajoryzowana, z konwencją $e^{-c\infty}=0$,
pozwala przejść do granicy w transformatach sum
częściowych. Zatem $\E e^{-c\zeta}=0$, co wymusza
$\Prob(\zeta=\infty)=1$. To jest bezpośredni dowód
przypadku kryterium eksplozji potrzebnego tutaj.

Dla przejść rozkładu $p_n(t)=\Prob(Z_t=n)$
uzasadnimy najpierw użycie równań w przód, chociaż przestrzeń
stanów jest nieskończona. Dla każdego ustalonego $K\ge2$
zastąpmy stany $K+1,K+2,\ldots$ jednym pochłaniającym
stanem $\partial$, nie zmieniając przejść z $1,\ldots,K$.
Do pierwszego osiągnięcia $K+1$ oba procesy są identyczne,
a ponieważ proces Yule'a wykonuje jedynie skoki $n\to n+1$,
dla $j\le K$ ich prawdopodobieństwa przebywania w stanie $j$
są równe. Proces obcięty ma skończoną przestrzeń stanów;
możemy więc zastosować wcześniej udowodnione równania
Kołmogorowa w przód (twierdzenie~\ref{thm:kolmogorov-eq}),
a otrzymane równania dla $p_j$, $j\le K$,
nie zależą od $K$. Dowodzi to dla wszystkich $n\ge1$ równań
$p_1'(t)=-\lambda p_1(t)$,
$p_n'(t)=-n\lambda p_n(t)+(n-1)\lambda p_{n-1}(t)$ dla $n\ge2$,
z warunkami $p_1(0)=1$ i $p_n(0)=0$ dla $n\ge2$.
Pierwsze równanie daje $p_1(t)=e^{-\lambda t}$.
Podstawiając kandydat
$p_n(t)=a(t)(1-a(t))^{n-1}$, gdzie $a(t)=e^{-\lambda t}$,
i korzystając z $a'=-\lambda a$, otrzymujemy po różniczkowaniu
$p_n'=-\lambda a(1-a)^{n-1}+\lambda(n-1)a^2(1-a)^{n-2}$.
Prawa strona równania Kołmogorowa wynosi
$-n\lambda a(1-a)^{n-1}+(n-1)\lambda a(1-a)^{n-2}$,
co po wyłączeniu $\lambda a(1-a)^{n-2}$ sprowadza się do
$-n(1-a)+(n-1)=na-1$,
tak samo jak pierwsza pochodna.
Jednoznaczność rozwiązań wynika indukcyjnie: każde równanie
jest liniowym równaniem niejednorodnym z wcześniej znaną
funkcją $p_{n-1}$. Ponieważ kandydat daje rozkład
prawdopodobieństwa ($\sum_{n\ge1}a(1-a)^{n-1}=1$),
odzyskujemy prawo procesu. Średnia rozkładu
geometrycznego na $\{1,2,\ldots\}$ wynosi $1/a=e^{\lambda t}$.
\end{proof}

\begin{uwaga}[Dwa odmienne mechanizmy wzrostu]
W procesie Galtona--Watsona wszyscy osobnicy rozmnażają się
w kolejnych \emph{pokoleniach}, a liczebność zmienia się
skokowo w chwilach całkowitych. W procesie Yule'a
urodzenia pojawiają się w \emph{czasie ciągłym}, przy czym
intensywność łączna rośnie proporcjonalnie do bieżącej liczby
osobników. Oba procesy mają strukturę rozgałęziającą,
ale inne trajektorie, rozkłady i narzędzia obliczeniowe.
\end{uwaga}


\section{Jak szybko zanika prawdopodobieństwo przetrwania?}
\begin{stwierdzenie}[Ogon czasu wymarcia w przypadku podkrytycznym]
Niech $Z$ będzie procesem Galtona--Watsona rozpoczynającym się
od $Z_0=1$ ze średnią potomstwa $m\in[0,1)$.
Niech $\tau=\inf\{n\ge0:Z_n=0\}$.
Wtedy dla każdego $n\ge0$
\[
\Prob(\tau>n)=\Prob(Z_n>0)\le m^n,
\qquad \E\tau\le\frac{1}{1-m}.
\]
\end{stwierdzenie}
\begin{proof}
Ponieważ zero jest stanem pochłaniającym,
$\{\tau>n\}=\{Z_n>0\}$. Dla nieujemnej zmiennej
całkowitoliczbowej nierówność Markowa daje
$\Prob(Z_n>0)=\Prob(Z_n\ge1)\le\E Z_n=m^n$.
Wzór ogonowy dla zmiennej o wartościach w $\N_0\cup\{\infty\}$,
wynikający z twierdzenia Tonellego, ma postać
$\E\tau=\sum_{n=0}^\infty\Prob(\tau>n)$.
Sumując oszacowania, otrzymujemy deklarowaną granicę.
\end{proof}

\begin{twierdzenie}[Asymptotyka Kołmogorowa w przypadku krytycznym]
\label{thm:gw-critical-survival}
Niech $Z_0=1$, $\E\xi=1$ i
$0<\sigma^2=\Var(\xi)<\infty$.
Wtedy
\[
\Prob(Z_n>0)\sim\frac{2}{\sigma^2 n}
\qquad(n\to\infty).
\]
W szczególności czas wymarcia $\tau$ jest skończony
prawie na pewno, lecz $\E\tau=\infty$.
\end{twierdzenie}
\begin{proof}
Warunki $m=1$ i $\sigma^2>0$ wymuszają $p_0>0$:
gdyby $\xi\ge1$ p.n., średnia $1$ wymuszałaby
$\xi=1$ p.n. Z kryterium wymierania wynika
$q_n=\Prob(Z_n=0)\uparrow1$.
Połóżmy $a_n=1-q_n>0$; dodatniość wynika
z istnienia dodatniego prawdopodobieństwa
co najmniej jednego dziecka w każdym z $n$
pokoleń (można śledzić pierwszą linię potomną).
Mamy $a_{n+1}=1-f(1-a_n)$.

Dla $0<s<1$ różniczkowanie szeregu potęgowego
daje $f''(s)=\sum_{k\ge2} k(k-1)p_k s^{k-2}$.
Skończoność drugiego momentu daje jednostajną
zbieżność szeregu pochodnych do rzędu drugiego
na $[0,1]$ (każdy wyraz jest ograniczony przez
$k(k-1)p_k$ dla drugiej pochodnej).
W szczególności $f'(1)=1$ oraz
$f''(1)=\E[\xi(\xi-1)]=\sigma^2$,
gdzie pochodne w $1$ są lewostronne.
Wzór Taylora daje zatem, dla $u\downarrow0$,
\[
1-f(1-u)=u-\tfrac12\sigma^2u^2+o(u^2).
\]
Oznaczmy $c=\sigma^2/2>0$. Ponieważ $a_n\to0$,
mamy $a_{n+1}/a_n\to1$ oraz
\[
\frac1{a_{n+1}}-\frac1{a_n}
=\frac{a_n-a_{n+1}}{a_na_{n+1}}\longrightarrow c.
\]
Sumując różnice i stosując twierdzenie o średnich
arytmetycznych ciągu zbieżnego, dostajemy
$1/(na_n)\to c$. Zatem $na_n\to1/c$.
Ze wzoru ogonowego $\E\tau=\sum_{n\ge0}a_n$
i rozbieżności szeregu harmonicznego wynika
ostatnia teza.
\end{proof}

\begin{twierdzenie}[Twierdzenie Yagloma]
\label{thm:gw-yaglom}
Przy założeniach twierdzenia~\ref{thm:gw-critical-survival}
warunkowy rozkład $Z_n/n$ przy $Z_n>0$ zbiega
według rozkładu do rozkładu wykładniczego
o średniej $\sigma^2/2$, czyli intensywności
$2/\sigma^2$:
\[
\Law\left(\frac{Z_n}{n}\,\middle|\,Z_n>0\right)
\longrightarrow\mathrm{Exp}(2/\sigma^2).
\]
\end{twierdzenie}
\begin{proof}
Niech $c=\sigma^2/2$ i
$G(u)=1-f(1-u)$ dla $0<u\le1$.
Dla $0<u<1$ mamy $0<G(u)\le u$,
ponieważ $f(1-u)\ge1-u$ w przypadku krytycznym,
a $f(1-u)<1$ dzięki $\Prob(\xi>0)>0$.
Z rozwinięcia użytego w poprzednim dowodzie
wynika
\[
D(u):=\frac1{G(u)}-\frac1u\longrightarrow c
\quad(u\downarrow0).
\]
Ustalmy $\theta>0$ i dla każdego $n\ge1$ połóżmy
$b_{n,0}=1-e^{-\theta/n}$ oraz
$b_{n,j+1}=G(b_{n,j})$ dla $0\le j<n$.
Mamy $0<b_{n,j}\le b_{n,0}\to0$ jednostajnie po $j$,
więc $\max_{0\le j<n}|D(b_{n,j})-c|\to0$.
Sumowanie różnic odwrotności daje
\[
\frac1{n b_{n,n}}
=\frac1{n b_{n,0}}+
\frac1n\sum_{j=0}^{n-1}D(b_{n,j})
\longrightarrow \frac1\theta+c .
\]
Zatem $n b_{n,n}\to\theta/(1+c\theta)$.
Ponieważ $b_{n,n}=1-f^{\circ n}(e^{-\theta/n})$
oraz $a_n=1-f^{\circ n}(0)$,
warunkowa transformata Laplace'a wynosi
\[
\E[e^{-\theta Z_n/n}\mid Z_n>0]
=\frac{f^{\circ n}(e^{-\theta/n})-f^{\circ n}(0)}
{a_n}
=1-\frac{b_{n,n}}{a_n}
\longrightarrow \frac1{1+c\theta}.
\]
Dla $\theta=0$ równość transformat jest automatyczna.
Ostatnia funkcja jest transformatą rozkładu
$\mathrm{Exp}(1/c)$.

Dla ścisłości wniosek o zbieżności rozkładów
można otrzymać bez dodatkowego twierdzenia
o transformatach. Średnie warunkowe
$\E[Z_n/n\mid Z_n>0]=1/(na_n)$ są ograniczone,
więc nierówność Markowa daje ciasność praw
na $[0,\infty)$. Każdy podciąg ma słabo zbieżny
podciąg (twierdzenie o wyborze Helly'ego dla
dystrybuant, z ciasnością). Jego granica ma
wskazaną transformatę, gdyż $x\mapsto e^{-\theta x}$
jest ograniczona i ciągła.
Jednoznaczność transformaty wynika już z wartości
dla $\theta=k\in\N_0$: po odwzorowaniu $x\mapsto e^{-x}$
otrzymujemy wszystkie momenty miary na $[0,1]$.
Wielomiany są jednostajnie gęste w $C([0,1])$
na mocy twierdzenia Weierstrassa, więc momenty
wyznaczają miarę, a odwzorowanie $x\mapsto e^{-x}$
jest różnowartościowe z borelowską odwrotnością
na $(0,1]$. Wszystkie granice podciągowe są
zatem identyczne, co kończy dowód.
\end{proof}

\begin{wniosek}[Średnia i wariancja pod warunkiem przetrwania]
Przy tych samych założeniach
\[
\E[Z_n/n\mid Z_n>0]\longrightarrow c,\qquad
\Var(Z_n/n\mid Z_n>0)\longrightarrow c^2,
\quad c=\sigma^2/2.
\]
\end{wniosek}
\begin{proof}
Ponieważ $\E Z_n=1$ oraz $\E Z_n^2=1+n\sigma^2$,
a oba momenty nie mają wkładu ze zdarzenia $Z_n=0$,
mamy
\[
\E[Z_n/n\mid Z_n>0]=\frac1{na_n},\qquad
\E[(Z_n/n)^2\mid Z_n>0]=\frac{1+n\sigma^2}{n^2a_n}.
\]
Asymptotyka $na_n\to1/c$ daje granice $c$ oraz
$\sigma^2c=2c^2$. Po odjęciu kwadratu średniej
otrzymujemy granicę wariancji.
\end{proof}

\begin{przyklad}[Krytyczne potomstwo zero lub dwa]
Dla $\Prob(\xi=0)=\Prob(\xi=2)=1/2$ mamy
$m=1$, $\sigma^2=1$ i
$a_{n+1}=a_n-a_n^2/2$, $a_0=1$.
Wymarcie następuje prawie na pewno, lecz
$\Prob(Z_n>0)\sim2/n$.
Warunkowo względem przetrwania liczebność jest
rzędu $n$, a $Z_n/n$ ma graniczny rozkład
wykładniczy o średniej $1/2$.
Nie ma sprzeczności z $\E Z_n=1$:
malejąca szansa przetrwania równoważy wzrost
średniej liczebności w ocalałych populacjach.
\end{przyklad}

\section{Urodzenia i śmierci w czasie ciągłym}
\begin{twierdzenie}[Prawdopodobieństwo wymarcia liniowego procesu urodzin i śmierci]
Niech $\lambda,\delta>0$. Rozważmy minimalny łańcuch Markowa
$Z$ na $\N_0$, rozpoczynający się od $Z_0=1$, z intensywnościami
$q_{n,n+1}=\lambda n$, $q_{n,n-1}=\delta n$ dla $n\ge1$,
$q_{nn}=-(\lambda+\delta)n$, zaś $0$ jest stanem pochłaniającym.
Wtedy proces jest nieeksplozywny, a prawdopodobieństwo
ostatecznego wymarcia wynosi
\[
q=\begin{cases}1,&\lambda\le\delta,\\
\delta/\lambda,&\lambda>\delta.
\end{cases}
\]
\end{twierdzenie}
\begin{proof}
Najpierw uzasadnimy brak eksplozji bez odwoływania się do
nieformalnego sprzężenia. Dla $K\ge2$ niech
$\tau_K=\inf\{t\ge0:Z_t\ge K\}$. Proces zatrzymany
$Z_{t\wedge\tau_K}$ przyjmuje wartości w skończonym
zbiorze $\{0,1,\ldots,K\}$. Dla funkcji $f(n)=n$
jego generator przed zatrzymaniem wynosi
$Qf(n)=(\lambda-\delta)n\le\lambda n$.
Wzór Dynkina (twierdzenie~\ref{thm:dynkin})
dla skończonego łańcucha z pochłaniającymi stanami
$0$ i $K$, który ma prawo procesu do zatrzymania, daje dla $u_K(t)=\E Z_{t\wedge\tau_K}$
\[
u_K(t)\le1+\lambda\int_0^tu_K(s)\,\mathrm d\lambda_1(s).
\]
Nierówność Gronwalla prowadzi do
$u_K(t)\le e^{\lambda t}$. Ponieważ
$Z_{t\wedge\tau_K}=K$ na $\{\tau_K\le t\}$,
$\Prob(\tau_K\le t)\le e^{\lambda t}/K$.
Eksplozja w skończonym czasie wymaga osiągnięcia
każdego poziomu $K$ poza zbiorem zerowym:
jeżeli łańcuch pozostaje w skończonym zbiorze stanów,
jego intensywności są ograniczone, a suma czasów
oczekiwania przy nieskończenie wielu skokach
jest nieskończona p.n. (twierdzenie~\ref{thm:bounded-rates},
lub sprzężenie z niezależnymi odstępami wykładniczymi
o wspólnej maksymalnej intensywności).
Ponieważ każdy skok zmienia liczebność o jeden,
ucieczka z każdego skończonego zbioru wymusza
osiągnięcie każdego $K$. Dla każdego $t<\infty$
$\Prob(\zeta\le t)\le\inf_{K\ge2}e^{\lambda t}/K=0$,
gdzie $\zeta$ oznacza czas eksplozji. Proces jest zatem
nieeksplozywny.

Rozważmy pierwsze zdarzenie po starcie od jednego osobnika.
Łączna intensywność wynosi $\lambda+\delta$, a jego rodzaj
jest niezależny od czasu oczekiwania: z prawdopodobieństwem
$\delta/(\lambda+\delta)$ osobnik umiera, z prawdopodobieństwem
$\lambda/(\lambda+\delta)$ rodzi drugiego osobnika.
Konstrukcję przez niezależne zegary uzasadniamy
następująco: każdemu osobnikowi przypisujemy niezależny
czas wykładniczy o intensywności $\lambda+\delta$
i niezależny typ zdarzenia (śmierć lub rozdzielenie
na dwóch osobników). Minimum $n$ zegarów ma
intensywność $n(\lambda+\delta)$, a typ zdarzenia
ma podane prawdopodobieństwa. Własność braku pamięci
daje po zdarzeniu niezależne nowe zegary pozostałych
osobników, więc konstrukcja ma dokładnie intensywności
i czasy przebywania z definicji minimalnego łańcucha.
Brak eksplozji pozwala prowadzić ją dla wszystkich czasów.
Po rozdzieleniu mamy dwa niezależne drzewa potomne,
każde o takim samym prawdopodobieństwie wymarcia $q$.
Wobec tego
\[
q=\frac{\delta}{\lambda+\delta}
+\frac{\lambda}{\lambda+\delta}q^2.
\]
Rozwiązaniami są $1$ i $\delta/\lambda$.
Uzasadnijmy wybór mniejszego pierwiastka. Każdy żyjący
osobnik ma niezależny czas do pierwszego zdarzenia
o rozkładzie $\mathrm{Exp}(\lambda+\delta)$; wtedy
z prawdopodobieństwem $\delta/(\lambda+\delta)$
umiera bez potomków, a z prawdopodobieństwem
$\lambda/(\lambda+\delta)$ rodzi nowego osobnika.
W drugim przypadku można potraktować dotychczasowego
osobnika (z uwagi na własność braku pamięci) i nowego
jako \emph{dwie niezależne linie potomne}. Drzewo
zdarzeń jest więc drzewem Galtona--Watsona z liczbą
dzieci $0$ lub $2$ i funkcją tworzącą
$f(s)=\delta/(\lambda+\delta)+\lambda s^2/(\lambda+\delta)$.
Zdarzenie, że drzewo ma skończoną liczbę węzłów,
jest sumą zdarzeń wymarcia do skończonej głębokości;
jego prawdopodobieństwo jest zatem granicą iteracji
$f^{\circ n}(0)$, czyli najmniejszym punktem stałym $f$.
Jeżeli drzewo jest skończone, wszystkie związane z nim
czasy oczekiwania są skończone prawie na pewno,
a więc proces rzeczywiście osiąga stan $0$ w skończonym
czasie. Odwrotnie, jeśli osiąga $0$ w skończonym czasie,
to brak eksplozji oznacza skończoną liczbę zdarzeń do tej
chwili, a zatem skończone drzewo genealogiczne.
Wymarcie procesu jest więc dokładnie wymarciem tego drzewa.
Stosując kryterium z wcześniejszej części rozdziału,
wybieramy najmniejszy pierwiastek: $q=1$ przy
$\lambda\le\delta$ i $q=\delta/\lambda$ przy
$\lambda>\delta$.
\end{proof}

\begin{zadanie}[Próg przeżycia populacji]
Niech każda żyjąca cząstka rodzi nową z intensywnością
$\lambda=3$ i ginie z intensywnością $\delta=2$,
niezależnie od pozostałych. Rozpoczynamy od jednej cząstki.
Jakie jest prawdopodobieństwo, że populacja kiedyś wymrze?

\emph{Rozwiązanie.} Pierwsze zdarzenie jest śmiercią
z prawdopodobieństwem $2/5$ lub narodzinami z prawdopodobieństwem
$3/5$. Stąd $q=2/5+(3/5)q^2$, więc
$(3q-2)(q-1)=0$.
Wybieramy mniejsze rozwiązanie: $q=2/3$.
Średni dodatni przyrost liczby osobników nie gwarantuje
przetrwania konkretnego drzewa genealogicznego.
\end{zadanie}


\begin{zadanie}[Wzrost procesu Yule'a]
Niech $Z$ będzie procesem Yule'a o intensywności $\lambda>0$
z jednym osobnikiem początkowym. Wyznaczyć
$\Prob(Z_t\ge k)$, $\Var(Z_t)$ i
$\Prob(Z_t\ge3\mid Z_t\ge2)$ dla $t>0$, $k\ge1$.

\emph{Rozwiązanie.} Wprowadzamy $a=e^{-\lambda t}$.
Z udowodnionego wzoru $\Prob(Z_t=n)=a(1-a)^{n-1}$
otrzymujemy przez zsumowanie szeregu geometrycznego
$\Prob(Z_t\ge k)=(1-a)^{k-1}$.
Dla rozkładu geometrycznego na $\{1,2,\dots\}$
$\E Z_t=1/a$ i $\Var(Z_t)=(1-a)/a^2$, a więc
$\Var(Z_t)=e^{2\lambda t}-e^{\lambda t}$.
Wreszcie
\[
\Prob(Z_t\ge3\mid Z_t\ge2)
=\frac{(1-a)^2}{1-a}=1-e^{-\lambda t}.
\]
To przykład odrębnej własności braku pamięci
\emph{rozkładu geometrycznego liczebności},
a nie wykładniczego rozkładu czasów między skokami.
\end{zadanie}

\section{Podsumowanie i kierunek dalszej teorii}
Łańcuch Markowa opisuje przejścia między stanami;
proces odnowy opisuje cykle; proces rozgałęziający
opisuje losową liczbę \emph{potomków}, którzy sami
stają się źródłem kolejnych losowań. Funkcje tworzące
zamieniają złożoną strukturę drzew na iterację funkcji,
a martyngały pozwalają śledzić zmienność populacji
w kolejnych pokoleniach. Dalszym krokiem będą martyngały w czasie ciągłym, lokalizacja,
zatrzymywanie i wariacja kwadratowa. Całka Itô i równania
stochastyczne zostaną skonstruowane dopiero w tomie III.


% ================================================================
% ROZDZIAŁ 11. MARTYNGAŁY W CZASIE CIĄGŁYM
% ================================================================

\chapter{Martyngały w czasie ciągłym: informacja, regularność i lokalizacja}
\label{ch:mart-cont}

\section{Dlaczego przejście od czasu dyskretnego nie jest automatyczne?}
W tomie I (rozdział~\ref{I-ch:nadmartyngaly})
badaliśmy martyngały w czasie dyskretnym.
Martyngał był ciągiem $(M_n)_{n\ge0}$, którego
warunkowa prognoza następnej wartości, przy znajomości przeszłości,
wynosiła dokładnie wartość obecną. W czasie ciągłym rozważamy
nieprzeliczalną rodzinę chwil. Pojawiają się dwa nowe pytania:
jak wybierać jedną wspólną wersję trajektorii oraz co znaczy
zatrzymać proces w chwili wybranej na podstawie obserwacji?
Później zobaczymy też, że proces ze skokami może być martyngałem,
mimo że jego trajektorie mają skończoną wariację całkowitą.

\begin{definicja}[Martyngał, podmartyngał i nadmartyngał]
Niech $(\Omega,\calF,\Prob)$ będzie przestrzenią probabilistyczną,
$(\calF_t)_{t\ge0}$ filtracją i $M=(M_t)_{t\ge0}$ procesem
rzeczywistym, adaptowanym, takim że $\E|M_t|<\infty$
dla każdego $t\ge0$. Proces $M$ jest \emph{martyngałem}, jeśli
$\E[M_t\mid\calF_s]=M_s$ prawie na pewno dla wszystkich
$0\le s\le t$. Jest \emph{podmartyngałem}, jeśli
$\E[M_t\mid\calF_s]\ge M_s$, a \emph{nadmartyngałem},
gdy zachodzi nierówność przeciwna. Filtracja spełnia
\emph{warunki zwykłe}, jeśli jest zupełna (każdy podzbiór zbioru z $\calF$ o prawdopodobieństwie
zero należy do $\calF_0$) i prawostronnie ciągła:
$\calF_t=\bigcap_{u>t}\calF_u$ dla każdego $t\ge0$.
\end{definicja}

Ważne, że równość dla każdej pary $s,t$ jest rozumiana osobno
prawie na pewno. Nie wolno bez dodatkowego argumentu zażądać
jednego zbioru miary jeden, na którym równości warunkowe
miałyby zachodzić jednocześnie dla nieprzeliczalnie wielu par.

\begin{przyklad}[Trzy modele uczciwej ewolucji]
Niech $W$ będzie standardowym procesem Wienera w jego
uzupełnionej filtracji naturalnej. Wówczas
$W_t$ jest martyngałem z powodu niezależnych wycentrowanych
przyrostów (rozdział~\ref{ch:wiener}). Uzupełnienie filtracji
o zbiory zerowe nie zmienia równości warunkowych. Proces $W_t^2-t$ również nim jest, ponieważ
\[
\E[W_t^2-t\mid\calF_s]
=W_s^2+\E[(W_t-W_s)^2]-t=W_s^2-s.
\]
Jeśli $N$ jest procesem Poissona o intensywności $\lambda>0$
w swojej filtracji naturalnej (rozdział~\ref{ch:poisson-process}),
to $N_t-\lambda t$ jest martyngałem, gdyż przyrost liczby
zdarzeń po chwili $s$ jest niezależny od przeszłości i ma
średnią $\lambda(t-s)$. Pierwsze dwa procesy mają trajektorie
ciągłe, ostatni ma skoki wysokości jeden.
\end{przyklad}

\begin{lemat}[Warunkowe wartości oczekiwane są jednostajnie całkowalne]
\label{lem:mart-cont-ui-cond}
Niech $X\in L^1(\Omega,\calF,\Prob)$. Rodzina
$\{\E[X\mid\calG]:\calG\subseteq\calF\text{ jest pod-}\sigma\text{-ciałem}\}$
jest jednostajnie całkowalna.
\end{lemat}
\begin{proof}
Ustalmy $K>0$ i rozłóżmy $X=U+V$, gdzie
$U=X\chi_{\{|X|\le K/2\}}$ oraz
$V=X\chi_{\{|X|>K/2\}}$. Dla dowolnego $\calG$
połóżmy $Y=\E[X\mid\calG]$. Wówczas
$|\E[U\mid\calG]|\le K/2$, skąd na zdarzeniu
$\{|Y|>K\}$ zachodzi $|\E[V\mid\calG]|>K/2$
i $|Y|\le K/2+|\E[V\mid\calG]|\le2|\E[V\mid\calG]|$.
Dlatego
\[
\E[|Y|\chi_{\{|Y|>K\}}]
\le2\E|\E[V\mid\calG]|
\le2\E|V|
=2\E[|X|\chi_{\{|X|>K/2\}}].
\]
Prawa strona nie zależy od $\calG$ i dąży do zera
przy $K\to\infty$ dzięki całkowalności $X$.
\end{proof}

\begin{wniosek}[Jednostajna całkowalność na skończonym horyzoncie]
Niech $M$ będzie martyngałem oraz $T<\infty$.
Wtedy rodzina $\{M_t:0\le t\le T\}$ jest jednostajnie
całkowalna, gdyż $M_t=\E[M_T\mid\calF_t]$.
\end{wniosek}

\section{Prawostronne modyfikacje: skąd bierze się regularność?}
\begin{definicja}[Trajektoria c\`adl\`ag]
Funkcja $f:[0,\infty)\to\R$ jest c\`adl\`ag, jeśli
jest prawostronnie ciągła i ma skończone granice lewostronne
we wszystkich chwilach $t>0$. Proces jest c\`adl\`ag,
gdy prawie wszystkie jego trajektorie mają tę własność.
\end{definicja}

\begin{lemat}[Zbieżność odwrotnych warunkowych wartości oczekiwanych]
Niech $(\calG_n)_{n\ge1}$ będzie malejącym ciągiem
pod-$\sigma$-ciał przestrzeni $(\Omega,\calF,\Prob)$,
$\calG_\infty=\bigcap_n\calG_n$, a $X\in L^1$.
Wtedy
$\E[X\mid\calG_n]\to\E[X\mid\calG_\infty]$ w $L^1$.
\end{lemat}
\begin{proof}
Najpierw niech $X\in L^2$ i $Y_n=\E[X\mid\calG_n]$.
Dla $m\ge n$ mamy $Y_m=\E[Y_n\mid\calG_m]$.
Ortogonalność rzutu warunkowego daje
$\|Y_n-Y_m\|_2^2=\|Y_n\|_2^2-\|Y_m\|_2^2$.
Normy $\|Y_n\|_2$ maleją i są nieujemne, zatem
$(Y_n)$ jest ciągiem Cauchy'ego w $L^2$. Niech $Y$
będzie jego granicą. Wybierzmy podciąg $(Y_{n_j})$ zbieżny prawie na pewno
do $Y$, co jest możliwe dzięki zbieżności w $L^2$.
Aby uzyskać jedną wersję mierzalną względem
przecięcia także bez założenia zupełności $\calG_k$,
definiujemy $\widehat Y$ jako granicę tego podciągu
tam, gdzie istnieje skończona granica, a poza tym
jako zero. Dla każdego $k$ definicja zależy wyłącznie
od ogona podciągu o indeksach $n_j\ge k$, dlatego
$\widehat Y$ jest $\calG_k$-mierzalna dla wszystkich $k$
i w konsekwencji $\calG_\infty$-mierzalna.
Ponadto $\widehat Y=Y$ p.n.; dalej używamy tej wersji.
Dla $A\in\calG_\infty$ i wszystkich $n$ mamy
$\E[Y_n\chi_A]=\E[X\chi_A]$; przejście do granicy
w $L^1$ identyfikuje $Y=\E[X\mid\calG_\infty]$.
Dla ogólnego $X\in L^1$ przybliżamy $X$ obcięciami
$X^{(K)}\in L^2$. Kontrakcja warunkowania w $L^1$
pozwala przenieść zbieżność z $X^{(K)}$ na $X$ po
przejściu $K\to\infty$.
\end{proof}

\begin{lemat}[Dlaczego z ograniczonej liczby przekroczeń wynikają granice?]
Niech $D$ będzie gęstym podzbiorem przedziału $[0,T]$,
a $f:D\to\R$ funkcją ograniczoną. Załóżmy, że dla każdej
pary liczb wymiernych $a<b$ funkcja $f$, oglądana w
rosnącym porządku chwil z $D$, ma skończenie wiele pełnych
przekroczeń przedziału $[a,b]$ w górę. Wtedy w każdej chwili
$t<T$ istnieje skończona prawa granica
$\lim_{r\downarrow t,\,r\in D}f(r)$, a w każdej chwili
$t>0$ istnieje skończona lewa granica wzdłuż $D$.
Funkcja zdefiniowana przez prawe granice (z wartością $f(T)$
w końcu przedziału, gdy $T\in D$) jest c\`adl\`ag.
\end{lemat}
\begin{proof}
Brak prawej granicy w chwili $t$ oznaczałby
$\liminf_{r\downarrow t,\,r\in D}f(r)
<\limsup_{r\downarrow t,\,r\in D}f(r)$.
Można wtedy wybrać wymierne $a<b$ pomiędzy tymi
liczbami oraz dowolnie długie skończone ciągi chwil
z $D\cap(t,t+\delta)$, w których wartości naprzemiennie
spadają poniżej $a$ i rosną powyżej $b$.
Po ustawieniu chwil rosnąco otrzymujemy dowolnie wiele
przekroczeń w górę, sprzecznie z założeniem.
To samo rozumowanie, z chwilami zbliżającymi się do $t$
od lewej, daje istnienie granicy lewostronnej.
Ograniczoność $f$ zapewnia skończoność obu granic.

Oznaczmy prawą granicę przez $g(t)$. Gdyby dla
$t_n\downarrow t$ zachodziło $g(t_n)\not\to g(t)$,
można by znaleźć $r_n\in D$ położone tuż za $t_n$,
tak by $r_n\downarrow t$ oraz
$|f(r_n)-g(t_n)|\to0$; byłaby to sprzeczność z
definicją $g(t)$. Analogicznie, dla $t_n\uparrow t$
wybieramy $r_n\in D$ między $t_n$ i $t$.
Istnienie lewej granicy $f$ w $t$ pociąga wówczas
istnienie tej samej lewej granicy dla $g$.
Zatem $g$ jest prawostronnie ciągła i ma granice
lewostronne.
\end{proof}

\begin{twierdzenie}[Regularizacja martyngału typu zamkniętego]
\label{thm:regularizacja}
Niech filtracja $(\calF_t)_{t\ge0}$ spełnia warunki zwykłe,
a $X\in L^1(\Omega,\calF,\Prob)$. Proces
$M_t=\E[X\mid\calF_t]$ ma modyfikację c\`adl\`ag.
Dokładniej, istnieje proces $\widetilde M$ o trajektoriach
c\`adl\`ag, taki że dla każdego $t\ge0$
$\widetilde M_t=M_t$ prawie na pewno.
\end{twierdzenie}
\begin{proof}
Ustalmy skończony horyzont $T>0$ i wybierzmy
$D_T=(\Q\cap[0,T])\cup\{T\}$.
Wybieramy rosnący ciąg skończonych podzbiorów $D_T$
zawierających $0,T$ i wyczerpujących $D_T$.
Dla dowolnych liczb
wymiernych $a<b$ i skończonego uporządkowanego podzbioru
$D_T$ zastosujmy nierówność przekroczeń do dyskretnego
martyngału otrzymanego przez próbkowanie $(M_t)$ w tych
chwilach. Oczekiwana liczba przekroczeń przedziału
$[a,b]$ jest ograniczona przez stałą zależną od
$a,b$ i $\E|X|$, a nie od liczby punktów próbkowania:
taki wniosek wynika z nierówności przekroczeń
Dooba z tomu I, rozdział~\ref{I-ch:nadmartyngaly}:
dla martyngału na skończonej siatce
$(b-a)\E U[a,b]\le\E[(M_T-a)^-]$
i prawa strona jest nie większa od $\E|X|+|a|$.
Tutaj $U[a,b]$ oznacza liczbę pełnych przekroczeń
w górę, a $y^-=\max(-y,0)$.
Przy dodawaniu punktów siatki liczba przekroczeń
nie maleje (definiujemy ją jako maksymalną
liczbę rozłącznych kolejnych par punktów).
Zbieżność monotoniczna zapewnia skończoną liczbę
przekroczeń wzdłuż $D_T$ prawie na pewno.
Ponadto $|M|$ jest podmartyngałem przez nierówność
Jensena. Słaba nierówność maksymalna Dooba
na każdej skończonej siatce daje dla $K>0$
$\Prob(\max |M_r|>K)\le \E|M_T|/K\le\E|X|/K$.
Przechodząc do całego $D_T$, a następnie z $K$
do nieskończoności, uzyskujemy
$\sup_{r\in D_T}|M_r|<\infty$ prawie na pewno.
Na wspólnym zbiorze pełnej miary, wyznaczonym
przez te warunki dla przeliczalnie wielu par
wymiernych $a<b$, stosujemy poprzedni lemat.
Dostajemy skończone prawe granice dla $0\le s<T$
i lewe dla $0<s\le T$.

Na tym zbiorze definiujemy
$\widetilde M_s=\lim_{r\downarrow s,r\in D_T}M_r$
dla $s<T$ oraz $\widetilde M_T=M_T$.
Na jego dopełnieniu przyjmujemy zerową trajektorię. Poprzedni lemat
o funkcjach rzeczywistych z granicami jednostronnymi
zapewnia, że $s\mapsto\widetilde M_s$ jest prawostronnie
ciągła i ma granice z lewej: gdyby prawostronna ciągłość
zawiodła, powstałyby dwie sekwencje punktów zbieżne do
tej samej chwili o różnych granicach, sprzeczne z istnieniem
wspólnej granicy prawostronnej próbek z $D_T$.

Aby sprawdzić, że to modyfikacja właściwego procesu, dla
każdego ustalonego $s<T$ wybieramy $r_n\in D_T$ malejące
do $s$. Z lematu o malejących $\sigma$-ciałach i prawostronnej
ciągłości filtracji
$M_{r_n}=\E[X\mid\calF_{r_n}]\to
\E[X\mid\calF_s]=M_s$ w $L^1$.
Z drugiej strony lewa strona zbiega prawie na pewno do
$\widetilde M_s$. Stąd $\widetilde M_s=M_s$ prawie na pewno.
Dla ustalonego $s$ granica wzdłuż malejących
$r_n$ jest $\bigcap_n\calF_{r_n}=\calF_s$-mierzalna:
można ją zdefiniować jako granicę ogona tam,
gdzie istnieje skończona granica, i zero w przeciwnym
razie. Zmiana na wspólnym zbiorze zerowym jest
$\calF_0$-mierzalna dzięki zupełności filtracji.
Zatem $\widetilde M$ jest adaptowany.

Łącząc konstrukcje na horyzontach $T=1,2,\ldots$,
bierzemy wspólny zbiór pełnej miary, na którym
istnieją wszystkie potrzebne granice oraz
w każdym całkowitym punkcie końcowym wartość
$M_T$ jest zgodna z prawą granicą uzyskaną
na dłuższym horyzoncie. Ten ostatni warunek
wynika z udowodnionej własności modyfikacji
i przeliczalności całkowitych $T$.
W punktach wewnętrznych zgodność jest bezpośrednia,
bo używamy tych samych wymiernych chwil.
Poza wspólnym zbiorem pełnej miary definiujemy
zerową trajektorię. Otrzymujemy adaptowany proces
o wszędzie c\`adl\`ag trajektoriach, który jest
modyfikacją $M$, a przez równości w ustalonych
chwilach również martyngałem.
Każde dwie prawostronnie ciągłe modyfikacje
są nierozróżnialne, gdyż zgadzają się p.n.
w każdej z przeliczalnie wielu chwil wymiernych,
a prawostronna ciągłość rozszerza zgodność
na wszystkie chwile jednocześnie.
\end{proof}

\begin{twierdzenie}[Ogólne twierdzenie regularizacyjne Dooba]
\label{thm:doob-general-regularization}
Niech $(\calF_t)$ spełnia warunki zwykłe, a $X$ będzie
rzeczywistym podmartyngałem. Istnieje adaptowana
modyfikacja c\`adl\`ag wtedy i tylko wtedy, gdy
$t\mapsto\E X_t$ jest prawostronnie ciągła.
Analogiczna teza zachodzi dla nadmartyngałów.
W szczególności każdy martyngał ma taką modyfikację.
\end{twierdzenie}
\begin{proof}
Najpierw wykażemy istnienie granic trajektorii
na chwilach wymiernych, niezależnie od założenia
o ciągłości średniej. Dla $T>0$ połóżmy
$D_T=(\Q\cap[0,T])\cup\{T\}$.
Na każdej skończonej siatce z końcem $T$
nierówność przekroczeń Dooba (tom I,
rozdział~\ref{I-ch:nadmartyngaly}) daje
$(b-a)\E U[a,b]\le \E[(X_T-a)^+]+
\E[(X_0-a)^-]$ (wystarcza to oszacowanie).
Liczba przekroczeń wzdłuż całego $D_T$
jest więc skończona p.n. dla każdej wymiernej
pary $a<b$.

Aby uzyskać także ograniczoność próbek, połóżmy
$C_u=\E[X_T\mid\calF_u]$ i $Y_u=C_u-X_u$
dla $0\le u\le T$. Z własności podmartyngału
$Y_u\ge0$ p.n., a $Y$ jest nadmartyngałem.
Na skończonej siatce maksymalna nierówność
dla nieujemnego nadmartyngału daje
$\Prob(\max Y_u>K)\le\E Y_0/K$:
zatrzymujemy na pierwszym przekroczeniu poziomu
$K$ (lub końcu siatki), stosujemy dyskretne
próbkowanie opcjonalne i nieujemność.
Nierówność maksymalna dla $|C|$ daje analogiczne
oszacowanie przez $\E|X_T|/K$.
Wyczerpując $D_T$ skończonymi siatkami, otrzymujemy
$\sup_{u\in D_T}|X_u|<\infty$ p.n.
Lemat o przekroczeniach pozwala zatem zdefiniować
prawostronne granice $\widetilde X_u$ wzdłuż $D_T$
dla $u<T$, z trajektoriami c\`adl\`ag po dołączeniu
wartości $X_T$ na końcu. Jak w dowodzie
twierdzenia~\ref{thm:regularizacja}, warunki
zwykłe zapewniają adaptację tej konstrukcji.

Potrzebujemy teraz zbieżności w $L^1$ przy
$u_j\downarrow u$. Oznaczmy $V_j=X_{u_j}$,
$\mathcal G_j=\calF_{u_j}$ i
$A_j=\E[V_j\mid\mathcal G_{j+1}]-V_{j+1}\ge0$
p.n. Z własności wieży dostajemy indukcyjnie
\[
V_n=\E\left[V_1-\sum_{j=1}^{n-1}A_j
\,\middle|\,\mathcal G_n\right].
\]
Istotnie $A_j$ jest $\mathcal G_{j+1}$-mierzalna,
a $\mathcal G_{j+1}\subseteq\mathcal G_j$;
kolejne warunkowanie powyższej równości i
odjęcie $A_n$ daje następny krok.
Ponieważ $\E V_n\ge\E X_u$, mamy
\[
\sum_{j\ge1}\E A_j
=\E V_1-\lim_n\E V_n
\le \E V_1-\E X_u<\infty.
\]
Zatem $A=\sum_{j\ge1}A_j$ jest całkowalna
i $V_n$ różni się od
$\E[V_1-A\mid\mathcal G_n]$ w normie $L^1$
o co najwyżej $\E\sum_{j\ge n}A_j\to0$.
Lemat o odwrotnych warunkowych wartościach
oczekiwanych daje zbieżność $V_n$ w $L^1$.
Jeżeli $u_j\in D_T$, jej granica jest
$\widetilde X_u$, ponieważ istnieje granica p.n.

W szczególności $\E\widetilde X_u=
\lim_{j\to\infty}\E X_{u_j}$ oraz
$\E[V_j\mid\calF_u]\ge X_u$.
Przechodząc do granicy w $L^1$ i korzystając
z $\calF_u$-mierzalności $\widetilde X_u$,
dostajemy $\widetilde X_u\ge X_u$ p.n.
Gdy średnia jest prawostronnie ciągła,
różnica ta ma średnią zero, więc
$\widetilde X_u=X_u$ p.n.
Łączenie horyzontów i wybór zerowej trajektorii
na wspólnym zbiorze zerowym odbywa się
jak w twierdzeniu~\ref{thm:regularizacja}.

Odwrotnie, jeśli istnieje modyfikacja
prawostronnie ciągła, dla każdego ciągu
$u_j\downarrow u$ mamy $X_{u_j}\to X_u$ p.n.
(zgodność wersji wystarczy w przeliczalnym
zbiorze tych chwil). Udowodniona powyżej
zbieżność w $L^1$ dla takiego ciągu pokazuje
$\E X_{u_j}\to\E X_u$. Ponieważ średnia
podmartyngału jest niemalejąca, to oznacza
jej prawostronną ciągłość.
Dla nadmartyngałów stosujemy wynik do $-X$,
a dla martyngału średnia jest stała.
\end{proof}

\begin{uwaga}[Dlaczego wymagana jest prawostronna ciągłość filtracji?]
Bez warunku $\calF_{t+}=\calF_t$ granica
$\E[X\mid\calF_r]$ przy $r\downarrow t$ może być
warunkową wartością względem większego $\sigma$-ciała
$\calF_{t+}$, a nie względem $\calF_t$.
Uzyskujemy wtedy proces prawostronnie ciągły, który w chwili
$t$ może wykorzystywać informację niedostępną w $\calF_t$.
To jest powód, dla którego warunki zwykłe są podstawowym
założeniem teorii w czasie ciągłym.
\end{uwaga}


\begin{przyklad}[Dlaczego nie wystarcza lewostronna informacja?]
Niech $\Omega=\{-1,1\}$, obydwa wyniki mają prawdopodobieństwo
$1/2$, a $X(\omega)=\omega$. Przyjmijmy $\calF_0=\{\varnothing,\Omega\}$
oraz $\calF_t=\sigma(X)$ dla $t>0$. Wtedy
$M_t=\E[X\mid\calF_t]$ ma postać $M_0=0$ i $M_t=X$ dla $t>0$.
Jest to martyngał, ale nie posiada modyfikacji prawostronnie ciągłej
w chwili zero. Istotnie, dla wszystkich wymiernych $r>0$
każda modyfikacja musiałaby mieć $\widetilde M_r=X$ poza jednym
wspólnym zbiorem zerowym, a więc prawa granica w zerze wynosiłaby
$X\ne0=\widetilde M_0$ prawie na pewno. Filtracja nie jest
prawostronnie ciągła w zerze, ponieważ
$\calF_{0+}=\bigcap_{t>0}\calF_t=\sigma(X)\supsetneq\calF_0$.
Ten przykład ujawnia, że regularizacja jest twierdzeniem o
\emph{procesie oraz informacji}, a nie wyłącznie o funkcjach czasu.
\end{przyklad}

W dalszej części rozdziału przyjmujemy warunki zwykłe
filtracji, chyba że wyraźnie zaznaczono inaczej.
Adaptowany proces o prawie na pewno c\`adl\`ag
trajektoriach można na mierzalnym zbiorze zerowym
zastąpić zerową trajektorią, zachowując adaptację
i nierozróżnialność. Dlatego w argumentach
trajektoryjnych możemy używać wszędzie c\`adl\`ag
wersji.

\section{Czasy zatrzymania i proces zatrzymany}

\begin{lemat}[Mierzalność wartości procesu zatrzymanego]
Niech $X=(X_t)_{t\ge0}$ będzie rzeczywistym procesem adaptowanym
o wszędzie prawostronnie ciągłych trajektoriach, a $\tau$ czasem zatrzymania.
Dla każdego $t\ge0$ zmienna $X_{t\wedge\tau}$ jest $\calF_t$-mierzalna.
\end{lemat}
\begin{proof}
Ustalmy $t>0$. Przybliżamy $\tau\wedge t$ od góry przez czasy
$\rho_n=2^{-n}\lceil2^n(\tau\wedge t)\rceil\wedge t$, które
przyjmują skończenie wiele wartości z $[0,t]$ i są czasami
zatrzymania. Dla każdego $n$ mamy
$X_{\rho_n}=\sum_{u\in D_n}X_u\chi_{\{\rho_n=u\}}$,
gdzie $D_n$ jest skończonym zbiorem wartości $\rho_n$.
Każdy składnik jest $\calF_t$-mierzalny (ponieważ $u\le t$),
a $\rho_n\downarrow\tau\wedge t$. Prawostronna ciągłość daje
$X_{\rho_n}\to X_{\tau\wedge t}$ punktowo, a więc
ostatnia zmienna jest mierzalna względem $\calF_t$
bez potrzeby uzupełniania filtracji.
Dla $t=0$ jest to adaptowana zmienna $X_0$.
Przy trajektoriach prawostronnie ciągłych jedynie
p.n. ten sam argument działa przy zupełnej filtracji.
\end{proof}


\begin{definicja}[Proces zatrzymany]
Niech $X$ będzie procesem adaptowanym c\`adl\`ag,
a $\tau$ czasem zatrzymania (dopuszczamy wartość $\infty$).
Dla $t\ge0$ definiujemy $X_t^\tau=X_{t\wedge\tau}$.
Gdy $\tau=\infty$, przyjmujemy $t\wedge\tau=t$.
\end{definicja}

\begin{twierdzenie}[Zatrzymanie w ustalonej chwili]
Niech $M$ będzie martyngałem c\`adl\`ag, a $\tau$
czasem zatrzymania. Wtedy $M^\tau$ jest martyngałem.
Jeżeli $X$ jest podmartyngałem (odpowiednio nadmartyngałem)
c\`adl\`ag \emph{klasy $(DL)$}, to $X^\tau$ pozostaje
podmartyngałem (odpowiednio nadmartyngałem).
Tutaj klasa $(DL)$ oznacza, że dla każdego $T<\infty$
rodzina $\{X_\sigma: \sigma\le T,\ \sigma
\text{ jest czasem zatrzymania}\}$ jest jednostajnie
całkowalna.
\end{twierdzenie}
\begin{proof}
Adaptacja wynika z poprzedniego lematu, a trajektorie
zatrzymane są c\`adl\`ag: przed zatrzymaniem
zgadzają się z pierwotną trajektorią, a od chwili
zatrzymania są stałe. Ustalmy $0\le s\le t$ i czas
$\rho=s\vee(t\wedge\tau)$. Ponieważ
$s\le\rho\le t$, twierdzenie o ograniczonym
próbkowaniu opcjonalnym (twierdzenie~\ref{thm:optional-cont})
daje całkowalność $M_\rho,M_{t\wedge\tau},M_{s\wedge\tau}$
oraz $\E[M_\rho\mid\calF_s]=M_s$.
Na zdarzeniu $\{\tau>s\}\in\calF_s$ mamy
$\rho=t\wedge\tau$, a na $\{\tau\le s\}$
proces zatrzymany jest już stały. Dlatego
\[
\begin{aligned}
\E[M_{t\wedge\tau}\mid\calF_s]
&=\chi_{\{\tau\le s\}}M_{s\wedge\tau}
 +\chi_{\{\tau>s\}}\E[M_\rho\mid\calF_s]\\
&=M_{s\wedge\tau}.
\end{aligned}
\]
Dla podmartyngału $X$ klasy $(DL)$ przybliżamy
$\rho$ od góry skończonymi czasami siatkowymi
$\rho_n=2^{-n}\lceil 2^n\rho\rceil\wedge t$,
do siatki włączając $s$ i $t$.
Mamy $s\le\rho_n\le t$. Próbkowanie opcjonalne
w czasie dyskretnym na tej siatce daje
$\E[X_{\rho_n}\mid\calF_s]\ge X_s$.
Po przejściu do granicy otrzymujemy
$\E[X_\rho\mid\calF_s]\ge X_s$.
Stosując ten sam rozkład na zdarzenia
$\{\tau\le s\}$ i $\{\tau>s\}$ co dla martyngału,
dostajemy
$\E[X_{t\wedge\tau}\mid\calF_s]\ge X_{s\wedge\tau}$.
Istotne jest przejście do granicy w $L^1$:
z klasy $(DL)$ otrzymujemy jednostajną całkowalność
wartości w przybliżających czasach zatrzymania,
a prawostronna ciągłość daje ich zbieżność prawie na pewno.
Nierówność przechodzi więc do granicy w warunkowych
wartościach oczekiwanych. Dla nadmartyngału zmieniamy
znak procesu i powtarzamy argument. Samo formalne
podstawienie losowej chwili w definicji podmartyngału
nie jest dowodem.
\end{proof}

\begin{definicja}[Klasy $(D)$ i $(DL)$]
Adaptowany proces c\`adl\`ag jest klasy $(D)$,
jeżeli rodzina wszystkich jego wartości w
ograniczonych czasach zatrzymania jest jednostajnie
całkowalna (bez wspólnego deterministycznego
ograniczenia tych czasów).
Jest klasy $(DL)$, jeżeli dla każdego $T<\infty$
rodzina wartości w czasach zatrzymania $\sigma\le T$
jest jednostajnie całkowalna.
\end{definicja}

\begin{twierdzenie}[Klasy całkowalności martyngału]
\label{thm:mart-cont-class-d}
Każdy martyngał c\`adl\`ag jest klasy $(DL)$.
Dla martyngału c\`adl\`ag równoważne są warunki:
\begin{enumerate}
\item jest klasy $(D)$;
\item rodzina $\{M_t:t\ge0\}$ jest jednostajnie całkowalna;
\item istnieje $X\in L^1$ mierzalna względem
$\calF_\infty=\sigma(\bigcup_{t\ge0}\calF_t)$
taka, że $M_t=\E[X\mid\calF_t]$ p.n. dla każdego $t$.
\end{enumerate}
W ostatnim przypadku dla każdego p.n. skończonego
czasu zatrzymania $\sigma$ zachodzi
$M_\sigma=\E[X\mid\calF_\sigma]$ p.n.
\end{twierdzenie}
\begin{proof}
Dla $\sigma\le T$ ograniczone próbkowanie
opcjonalne (twierdzenie~\ref{thm:optional-cont})
daje $M_\sigma=\E[M_T\mid\calF_\sigma]$.
Lemat~\ref{lem:mart-cont-ui-cond} zapewnia
jednostajną całkowalność na ustalonym horyzoncie.
Z $(D)$ wynika drugi warunek przez wybór czasów
deterministycznych.
Przy drugim warunku dyskretny martyngał $(M_n)$
jest jednostajnie całkowalny, więc zbiega w $L^1$
do $X\in L^1(\calF_\infty)$
(tom I, rozdział~\ref{I-ch:nadmartyngaly}).
Dla każdego ustalonego $t$ i całkowitego $n\ge t$
mamy $M_t=\E[M_n\mid\calF_t]$.
Przejście do granicy w $L^1$ daje trzeci warunek.
Jeśli zachodzi trzeci warunek, dla ograniczonego
$\sigma\le T$ z własności wieży i próbkowania
opcjonalnego
$M_\sigma=\E[M_T\mid\calF_\sigma]
=\E[X\mid\calF_\sigma]$.
Stąd $(D)$ z tego samego lematu.

Dla p.n. skończonego $\sigma$ mamy
$M_{\sigma\wedge n}\to M_\sigma$ p.n.
(ostatecznie $\sigma\wedge n=\sigma$),
a klasa $(D)$ daje również zbieżność w $L^1$.
Dla $A\in\calF_\sigma$ zdarzenie
$A_n=A\cap\{\sigma\le n\}$ należy do
$\calF_{\sigma\wedge n}$, co wynika bezpośrednio
z definicji $\sigma$-ciała w czasie zatrzymania.
Zatem
$\E[M_{\sigma\wedge n}\chi_{A_n}]=\E[X\chi_{A_n}]$.
Przejście do granicy i mierzalność $M_\sigma$
(twierdzenie~\ref{thm:sampled}) identyfikują
żądaną warunkową wartość oczekiwaną.
\end{proof}

\begin{uwaga}[Nie mylić zatrzymania z granicą $M_\tau$]
Powyższe twierdzenie dotyczy procesu $M_{t\wedge\tau}$
przy każdym skończonym $t$. Nie wolno z niego wnioskować
$\E M_\tau=\E M_0$ dla nieograniczonego czasu $\tau$:
przejście $t\to\infty$ wymaga dodatkowego warunku
jednostajnej całkowalności.
\end{uwaga}

\begin{zadanie}[Co odróżnia martyngał od procesu o niemalejącej średniej?]
Niech $W$ będzie standardowym procesem Wienera w jego
filtracji naturalnej. Sprawdzić bez rachunku It\^o, które
z procesów $W_t$, $W_t^2$ i $W_t^2-t$ są martyngałami.

\emph{Rozwiązanie.} Dla $s<t$ zapiszmy
$W_t=W_s+U$, gdzie $U\sim\Normal(0,t-s)$
i $U$ jest niezależne od $\calF_s$.
Stąd $\E[W_t\mid\calF_s]=W_s$.
Po podniesieniu do kwadratu i warunkowaniu otrzymujemy
\[
\E[W_t^2\mid\calF_s]
=W_s^2+2W_s\E U+\E U^2=W_s^2+(t-s).
\]
Proces $W_t^2$ jest zatem podmartyngałem, ale nie
martyngałem. Odejmując od niego deterministyczny
proces $t$, dostajemy
$\E[W_t^2-t\mid\calF_s]=W_s^2-s$.
Wyraz $t$ jest dokładnie kompensacją narastającej
wariancji, którą w kolejnym rozdziale opiszemy za
pomocą nawiasu procesu.
\end{zadanie}

\section{Martyngały lokalne}
\begin{definicja}[Lokalny martyngał i lokalizacja]
Niech $M$ będzie adaptowanym procesem c\`adl\`ag o
wartościach rzeczywistych. Nazywamy go \emph{lokalnym
martyngałem}, jeżeli istnieją czasy zatrzymania
$\tau_k\uparrow\infty$ prawie na pewno, takie że dla
każdego $k$ proces $M^{\tau_k}$ jest martyngałem.
Ciąg $(\tau_k)$ nazywamy \emph{ciągiem lokalizującym}.
W szczególności $M_0\in L^1$, ponieważ
$M^{\tau_k}_0=M_0$. Można zawsze wymagać
$\tau_k\le k$: zastąpienie $\tau_k$ przez
$\tau_k\wedge k$ zachowuje własność martyngałową
na mocy twierdzenia o zatrzymywaniu i nadal daje
ciąg rosnący do nieskończoności p.n.
\end{definicja}

\begin{twierdzenie}[Zatrzymanie ograniczonego procesu lokalnego]
Niech $M$ będzie ciągłym lokalnym martyngałem z $M_0=0$,
a $\tau_K=\inf\{t\ge0:|M_t|\ge K\}$ dla $K>0$.
Wówczas $M^{\tau_K}$ jest martyngałem ograniczonym
przez $K$. Czasy $\tau_K\uparrow\infty$ prawie na pewno,
gdy $K\to\infty$.
\end{twierdzenie}
\begin{proof}
Czasy $\tau_K$ są czasami zatrzymania, ponieważ
dla ciągłej trajektorii i $t\ge0$
$\{\tau_K\le t\}=\{\sup_{0\le u\le t}|M_u|\ge K\}$,
a supremum można obliczyć po chwilach wymiernych
z $[0,t]$ wraz z punktem $t$.
Niech $\sigma_n$ lokalizuje $M$. Proces
$M^{\tau_K\wedge\sigma_n}$ jest martyngałem na mocy
próbkowania opcjonalnego zastosowanego do
$M^{\sigma_n}$: zatrzymujemy go dodatkowo w $\tau_K$.
Ciągłość $M$ daje $|M_{t\wedge\tau_K\wedge\sigma_n}|\le K$.
Dla ustalonego $t$ mamy
$M_{t\wedge\tau_K\wedge\sigma_n}\to M_{t\wedge\tau_K}$
prawie na pewno, a wspólne ograniczenie przez $K$
pozwala przejść do granicy w $L^1$.
Dla $s\le t$ przechodzimy do granicy w warunkowych
równościach martyngałowych, używając kontrakcji warunkowania
w $L^1$. Otrzymujemy martyngał $M^{\tau_K}$.
Na każdym ograniczonym przedziale trajektoria ciągła jest
ograniczona, więc dla dużego $K$ nie przekracza poziomu
$K$ na tym przedziale. To dowodzi $\tau_K\uparrow\infty$.
\end{proof}

\begin{zadanie}[Kompensowany proces Poissona i momenty]
Niech $N$ będzie procesem Poissona o intensywności
$\lambda>0$ oraz $M_t=N_t-\lambda t$. Wykazać,
że $M$ jest martyngałem, i wyznaczyć jego wariancję.

\emph{Rozwiązanie.} Dla $s<t$ przyrost
$N_t-N_s$ ma rozkład Poissona z parametrem
$\lambda(t-s)$ i jest niezależny od $\calF_s$.
Stąd $\E[M_t-M_s\mid\calF_s]=0$. Ponadto
$\Var(M_t)=\Var(N_t)=\lambda t$.
W przeciwieństwie do $W_t$, trajektoria $M$ jest
odcinkami liniowa i ma skoki wielkości $1$.
\end{zadanie}


\begin{twierdzenie}[Nieujemny lokalny martyngał jest nadmartyngałem]
\label{thm:positive-local-supermart}
Niech $M=(M_t)_{t\ge0}$ będzie rzeczywistym lokalnym
martyngałem c\`adl\`ag względem filtracji $(\calF_t)$,
$M_t\ge0$ dla każdego $t$ poza jednym zbiorem zerowym,
oraz $\E M_0<\infty$. Wówczas $M_t\in L^1$ dla wszystkich $t$
i dla $0\le s\le t$ zachodzi
$\E[M_t\mid\calF_s]\le M_s$ prawie na pewno.
W szczególności $\E M_t\le\E M_0$.
\end{twierdzenie}
\begin{proof}
Niech $(\tau_k)$ lokalizuje $M$. Dla każdego $k$ proces
$M^{\tau_k}$ jest martyngałem, toteż
$\E[M_{t\wedge\tau_k}\mid\calF_s]=M_{s\wedge\tau_k}$.
Z warunkowego lematu Fatou dla nieujemnych zmiennych dostajemy
\[
\E[M_t\mid\calF_s]
\le\liminf_{k\to\infty}\E[M_{t\wedge\tau_k}\mid\calF_s]
=\lim_{k\to\infty}M_{s\wedge\tau_k}=M_s.
\]
W tym rachunku lewą stronę początkowo rozumiemy jako rozszerzoną
warunkową wartość oczekiwaną. Niezależnie, z bezwarunkowego
lematu Fatou i $\E M_{t\wedge\tau_k}=\E M_0$ otrzymujemy
$\E M_t\le\E M_0<\infty$ dla każdego $t$.
Wobec tego wszystkie występujące warunkowe wartości
oczekiwane są zwykłymi zmiennymi całkowalnymi, a poprzednia
nierówność dowodzi własności nadmartyngału.
Nierówność może być ostra: lokalizacja zapewnia uczciwość
\emph{przed zatrzymaniem}, ale nie gwarantuje jednostajnej
całkowalności potrzebnej do zachowania średniej po usunięciu
lokalizacji.
\end{proof}

\begin{twierdzenie}[Kryteria dla martyngałów lokalnych]
\label{thm:local-mart-criteria}
Niech $M$ będzie lokalnym martyngałem c\`adl\`ag.
\begin{enumerate}
\item Dla każdego czasu zatrzymania $\rho$
proces $M^\rho$ jest lokalnym martyngałem.
\item Suma i kombinacje liniowe lokalnych
martyngałów są lokalnymi martyngałami.
\item Jeżeli $M$ jest klasy $(DL)$, to jest martyngałem.
Wystarczy również, że dla każdego $T<\infty$
zmienna $\sup_{t\le T}|M_t|$ jest całkowalna.
\item Jeżeli $M\ge0$, to jest martyngałem wtedy
i tylko wtedy, gdy $\E M_t=\E M_0$ dla każdego $t$.
\end{enumerate}
\end{twierdzenie}
\begin{proof}
Niech $(\tau_n)$ lokalizuje $M$.
Zatrzymując martyngał $M^{\tau_n}$ dodatkowo
w $\rho$, dostajemy martyngał
$(M^\rho)^{\tau_n}$, co dowodzi pierwszej tezy.
Dla dwóch procesów z ciągami lokalizującymi
$(\tau_n)$ i $(\sigma_n)$ wspólny ciąg
$(\tau_n\wedge\sigma_n)$ nadal rośnie do
nieskończoności p.n. i lokalizuje ich kombinację
liniową.
Przy klasie $(DL)$ dla ustalonego $t$
rodzina $M_{t\wedge\tau_n}$ jest jednostajnie
całkowalna. Zbieżność p.n. do $M_t$ (ostateczna
równość) daje zbieżność w $L^1$, a stąd
równości martyngałowe przez kontrakcję warunkowania.
Całkowalny supremum na każdym skończonym
horyzoncie daje klasę $(DL)$ przez dominację.
Jeżeli $M\ge0$, to z twierdzenia
\ref{thm:positive-local-supermart}
$\E[M_t\mid\calF_s]\le M_s$.
Przy stałych średnich nieujemna różnica
$M_s-\E[M_t\mid\calF_s]$ ma średnią zero,
więc znika p.n. Odwrotna implikacja jest
bezpośrednia.
\end{proof}

\begin{uwaga}[Różnica między lokalnym martyngałem a martyngałem]
Każdy martyngał c\`adl\`ag jest lokalnym martyngałem
(można wziąć $\tau_k=k$), ale implikacja odwrotna nie jest
ogólnie prawdziwa. Poprzednie twierdzenie pokazuje jedynie,
że nieujemny lokalny martyngał jest nadmartyngałem;
aby stwierdzić równość warunkowych wartości oczekiwanych,
trzeba jeszcze kontrolować brak utraty masy przy lokalizacji.
\end{uwaga}


\begin{twierdzenie}[Odwrotność trójwymiarowego procesu Bessela]
\label{thm:inverse-bessel-strict-local}
Niech $B$ będzie standardowym ruchem Browna
w $\R^3$, $x\ne0$, $r=|x|$ i
$R_t=|x+B_t|$. Proces $R$ nazywamy
trójwymiarowym procesem Bessela rozpoczętym od $r$.
W jego filtracji naturalnej (uzupełnionej)
oraz w filtracji ruchu Browna spełniającej
warunki zwykłe i zachowującej niezależność
przyrostów proces $M_t=1/R_t$ jest dodatnim
ciągłym lokalnym martyngałem, ale nie jest
martyngałem. Dla $t>0$
\[
\E M_t=\frac{2\Phi(r/\sqrt t)-1}{r}
<\frac1r=M_0,
\]
gdzie $\Phi$ jest dystrybuantą standardowego
rozkładu normalnego.
\end{twierdzenie}
\begin{proof}
Z twierdzenia o trafianiu w punkty przez ruch
Browna w rozdziale~\ref{ch:brown-multi}
wynika, że $x+B_t$ prawie na pewno nigdy nie
osiąga zera. Na wspólnym zbiorze zerowym
możemy przyjąć stałą trajektorię $R_t=r$.
Dla $n$ tak dużego, że $1/n<r<n$, połóżmy
\[
\tau_n=\inf\{t\ge0:R_t\notin(1/n,n)\}.
\]
Są to czasy zatrzymania względem obu filtracji,
przez ciągłość trajektorii i twierdzenie
o trafianiu w zbiory domknięte
(twierdzenie~\ref{thm:hitting-closed}).
Na każdym zwartym przedziale $R$ ma dodatnie
minimum i skończone maksimum, więc
$\tau_n\uparrow\infty$ p.n.

Funkcja $h(y)=|y|^{-1}$ jest harmoniczna
na $\R^3\setminus\{0\}$: dla $u=|y|>0$
jej radialny laplasjan wynosi
$h''(u)+2h'(u)/u=2/u^3-2/u^3=0$
(por. twierdzenie~\ref{thm:brown-annulus}).
Wybieramy gładką funkcję o zwartym nośniku
$h_n$ równą $h$ na otoczeniu domkniętego
pierścienia $\{1/n\le |y|\le n\}$,
mnożąc $h$ przez gładką funkcję odcinającą
równą jeden na tym otoczeniu i zero blisko
zera oraz poza dostatecznie dużą kulą.
Lemat~\ref{lem:brown-generator-martingale}
daje ciągły martyngał
\[
h_n(x+B_t)-h_n(x)
-\frac12\int_0^t\Delta h_n(x+B_v)\,\mathrm d\lambda_1(v).
\]
Po zatrzymaniu w $\tau_n$ całka znika,
gdyż $\Delta h_n=0$ na pierścieniu,
a pozostała różnica wynosi
$M_{t\wedge\tau_n}-1/r$.
Zatem $M^{\tau_n}$ jest martyngałem
w filtracji ruchu Browna. W mniejszej filtracji
$R$ pozostaje nim przez własność wieży
i adaptację. Dowód lokalności nie wymaga
rachunku It\^o.

Obliczymy teraz średnią bez lokalizacji.
Dzięki symetrii obrotowej możemy przyjąć
$x=(0,0,r)$. Dla $t>0$ gęstość $x+B_t$
w punkcie $y$ wynosi
$(2\pi t)^{-3/2}\exp(-|y-x|^2/(2t))$.
We współrzędnych sferycznych
$\rho=|y|>0$, $\alpha$ jest kątem względem osi
$x$, a element objętości ma postać
$\rho^2\sin\alpha\,\mathrm d\lambda_3(\rho,\alpha,\beta)$.
Po całkowaniu po kącie azymutalnym $\beta$
oraz podstawieniu $v=\cos\alpha$ otrzymujemy
\[
\begin{aligned}
\E M_t
&=\frac{2\pi}{(2\pi t)^{3/2}}\int_0^\infty
\rho e^{-(\rho^2+r^2)/(2t)}
\left(\int_{-1}^1 e^{\rho rv/t}\,\mathrm d\lambda_1(v)\right)
\,\mathrm d\lambda_1(\rho)\\
&=\frac1{r\sqrt{2\pi t}}\int_0^\infty
\left(e^{-(\rho-r)^2/(2t)}-
e^{-(\rho+r)^2/(2t)}\right)
\,\mathrm d\lambda_1(\rho)\\
&=\frac1r\left(\Phi(r/\sqrt t)-\Phi(-r/\sqrt t)\right).
\end{aligned}
\]
W drugiej równości całkujemy funkcję wykładniczą
po $v$, w ostatniej podstawiamy
$(\rho-r)/\sqrt t$ oraz $(\rho+r)/\sqrt t$.
Symetria rozkładu normalnego daje żądany wzór.
Dla każdego $t>0$ średnia jest ściśle mniejsza
niż $1/r$, więc proces nie jest martyngałem.
Jest to \emph{ścisły lokalny martyngał},
w odróżnieniu od prawdziwego martyngału,
który traci średnią dopiero w granicy
$t\to\infty$.
\end{proof}

\section{Wykładnik gaussowski: martyngał, który traci średnią w granicy}
\begin{twierdzenie}[Wykładnik procesu Wienera]
Niech $W$ będzie standardowym procesem Wienera, a
$\theta\in\R$. Wtedy
$Z_t=\exp(\theta W_t-\tfrac12\theta^2t)$
jest dodatnim, ciągłym martyngałem o $\E Z_t=1$
dla każdego $t\ge0$. Gdy $\theta\ne0$,
$Z_t\to0$ prawie na pewno przy $t\to\infty$.
W szczególności rodzina $\{Z_t:t\ge0\}$ nie jest
jednostajnie całkowalna. Jest to konkretny przykład
martyngału ograniczonego w $L^1$, który nie zbiega w $L^1$.
\end{twierdzenie}
\begin{proof}
Niech $s<t$. Z niezależności przyrostów
$W_t-W_s\sim\Normal(0,t-s)$ i z gaussowskiej
funkcji tworzącej momenty mamy
\[
\E[Z_t\mid\calF_s]
=Z_s\E\exp\bigl(\theta(W_t-W_s)-\tfrac12\theta^2(t-s)\bigr)
=Z_s.
\]
Dla liczb całkowitych $n$, przyrosty
$W_n-W_{n-1}$ są niezależne i mają ten sam rozkład
$\Normal(0,1)$; mocne prawo wielkich liczb daje
$W_n/n\to0$ prawie na pewno.
Dla $u>0$ zasada odbicia (twierdzenie~\ref{thm:wiener-reflection})
i symetria dają
$\Prob(\sup_{0\le s\le1}|W_s|>u)\le
4(1-\Phi(u))\le C e^{-u^2/2}$ dla $u\ge1$.
Dlatego dla każdego $\varepsilon>0$ suma
$\sum_n\Prob(\sup_{n\le t\le n+1}|W_t-W_n|>\varepsilon n)$
jest skończona. Lemat Borela--Cantellego pokazuje, że
$\sup_{n\le t\le n+1}|W_t-W_n|/n\to0$ prawie na pewno.
Stąd $W_t/t\to0$ dla $t\to\infty$, a więc
$t^{-1}\log Z_t=\theta W_t/t-\theta^2/2\to-\theta^2/2<0$.
To wymusza $Z_t\to0$ prawie na pewno.
Gdyby rodzina $(Z_t)$ była jednostajnie całkowalna,
zbieżność prawie na pewno pociągałaby zbieżność w $L^1$,
a zatem $\E Z_t\to0$. Przeczy to $\E Z_t=1$.
\end{proof}

\section{Przypadek zatrzymania na pierwszym trafieniu}
\begin{zadanie}[Dlaczego $\E M_\tau$ może różnić się od $\E M_0$?]
Niech $W$ będzie procesem Wienera, a
$\tau=\inf\{t\ge0:W_t=1\}$. Wykazać, że
$\tau<\infty$ prawie na pewno, lecz
$\E W_\tau=1\ne0=\E W_0$.

\emph{Rozwiązanie.} Z zasady odbicia (twierdzenie~\ref{thm:wiener-reflection})
$\Prob(\tau\le t)=2(1-\Phi(1/\sqrt t))\to1$.
Z ciągłości trajektorii $W_\tau=1$ na $\{\tau<\infty\}$.
Dla każdego skończonego $T$ proces $W_{t\wedge\tau}$
jest martyngałem i $\E W_{T\wedge\tau}=0$.
Nie można jednak przejść do granicy pod wartością oczekiwaną:
rodzina $\{W_{T\wedge\tau}:T\ge0\}$ nie jest
jednostajnie całkowalna, skoro granica prawie na pewno
ma wartość oczekiwaną $1$, a wszystkie zmienne mają
średnią $0$. To konkretny przykład nieuprawnionego
stosowania próbkowania opcjonalnego.
\end{zadanie}


% ================================================================
% ROZDZIAŁ 12. PRÓBKOWANIE, NIERÓWNOŚCI, ZBIEŻNOŚĆ
% ================================================================

\chapter{Zatrzymywanie, nierówności maksymalne i zbieżność martyngałów}
\label{ch:stop-cont}

W całym rozdziale zakładamy warunki zwykłe filtracji,
jak w rozdziale~\ref{ch:mart-cont}. Dzięki zupełności
możemy wybierać wszędzie c\`adl\`ag trajektorie
oraz poprawiać czasy zatrzymania na zbiorach zerowych.

\section{Co wolno zrobić w losowej chwili?}
Jeżeli znamy przyszłość, możemy zawsze „wybrać” korzystną
chwilę sprzedaży. Twierdzenia o zatrzymywaniu wyjaśniają,
dlaczego nie uzyskujemy takiego przywileju, gdy chwilę
zatrzymania wybieramy bez zaglądania w przyszłość.
Samo pojęcie czasu zatrzymania nie wystarcza jednak do
przeniesienia równości wartości oczekiwanych na nieograniczony
horyzont. W każdym zastosowaniu trzeba kontrolować
całkowalność rodziny zatrzymanych wartości.

\begin{twierdzenie}[Próbkowanie opcjonalne w ograniczonym horyzoncie]
\label{thm:ost-cont}
Niech $(\calF_t)$ będzie filtracją spełniającą
warunki zwykłe, $M$ martyngałem
c\`adl\`ag i $\sigma,\tau$ czasami zatrzymania takimi,
że $0\le\sigma\le\tau\le T$ prawie na pewno dla ustalonego
$T<\infty$. Wtedy $M_\sigma,M_\tau\in L^1$ i
\[
\E[M_\tau\mid\calF_\sigma]=M_\sigma
\qquad\text{prawie na pewno}.
\]
Dla podmartyngału (odpowiednio nadmartyngału) c\`adl\`ag
ten sam wniosek zachodzi z odpowiednią nierównością,
jeżeli $\{M_\rho:\rho\le T\text{ jest czasem zatrzymania}\}$
jest rodziną jednostajnie całkowalną (w szczególności
gdy proces jest klasy $(DL)$).
\end{twierdzenie}
\begin{proof}
Przypadek $T=0$ jest natychmiastowy.
Gdy ograniczenia i uporządkowanie czasów zachodzą
tylko p.n., na wspólnym zbiorze zerowym przyjmujemy
oba czasy równe zero. Zbiór ten należy do $\calF_0$,
więc zmienione czasy nadal są czasami zatrzymania.
Ich $\sigma$-ciała są identyczne z pierwotnymi:
zdarzenia testujące definicję $\calF_\sigma$
zmieniają się tylko o podzbiory tego zbioru zerowego,
należące do każdej $\calF_t$.
Wartości w czasach zatrzymania są mierzalne
w odpowiednich $\sigma$-ciałach na mocy
twierdzenia~\ref{thm:sampled}.
Możemy zatem zakładać nierówności punktowe.
Podzielmy odcinek $[0,T]$ na $2^n$ równych części
o końcach $t_k^n=kT/2^n$. Niech $\sigma_n,\tau_n$
będą najmniejszymi punktami tej siatki, nie mniejszymi
od odpowiednich czasów (dla wartości $T$ przyjmujemy $T$).
Wówczas $\sigma_n,\tau_n$ są czasami zatrzymania,
$\sigma_n\le\tau_n\le T$ i
$\sigma_n\downarrow\sigma$, $\tau_n\downarrow\tau$.
Istotnie zdarzenie $\{\sigma_n\le t\}$ jest
$\{\sigma\le t_k^n\}$, gdzie $t_k^n$ jest ostatnim
punktem siatki nie większym od $t$ (dla $t\ge T$
zdarzenie jest całym $\Omega$); analogicznie dla $\tau_n$.
Zagnieżdżenie siatek daje monotoniczność przybliżeń.
W skończonym schemacie próbkowania $\calF_{\sigma_n}$
jest dokładnie $\sigma$-ciałem dyskretnego czasu:
dla każdej wartości $t_k^n$ zdarzenie
$A\cap\{\sigma_n=t_k^n\}$ musi należeć
do $\calF_{t_k^n}$. Równoważność tej charakteryzacji
z definicją wynika przez rozkład na skończenie wiele
wartości czasu.
Wartości $M_{t_k^n}$ tworzą martyngał dyskretny, więc
jego twierdzenie o zatrzymaniu daje dla
$A\in\calF_\sigma\subseteq\calF_{\sigma_n}$
\[
\E[M_{\tau_n}\chi_A]=\E[M_{\sigma_n}\chi_A].
\]
Ciągłość prawostronna trajektorii zapewnia zbieżność
$M_{\tau_n}\to M_\tau$ i $M_{\sigma_n}\to M_\sigma$
prawie na pewno. Żeby uzasadnić przejście do wartości
oczekiwanych, zauważmy, że w skończonym schemacie próbkowania
$M_{\tau_n}=\E[M_T\mid\calF_{\tau_n}]$ oraz
$M_{\sigma_n}=\E[M_T\mid\calF_{\sigma_n}]$.
Rodziny te są jednostajnie całkowalne na mocy
lematu~\ref{lem:mart-cont-ui-cond}. Ich zbieżność prawie na pewno
jest zatem także zbieżnością w $L^1$.
Przejście $n\to\infty$ daje
$\E[M_\tau\chi_A]=\E[M_\sigma\chi_A]$ dla każdego
$A\in\calF_\sigma$, czyli żądaną równość warunkową.
W przypadku podmartyngałów i nadmartyngałów
zakładana jednostajna całkowalność wartości w ograniczonych
czasach zatrzymania pozwala przejść do granicy w $L^1$
dokładnie w ten sam sposób, przy zachowaniu właściwego
zwrotu nierówności. Dla ogólnych procesów bez tej
kontroli nie uzasadniamy przejścia do granicy przez
samą zbieżność prawie na pewno.
\end{proof}

\begin{uwaga}[Związek z procesem zatrzymanym]
Każdy martyngał c\`adl\`ag jest klasy $(DL)$
(twierdzenie~\ref{thm:mart-cont-class-d}).
Dla podmartyngału założenie o jednostajnej
całkowalności wartości w czasach zatrzymania
jest warunkiem dostatecznym powyższego dowodu,
a nie definicją podmartyngału.
Twierdzenie o zatrzymywaniu procesu
(twierdzenie~\ref{thm:mart-cont-stopping})
jest zastosowaniem ograniczonego próbkowania
opcjonalnego; nie powtarzamy jego dowodu.
\end{uwaga}

\begin{zadanie}[Nie wolno podejmować decyzji na podstawie przyszłości]
Niech $W$ będzie standardowym procesem Wienera
w uzupełnionej filtracji naturalnej i określmy
losową chwilę $\rho=0$ na $\{W_1\ge0\}$ oraz $\rho=1$
na $\{W_1<0\}$. Czy $\rho$ jest czasem zatrzymania?
Ile wynosi $\E W_\rho$?

\emph{Rozwiązanie.} Zdarzenie $\{\rho\le0\}=\{W_1\ge0\}$
nie należy do $\calF_0$ uzupełnionej filtracji naturalnej
ruchu Browna: ma prawdopodobieństwo $1/2$ i zależy od przyszłej
obserwacji. Zatem $\rho$ nie jest czasem zatrzymania.
Ponieważ $W_0=0$ i $W_1\sim\Normal(0,1)$,
\[
\E W_\rho=\E[W_1\chi_{\{W_1<0\}}]
=-\frac1{\sqrt{2\pi}},
\qquad
\text{gdyż }\int_{-\infty}^0 x e^{-x^2/2}\,
\mathrm d\lambda_1(x)=-1.
\]
Nierówność średnich nie przeczy twierdzeniu o próbkowaniu
opcjonalnym: jego podstawowe założenie, czyli brak
korzystania z przyszłości, nie jest spełnione.
\end{zadanie}


\section{Maksimum procesu a jego wartość końcowa}
\begin{twierdzenie}[Nierówność maksymalna Dooba]
\label{thm:doob-cont}
Niech $X$ będzie nieujemnym podmartyngałem c\`adl\`ag
na $[0,T]$ i $X_T^*=\sup_{0\le t\le T}X_t$.
Dla $a>0$ zachodzi nierówność ogonowa
\[
a\Prob(X_T^*\ge a)\le
\E[X_T\chi_{\{X_T^*\ge a\}}]\le\E X_T.
\]
Jeżeli $p>1$ i $X_T\in L^p$, to
\[
\E[(X_T^*)^p]\le
\left(\frac{p}{p-1}\right)^p\E X_T^p.
\]
W szczególności dla martyngału c\`adl\`ag $M$
można zastosować obie nierówności do $X_t=|M_t|$.
\end{twierdzenie}
\begin{proof}
Ustalmy skończony zbiór chwil
$0=t_0<\cdots<t_n=T$ i połóżmy $Y_k=X_{t_k}$.
Jest to nieujemny podmartyngał dyskretny.
Niech $Y=\max_{0\le k\le n}Y_k$.
Jeżeli $\rho=\min\{k:Y_k\ge a\}\wedge n$,
z konwencją $\inf\varnothing=\infty$,
to na zdarzeniu $\{Y\ge a\}$ mamy $Y_\rho\ge a$.
Próbkowanie podmartyngału daje
$a\Prob(Y\ge a)\le\E[Y_n\chi_{\{Y\ge a\}}]$:
dokładniej sumujemy nierówności warunkowe po rozłącznych
zdarzeniach $\{\rho=k,Y\ge a\}$, które należą
do $\calF_{t_k}$: dla $k<n$ jest to $\{\rho=k\}$,
a dla $k=n$ przecinamy $\{\rho=n\}$ z $\{Y_n\ge a\}$.
Dla $p>1$ z $0\le Y_k\le\E[X_T\mid\calF_{t_k}]$
p.n. i warunkowej nierówności Jensena wynika
$\E Y_k^p\le\E X_T^p$. W szczególności
$\E Y^p\le\sum_{k=0}^n\E Y_k^p<\infty$.
Całkujemy nierówność ogonową względem
$p a^{p-2}\,\mathrm d\lambda_1(a)$ po $a\in(0,\infty)$.
Twierdzenie Tonellego daje
\[
\E Y^p\le\frac{p}{p-1}\E[Y_nY^{p-1}].
\]
Z nierówności H\"oldera prawa strona jest nie większa od
$\frac{p}{p-1}\|Y_n\|_p\|Y\|_p^{p-1}$.
Jeśli $\E Y^p=0$, teza jest oczywista;
w przeciwnym razie dzielimy przez odpowiednią potęgę
i otrzymujemy $\|Y\|_p\le p(p-1)^{-1}\|Y_n\|_p$.
Przechodzimy do rosnącego ciągu skończonych siatek,
których suma jest gęsta i zawiera $T$.
Dzięki prawostronnej ciągłości trajektorii
$Y$ rośnie do $X_T^*$, który jest zatem mierzalny.
Zbieżność monotoniczna daje oszacowanie w $L^p$.
W wersji ogonowej supremum funkcji c\`adl\`ag
może nie być osiągane. Dla $0<\varepsilon<a$
zdarzenia $\{Y\ge a-\varepsilon\}$ rosną wraz
z zagęszczaniem siatek do zdarzenia $A_\varepsilon$,
które zawiera $\{X_T^*\ge a\}$ i jest zawarte
w $\{X_T^*\ge a-\varepsilon\}$.
Dyskretna nierówność i zbieżność monotoniczna
dają $(a-\varepsilon)\Prob(A_\varepsilon)
\le\E[X_T\chi_{A_\varepsilon}]$.
Dla ciągu $\varepsilon\downarrow0$ zdarzenia
$A_\varepsilon$ maleją do $\{X_T^*\ge a\}$,
więc przejście do granicy daje nierówność ogonową.
Dla $X=|M|$ używamy warunkowej nierówności Jensena,
która dowodzi własności podmartyngałowej $X$.
\end{proof}

\begin{wniosek}[Maksimum ruchu Browna w $L^2$]
Dla standardowego procesu Wienera $W$ oraz $T>0$
\[
\E\sup_{0\le t\le T}|W_t|^2\le4T.
\]
Dla standardowego $d$-wymiarowego ruchu Browna $B$
\[
\E\sup_{0\le t\le T}|B_t|^2\le4dT.
\]
\end{wniosek}
\begin{proof}
Pierwsza nierówność wynika z twierdzenia dla $p=2$
i $\E W_T^2=T$. W drugiej używamy
$|B_t|^2=\sum_i(B_t^{(i)})^2$ oraz
$\sup_t\sum_i a_i(t)\le\sum_i\sup_t a_i(t)$.
\end{proof}

\section{Nierówność przekroczeń i istnienie granic}
\begin{definicja}[Przekroczenie przedziału]
Dla funkcji rzeczywistej $x$ i liczb $a<b$ \emph{
pełne przekroczenie w górę przedziału $[a,b]$} to
para chwil $u<v$ takich, że $x(u)\le a$
oraz $x(v)\ge b$, przy czym po osiągnięciu $b$
rozpoczynamy poszukiwanie następnego przekroczenia
od ponownego zejścia do poziomu $a$. Maksymalną liczbę
takich rozłącznych przekroczeń do czasu $T$ oznaczamy
$U_{a,b}(T)$ (supremum liczby par w skończonych
ciągach $u_1<v_1<u_2<v_2<\cdots\le T$).
W dowodach używamy liczby przekroczeń na ustalonych
przeliczalnych zbiorach chwil, mierzalnej jako
supremum po skończonych podzbiorach.
\end{definicja}

\begin{lemat}[Przekroczenia dla martyngału dyskretnego]
Niech $(M_k)_{k=0}^n$ będzie martyngałem całkowalnym,
a $a<b$. Wtedy
\[
(b-a)\E U_{a,b}(n)\le\E[(a-M_n)^+].
\]
\end{lemat}
\begin{proof}
Budujemy przewidywalną strategię: nabywamy jedną jednostkę
aktywa, gdy $M_k\le a$, sprzedajemy przy pierwszym
$M_l\ge b$ i ponawiamy ten schemat. Pozycję otwartą
w chwili $n$ wyceniamy według $M_n$.
Dla każdego pełnego cyklu zysk jest co najmniej $b-a$.
Ewentualna niedokończona pozycja otwarta po cenie
co najwyżej $a$ ma zysk większy lub równy
$M_n-a\ge-(a-M_n)^+$. Łączny zysk $G_n$ spełnia
$G_n\ge(b-a)U_{a,b}(n)-(a-M_n)^+$.
Aby uzasadnić przewidywalność, definiujemy $H_k$ w chwili $k$
wyłącznie na podstawie $M_0,\ldots,M_k$: po zakupie utrzymujemy
$H_k=1$ do pierwszego osiągnięcia poziomu $b$, a po sprzedaży
przyjmujemy $H_k=0$ aż do kolejnego osiągnięcia poziomu $a$.
W szczególności $H_k\in\{0,1\}$ jest $\calF_k$-mierzalne.
Ponieważ
$G_n=\sum_{k=0}^{n-1}H_k(M_{k+1}-M_k)$ jest sumą skończoną,
$\E|G_n|<\infty$ i
$\E G_n=\sum_k\E[H_k\E(M_{k+1}-M_k\mid\calF_k)]=0$. Biorąc wartości
oczekiwane w nierówności, otrzymujemy tezę.
\end{proof}

\begin{twierdzenie}[Zbieżność przy ograniczeniu w $L^1$]
\label{thm:mart-cont-convergence}
Niech $M$ będzie podmartyngałem c\`adl\`ag na
$[0,\infty)$ takim, że $\sup_{t\ge0}\E M_t^+<\infty$.
Wtedy $\sup_t\E|M_t|<\infty$ oraz istnieje zmienna
losowa $M_\infty\in L^1$ taka, że
$M_t\to M_\infty$ prawie na pewno przy $t\to\infty$.
W szczególności wynik dotyczy martyngałów ograniczonych
w $L^1$. Dla nadmartyngału wystarcza
$\sup_t\E M_t^-<\infty$.
Nie wnioskujemy z tych założeń zbieżności w $L^1$.
\end{twierdzenie}
\begin{proof}
Z monotoniczności średnich podmartyngału mamy
$\E|M_t|=2\E M_t^+-\E M_t\le2\sup_u\E M_u^+-\E M_0$.
Wybierzmy przeliczalny gęsty zbiór chwil
$D=\Q\cap[0,\infty)$ oraz rosnące skończone
podzbiory zawierające zero i wyczerpujące $D$.
Na każdej siatce próbkowany proces jest
podmartyngałem dyskretnym. Nierówność przekroczeń
Dooba (tom I, rozdział~\ref{I-ch:nadmartyngaly})
daje dla ostatniej chwili $v$ siatki
$(b-a)\E U\le\E[(M_v-a)^+]-\E[(M_0-a)^+]
\le\sup_t\E M_t^++|a|$.
Zbieżność monotoniczna oraz przeliczalność
par wymiernych $a<b$ dają wspólny zbiór pełnej
miary, na którym liczba przekroczeń każdej
takiej pary na $D$ jest skończona.
Gdyby wartości $M_t$ dla $t\in D$ nie miały granicy
w rozszerzonej prostej rzeczywistej przy $t\to\infty$,
ich dolna granica byłaby ściśle mniejsza od górnej.
Wybierając pomiędzy nimi wymierne $a<b$,
moglibyśmy naprzemiennie wybierać coraz późniejsze
chwile z wartościami poniżej $a$ i powyżej $b$,
uzyskując nieskończenie wiele przekroczeń.
Zatem granica istnieje.
Stosując lemat Fatou wzdłuż całkowitych chwil
$n\to\infty$, dostajemy
$\E\liminf_n |M_n|\le\sup_t\E|M_t|<\infty$.
Ponieważ granica wzdłuż $D$ już istnieje p.n.,
jest skończona i należy do $L^1$.
Jako mierzalną wersję $M_\infty$ bierzemy granicę
$(M_n)$ tam, gdzie jest skończona, i zero poza tym;
jest ona $\calF_\infty$-mierzalna, gdzie
$\calF_\infty=\sigma(\bigcup_t\calF_t)$.
Prawostronna ciągłość trajektorii rozszerza wynik
na wszystkie chwile: dla każdej chwili $t$
można dobrać $r\in D$ tuż za $t$ z
$|M_r-M_t|$ dowolnie małym. Gdyby pewien ciąg
$t_n\to\infty$ nie dawał tej samej granicy,
dobierając $r_n\ge t_n$ z błędem nie większym niż
$1/n$, otrzymalibyśmy sprzeczność ze zbieżnością
próbek wymiernych.
Dla nadmartyngału stosujemy dowód do $-M$.
\end{proof}

\begin{twierdzenie}[Zbieżność jednostajnie całkowalnego martyngału]
\label{thm:mart-cont-ui-convergence}
Niech $M$ będzie martyngałem c\`adl\`ag na
$[0,\infty)$, a rodzina $\{M_t:t\ge0\}$
jednostajnie całkowalna. Wtedy istnieje
$M_\infty\in L^1$ takie, że $M_t\to M_\infty$
prawie na pewno i w $L^1$ oraz
\[
M_t=\E[M_\infty\mid\calF_t]
\qquad (t\ge0).
\]
Odwrotnie, każda taka reprezentacja z $M_\infty\in L^1$
tworzy jednostajnie całkowalny martyngał.
\end{twierdzenie}
\begin{proof}
Jednostajna całkowalność zapewnia ograniczenie w $L^1$,
więc istnieje granica prawie pewna z poprzedniego
twierdzenia. Kryterium Vitaliego daje też zbieżność
w $L^1$. Dla każdego ustalonego $s$ mamy
$M_s=\E[M_t\mid\calF_s]$ przy $t\ge s$.
Kontrakcja warunkowania w $L^1$ pozwala przejść
do granicy $t\to\infty$ po prawej stronie.
Wtedy $M_s=\E[M_\infty\mid\calF_s]$.
Implikacja odwrotna wynika z lematu~\ref{lem:mart-cont-ui-cond}
oraz własności wieży.
Związek z klasą $(D)$ podaje
twierdzenie~\ref{thm:mart-cont-class-d}.
\end{proof}

\begin{wniosek}[Zbieżność martyngałów w $L^p$]
\label{cor:mart-cont-lp-convergence}
Niech $p>1$, a $M$ będzie martyngałem c\`adl\`ag
spełniającym $\sup_{t\ge0}\E|M_t|^p<\infty$.
Wtedy istnieje $M_\infty\in L^p$ takie, że
$M_t\to M_\infty$ prawie na pewno i w $L^p$.
\end{wniosek}
\begin{proof}
Ograniczenie w $L^p$ implikuje jednostajną całkowalność
rodziny $(M_t)$, więc otrzymujemy granicę w $L^1$
i reprezentację $M_t=\E[M_\infty\mid\calF_t]$.
Lemat Fatou daje $M_\infty\in L^p$.
Warunkowa nierówność Jensena implikuje
$|M_t|^p\le\E[|M_\infty|^p\mid\calF_t]$.
Rodzina po prawej stronie jest jednostajnie całkowalna,
a więc taka jest również $\{|M_t|^p:t\ge0\}$.
Ponieważ $|M_t-M_\infty|^p\le2^{p-1}
(|M_t|^p+|M_\infty|^p)$, rodzina tych różnic do
potęgi $p$ też jest jednostajnie całkowalna.
Zbieżność prawie na pewno i kryterium Vitaliego dają
$\E|M_t-M_\infty|^p\to0$.
\end{proof}

\section{Kiedy można zatrzymać martyngał bez ograniczenia czasu?}
\begin{twierdzenie}[Próbkowanie dla jednostajnie całkowalnego martyngału]
\label{thm:mart-cont-ui-ost}
Niech $M$ będzie jednostajnie całkowalnym martyngałem
c\`adl\`ag, $\tau$ skończonym prawie na pewno czasem
zatrzymania. Wtedy $M_\tau\in L^1$ i
$\E M_\tau=\E M_0$. Ponadto, jeśli $\sigma\le\tau$
jest innym skończonym prawie na pewno czasem zatrzymania,
to $\E[M_\tau\mid\calF_\sigma]=M_\sigma$.
\end{twierdzenie}
\begin{proof}
Na mocy twierdzeń~\ref{thm:mart-cont-ui-convergence}
i~\ref{thm:mart-cont-class-d} dla każdego p.n.
skończonego czasu zatrzymania $\rho$ mamy
$M_\rho=\E[M_\infty\mid\calF_\rho]$.
Jeżeli $\sigma\le\tau$ p.n., po zmianie na zbiorze
zerowym możemy przyjąć uporządkowanie punktowe.
Wtedy $\calF_\sigma\subseteq\calF_\tau$,
a własność wieży daje
\[
\E[M_\tau\mid\calF_\sigma]
=\E[M_\infty\mid\calF_\sigma]=M_\sigma.
\]
Wybór $\sigma=0$ i wzięcie wartości oczekiwanej
dają $\E M_\tau=\E M_0$.
\end{proof}

\begin{zadanie}[Oczekiwanie na pierwszy sukces a uczciwa gra]
Niech $W$ będzie standardowym procesem Wienera,
$\tau_a=\inf\{t\ge0:W_t=a\}$, $a>0$. Wyjaśnić,
dlaczego nie możemy zastosować twierdzenia o
nieograniczonym zatrzymywaniu do $M=W$ i $\tau_a$.

\emph{Rozwiązanie.} Chociaż $\tau_a<\infty$ prawie
na pewno i $W_{\tau_a}=a$, martyngał $W$ nie jest
jednostajnie całkowalny na całej półprostej
($\E|W_t|$ rośnie jak $\sqrt t$).
Równość $\E W_{\tau_a}=\E W_0$ byłaby zatem
nieuprawniona. Ograniczone czasy $\tau_a\wedge T$
są dozwolone i dają $\E W_{\tau_a\wedge T}=0$;
brakuje jednak jednostajnej całkowalności przy
$T\to\infty$. Na tym właśnie polega różnica
między lokalną a globalną uczciwością.
\end{zadanie}

\begin{zadanie}[Ruina z dwiema granicami w czasie ciągłym]
Niech $W$ będzie standardowym procesem Wienera,
$a,b>0$ i
$\tau=\inf\{t\ge0:W_t\notin(-a,b)\}$.
Obliczyć prawdopodobieństwo wyjścia górą i wartość
oczekiwaną $\tau$, stosując tylko martyngały
$W_t$ oraz $W_t^2-t$.

\emph{Rozwiązanie.} Zatrzymujemy drugi martyngał w
$\tau\wedge n$:
$\E W_{\tau\wedge n}^2=\E(\tau\wedge n)$.
Ponieważ $|W_{\tau\wedge n}|\le\max(a,b)$,
$\E(\tau\wedge n)\le\max(a,b)^2$.
Zbieżność monotoniczna daje $\E\tau<\infty$, więc
$\tau<\infty$ prawie na pewno. Z ciągłości
$W_\tau\in\{-a,b\}$.
Proces $W^\tau$ jest ograniczony, a zatem
jednostajnie całkowalny. Z próbkowania opcjonalnego
$\E W_\tau=0$. Jeśli $p=\Prob(W_\tau=b)$, to
$bp-a(1-p)=0$, skąd
\[
\Prob(W_\tau=b)=\frac{a}{a+b},\qquad
\Prob(W_\tau=-a)=\frac{b}{a+b}.
\]
Ponownie używając $W_t^2-t$ i zbieżności zmajoryzowanej
po stronie kwadratu otrzymujemy
$\E\tau=\E W_\tau^2=b^2p+a^2(1-p)=ab$.
Rachunek ten pokazuje, jak zamienić zadanie o czasie
wyjścia na dwa równania dla momentów zatrzymanego
martyngału, bez rozwiązywania równania różniczkowego.
\end{zadanie}

\begin{zadanie}[Transformata Laplace'a czasu pierwszego trafienia]
Niech $a>0$ oraz
$\tau_a=\inf\{t\ge0:W_t=a\}$. Dla $u>0$
obliczyć $\E e^{-u\tau_a}$ bez posługiwania się
gęstością czasu trafienia.

\emph{Rozwiązanie.} Dla $\theta>0$ proces
$Z_t=\exp(\theta W_t-\theta^2t/2)$ jest
martyngałem. Dla każdego $T$ jego zatrzymana
wartość spełnia $0<Z_{T\wedge\tau_a}\le e^{\theta a}$:
przed trafieniem $W_s<a$, a składnik
$-\theta^2s/2$ jest niedodatni. Zatem rodzina
zatrzymanych wartości jest jednostajnie całkowalna.
Ponieważ $\tau_a<\infty$ prawie na pewno,
z próbkowania opcjonalnego i zbieżności zmajoryzowanej
\[
1=\E Z_{\tau_a}
=e^{\theta a}\E e^{-\theta^2\tau_a/2}.
\]
Po podstawieniu $\theta=\sqrt{2u}$ dostajemy
$\E e^{-u\tau_a}=e^{-a\sqrt{2u}}$.
Zauważmy, że tego zatrzymania \emph{nie} wolno
zastosować wprost do $W$, ponieważ $W_{\tau_a}$
ma średnią $a$, a $W_0$ średnią zero. Wykładnik
jest odpowiednim narzędziem właśnie dlatego, że
po zatrzymaniu posiada deterministyczne ograniczenie.
\end{zadanie}


% ================================================================
% ROZDZIAŁ 13. WARIACJA KWADRATOWA
% ================================================================

\chapter{Wariacja kwadratowa i kowariacja ciągłych martyngałów}
\label{ch:quadratic-cont}
\paragraph{Odsyłacze międzytomowe.} Konstrukcja całki Itô wykorzystująca wariację kwadratową znajduje się w tomie III, rozdziale~\ref{III-ch:ito-construction}.


W całym rozdziale zakładamy warunki zwykłe filtracji
(rozdział~\ref{ch:mart-cont}). Procesy o prawie na pewno
ciągłych trajektoriach zastępujemy w razie potrzeby
nierozróżnialnymi wersjami o wszędzie ciągłych
trajektoriach, zmieniając je na zbiorze zerowym z $\calF_0$.

\section{Dlaczego kwadraty przyrostów są inne niż zwykła wariacja?}
Dla gładkiej funkcji $f$ na $[0,T]$ suma
$\sum_i(f(t_{i+1})-f(t_i))^2$ dąży do zera, gdy
średnica podziału dąży do zera: jest nie większa od
$\max_i|\Delta_i f|\sum_i|\Delta_i f|$, a pierwsza
wielkość maleje do zera, podczas gdy druga pozostaje
ograniczona przez wariację całkowitą. Dla ruchu Browna
typowy przyrost jest rzędu $\sqrt{\Delta t}$, zatem
jego kwadrat ma rząd $\Delta t$. Suma kwadratów
nie musi już zanikać. Właśnie ta obserwacja motywuje
rachunek Itô, który pojawi się w tomie III.

\begin{definicja}[Sumy kwadratów przyrostów]
Niech $X$ będzie rzeczywistym procesem o trajektoriach
ciągłych i $T>0$. Dla deterministycznego podziału
$\pi=\{0=t_0<t_1<\cdots<t_m=T\}$, $t\in[0,T]$,
połóżmy
\[
V_t^\pi(X)=\sum_{i=0}^{m-1}
\bigl(X_{t\wedge t_{i+1}}-X_{t\wedge t_i}\bigr)^2.
\]
Jest to ciągły proces nieujemny, choć dla ustalonej
siatki nie musi być niemalejący pomiędzy jej węzłami.
Średnicą podziału nazywamy $|\pi|=\max_i(t_{i+1}-t_i)$.
Gdy istnieje ciągły proces $A$ taki, że dla każdego
$T>0$ i każdego ciągu deterministycznych podziałów
$\pi_n$ przedziału $[0,T]$ o średnicach $|\pi_n|\to0$
mamy
$\sup_{t\le T}|V_t^{\pi_n}(X)-A_t|\to0$
według prawdopodobieństwa, proces $A$ nazywamy
\emph{wariacją kwadratową} $X$ i oznaczamy $[X]$.
Skrót \emph{ucp} oznacza właśnie zbieżność jednostajną
na przedziałach zwartych według prawdopodobieństwa.
Taka granica jest jednoznaczna z dokładnością do
nierozróżnialności: dla dwóch granic $A,A'$ na każdym
$[0,T]$ nierówność trójkąta daje
$\Prob(\sup_{t\le T}|A_t-A'_t|>\varepsilon)=0$
dla każdego $\varepsilon>0$. Wystarczy następnie
wziąć przeliczalnie wiele dodatnich całkowitych $T$
i $\varepsilon=1/n$.
Jeżeli $X$ jest adaptowany, to $V^\pi(X)$ również;
ciągłą granicę można wybrać adaptowaną, korzystając
z zupełności filtracji, jak w dowodzie poniżej.
\end{definicja}

\begin{lemat}[Znikanie wariacji kwadratowej dla trajektorii o skończonej wariacji]
\label{lem:finite-var-zero-quadratic}
Niech $f:[0,T]\to\R$ będzie funkcją ciągłą o
skończonej wariacji całkowitej $\operatorname{Var}_{[0,T]}(f)$.
Dla dowolnego ciągu podziałów o średnicach malejących
do zera
$\sup_{t\le T}V_t^{\pi_n}(f)\to0$.
\end{lemat}
\begin{proof}
Dla każdego $t\le T$ mamy
\[
V_t^\pi(f)\le
\Bigl(\max_i\sup_{u,v\in[t_i,t_{i+1}]}|f(u)-f(v)|\Bigr)
\operatorname{Var}_{[0,T]}(f).
\]
Z jednostajnej ciągłości $f$ na zwartym przedziale
pierwszy czynnik dąży do zera, jednostajnie w $t$.
Drugi czynnik nie zależy od podziału ani od $t$.
\end{proof}

\begin{twierdzenie}[Wariacja kwadratowa procesu Wienera]
\label{thm:brownian-bracket}
Niech $W$ będzie standardowym procesem Wienera,
a $T>0$. Dla dowolnych deterministycznych podziałów
$\pi_n$ o $|\pi_n|\to0$ zachodzi
$V_T^{\pi_n}(W)\to T$ w $L^2$, a nawet
\[
\sup_{0\le t\le T}|V_t^{\pi_n}(W)-t|
\longrightarrow0
\quad\text{według prawdopodobieństwa}.
\]
Ponadto dla każdego ciągu siatek spełniającego
$\sum_n|\pi_n|<\infty$ zbieżność jednostajna
zachodzi prawie na pewno. Dotyczy to w szczególności
siatek równomiernych dyadycznych.
Wynik jest zgodny z konstrukcją w rozdziale~\ref{ch:wiener}.
\end{twierdzenie}
\begin{proof}
Z niezależności przyrostów $\Delta_iW\sim
\Normal(0,h_i)$, gdzie $h_i=t_{i+1}-t_i$,
wynika $\E(\Delta_iW)^2=h_i$ i
$\Var((\Delta_iW)^2)=2h_i^2$.
Dlatego
\[
\E V_T^\pi(W)=T,\qquad
\E|V_T^\pi(W)-T|^2=2\sum_i h_i^2
\le2T|\pi|.
\]
To daje zbieżność w $L^2$ dla wszystkich siatek.
Aby uzyskać zbieżność jednostajną, w węzłach podziału
rozważmy dyskretny martyngał
$D_k=\sum_{i<k}((\Delta_iW)^2-h_i)$.
Nierówność maksymalna Dooba w $L^2$
(tom I, rozdział~\ref{I-ch:nadmartyngaly}) daje
\[
\E\max_{k\le m}|D_k|^2\le4\E|D_m|^2
\le8T|\pi|.
\]
Między węzłami różnica $V_t^\pi(W)-t$ może
zmienić się względem lewego węzła najwyżej o
$\max_i\sup_{u,v\in[t_i,t_{i+1}]}|W_u-W_v|^2+|\pi|$.
Ta wielkość dąży prawie na pewno do zera dzięki
jednostajnej ciągłości trajektorii. Stąd teza ucp.
Jeśli $\sum_n|\pi_n|<\infty$, to
$\sum_n\E\max_k|D_k|^2<\infty$
(w szczególności dla siatek dyadycznych $|\pi_n|=T2^{-n}$).
Nierówność Markowa i lemat Borela--Cantellego dają
zbieżność $\max_k|D_k|\to0$ prawie na pewno.
Oscylacje wewnątrz komórek także zanikają prawie na pewno,
więc zbieżność jest jednostajna prawie na pewno.
\end{proof}

\begin{center}
\begin{tikzpicture}[x=1.5cm,y=0.85cm,>=Latex]
\draw[->] (0,0)--(4.4,0) node[right] {$t$};
\draw[->] (0,0)--(0,3.8) node[above] {$x(t)$};
\foreach \x in {1,2,3,4}{\draw[gray,densely dotted] (\x,0)--(\x,3.3);}
\draw[blue!75!black,thick] (0,1)--(0.55,2.4)--(1,1.8)--(1.6,0.6)--(2,1.2)--(2.5,2.8)--(3,2.2)--(3.5,1.4)--(4,2.8);
\foreach \x/\y in {0/1,1/1.8,2/1.2,3/2.2,4/2.8}{\fill[blue!75!black] (\x,\y) circle (1.6pt);}
\end{tikzpicture}

\begin{tikzpicture}[x=1.5cm,y=0.85cm,>=Latex]
\draw[->] (0,0)--(4.4,0) node[right] {$t$};
\draw[->] (0,0)--(0,4.2) node[above] {$V_t^\pi(x)$};
\foreach \x in {1,2,3,4}{\draw[gray,densely dotted] (\x,0)--(\x,3.85);}
\draw[blue!75!black,thick,domain=0:0.55,samples=25] plot(\x,{(1.4*\x/0.55)^2});
\draw[blue!75!black,thick,domain=0.55:1,samples=25] plot(\x,{(1.4-0.6*(\x-0.55)/0.45)^2});
\draw[blue!75!black,thick,domain=1:1.6,samples=25] plot(\x,{0.64+(1.2*(\x-1)/0.6)^2});
\draw[blue!75!black,thick,domain=1.6:2,samples=25] plot(\x,{0.64+(1.2-0.6*(\x-1.6)/0.4)^2});
\draw[blue!75!black,thick,domain=2:2.5,samples=25] plot(\x,{1+(1.6*(\x-2)/0.5)^2});
\draw[blue!75!black,thick,domain=2.5:3,samples=25] plot(\x,{1+(1.6-0.6*(\x-2.5)/0.5)^2});
\draw[blue!75!black,thick,domain=3:3.5,samples=25] plot(\x,{2+(0.8*(\x-3)/0.5)^2});
\draw[blue!75!black,thick,domain=3.5:4,samples=30] plot(\x,{2+(-0.8+1.4*(\x-3.5)/0.5)^2});
\draw[orange!80!black,thick,dashed] (0,0)--(1,0)--(1,0.64)--(2,0.64)--(2,1)--(3,1)--(3,2)--(4,2)--(4,2.36);
\end{tikzpicture}
\par\smallskip\small Góra: przykładowa ciągła trajektoria
odcinkami liniowa i siatka $\{0,1,2,3,4\}$.
Dół: suma kwadratów z niedokończonym przyrostem (linia ciągła)
oraz suma tylko pełnych przyrostów (linia przerywana).
Pierwsza może lokalnie maleć, druga jest niemalejąca.
Nie jest to symulacja ruchu Browna.
\end{center}

\section{Dyskretny lemat izometrii: fundament konstrukcji}
Zanim skonstruujemy wariację kwadratową \emph{dowolnego}
ciągłego martyngału, potrzebujemy tylko elementarnego
faktu o skończonych sumach. Nie konstruujemy jeszcze
ogólnej całki Itô.

\begin{lemat}[Izometria dla skończonej sumy przewidywalnej]
\label{lem:discrete-isometry}
Niech $0=t_0<\cdots<t_m=T$ oraz
$M$ będzie martyngałem ciągłym i kwadratowo całkowalnym
na $[0,T]$. Jeżeli $H_i$ jest ograniczoną
$\calF_{t_i}$-mierzalną zmienną losową, to
\[
S_k:=\sum_{i=0}^{k-1}H_i(M_{t_{i+1}}-M_{t_i}),
\qquad k=0,\ldots,m,
\]
jest martyngałem dyskretnym, a
\[
\E S_m^2
=\sum_{i=0}^{m-1}\E[H_i^2(M_{t_{i+1}}-M_{t_i})^2]
=\E\sum_{i=0}^{m-1}H_i^2
\E[(M_{t_{i+1}}-M_{t_i})^2\mid\calF_{t_i}].
\]
\end{lemat}
\begin{proof}
Warunkowanie daje
$\E[H_i(M_{t_{i+1}}-M_{t_i})\mid\calF_{t_i}]=0$,
więc $S$ jest martyngałem. Po rozwinięciu kwadratu
$S_m^2$ wyrazy mieszane mają średnią zero:
jeśli $i<j$, to iloczyn
$H_i\Delta_iM H_j$ jest $\calF_{t_j}$-mierzalny,
a $\E[\Delta_jM\mid\calF_{t_j}]=0$.
Pozostają wyłącznie wyrazy diagonalne; ich druga
postać to zwykła własność warunkowej wartości oczekiwanej.
\end{proof}

\section{Konstrukcja nawiasu ciągłego martyngału bez całki Itô}
\begin{twierdzenie}[Istnienie wariacji kwadratowej]
\label{thm:general-bracket}
Niech $(\calF_t)_{t\ge0}$ będzie filtracją
spełniającą warunki zwykłe,
a $M$ ciągłym lokalnym martyngałem rzeczywistym,
$M_0=0$. Istnieje, z dokładnością do nierozróżnialności, dokładnie jeden
ciągły, adaptowany, niemalejący proces $[M]$, rozpoczynający się od zera,
taki że $M_t^2-[M]_t$ jest lokalnym martyngałem.
Dla każdego $T>0$ i dowolnej rodziny deterministycznych
podziałów $\pi_n$ przedziału $[0,T]$ o
$|\pi_n|\to0$ zachodzi
\[
\sup_{0\le t\le T}|V_t^{\pi_n}(M)-[M]_t|
\xrightarrow{\Prob}0.
\]
Dla ciągłego lokalnego martyngału o dowolnej
wartości początkowej definiujemy $[M]=[M-M_0]$;
przyrosty obu procesów są identyczne.
Wówczas $(M-M_0)^2-[M]$ jest lokalnym martyngałem.
Nie twierdzimy bez dodatkowego założenia,
że $M^2-[M]$ jest lokalnym martyngałem,
gdyż $M_0^2$ może nie być całkowalna.
\end{twierdzenie}
\begin{proof}
\emph{Krok 1: martyngał ograniczony na $[0,T]$.}
Załóżmy najpierw, że $|M_t|\le K$ dla $t\le T$.
Dla podziału $\pi$ definiujemy
\[
I_t^\pi=\sum_i M_{t_i}
(M_{t\wedge t_{i+1}}-M_{t\wedge t_i}).
\]
Dla ustalonego podziału $I^\pi$ jest ciągłym martyngałem:
każdy składnik jest zerowy przed $t_i$, a dla
$t\ge t_i$ jest martyngałem $M_{t_i}(M_{t\wedge t_{i+1}}-M_{t_i})$.
Dla $s<t_i<t$ warunkujemy najpierw względem
$\calF_{t_i}$, uzyskując zero, a następnie względem
$\calF_s$. Dla $t_i\le s\le t$ wystarcza
własność martyngałowa zatrzymanego procesu i
$\calF_s$-mierzalność współczynnika.
Ograniczenie $M$ zapewnia całkowalność.
Z dyskretnej izometrii i ortogonalności przyrostów
$\E|I_t^\pi|^2\le K^2\E M_t^2$.

Pokażemy, że $(I^\pi)$ jest ciągiem Cauchy'ego w
ucp przy malejącej średnicy podziałów. Ustalmy
$\varepsilon>0$ i dwa podziały $\pi,\rho$.
Niech $\sigma$ będzie pierwszą chwilą, w której
oscylacja $M$ od lewego końca komórki $\pi$ lub
$\rho$ osiąga poziom $\varepsilon$, oraz $T$,
jeśli nie nastąpi to wcześniej. Formalnie dla każdej komórki $[u,v]$ obu siatek
bierzemy pierwszy czas $q\in[u,v]$ taki, że
$|M_q-M_u|\ge\varepsilon$ (lub $\infty$, gdy go nie ma).
Zdarzenie trafienia do chwili $t\ge u$ wyraża się przez
$\sup_{u\le q\le v\wedge t}|M_q-M_u|\ge\varepsilon$,
które jest $\calF_t$-mierzalne dzięki ciągłości
i możliwości użycia przeliczalnego gęstego zbioru
chwil wraz z końcami przedziału.
Minimum tych czasów i $T$ jest więc czasem zatrzymania.
Z jednostajnej ciągłości trajektorii wynika
$\Prob(\sigma<T)\to0$ przy
$|\pi|\vee|\rho|\to0$.
Dla zatrzymanego martyngału $M^\sigma$ przyrosty
w każdej komórce obu siatek mają moduł co najwyżej
$\varepsilon$ (również po zatrzymaniu).
Na wspólnym rozdrobnieniu siatek różnicę
$I^\pi(M^\sigma)-I^\rho(M^\sigma)$ zapisujemy jako
skończoną sumę przewidywalną z współczynnikami
$H_j=M^\sigma_{\text{lewy }\pi}-
M^\sigma_{\text{lewy }\rho}$ o $|H_j|\le2\varepsilon$.
Z lematu~\ref{lem:discrete-isometry} oraz nierówności
maksymalnej Dooba w $L^2$ otrzymujemy
\[
\E\sup_{t\le T}|I_t^\pi(M^\sigma)-
I_t^\rho(M^\sigma)|^2
\le16\varepsilon^2\E (M^\sigma_T)^2
\le16\varepsilon^2\E M_T^2.
\]
Pierwsze oszacowanie dla supremum można otrzymać także
bezpośrednio z ciągłoczasowej nierówności Dooba;
ostatnią nierówność zapewnia próbkowanie opcjonalne
dla podmartyngału $M^2$.
Na $\{\sigma=T\}$ oba procesy zatrzymane są zgodne
z niezatrzymanym na $[0,T]$. Stąd dla $\eta>0$
\[
\Prob\left(\sup_{t\le T}|I_t^\pi-I_t^\rho|>\eta\right)
\le\Prob(\sigma<T)+16\varepsilon^2\E M_T^2/\eta^2.
\]
Najpierw wybieramy małe $\varepsilon$, potem drobne
siatki: prawa strona może być dowolnie mała.
Aby skonstruować granicę, wybieramy jeden ciąg siatek
o średnicach zbiegających do zera, np. dyadycznych.
Z własności Cauchy'ego wybieramy podciąg $(\pi_{n_k})$
tak, aby
$\Prob(\sup_{t\le T}|I^{\pi_{n_{k+1}}}_t-I^{\pi_{n_k}}_t|>2^{-k})
\le2^{-k}$.
Lemat Borela--Cantellego i zupełność przestrzeni
$C([0,T])$ z normą supremum dają jednostajną
zbieżność podciągu p.n. do ciągłej trajektorii $I$.
Poza zbiorem pełnej miary przyjmujemy zerową
trajektorię; dzięki warunkom zwykłym $I_t$ jest
$\calF_t$-mierzalna. Oszacowanie Cauchy'ego i
nierówność trójkąta pokazują, że dla każdego
innego ciągu siatek o zanikającej średnicy
granica ucp jest tym samym $I$.
Ponieważ zmienne $I_t^\pi$ są jednostajnie ograniczone
w $L^2$, zbieżność według prawdopodobieństwa implikuje
zbieżność w $L^1$ w każdej ustalonej chwili.
Przechodząc do granicy w warunkowych równościach,
stwierdzamy, że $I$ jest martyngałem.


Oszacowanie $\E|I_t^\pi|^2\le K^2\E M_t^2$
stosuje się również dla $t$ wewnątrz komórki,
ponieważ można zakończyć podział w $t$.


\emph{Krok 2: konstrukcja procesu niemalejącego.}
Z elementarnej tożsamości
$(b^2-a^2)=2a(b-a)+(b-a)^2$
na każdym odcinku podziału otrzymujemy, również
w chwilach leżących wewnątrz jego komórek,
\[
M_t^2-M_0^2=2I_t^\pi+V_t^\pi(M).
\]
Zatem $V^\pi(M)\to A:=M^2-M_0^2-2I$ w ucp.
Z ucp wybieramy podciąg podziałów, dla którego
zbieżność jest jednostajna prawie na pewno.
Proces $A$ ma wtedy trajektorie ciągłe, $A_0=0$
i $A_t\ge0$. Pokażemy, że jest niemalejący:
dla dowolnych $s<t$ proces $V^\pi$ może się obniżyć
między $s$ i $t$ tylko wskutek zmiany niedokończonego
kwadratu w komórce zawierającej $s$. Dlatego
\[
V_t^\pi(M)-V_s^\pi(M)
\ge-\max_i\sup_{u,v\in[t_i,t_{i+1}]}
|M_u-M_v|^2.
\]
Prawa strona dąży prawie na pewno do zera dzięki
jednostajnej ciągłości trajektorii. Po przejściu
do granicy $A_t\ge A_s$, najpierw dla wymiernych
$s,t$, a potem dla wszystkich dzięki ciągłości.
Ponieważ $M^2-A=2I$ jest martyngałem na $[0,T]$,
otrzymaliśmy żądany proces w przypadku ograniczonym.

\emph{Krok 3: zgodność konstrukcji na horyzontach i lokalizacja.}
Najpierw dla globalnie ograniczonego martyngału
łączymy wyniki z horyzontów $T=1,2,\ldots$.
Dla $S<T$ całkowitych porównujemy podziały $[0,T]$
zawierające $S$, których ograniczenia do $[0,S]$
mają zanikającą średnicę. Sumy na $[0,S]$ są
identyczne, więc jednoznaczność granicy ucp daje
zgodność procesów skonstruowanych na obu horyzontach.
Przeliczalność horyzontów pozwala połączyć je poza
jednym zbiorem zerowym w ciągły, adaptowany,
niemalejący proces $A$ na całej półprostej.
Proces $M^2-A$ jest martyngałem.

Dla dowolnego ciągłego lokalnego martyngału o $M_0=0$
niech $\tau_K=\inf\{t:|M_t|\ge K\}$, $K\in\mathbb N$,
$K\ge1$. Na mocy rozdziału~\ref{ch:mart-cont}
proces $M^{\tau_K}$ jest ograniczonym martyngałem,
a $\tau_K\uparrow\infty$ p.n. Oznaczmy przez $A^K$
jego skonstruowaną wariację kwadratową.
Dla każdego podziału mamy dokładną tożsamość
\[
V_t^\pi(M^{\tau_K})
=V_{t\wedge\tau_K}^\pi(M^{\tau_{K+1}}).
\]
Zbieżność ucp po prawej stronie jest zachowana
po zastąpieniu $t$ przez $t\wedge\tau_K$, ponieważ
supremum błędu nie może się zwiększyć.
Z jednoznaczności granicy ucp wynika więc
$A^K_t=A^{K+1}_{t\wedge\tau_K}$ jednocześnie
dla wszystkich $t$, poza jednym zbiorem zerowym.
Na wspólnym zbiorze pełnej miary dla wszystkich $K$
definiujemy $[M]_t$ jako stabilizującą się wartość
$A^K_t$ dla dostatecznie dużego $K$ (gdy $t\le\tau_K$).
Na dopełnieniu przyjmujemy zerową trajektorię.
Na każdym zwartym przedziale $[M]$ zgadza się z pewnym
$A^K$, więc jest ciągły, niemalejący i skończony.
Adaptacja wynika z tego, że w każdej ustalonej chwili
jest granicą $\calF_t$-mierzalnych $A^K_t$,
ze zmianą na zbiorze z $\calF_0$.
Ponadto $[M]_{t\wedge\tau_K}=A^K_t$.
W konsekwencji
\[
(M^2-[M])_{t\wedge\tau_K}
=(M^{\tau_K}_t)^2-A^K_t
\]
jest martyngałem. To dowodzi, że $M^2-[M]$
jest lokalnym martyngałem.
Dla każdego $T>0$ i $\eta>0$ mamy
\[
\Prob\left(\sup_{t\le T}|V_t^\pi(M)-[M]_t|>\eta\right)
\le \Prob(\tau_K<T)+
\Prob\left(\sup_{t\le T}|V_t^\pi(M^{\tau_K})-A^K_t|>\eta\right).
\]
Najpierw przechodzimy z $|\pi|$ do zera przy ustalonym $K$,
następnie z $K$ do nieskończoności.
Otrzymujemy zbieżność ucp dla dowolnych ciągów podziałów.

\emph{Krok 4: jednoznaczność.}
Jeśli $A$ i $A'$ spełniają tezę, to $D=A-A'$ jest
ciągłym lokalnym martyngałem (różnicą dwóch lokalnych
martyngałów po wspólnej lokalizacji)
o skończonej wariacji na przedziałach zwartych,
gdyż $A,A'$ są niemalejące i skończone.
Z twierdzenia~\ref{thm:continuous-local-finite-variation}
wynika, że $D$ jest nierozróżnialny od stałego procesu
$D_0=0$. Stąd $A=A'$ z dokładnością do nierozróżnialności.
\end{proof}

\begin{uwaga}[Co dokładnie zostało udowodnione?]
Istnienie i jednoznaczność nawiasu dla \emph{ciągłych}
lokalnych martyngałów nie wymagają w tym dowodzie
całki Itô ani pełnego twierdzenia Dooba--Meyera.
Konstrukcja używa jedynie skończonych sum przewidywalnych,
ortogonalności przyrostów, zatrzymywania i lokalizacji.
Dla procesów ze skokami trzeba osobno odróżnić
wariację kwadratową (zawierającą kwadraty skoków)
od jej przewidywalnej kompensacji; por. rozdział~\ref{ch:jump-mart}.
\end{uwaga}


\begin{twierdzenie}[Dodanie procesu o skończonej wariacji]
\label{thm:finite-var-drift-bracket}
Niech $X$ będzie procesem ciągłym mającym
wariację kwadratową w sensie ucp, a $A$
procesem ciągłym o trajektoriach lokalnie
skończonej wariacji całkowitej p.n.
Wtedy
\[
[X+A]=[X],\qquad [A]=0,
\]
a sumy mieszane $\sum_i\Delta_i^tX\Delta_i^tA$
zbiegają do zera w ucp. W szczególności wynik
dotyczy $X$ będącego ciągłym lokalnym martyngałem.
\end{twierdzenie}
\begin{proof}
Na każdym $[0,T]$ lemat
\ref{lem:finite-var-zero-quadratic} daje
$\sup_{t\le T}V_t^\pi(A)\to0$ p.n. wzdłuż
dowolnego ciągu podziałów o zanikającej średnicy.
Z nierówności Cauchy'ego--Schwarza
\[
\sup_{t\le T}\left|\sum_i\Delta_i^tX\Delta_i^tA\right|
\le
\sqrt{\sup_{t\le T}V_t^\pi(X)}
\sqrt{\sup_{t\le T}V_t^\pi(A)}.
\]
Pierwszy czynnik jest ograniczony według
prawdopodobieństwa: z ucp supremum $V^\pi(X)$
zbiega według prawdopodobieństwa do skończonego
supremum ciągłego procesu $[X]$.
Iloczyn z czynnikiem zbiegającym do zera
według prawdopodobieństwa zbiega do zera
według prawdopodobieństwa (najpierw ograniczamy
pierwszy czynnik przez dużą stałą z dużym
prawdopodobieństwem, a następnie przechodzimy
do granicy dla drugiego).
Rozwinięcie kwadratu przyrostu $X+A$ kończy dowód.
\end{proof}

\begin{przyklad}[Dodanie deterministycznego dryfu nie zmienia wariacji]
Niech $W$ będzie procesem Wienera i $b\in\R$, a
$X_t=W_t+bt$. Wtedy dla każdej siatki $\pi$ na $[0,T]$
\[
V_t^\pi(X)=V_t^\pi(W)
 +2b\sum_i\Delta_i^tW\,\Delta_i^t s
 +b^2\sum_i(\Delta_i^t s)^2,
\]
gdzie $\Delta_i^tW=W_{t\wedge t_{i+1}}-W_{t\wedge t_i}$
oraz $\Delta_i^t s=(t\wedge t_{i+1})-(t\wedge t_i)$.
Nierówność Cauchy'ego--Schwarza daje, jednostajnie w $t$,
\[
\left|\sum_i\Delta_i^tW\,\Delta_i^t s\right|
\le \sqrt{V_t^\pi(W)}\,\sqrt{T|\pi|},\qquad
\sum_i(\Delta_i^t s)^2\le T|\pi|.
\]
Z twierdzenia~\ref{thm:brownian-bracket} supremum $V_t^\pi(W)$
jest ograniczone według prawdopodobieństwa, toteż oba dodatkowe
składniki znikają w ucp. W konsekwencji $[W+bt]_t=t$.
Dryf wpływa na średnią $X_t$, ale nie na jego kwadratową
fluktuację w drobnej skali. Jest to jedno z podstawowych
rozróżnień wykorzystywanych później w rachunku It\^o.
\end{przyklad}

\begin{uwaga}[Dla ciągłego martyngału oba nawiasy się pokrywają]
Proces $[M]$ dla ciągłego lokalnego martyngału jest adaptowany
i ciągły, a zatem przewidywalny (rozdział~\ref{ch:filtracje}). Istotnie, na każdym przedziale
skończonym można go aproksymować punktowo procesami prostymi,
które na $(k2^{-n},(k+1)2^{-n}]$ przyjmują znaną już w chwili
$k2^{-n}$ wartość $[M]_{k2^{-n}}$; ich granica jest
przewidywalna jako granica funkcji mierzalnych względem
$\sigma$-ciała przewidywalnego. Różnica $(M-M_0)^2-[M]$
jest lokalnym martyngałem. W przyszłej ogólnej definicji
przewidywalnego nawiasu, opartej na kompensatorze $[M]$,
otrzymamy więc $\langle M\rangle=[M]$.
Rozbieżność tych nawiasów zaczyna mieć znaczenie przy skokach.
\end{uwaga}


\section{Kowariacja i iloczyn martyngałów}
\begin{definicja}[Kowariacja]
Dla ciągłych lokalnych martyngałów $M,N$ definiujemy
\[
[M,N]_t=\frac14\bigl([M+N]_t-[M-N]_t\bigr).
\]
Jest to ciągły proces o skończonej wariacji lokalnej.
\end{definicja}

\begin{twierdzenie}[Wzór iloczynowy dla nawiasu]
\label{thm:cross-bracket}
Niech $M,N$ będą ciągłymi lokalnymi martyngałami
z $M_0=N_0=0$. Wtedy
$M_tN_t-[M,N]_t$ jest lokalnym martyngałem oraz
$[M,N]$ jest granicą ucp sum mieszanych
$\sum_i\Delta_i^tM\Delta_i^tN$ na dowolnych
ciągach podziałów deterministycznych o średnicach
zbiegających do zera, gdzie
$\Delta_i^tM=M_{t\wedge t_{i+1}}-M_{t\wedge t_i}$
i analogicznie dla $N$.
Ponadto $[M,M]=[M]$ i
$[aM+bN,K]=a[M,K]+b[N,K]$ dla stałych $a,b$.
\end{twierdzenie}
\begin{proof}
Tożsamość
$4M_tN_t=(M_t+N_t)^2-(M_t-N_t)^2$
po odjęciu analogicznej tożsamości dla nawiasów
daje różnicę dwóch lokalnych martyngałów;
po wspólnej lokalizacji jest to martyngał.
Również dla sum po podziałach mamy
$4\sum\Delta M\Delta N=
\sum(\Delta M+\Delta N)^2-
\sum(\Delta M-\Delta N)^2$.
Zbieżność ucp obu sum kwadratów daje drugą tezę.
Dwuliniowość wynika z tej granicy, ponieważ
sumy mieszane są symetryczne i dwuliniowe dla każdego podziału,
a zatem ich granice ucp również (każda tożsamość zachodzi
z dokładnością do nierozróżnialności).
Dla $M=N$ sumy mieszane są sumami kwadratów,
co daje $[M,M]=[M]$. Wszystkie tożsamości dotyczące
przyrostów i kowariacji pozostają prawdziwe dla
dowolnych wartości początkowych, po odjęciu $M_0,N_0$.
\end{proof}

\begin{twierdzenie}[Macierz kowariacji i nierówność Kunity--Watanabego]
\label{thm:kunita-watanabe}
Niech $M=(M^{(1)},\ldots,M^{(d)})$ będzie
wektorem ciągłych lokalnych martyngałów.
Macierz $C_t=([M^{(i)},M^{(j)}]_t)_{i,j=1}^d$
jest symetryczna i jej przyrosty $C_t-C_s$
dla $s\le t$ są dodatnio półokreślone,
jednocześnie dla wszystkich $s,t$ p.n.
Dla dwóch składowych, oznaczonych $U,V$, oraz
dowolnych funkcji borelowskich $f,g$ na $[0,T]$
zachodzi trajektoryjna nierówność
\[
\int_{[0,T]} |fg|\,\mathrm d|[U,V]|
\le
\left(\int_{[0,T]} |f|^2\,\mathrm d[U]\right)^{1/2}
\left(\int_{[0,T]} |g|^2\,\mathrm d[V]\right)^{1/2},
\]
jeżeli obie całki po prawej są skończone.
Zapis $\mathrm d[U]$ oznacza miarę Lebesgue'a--Stieltjesa
trajektorii $[U]$, z zerową masą w zerze, a
$\mathrm d|[U,V]|$ oznacza wariację całkowitą
miary znakowanej indukowanej przez $[U,V]$.
W szczególności
\[
\operatorname{Var}_{[s,t]}([U,V])
\le\sqrt{([U]_t-[U]_s)([V]_t-[V]_s)}.
\]
\end{twierdzenie}
\begin{proof}
Dla $a\in\Q^d$ z dwuliniowości nawiasu mamy
$a^\top C_t a=[\sum_i a_i M^{(i)}]_t$.
Przecinając zbiory pełnej miary dla przeliczalnie
wielu $a$, otrzymujemy wszystkie te tożsamości
jednocześnie w czasie. Prawa strona jest
niemalejąca. Przez gęstość $\Q^d$ i ciągłość
formy kwadratowej w $a$ przyrosty $C_t-C_s$
są dodatnio półokreślone także dla dowolnego
$a\in\R^d$.

Ustalmy trajektorię z tego zbioru pełnej miary,
a na $[0,T]$ oznaczmy przez $\mu,\nu,\gamma$
miary Lebesgue'a--Stieltjesa odpowiednio
$[U],[V],[U,V]$. Dwie pierwsze są nieujemne
i skończone, trzecia znakowana i skończonej
wariacji. Dla dowolnych $a,b\in\R$ miara
$a^2\mu+2ab\gamma+b^2\nu$ jest nieujemna,
gdyż odpowiada niemalejącej funkcji $[aU+bV]$.
Stąd dla każdego zbioru borelowskiego $E$
macierz z elementami $\mu(E),\gamma(E),\nu(E)$
jest dodatnio półokreślona, a więc
$|\gamma(E)|^2\le\mu(E)\nu(E)$.
W szczególności $\gamma\ll\kappa=\mu+\nu$.
Z twierdzenia Radona--Nikodyma oznaczmy
$p=\mathrm d\mu/\mathrm d\kappa$,
$q=\mathrm d\nu/\mathrm d\kappa$,
$h=\mathrm d\gamma/\mathrm d\kappa$.
Mamy $p,q\ge0$, $p+q=1$ $\kappa$-p.w.
Dla wymiernych $a,b$ nieujemność powyższej
miary daje $a^2p+2abh+b^2q\ge0$ $\kappa$-p.w.
Po przecięciu przeliczalnej rodziny zbiorów
pełnej miary i przejściu do rzeczywistych $a,b$
dostajemy $|h|^2\le pq$.
Wariacja całkowita $\gamma$ ma gęstość $|h|$
względem $\kappa$ (wynika to z rozkładu na
dodatnią i ujemną część gęstości).
Zatem
\[
\int|fg|\,\mathrm d|\gamma|
\le\int |f|\sqrt p\,|g|\sqrt q\,\mathrm d\kappa
\le
\left(\int|f|^2\,\mathrm d\mu\right)^{1/2}
\left(\int|g|^2\,\mathrm d\nu\right)^{1/2}
\]
przez nierówność Cauchy'ego--Schwarza.
Dla $f=g=\chi_{(s,t]}$ otrzymujemy oszacowanie
wariacji na $[s,t]$, ponieważ miary nie mają atomów
z uwagi na ciągłość funkcji.
Argument jest trajektoryjny, więc obowiązuje
także dla funkcji zależnych od $\omega$,
o ile dla danej trajektorii są borelowskie.
\end{proof}

\begin{wniosek}[Kowariacja ruchu Browna]
Dla standardowego ruchu Browna
$B=(B^{(1)},\ldots,B^{(d)})$ zachodzi
$[B^{(i)},B^{(j)}]_t=\delta_{ij}t$.
Jeżeli $X_t=AB_t$ dla deterministycznej macierzy $A$,
to $[X^{(i)},X^{(j)}]_t=(AA^\top)_{ij}t$.
\end{wniosek}
\begin{proof}
Dla $i=j$ stosujemy twierdzenie
\ref{thm:brownian-bracket}. Dla $i\ne j$ niezależność
współrzędnych daje, dla ustalonego podziału,
\[
\E\left|\sum_k\Delta_kB^{(i)}\Delta_kB^{(j)}\right|^2
=\sum_k (t_{k+1}-t_k)^2\le T|\pi|\to0.
\]
Granica w $L^2$ jest zero, a z poprzedniego
twierdzenia musi być równa kowariacji. Zależność
dla $AB$ wynika z dwuliniowości.
\end{proof}


\section{Warunkowa wariancja przyrostu i elementarna izometria}
\begin{twierdzenie}[Wariancja warunkowa przez nawias]
\label{thm:conditional-variance-bracket}
Niech $M$ będzie ciągłym martyngałem na $[0,T]$
z $M_T\in L^2$ oraz $M_0=0$.
Wtedy $[M]_T\in L^1$, a $M^2-[M]$ jest
martyngałem na $[0,T]$.
Dla każdych $0\le s\le t\le T$ zachodzi
\[
\E[(M_t-M_s)^2\mid\calF_s]
=\E[[M]_t-[M]_s\mid\calF_s].
\]
W szczególności $\E M_t^2=\E\bigl([M]_t\bigr)$.
\end{twierdzenie}
\begin{proof}
Dla $\tau_K=\inf\{t\ge0:|M_t|\ge K\}$
(z konwencją $\inf\varnothing=\infty$ na $[0,T]$)
proces $M^{\tau_K}$ jest ograniczonym martyngałem.
Z twierdzenia~\ref{thm:general-bracket} i własności
zatrzymywania nawiasu (twierdzenie~\ref{thm:stopped-bracket})
wynika, że
$M_{t\wedge\tau_K}^2-[M]_{t\wedge\tau_K}$
jest martyngałem na $[0,T]$, rozpoczynającym się
od zera. W szczególności
$\E[M]_{t\wedge\tau_K}=\E M_{t\wedge\tau_K}^2$.
Z nierówności Dooba (twierdzenie~\ref{thm:doob-cont})
$\E\sup_{u\le T}|M_u|^2\le4\E M_T^2<\infty$.
Kwadraty zatrzymanych wartości zbiegają więc w $L^1$
do $M_t^2$ przez zbieżność zmajoryzowaną.
Ponieważ $[M]$ jest niemalejący i $\tau_K$ rosną,
zbieżność monotoniczna daje
\[
\E[M]_t=\lim_K\E[M]_{t\wedge\tau_K}=\E M_t^2<\infty.
\]
W szczególności $[M]_{t\wedge\tau_K}\to[M]_t$
w $L^1$. Przechodząc do granicy w warunkowych
równościach dla zatrzymanych martyngałów,
uzyskujemy własność martyngałową $M^2-[M]$.
Zatem
$\E[M_t^2-M_s^2\mid\calF_s]
=\E[[M]_t-[M]_s\mid\calF_s]$.
Ponieważ $M_t-M_s$ jest warunkowo wycentrowany,
$\E[2M_s(M_t-M_s)\mid\calF_s]=0$.
Rozwinięcie $(M_t-M_s)^2=M_t^2-M_s^2-2M_s(M_t-M_s)$
daje tezę. Biorąc $s=0$, otrzymujemy równość średnich.
\end{proof}

\begin{wniosek}[Warunkowa kowariancja przyrostów]
\label{cor:conditional-cross-bracket}
Dla ciągłych martyngałów $M,N$ kwadratowo
całkowalnych na $[0,T]$ z zerowymi wartościami
początkowymi proces $MN-[M,N]$ jest martyngałem i
\[
\E[(M_t-M_s)(N_t-N_s)\mid\calF_s]
=\E[[M,N]_t-[M,N]_s\mid\calF_s]
\]
dla $0\le s\le t\le T$. W szczególności
$\E M_tN_t=\E[M,N]_t$.
\end{wniosek}
\begin{proof}
Stosujemy twierdzenie
\ref{thm:conditional-variance-bracket}
do $M+N$ i $M-N$ oraz odejmujemy
odpowiednie równości, dzieląc przez cztery.
Całkowalność wynika również z nierówności
Cauchy'ego--Schwarza dla iloczynu i z
nierówności~\ref{thm:kunita-watanabe} dla wariacji
kowariacji, gdyż $\E[M]_T,\E[N]_T<\infty$.
\end{proof}

\begin{wniosek}[Izometria elementarnych strategii]
\label{cor:elementary-bracket-isometry}
Przy założeniach poprzedniego twierdzenia niech
$0=t_0<\cdots<t_n=T$ i $H_i$ będą ograniczonymi
$\calF_{t_i}$-mierzalnymi zmiennymi. Wtedy
\[
\E\left(\sum_{i=0}^{n-1}H_i(M_{t_{i+1}}-M_{t_i})\right)^2
=\E\sum_{i=0}^{n-1}H_i^2
([M]_{t_{i+1}}-[M]_{t_i}).
\]
\end{wniosek}
\begin{proof}
Stosujemy dyskretną izometrię z lematu
\ref{lem:discrete-isometry}. Dla każdej komórki
zastępujemy warunkowy drugi moment przyrostu przez
warunkowy przyrost $[M]$, zgodnie z poprzednim
twierdzeniem. Ponieważ $H_i^2$ jest
$\calF_{t_i}$-mierzalne, można je wyciągnąć przed
warunkowanie, a następnie wziąć wartość oczekiwaną.
\end{proof}

\begin{wniosek}[Izometria mieszana elementarnych strategii]
\label{cor:elementary-cross-isometry}
Przy założeniach wniosku~\ref{cor:conditional-cross-bracket}
niech $H_i,K_i$ będą ograniczonymi
$\calF_{t_i}$-mierzalnymi zmiennymi.
Wówczas
\[
\E\left[
\left(\sum_i H_i\Delta_i M\right)
\left(\sum_i K_i\Delta_i N\right)
\right]
=\E\sum_i H_iK_i
([M,N]_{t_{i+1}}-[M,N]_{t_i}).
\]
\end{wniosek}
\begin{proof}
Po rozwinięciu iloczynu wyrazy dla różnych
indeksów mają średnią zero, warunkując względem
początku późniejszego przyrostu.
Wyrazy diagonalne przekształcamy przez
wniosek~\ref{cor:conditional-cross-bracket}
i wyciągnięcie $H_iK_i$ przed warunkowanie.
Całkowalność wszystkich wyrazów zapewniają
ograniczenie współczynników, kwadratowa
całkowalność martyngałów oraz oszacowanie
wariacji z twierdzenia~\ref{thm:kunita-watanabe}.
\end{proof}

\begin{uwaga}[Gdzie w tej formule ukrywa się całka Itô?]
Lewa strona opisuje skończoną sumę przewidywalnych
zakładów, których stawka $H_i$ jest znana przed
przyrostem $M_{t_{i+1}}-M_{t_i}$. Prawa strona
mierzy ich kwadratowy koszt, ważony przyrostami
wariacji kwadratowej. W tomie III rozszerzymy tę
równość z procesów prostych na odpowiednie klasy
procesów przewidywalnych przez domknięcie podprzestrzeni w odpowiedniej przestrzeni
Hilberta; właśnie wtedy pojawi się całka stochastyczna.
\end{uwaga}

\begin{uwaga}[Dlaczego zwykła wariacja i nawias opisują inne zjawiska?]
Dla funkcji ciągłej $f$ o skończonej wariacji całkowitej
suma $\sum_i(f(t_{i+1})-f(t_i))^2$ zanika przy
zagęszczaniu podziału, niezależnie od tego, jak duża
jest całkowita zmiana $f$. Dla ruchu Browna te same
sumy dążą do długości przedziału. Nawias jest więc
miarą fluktuacji drugiego rzędu, a nie sumą bezwzględnych
przyrostów. W szczególności ciągły lokalny martyngał
może mieć skończony nawias na $[0,T]$, mimo że jego
trajektorie mają nieskończoną wariację całkowitą.
\end{uwaga}

\section{Wariacja kwadratowa a czas zatrzymania}
\begin{twierdzenie}[Wariacja kwadratowa procesu zatrzymanego]
\label{thm:stopped-bracket}
Jeśli $M$ jest ciągłym lokalnym martyngałem i
$\tau$ czasem zatrzymania, to
$[M^\tau]_t=[M]_{t\wedge\tau}$ dla $t\ge0$.
W szczególności dla procesu Wienera
$[W^\tau]_t=t\wedge\tau$.
\end{twierdzenie}
\begin{proof}
Dla każdego deterministycznego podziału $\pi$ i każdej
chwili $t$ zachodzi dokładna równość
$V_t^\pi(M^\tau)=V_{t\wedge\tau}^\pi(M)$.
Wynika to bezpośrednio z definicji: po chwili $\tau$
proces zatrzymany jest stały, a w komórce zawierającej
$\tau$ ostatni przyrost kończy się w $\tau$.
Z twierdzenia~\ref{thm:local-mart-criteria} proces
$M^\tau$ pozostaje lokalnym martyngałem.
Z twierdzenia~\ref{thm:general-bracket} obie strony
mają granice ucp, odpowiednio $[M^\tau]_t$ oraz
$[M]_{t\wedge\tau}$: dla prawej strony supremum
błędu po zatrzymaniu nie przekracza supremum
błędu na całym $[0,T]$. Jednoznaczność granicy według
prawdopodobieństwa daje równość procesów.
Przecięcie zdarzeń pełnej miary dla chwil wymiernych
i ciągłość nawiasów zapewniają jednoczesność dla wszystkich $t$.
\end{proof}

\begin{zadanie}[Brownowski martyngał i czas wyjścia z przedziału]
Niech $W$ będzie procesem Wienera,
$\tau=\inf\{t\ge0:|W_t|=a\}$ dla $a>0$,
a $M_t=W_{t\wedge\tau}$.
Wyznaczyć $[M]_t$ oraz $\E[M_t^2]$.

\emph{Rozwiązanie.} Z twierdzenia o zatrzymaniu nawiasu
$[M]_t=t\wedge\tau$. Ponieważ $|M_t|\le a$,
$M$ jest ograniczonym martyngałem, a
$M_t^2-[M]_t$ martyngałem o średniej zero przy $t=0$.
Dlatego $\E M_t^2=\E(t\wedge\tau)$.
Gdy $t\to\infty$, lewa strona dąży do $a^2$,
a prawa do $\E\tau$ dzięki zbieżności monotonicznej.
Otrzymujemy ponownie $\E\tau=a^2$, bez używania
całki Itô.
\end{zadanie}


% ================================================================
% ROZDZIAŁ 14. SKOKI, KOMPENSATORY, PODSUMOWANIE TOMU
% ================================================================

\chapter{Martyngały ze skokami, kompensacja i granica zakresu tomu II}
\label{ch:jump-mart}

W całym rozdziale przyjmujemy warunki zwykłe filtracji
(rozdział~\ref{ch:mart-cont}). Procesy c\`adl\`ag
rozpatrujemy w wersjach o takich trajektoriach
poza jednym zbiorem zerowym; na tym zbiorze
można przyjąć zerową trajektorię dzięki zupełności.
Konstrukcję ograniczamy do przypadków, w których
można uzasadnić wszystkie własności bez ogólnej
teorii całkowania względem półmartyngałów.

\section{Dlaczego trzeba oddzielić dwa rodzaje wariacji kwadratowej?}
Dla ciągłego lokalnego martyngału jego wariacja kwadratowa
jest ciągłym procesem rosnącym i opisuje akumulację
losowych fluktuacji. Przy skokach pojawia się dodatkowe
zjawisko: pojedynczy skok o wysokości $\Delta X_t$
wnosi do wariacji kwadratowej dokładnie $(\Delta X_t)^2$.
Dla martyngałów z częścią skokową trzeba odróżniać
rzeczywistą wariację kwadratową od jej przewidywalnej
kompensacji. Nie jest ona na ogół warunkową średnią
całego nawiasu w jednej ustalonej chwili. Obie wielkości są użyteczne, ale w ogólności
nie są sobie równe.

\begin{definicja}[Skok i suma kwadratów skoków]
Dla procesu c\`adl\`ag $X$ definiujemy skok
$\Delta X_t=X_t-X_{t-}$, $t>0$. Dla trajektorii lokalnie skończonej wariacji
całkowitej poniżej udowodnimy, że jej wariacja
kwadratowa wynosi
\[
[X]_t=\sum_{0<s\le t}(\Delta X_s)^2.
\]
Dla szerszych klas procesów granica sum kwadratów
przyrostów wymaga dodatkowej konstrukcji, stanowiącej
część teorii półmartyngałów rozwijanej w tomie III.
\end{definicja}


\begin{definicja}[Wariacja kwadratowa w klasie procesów z niniejszego rozdziału]
Dla rzeczywistego procesu c\`adl\`ag $X$ i podziału
$\pi=\{0=t_0<\cdots<t_m=T\}$ definiujemy
$V_t^\pi(X)=\sum_i(X_{t\wedge t_{i+1}}-X_{t\wedge t_i})^2$.
Jeśli dla każdego ciągu deterministycznych podziałów o
średnicach dążących do zera procesy $V^\pi(X)$ zbiegają
jednostajnie na $[0,T]$ według prawdopodobieństwa do
procesu c\`adl\`ag $A$, piszemy $[X]=A$ na $[0,T]$.
Jednoznaczność z dokładnością do nierozróżnialności
wynika z nierówności trójkąta dla supremum i
przeliczalnej rodziny horyzontów, jak w rozdziale
\ref{ch:quadratic-cont}.
Istnienie granicy wykażemy dla procesów o lokalnie
skończonej wariacji i ich sum z ciągłymi procesami
o istniejącej wariacji kwadratowej. Ogólną teorię
półmartyngałów pozostawiamy tomowi III.
Ważne jest użycie niedokończonego przyrostu:
gdy $t$ jest chwilą skoku, jego kwadrat jest
uwzględniony już w chwili $t$, a nie dopiero
w następnym węźle podziału.
\end{definicja}

\begin{lemat}[Wariacja kwadratowa funkcji o skończonej wariacji]
\label{lem:cadlag-finite-variation-qv}
Niech $f:[0,T]\to\R$ będzie funkcją c\`adl\`ag
o skończonej wariacji całkowitej.
Dopuszczamy nieskończenie wiele skoków. Dla dowolnych deterministycznych podziałów
o średnicach malejących do zera zachodzi
\[
\sup_{t\le T}\left|V_t^\pi(f)-
\sum_{0<s\le t}(\Delta f_s)^2\right|\longrightarrow0.
\]
W szczególności zbieżność zachodzi również w chwilach
skoku, z prawostronną konwencją $f_t$ jako wartości
trajektorii już po skoku.
\end{lemat}
\begin{proof}
Najpierw zakładamy skończoną liczbę skoków.
Rozłóżmy $f=g+j$, gdzie
$j(t)=\sum_{0<s\le t}\Delta f_s$,
a $g=f-j$ jest ciągła i ma skończoną wariację
całkowitą. Dla dostatecznie drobnych podziałów
każda półotwarta komórka $(t_i,t_{i+1}]$
zawiera co najwyżej jeden skok.
Suma $(\Delta j)^2$ po pełnych komórkach oraz
aktualnej komórce jest dokładnie sumą kwadratów
skoków, które wystąpiły nie później niż $t$. Suma $(\Delta g)^2$ dąży do zera
z lematu~\ref{lem:finite-var-zero-quadratic}. Wyrazy mieszane są
ograniczone przez nierówność Cauchy'ego--Schwarza:
\[
\left|2\sum\Delta g\Delta j\right|
\le2\Bigl(\sum(\Delta g)^2\Bigr)^{1/2}
\Bigl(\sum(\Delta j)^2\Bigr)^{1/2}\longrightarrow0.
\]
W chwili skoku $t$ wartość $f_t$ zawiera już skok,
a więc odpowiadający mu kwadrat pojawia się również
w sumie $V_t^\pi$. Oszacowania są jednostajne w $t$,
ponieważ liczba skoków jest skończona, a funkcja $g$
jest jednostajnie ciągła.

W ogólnym przypadku dla każdego skończonego zbioru
chwil skoku suma modułów skoków jest nie większa
od $\operatorname{Var}_{[0,T]}(f)$: dobieramy rozłączne
przedziały kończące się w tych chwilach i przechodzimy
z lewymi końcami do chwil skoków, używając lewych granic.
W szczególności dla każdego $\varepsilon>0$ tylko
skończenie wiele skoków ma moduł większy od
$\varepsilon$, wszystkie niezerowe skoki tworzą
zbiór przeliczalny oraz
$\sum_{0<s\le T}|\Delta f_s|\le\operatorname{Var}_{[0,T]}(f)$.
Szereg $j(t)=\sum_{0<s\le t}\Delta f_s$
jest jednostajnie zbieżny i określa funkcję
c\`adl\`ag o skończonej wariacji z tymi samymi
skokami co $f$. Istotnie sumy skończone mają
te własności, a jednostajna granica zachowuje
prawostronną ciągłość i lewe granice.
Wariacja ogona jest nie większa niż suma modułów
odrzuconych skoków. Funkcja $g=f-j$ jest ciągła.
Wybierzmy skończone sumy skokowe $j_n$ z
$\operatorname{Var}(j-j_n)\to0$.
Dla każdego podziału i $t\le T$ nierówność
trójkąta dla normy euklidesowej daje
\[
\left|\sqrt{V_t^\pi(f)}-\sqrt{V_t^\pi(g+j_n)}\right|
\le\sqrt{V_t^\pi(j-j_n)}
\le\operatorname{Var}_{[0,T]}(j-j_n).
\]
Obie pierwiastkowane sumy są jednostajnie
ograniczone przez
$\operatorname{Var}(g)+\operatorname{Var}(j)+
\sum_{0<s\le T}|\Delta f_s|$.
Zatem także różnica samych kwadratów jest
jednostajnie mała przy dużym $n$, niezależnie
od siatki i $t$. Suma kwadratów skoków
$j_n$ zbiega jednostajnie do sumy kwadratów
wszystkich skoków, gdyż szereg kwadratów
jest zbieżny. Stosujemy przypadek skończony
przy ustalonym $n$, a następnie przechodzimy
z $n$ do nieskończoności.
\end{proof}

\begin{definicja}[Kompensator i przewidywalny nawias]
\label{def:jump-predictable-bracket}
Dla adaptowanego niemalejącego procesu c\`adl\`ag
$V$ z $V_0=0$ kompensatorem nazywamy przewidywalny,
niemalejący proces c\`adl\`ag $A$ z $A_0=0$,
taki że $V-A$ jest lokalnym martyngałem.
W przypadku ciągłych lokalnych martyngałów
$[M]$ jest już przewidywalny i sam jest swoim
kompensatorem. Dla lokalnie kwadratowo
całkowalnego martyngału $M$ z $M_0=0$
przewidywalny nawias $\langle M\rangle$,
gdy istnieje, jest przewidywalnym niemalejącym
procesem z wartością początkową zero, takim że
$M^2-\langle M\rangle$ jest lokalnym martyngałem.
W rozważanych niżej przykładach sprawdzimy
bezpośrednio, że również $M^2-[M]$ jest
lokalnym martyngałem, a zatem $\langle M\rangle$
jest kompensatorem $[M]$.
Jednoznaczność udowodnimy poniżej. Ogólne
twierdzenie o istnieniu kompensatora wymaga
dodatkowej teorii (twierdzenia Dooba--Meyera)
i nie jest tutaj zakładane.
\end{definicja}

\begin{twierdzenie}[Nawias i kompensator procesu Poissona]
\label{thm:poisson-brackets}
Niech $N$ będzie procesem Poissona o intensywności
$\lambda>0$, w jego uzupełnionej filtracji naturalnej,
a $M_t=N_t-\lambda t$. Wtedy $M$ jest kwadratowo
całkowalnym martyngałem z
\[
[M]_t=N_t,\qquad \langle M\rangle_t=\lambda t.
\]
Pierwszy proces jest wariacją kwadratową trajektorii,
a drugi \emph{przewidywalnym nawiasem}: jest ciągły,
niemalejący i przewidywalny, a $M_t^2-\langle M\rangle_t$
jest martyngałem. Jednoznaczność tej kompensacji w klasie procesów
przewidywalnych wynika z wniosku
\ref{cor:predictable-compensator-unique}.
Ponadto $\E M_t^2=\lambda t$.
\end{twierdzenie}
\begin{proof}
Proces $M$ jest kwadratowo całkowalnym martyngałem,
gdyż przyrosty $N$ są niezależne od przeszłości
i mają średnią $\lambda(t-s)$ oraz skończony
drugi moment (rozdział~\ref{ch:poisson-process}).
Trajektorie $M$ mają skoki równe $1$ w chwilach
zdarzeń Poissona i liniowe odcinki o nachyleniu
$-\lambda$ pomiędzy nimi. Na każdym skończonym
przedziale liczba skoków jest skończona prawie na pewno,
więc poprzedni lemat daje
$[M]_t=\sum_{0<s\le t}1=N_t$.
Dla $s<t$, pisząc $U=N_t-N_s$, mamy
$M_t=M_s+(U-\lambda(t-s))$.
Zmienna $U$ jest niezależna od $\calF_s$, a jej
średnia i wariancja wynoszą $\lambda(t-s)$.
Stąd
\[
\E[M_t^2\mid\calF_s]
=M_s^2+\lambda(t-s),
\]
co pokazuje, że $M_t^2-\lambda t$ jest martyngałem.
Deterministyczna funkcja $t\mapsto\lambda t$ jest
przewidywalna. Równanie przy $s=0$ daje też
$\E M_t^2=\lambda t$.
Ponieważ $[M]-\lambda t=N-\lambda t=M$
jest martyngałem, także $M^2-[M]$
jest martyngałem jako różnica dwóch martyngałów.
\end{proof}


\begin{lemat}[Jednoznaczność ciągłego kompensatora kwadratu]
\label{lem:continuous-compensator-unique}
Niech $M$ będzie lokalnym martyngałem c\`adl\`ag, a $A$ i $B$
ciągłymi, adaptowanymi procesami o lokalnie skończonej wariacji,
$A_0=B_0=0$. Jeżeli zarówno $M^2-A$, jak i $M^2-B$
są lokalnymi martyngałami, to $A$ i $B$ są nierozróżnialne.
\end{lemat}
\begin{proof}
Różnica $D=A-B=(M^2-B)-(M^2-A)$ jest ciągłym
lokalnym martyngałem o skończonej wariacji na przedziałach
zwartych. Z twierdzenia~\ref{thm:continuous-local-finite-variation}
wynika, że każdy taki proces jest stały. Ponieważ $D_0=0$,
proces $D$ jest nierozróżnialny od zera.
\end{proof}

\begin{lemat}[Całkowanie względem martyngału o całkowalnej wariacji]
\label{lem:predictable-fv-mean-zero}
Niech $D$ będzie rzeczywistym martyngałem c\`adl\`ag
z $D_0=0$ oraz $\E\operatorname{Var}_{[0,T]}(D)<\infty$.
Dla każdego ograniczonego procesu przewidywalnego $H$
zachodzi
\[
\E\int_{(0,T]} H_t\,\mathrm dD_t=0.
\]
Całka jest tutaj zwykłą trajektoryjną całką
względem miary znakowanej Lebesgue'a--Stieltjesa,
a nie całką It\^o.
\end{lemat}
\begin{proof}
Dla generatorów $\chi_F(\omega)\chi_{(a,b]}(t)$,
$F\in\calF_a$ i $0\le a<b\le T$, lewa strona
wynosi $\E[\chi_F(D_b-D_a)]=0$.
Rodzina tych prostokątów, wraz ze zbiorem pustym,
jest zamknięta na skończone przecięcia i generuje
$\sigma$-ciało przewidywalne ograniczone do
$(0,T]\times\Omega$ (rozdział~\ref{ch:filtracje}).
Dla każdego zbioru z tego $\sigma$-ciała całka
jego funkcji charakterystycznej jest mierzalna:
najpierw dla prostokątów, a następnie przez
twierdzenie o klasie monotonicznej i trajektoryjną
zbieżność zmajoryzowaną względem miary wariacji.
Funkcja zbioru otrzymana przez wzięcie wartości
oczekiwanej jest skończoną miarą znakowaną,
ponieważ całkowalna wariacja majoruje moduły
całek i pozwala przechodzić do granicy dla
sum rozłącznych. Jej wartość na całej przestrzeni
wynosi $\E D_T=0$. Z jednoznaczności miar
na układzie generującym (lub twierdzenia
$\pi$--$\lambda$ dla zbiorów, na których
miara ma wartość zero) jest identycznie zerowa.
Aproksymacja ograniczonego $H$ funkcjami
prostymi kończy dowód.
\end{proof}

\begin{twierdzenie}[Przewidywalny lokalny martyngał o skończonej wariacji]
\label{thm:predictable-fv-local-constant}
Jeżeli $D$ jest przewidywalnym lokalnym martyngałem
c\`adl\`ag o lokalnie skończonej wariacji całkowitej,
to $D$ jest nierozróżnialny od procesu stałego $D_0$.
\end{twierdzenie}
\begin{proof}
Odejmujemy $D_0$ i niech
$V_t=\operatorname{Var}_{[0,t]}(D)$.
Proces $V$ jest adaptowany, niemalejący i c\`adl\`ag:
wariację można obliczać po podziałach z wymiernymi
punktami wewnętrznymi oraz końcem $t$ dzięki
prawostronnej ciągłości; miara wariacji trajektorii
ma atom w $t>0$ o masie $|\Delta D_t|$.
Zatem $\Delta V_t=|\Delta D_t|$.
Niech $(\sigma_n)$ będzie ciągiem lokalizującym $D$
i połóżmy
\[
\tau_n=\sigma_n\wedge n\wedge
\inf\{t\ge0: V_t\ge n\}.
\]
Są to czasy zatrzymania, gdyż dla niemalejącego
prawostronnie ciągłego $V$ zdarzenie trafienia
do poziomu $n$ do chwili $t$ jest równe
$\{V_t\ge n\}$. Mamy $\tau_n\uparrow\infty$ p.n.
Proces $D^{\tau_n}$ jest martyngałem na mocy
twierdzenia~\ref{thm:mart-cont-stopping}
zastosowanego do $D^{\sigma_n}$.
Ponadto jest przewidywalny: dla dowolnego
czasu zatrzymania $\tau$ rozkład
\[
D_{t\wedge\tau}=D_t\chi_{\{t\le\tau\}}+
D_\tau\chi_{\{\tau<t\}}
\]
składa się z procesów przewidywalnych
(w drugim składniku przyjmujemy zero przed
$\tau$, także gdy $\tau=\infty$).
Wskaźnik pierwszego zdarzenia jest lewostronnie
ciągły i adaptowany, podobnie drugi składnik
jest lewostronnie ciągłym procesem adaptowanym.
Wartość w zerze jest mierzalna względem $\calF_0$.

Przed $\tau_n$ wariacja jest ograniczona przez $n$,
więc $V_{\tau_n}\le n+|\Delta D_{\tau_n}|$
(przy $\tau_n=0$ przyjmujemy skok równy zero).
Ponieważ $|D_{\tau_n-}|\le n$ i $D_{\tau_n}$
jest całkowalna na mocy ograniczonego próbkowania
opcjonalnego dla $D^{\sigma_n}$, otrzymujemy
$\E V_{\tau_n}<\infty$.
Zatrzymany proces ma zatem całkowalną wariację
na każdym horyzoncie.
W dalszej części oznaczmy go przez $L$.

Proces $L_-$ jest przewidywalny jako adaptowany
proces o lewostronnie ciągłych trajektoriach
dla $t>0$. Stąd
$H_t=\operatorname{sgn}(\Delta L_t)$
jest ograniczony i przewidywalny
(przyjmujemy $H_0=0$ oraz $\operatorname{sgn}0=0$).
Każda trajektoria o skończonej wariacji
ma co najwyżej przeliczalnie wiele skoków.
Ponieważ $H$ znika poza nimi, mamy
\[
\int_{(0,T]}H_t\,\mathrm dL_t
=\sum_{0<t\le T}|\Delta L_t|.
\]
Lemat~\ref{lem:predictable-fv-mean-zero}
daje średnią zero tej nieujemnej sumy,
więc $L$ nie ma skoków p.n. na $[0,T]$.
Po przecięciu zdarzeń pełnej miary dla
całkowitych $T$ proces $L$ jest ciągły.
Twierdzenie~\ref{thm:continuous-local-finite-variation}
daje $L=0$ z dokładnością do nierozróżnialności.
Przeliczalność lokalizacji i $\tau_n\uparrow\infty$
kończą dowód dla $D$.
\end{proof}

\begin{wniosek}[Jednoznaczność kompensatorów przewidywalnych]
\label{cor:predictable-compensator-unique}
Jeśli $A,B$ są kompensatorami tego samego
procesu $V$ z definicji
\ref{def:jump-predictable-bracket}, to są
nierozróżnialne. Także przewidywalny nawias,
jeśli istnieje, jest jednoznaczny.
\end{wniosek}
\begin{proof}
Różnica $A-B$ jest przewidywalnym lokalnym
martyngałem o lokalnie skończonej wariacji
z wartością początkową zero, zarówno gdy
odejmujemy dwa kompensatory, jak i dwa nawiasy
przewidywalne. Stosujemy twierdzenie
\ref{thm:predictable-fv-local-constant}.
\end{proof}

\begin{uwaga}[Co oznacza kompensator?]
Wzór $\langle M\rangle_t=\lambda t$ nie oznacza,
że losowa liczba skoków $N_t$ jest równa $\lambda t$.
Oznacza, że $[M]_t-\langle M\rangle_t=N_t-\lambda t$
jest martyngałem: po odjęciu przewidywalnego tempa
wzrostu pozostaje wycentrowana fluktuacja.
Przewidywalne kompensatory mogą też mieć skoki.
Ich jednoznaczność wynika z wniosku
\ref{cor:predictable-compensator-unique};
ogólną konstrukcję pozostawiamy tomowi III.
\end{uwaga}


\begin{uwaga}[Różnica między nawiasami jest rzeczywista]
W chwili skoku $N$ wariacja kwadratowa $[M]$ skacze
o jeden, a przewidywalny nawias $\langle M\rangle$
pozostaje ciągły. Różnica $N_t-\lambda t=M_t$
sama jest martyngałem. Nie należy zatem utożsamiać
$[M]$ i $\langle M\rangle$, chociaż ich wartości
oczekiwane są równe: $\E\bigl([M]_t\bigr)=\E\bigl(\langle M\rangle_t\bigr)=\lambda t$.
\end{uwaga}

\begin{przyklad}[Przewidywalny skok nawiasu]
\label{ex:predictable-bracket-jump}
Niech $Y$ przyjmuje wartości $-1,1$ z równymi
prawdopodobieństwami, a filtracja będzie trywialna
przed chwilą $1$ i równa $\sigma(Y)$ od chwili $1$
(z uzupełnieniem). Jest to filtracja prawostronnie
ciągła. Proces $M_t=Y\chi_{[1,\infty)}(t)$
jest martyngałem: przed ujawnieniem $Y$ jego
warunkowa średnia jest zerowa, a potem proces
jest stały.
W tym przypadku
\[
[M]_t=Y^2\chi_{[1,\infty)}(t)
=\chi_{[1,\infty)}(t)=\langle M\rangle_t.
\]
Proces deterministyczny $\chi_{[1,\infty)}$
jest przewidywalny, gdyż zbiór $[1,\infty)$
jest przeliczalnym przecięciem przedziałów
$(1-1/n,\infty)$, mierzalnych w
$\sigma$-ciele przewidywalnym (rozdział
\ref{ch:filtracje}). Różnica
$M^2-\langle M\rangle$ jest zerowa.
Nawias przewidywalny może więc mieć skok,
a procesy o skokach nie muszą mieć różnych
nawiasów. W tym przykładzie nie istnieje ciągły
kompensator kwadratu nawet w sensie lokalnym:
każdy ciągły adaptowany kandydat jest
przewidywalny i z jednoznaczności musiałby
być nierozróżnialny od $\chi_{[1,\infty)}$.
\end{przyklad}

\begin{center}
\begin{tikzpicture}[x=1.5cm,y=0.8cm,>=Latex]
\draw[->] (0,0)--(4.5,0) node[right] {$t$};
\draw[->] (0,0)--(0,4.5) node[above] {nawias};
\draw[blue!75!black,very thick] (0,0)--(.7,0) (.7,1)--(1.6,1) (1.6,2)--(3.2,2) (3.2,3)--(4,3);
\foreach \x/\y in {.7/1,1.6/2,3.2/3}{\fill[blue!75!black] (\x,\y) circle (1.7pt); \draw[blue!75!black,fill=white] (\x,{\y-1}) circle (1.7pt);}
\draw[orange!80!black,thick,dashed] (0,0)--(4,4);
\node[blue!75!black,below] at (3.7,3) {$[M]_t=N_t$};
\node[orange!80!black,above left] at (3.6,3.6) {$\langle M\rangle_t=t$};
\end{tikzpicture}
\par\smallskip\small Przykładowa trajektoria nawiasu
kompensowanego procesu Poissona o intensywności $1$.
Nawias przewidywalny jest deterministyczny, choć
zrealizowana liczba skoków jest losowa.
\end{center}

\section{Czasy przewidywalne i skoki procesów zliczających}

\begin{definicja}[Czas przewidywalny i całkowicie nieosiągalny]
Czas zatrzymania $\tau$ nazywamy przewidywalnym,
jeżeli istnieje niemalejący ciąg czasów zatrzymania
$\tau_n\uparrow\tau$ p.n. oraz $\tau_n<\tau$
na $\{\tau>0\}$ p.n. (dla $\tau=\infty$
żądamy $\tau_n<\infty$).
Czasy $\tau_n$ nazywamy ciągiem zapowiadającym.
Czas zatrzymania $S$ nazywamy całkowicie
nieosiągalnym, jeżeli dla każdego
przewidywalnego czasu zatrzymania $\tau$
\[
\Prob(0<S=\tau<\infty)=0
\]
oraz $\Prob(S=0)=0$.
Dodatnia deterministyczna chwila jest
przewidywalna, np. przez ciąg $(t-1/n)^+$.
\end{definicja}

\begin{lemat}[Wykres czasu przewidywalnego]
\label{lem:predictable-time-graph}
Jeżeli $\tau$ jest przewidywalnym czasem zatrzymania,
to proces $\chi_{\{0<t=\tau<\infty\}}$
jest przewidywalny.
\end{lemat}
\begin{proof}
Dla dowolnych czasów zatrzymania $\alpha,\beta$
proces $\chi_{\{\alpha<t\le\beta\}}$
jest lewostronnie ciągły i adaptowany
na $(0,\infty)$, zatem przewidywalny.
Dla ciągu zapowiadającego mamy punktowo
poza wspólnym zbiorem zerowym
\[
\{0<t=\tau<\infty\}
=\bigcap_{n\ge1}\{\tau_n<t\le\tau\},
\qquad t>0.
\]
Przy $\tau=\infty$ żaden skończony $t$
nie należy do wszystkich zbiorów po prawej.
Zupełność filtracji umożliwia zmianę
na wspólnym zbiorze zerowym.
\end{proof}

\begin{twierdzenie}[Tożsamość kompensacyjna]
\label{thm:counting-compensation}
Niech $N$ będzie adaptowanym procesem zliczającym
z $N_0=0$ i jednostkowymi skokami, nieeksplodującym,
a $A$ przewidywalnym niemalejącym procesem
c\`adl\`ag z $A_0=0$. Załóżmy, że $N-A$
jest martyngałem oraz $\E N_T+\E A_T<\infty$
dla każdego $T<\infty$. Wtedy dla każdego
nieujemnego procesu przewidywalnego $H$
\[
\E\sum_{0<s\le T}H_s\Delta N_s
=\E\int_{(0,T]}H_s\,\mathrm dA_s
\]
z dopuszczeniem wartości nieskończonej.
Dla procesu Poissona o intensywności $\lambda$
w jego uzupełnionej filtracji naturalnej
prawa strona jest równa
$\lambda\E\int_0^T H_s\,\mathrm d\lambda_1(s)$.
Jeśli $A$ jest ciągły, wszystkie kolejne czasy
skoków $N$ są całkowicie nieosiągalne.
\end{twierdzenie}
\begin{proof}
Dla ograniczonego $H$ stosujemy lemat
\ref{lem:predictable-fv-mean-zero} do $D=N-A$,
którego wariacja na $[0,T]$ jest
ograniczona przez $N_T+A_T$.
Całka względem $N$ jest sumą po chwilach
jego skoków. Dla dowolnego nieujemnego $H$
przechodzimy od $H\wedge n$ przez zbieżność
monotoniczną po obu stronach.
Przypadek Poissona wynika z twierdzenia
\ref{thm:poisson-brackets}.

Dla przewidywalnego $\tau$ stosujemy tożsamość
do procesu $H_t=\chi_{\{0<t=\tau<\infty\}}$,
przewidywalnego z lematu
\ref{lem:predictable-time-graph}.
Przy ciągłym $A$ jego miara nie ma atomów,
więc prawa strona wynosi zero.
Otrzymujemy $\Prob(0<\tau\le T,\Delta N_\tau=1)=0$.
Przechodząc po całkowitych $T$ do nieskończoności,
widzimy, że żaden kolejny czas skoku
$S_k=\inf\{t:N_t\ge k\}$ nie może z dodatnim
prawdopodobieństwem pokryć się w skończonej
chwili z czasem przewidywalnym.
Czasy $S_k$ są czasami zatrzymania przez
adaptację, monotoniczność i prawostronną
ciągłość $N$, a $S_k>0$ p.n. dzięki $N_0=0$
i prawostronnej ciągłości. To kończy dowód.
\end{proof}

\begin{twierdzenie}[Warunkowe momenty skoku w czasie przewidywalnym]
\label{thm:predictable-jump-moments}
Niech $M$ będzie martyngałem c\`adl\`ag
z $M_0=0$ i $M_T\in L^2$, a $A=\langle M\rangle$
niech istnieje i spełnia $\E A_T<\infty$.
Niech $\tau$ będzie przewidywalnym czasem
zatrzymania, $0<\tau\le T$ p.n.
Definiujemy informację bezpośrednio przed $\tau$ przez
\[
\calF_{\tau-}=\sigma\bigl(\calF_0,\,
\{F\cap\{u<\tau\}:u\ge0,\ F\in\calF_u\}\bigr).
\]
Wówczas
\[
\E[\Delta M_\tau\mid\calF_{\tau-}]=0,\qquad
\E[(\Delta M_\tau)^2\mid\calF_{\tau-}]=\Delta A_\tau.
\]
\end{twierdzenie}
\begin{proof}
Niech $\tau_n$ będzie ciągiem zapowiadającym
$\tau$. Możemy przyjąć $0\le\tau_n<\tau\le T$
i $\tau_n\uparrow\tau$ poza jednym zbiorem
zerowym. Z dokładnością do uzupełnienia
\[
\calF_{\tau-}=\sigma\left(\bigcup_n\calF_{\tau_n}\right).
\]
Istotnie, dla $F\in\calF_{\tau_n}$ mamy
$F=\bigcup_{q\in\Q_{\ge0}}(F\cap\{\tau_n\le q\}
\cap\{q<\tau\})$ p.n., co daje jedną inkluzję.
Odwrotnie, dla $F\in\calF_u$ zdarzenie
$F\cap\{u<\tau\}$ jest sumą zbiorów
$F\cap\{u<\tau_n\}\in\calF_{\tau_n}$.
Ostatnia mierzalność wynika bezpośrednio
z testu definicji $\calF_{\tau_n}$.
Ponadto $\calF_0\subseteq\calF_{\tau_n}$.

Dla każdego przewidywalnego procesu $H$
zmienna $H_\tau$ jest $\calF_{\tau-}$-mierzalna.
Wystarczy sprawdzić to dla funkcji
charakterystycznych generatorów
$F\times(a,b]$, $F\in\calF_a$:
ich wartość w $\tau$ jest wskaźnikiem
$F\cap\{a<\tau\}\cap\{\tau\le b\}$.
Oba zdarzenia są $\calF_{\tau-}$-mierzalne.
Własność rozszerza się na całe $\sigma$-ciało
przewidywalne przez domknięcie na dopełnienia
i przeliczalne sumy. W szczególności
$A_\tau$ i $A_{\tau-}$ są mierzalne względem
$\calF_{\tau-}$ (proces $A_-$ jest przewidywalny),
podobnie $M_{\tau-}$.

Proces $M^2-A$ jest prawdziwym martyngałem
na $[0,T]$, ponieważ moduły jego wartości
w czasach zatrzymania są zmajoryzowane przez
$\sup_{s\le T}|M_s|^2+A_T\in L^1$
(twierdzenia~\ref{thm:doob-cont}
i~\ref{thm:local-mart-criteria}).
Dla $Y=M$ albo $Y=M^2-A$ ograniczone
próbkowanie opcjonalne daje
$\E[Y_\tau\mid\calF_{\tau_n}]=Y_{\tau_n}$.
Prawa strona zbiega p.n. i w $L^1$
do $Y_{\tau-}$ przez powyższą majorantę
(dla $Y=M$ wystarcza $\sup_{s\le T}|M_s|$).
Dla $F\in\bigcup_n\calF_{\tau_n}$
przechodzimy do granicy w równości
$\E[\chi_FY_\tau]=\E[\chi_FY_{\tau_n}]$
dla dostatecznie dużego $n$.
Twierdzenie o klasie monotonicznej rozszerza
równość na $\calF_{\tau-}$.
Stąd $\E[\Delta M_\tau\mid\calF_{\tau-}]=0$
i $\E[\Delta(M^2)_\tau-\Delta A_\tau
\mid\calF_{\tau-}]=0$.
Ponieważ
$\Delta(M^2)_\tau=2M_{\tau-}\Delta M_\tau+
(\Delta M_\tau)^2$,
a iloczyn jest całkowalny przez nierówność
Cauchy'ego--Schwarza, pierwsza równość
zeruje wyraz mieszany po warunkowaniu
(można najpierw obciąć $M_{\tau-}$).
Mierzalność $\Delta A_\tau$ daje drugą tezę.
\end{proof}

\begin{uwaga}[Intensywność nie musi być deterministyczna]
\label{rem:stochastic-counting-intensity}
Jeżeli kompensator procesu zliczającego ma postać
$A_t=\int_0^t a_s\,\mathrm d\lambda_1(s)$,
gdzie $a$ jest nieujemnym procesem przewidywalnym,
nazywamy $a$ intensywnością stochastyczną.
Tożsamość~\ref{thm:counting-compensation} daje
\[
\E\sum_{0<s\le T}H_s\Delta N_s
=\E\int_0^T H_s a_s\,\mathrm d\lambda_1(s).
\]
Intensywność jest jednoznaczna z dokładnością
do równości prawie wszędzie względem
$\lambda_1\otimes\Prob$: jednoznaczność
kompensatora daje równość wszystkich całek
do chwili $t$ na wspólnym zbiorze pełnej miary,
a więc równość gęstości prawie wszędzie
dla każdej takiej trajektorii i przez twierdzenie
Fubiniego. Nie wynika stąd niezależność
przyrostów procesu zliczającego.
\end{uwaga}

\begin{przyklad}[Losowa stała intensywność]
\label{ex:random-rate-counting}
Niech $\Lambda$ przyjmuje wartości $1,2$ z równymi
prawdopodobieństwami, a $P$ będzie niezależnym
procesem Poissona o intensywności $1$.
Definiujemy $N_t=P_{\Lambda t}$ i filtrację
generowaną przez $\Lambda$ oraz obserwacje
$N_s$ dla $s\le t$, z warunkami zwykłymi.
Warunkowo względem $\Lambda=\ell$ jest to
proces Poissona o intensywności $\ell$.
Niezależność przyrostów warunkowo względem
tej zmiennej daje
$\E[N_t-N_s\mid\calF_s]=\Lambda(t-s)$
oraz warunkową wariancję $\Lambda(t-s)$.
Równości zachowują się po zwykłym uzupełnieniu
filtracji: warunkujemy najpierw dla
$s_n\downarrow s$ i stosujemy odwrotne
warunkowe wartości oczekiwane oraz
prawostronną ciągłość procesów.
Stąd $N_t-\Lambda t$ i
$(N_t-\Lambda t)^2-\Lambda t$ są martyngałami.
Zatem kompensator $A_t=\Lambda t$ jest losowy.
Natomiast dla $h,k>0$, z prawa kowariancji
całkowitej,
\[
\operatorname{Cov}(N_h,N_{h+k}-N_h)
=hk\operatorname{Var}(\Lambda)=hk/4>0.
\]
Przyrosty na rozłącznych przedziałach
nie są więc niezależne bez warunkowania.
\end{przyklad}

\section{Złożony proces Poissona i jego fluktuacje}
\begin{twierdzenie}[Kompensowane sumy losowych skoków]
\label{thm:compound-poisson-brackets}
Niech $N$ będzie procesem Poissona o intensywności
$\lambda>0$, a $(Y_k)_{k\ge1}$ ciągiem niezależnych
zmiennych o jednakowym rozkładzie, niezależnym od $N$.
Załóżmy $\E Y_1^2<\infty$, oznaczmy
$\mu=\E Y_1$, $m_2=\E Y_1^2$ i zdefiniujmy
\[
J_t=\sum_{k=1}^{N_t}Y_k,\qquad
M_t=J_t-\lambda\mu t.
\]
Wtedy $M$ jest martyngałem w naturalnej filtracji
obserwacji zdarzeń i ich oznaczeń, z
\[
\E M_t^2=\lambda m_2t,\qquad
[M]_t=\sum_{k=1}^{N_t}Y_k^2,\qquad
\langle M\rangle_t=\lambda m_2t.
\]
\end{twierdzenie}
\begin{proof}
Stosujemy własność niezależnych przyrostów
w filtracji obserwowanych oznaczeń z rozdziału
\ref{ch:poisson-process}. W szczególności nie
ujawniamy z góry całego ciągu $(Y_k)$.
Przyrost $J_t-J_s$ jest niezależny od informacji
do chwili $s$. Warunkowo względem $N_t-N_s=n$ jest sumą
$n$ niezależnych kopii $Y_1$. Zatem
$\E(J_t-J_s)=\lambda(t-s)\mu$ i $M$ jest martyngałem.
Z prawa wariancji całkowitej
\[
\begin{aligned}
\Var(J_t-J_s)
&=\E[(N_t-N_s)\Var(Y_1)]
 +\Var((N_t-N_s)\mu)\\
&=\lambda(t-s)(\Var Y_1+\mu^2)
=\lambda(t-s)m_2.
\end{aligned}
\]
Stąd $\E[M_t^2\mid\calF_s]
=M_s^2+\lambda m_2(t-s)$, a więc
$M_t^2-\lambda m_2t$ jest martyngałem.
Trajektoria $M$ ma pomiędzy skokami wariację skończoną,
a jej przyrosty w chwilach zdarzeń są równe
kolejnym wartościom $Y_k$ (oznaczenie zerowe
nie daje niezerowego skoku).
Z lematu~\ref{lem:cadlag-finite-variation-qv}
mamy $[M]_t=\sum_{k\le N_t}Y_k^2$.
Stosując ten sam rachunek warunkowej średniej
do oznaczeń $Y_k^2$ (wystarcza ich całkowalność),
otrzymujemy martyngał $[M]_t-\lambda m_2t$.
Stąd również $M^2-[M]$ jest martyngałem.
Deterministyczny proces $\lambda m_2t$ jest
ciągły, niemalejący i przewidywalny, więc spełnia
definicję przewidywalnego nawiasu $M$.
\end{proof}

\begin{twierdzenie}[Tożsamość kompensacyjna dla oznaczeń]
\label{thm:marked-poisson-compensation}
Niech $S_k$ będą czasami zdarzeń procesu Poissona
o intensywności $\lambda$, a $Y_k$ niezależnymi
oznaczeniami o wspólnym rozkładzie $\nu$ na $\R$,
niezależnymi od procesu Poissona. Filtracja
ujawnia oznaczenie $Y_k$ dopiero w $S_k$.
Dla nieujemnej funkcji $H(t,\omega,y)$
mierzalnej względem iloczynu $\sigma$-ciała
przewidywalnego i $\mathcal B(\R)$ zachodzi
\[
\E\sum_{k:S_k\le T}H(S_k,\omega,Y_k)
=\lambda\E\int_0^T\int_{\R}H(t,\omega,y)
\,\mathrm d\nu(y)\,\mathrm d\lambda_1(t).
\]
\end{twierdzenie}
\begin{proof}
Najpierw bierzemy
$H(t,\omega,y)=\chi_F(\omega)\chi_{(a,b]}(t)
\chi_B(y)$, $F\in\calF_a$, $0\le a<b\le T$,
$B\in\mathcal B(\R)$.
Z własności oznaczonych przyrostów
(rozdział~\ref{ch:poisson-process})
warunkowa średnia liczby przyszłych zdarzeń
z oznaczeniem w $B$ wynosi
$\lambda(b-a)\nu(B)$: warunkowo względem
liczby przyszłych zdarzeń oznaczenia są
niezależnymi kopiami o prawie $\nu$.
Po pomnożeniu przez $\chi_F$ i wzięciu
wartości oczekiwanej otrzymujemy tezę.
Te prostokąty tworzą układ zamknięty na
skończone przecięcia, generujący rozważane
$\sigma$-ciało dla $0<t\le T$.
Obie strony dla funkcji charakterystycznych
są skończonymi miarami (ich całkowita masa
wynosi $\lambda T$); lewa strona jest
mierzalna najpierw dla generatorów,
a następnie przez klasę monotoniczną.
Z jednoznaczności miar równość zachodzi
dla wszystkich zbiorów mierzalnych.
Aproksymacja funkcjami prostymi i zbieżność
monotoniczna dają wynik.
\end{proof}

\begin{zadanie}[Ryzyko portfela ze skokami]
Niech $N$ będzie procesem Poissona o intensywności $2$,
a skoki $Y_k$ przyjmują wartości $+3$ i $-1$,
każdą z prawdopodobieństwem $1/2$.
Wyznaczyć wartość oczekiwaną i wariancję $J_t$,
a także $\E[J_t\mid N_t=3]$.

\emph{Rozwiązanie.} Mamy
$\mu=(3-1)/2=1$, $m_2=(9+1)/2=5$.
Stąd $\E J_t=2t$, $\Var(J_t)=2t\cdot5=10t$.
Przy warunku $N_t=3$ suma zawiera dokładnie
trzy niezależne oznaczenia, zatem
$\E[J_t\mid N_t=3]=3\mu=3$.
To nie jest to samo co $\E J_t=2t$:
warunkowanie ujawnia informację o liczbie skoków.
\end{zadanie}


\begin{zadanie}[Rzeczywista i przewidywalna wariacja --- dwa różne wyniki]
Niech $N$ będzie procesem Poissona o intensywności $2$,
a $M_t=N_t-2t$. Obliczyć prawdopodobieństwo zdarzenia
$\{[M]_1=0\}$, wartości oczekiwane $[M]_1$ oraz
$\langle M\rangle_1$, i wyjaśnić, dlaczego na tym
zdarzeniu oba nawiasy są różne.

\emph{Rozwiązanie.} Ponieważ $[M]_1=N_1$,
$\Prob([M]_1=0)=\Prob(N_1=0)=e^{-2}$.
Natomiast $\langle M\rangle_1=2$ deterministycznie,
także gdy do chwili $1$ nie wystąpił żaden skok.
Oba nawiasy mają jednak tę samą wartość oczekiwaną:
$\E\bigl([M]_1\bigr)=\E N_1=2
=\E\bigl(\langle M\rangle_1\bigr)$.
Natomiast $\E M_1=0$; są to trzy różne wielkości.
Jest to najprostsza ilustracja różnicy między wielkością
zrealizowaną a jej przewidywalną kompensacją.
\end{zadanie}


\section{Skończone sumy przewidywalne dla procesów ze skokami}
\begin{twierdzenie}[Izometria kompensowanej sumy]
\label{thm:compound-poisson-isometry}
Niech $M$ będzie kompensowanym procesem z twierdzenia
\ref{thm:compound-poisson-brackets}, a $0=t_0<\cdots<t_n=T$.
Jeśli $H_i\in L^\infty(\calF_{t_i})$, to
$S=\sum_{i=0}^{n-1}H_i(M_{t_{i+1}}-M_{t_i})$
spełnia
\[
\E S=0,\qquad
\E S^2=\lambda m_2\sum_{i=0}^{n-1}
(t_{i+1}-t_i)\E H_i^2.
\]
\end{twierdzenie}
\begin{proof}
Pierwsza równość i zanik wyrazów mieszanych w
$\E S^2$ wynikają z niezależnych wycentrowanych
przyrostów, dokładnie jak w lemacie o skończonych
sumach przewidywalnych z rozdziału
\ref{ch:quadratic-cont}. Dla każdego $i$
\[
\E[H_i^2(M_{t_{i+1}}-M_{t_i})^2]
=\E H_i^2\,\lambda m_2(t_{i+1}-t_i),
\]
ponieważ kwadrat przyrostu jest niezależny od
$\calF_{t_i}$. Po zsumowaniu dostajemy tezę.
\end{proof}

\begin{twierdzenie}[Izometria dla kompensowanego procesu zliczającego]
\label{thm:counting-compensated-isometry}
Przy założeniach twierdzenia
\ref{thm:counting-compensation} załóżmy dodatkowo,
że $A$ jest ciągły oraz $M=N-A$ jest martyngałem
kwadratowo całkowalnym na $[0,T]$.
Wtedy
\[
[M]_t=N_t,\qquad \langle M\rangle_t=A_t.
\]
Dla $0=t_0<\cdots<t_m=T$ i ograniczonych
$\calF_{t_i}$-mierzalnych $H_i$ zachodzi
\[
\E\left(\sum_i H_i(M_{t_{i+1}}-M_{t_i})\right)^2
=\E\sum_i H_i^2(A_{t_{i+1}}-A_{t_i}).
\]
\end{twierdzenie}
\begin{proof}
Z lematu~\ref{lem:cadlag-finite-variation-qv}
mamy $[M]=N$, gdyż $M$ ma lokalnie skończoną
wariację, a ciągłość $A$ daje $\Delta M=\Delta N$.
Udowodnimy bez całki It\^o, że $M^2-N$ jest
lokalnym martyngałem. Dla trajektorii o skończonej
wariacji zachodzi tożsamość
\[
M_t^2=2\int_{(0,t]}M_{s-}\,\mathrm dM_s+
\sum_{0<s\le t}(\Delta M_s)^2 .
\]
Wynika ona z teleskopowania kwadratów na
podziałach: sumy kwadratów zbiegają do $[M]_t$,
a sumy z lewym końcem przyrostu zbiegają
do całki względem miary znakowanej $\mathrm dM$.
Istotnie, na $(t_i,t_{i+1}]$ funkcja
podcałkowa jest równa $M_{t_i}$; dla każdego
$s>0$ zbiega ona do $M_{s-}$ przy zanikającej
średnicy, gdyż $t_i<s$. Zbieżność zmajoryzowana
względem skończonej miary wariacji daje tezę,
także dla podziału kończącego się w $t$.

Niech $\rho_k=\inf\{t:N_t+A_t\ge k\}\wedge k$.
Proces $N+A$ jest niemalejący i adaptowany,
jego skoki mają rozmiar jeden, zatem
$(N+A)_{t\wedge\rho_k}\le k+1$.
Czasy $\rho_k\uparrow\infty$ p.n.
Zatrzymany $M^{\rho_k}$ jest ograniczonym
martyngałem i ma wariację nie większą niż $k+1$.
Dla $0\le s<t$ i $F\in\calF_s$ stosujemy
lemat~\ref{lem:predictable-fv-mean-zero}
do $D=M^{\rho_k}$ i ograniczonego
procesu przewidywalnego
$\chi_F\chi_{(s,t]}D_-$.
Otrzymujemy, że trajektoryjna całka
$\int_{(0,t]}D_{u-}\,\mathrm dD_u$
jest martyngałem (adaptacja wynika też
z granic skończonych sum).
Tożsamość kwadratowa daje martyngał
$(M^{\rho_k})^2-N^{\rho_k}$.
Zatem $M^2-N$ jest lokalnym martyngałem,
a dodając $N-A$ otrzymujemy lokalny
martyngał $M^2-A$.

Na $[0,T]$ rodzina wartości $M^2-A$ w
ograniczonych czasach zatrzymania jest
jednostajnie całkowalna, ponieważ jej moduły
są ograniczone przez
$\sup_{s\le T}|M_s|^2+A_T\in L^1$
(nierówność Dooba, twierdzenie
\ref{thm:doob-cont}). Jest więc prawdziwym
martyngałem na $[0,T]$ na mocy kryterium
lokalnego martyngału z rozdziału
\ref{ch:mart-cont}.
Proces $A$ spełnia definicję nawiasu
przewidywalnego, a jednoznaczność daje
wniosek~\ref{cor:predictable-compensator-unique}.
Rozwijając kwadrat przyrostu i warunkując
względem jego początku otrzymujemy
$\E[(M_v-M_u)^2\mid\calF_u]
=\E[A_v-A_u\mid\calF_u]$.
Ortogonalność różnych przyrostów
z przewidywalnymi współczynnikami
oraz własność wieży dają żądaną izometrię.
Nie potrzeba niezależności przyrostów $N$.
\end{proof}

\section{Wykładnik procesu Poissona i oszacowania ogonów}
\begin{twierdzenie}[Martyngał wykładniczy procesu Poissona]
\label{thm:poisson-exponential-mart}
Niech $N$ będzie procesem Poissona o intensywności
$\lambda>0$ w uzupełnionej filtracji naturalnej.
Dla każdej $\theta\in\R$ proces
\[
Z_t(\theta)=\exp\{\theta N_t-\lambda t(e^\theta-1)\}
\]
jest dodatnim martyngałem o wartości oczekiwanej $1$.
Dla $T>0$ oraz $a>\lambda T$ zachodzi
\[
\Prob(N_T\ge a)\le
\exp\left(-a\log\frac{a}{\lambda T}+a-\lambda T\right).
\]
\end{twierdzenie}
\begin{proof}
Dla $s<t$ niezależność przyrostu
$U=N_t-N_s\sim\mathrm{Poisson}(\lambda(t-s))$
i jego funkcja tworząca
$\E e^{\theta U}=\exp(\lambda(t-s)(e^\theta-1))$
dają $\E[Z_t(\theta)\mid\calF_s]=Z_s(\theta)$.
Dla $\theta>0$ nierówność Markowa zastosowana do
$e^{\theta N_T}$ prowadzi do
$\Prob(N_T\ge a)\le
\exp(-\theta a+\lambda T(e^\theta-1))$.
Minimalizujemy wykładnik względem $\theta>0$:
pochodna wynosi $-a+\lambda Te^\theta$, więc
minimum przypada przy $\theta=\log(a/(\lambda T))$.
Podstawienie daje tezę. Jest to ścisły przykład tego,
jak martyngał wykładniczy prowadzi do nierówności
koncentracyjnej, a nie tylko do wzorów na momenty.
\end{proof}

\begin{twierdzenie}[Martyngał wykładniczy przy losowym kompensatorze]
\label{thm:counting-exponential-bound}
Niech $N,A$ spełniają założenia twierdzenia
\ref{thm:counting-compensation}, a $A$ będzie
ciągły. Dla $\theta\in\R$ proces
\[
Z_t=\exp\{\theta N_t-(e^\theta-1)A_t\}
\]
jest dodatnim lokalnym martyngałem, zatem
nadmartyngałem o $\E Z_t\le1$.
Dla $M=N-A$, $x,v>0$ i $T>0$ zachodzi
\[
\Prob\left(\sup_{t\le T\wedge\eta_v}M_t\ge x\right)
\le \exp\{-v h(x/v)\},
\qquad \eta_v=\inf\{t:A_t>v\},
\qquad h(r)=(1+r)\log(1+r)-r .
\]
\end{twierdzenie}
\begin{proof}
Trajektorie $Z$ mają lokalnie skończoną
wariację, ponieważ $N$ i $A$ mają tę własność.
Na odcinkach bez zdarzeń $N$ reguła łańcuchowa
dla ciągłej funkcji niemalejącej $A$ daje
$\mathrm dZ=-(e^\theta-1)Z\,\mathrm dA$.
Reguła ta nie wymaga różniczkowalności $A$:
teleskopujemy funkcję wykładniczą po
podziałach i szacujemy błąd liniowego
przybliżenia przez stałą razy
$\max_i|\Delta_i A|\sum_i|\Delta_i A|$,
który dąży do zera przez ciągłość $A$.
W chwili skoku $\Delta N=1$ mamy
$\Delta Z=(e^\theta-1)Z_-$.
Miara $\mathrm dA$ nie ma atomów,
więc $Z$ i $Z_-$ są równe względem tej miary
prawie wszędzie. Sumując części ciągłe
i skoki otrzymujemy trajektoryjnie
\[
Z_t=1+(e^\theta-1)\int_{(0,t]}Z_{s-}\,\mathrm d(N-A)_s .
\]
Nie używamy tutaj całki It\^o.
Lokalizujemy czasami
$\rho_k=k\wedge\inf\{t:N_t+A_t\ge k\}$.
Procesy $(N-A)^{\rho_k}$ są martyngałami
o wariacji ograniczonej przez $k+1$.
Wartości $Z_{s-}$ przed zatrzymaniem
są ograniczone przez stałą zależną od $k,\theta$.
Zatrzymana tożsamość jest zatem całką
względem $(N-A)^{\rho_k}$ z ograniczonym
przewidywalnym współczynnikiem
$(Z^{\rho_k})_-$.
Lemat~\ref{lem:predictable-fv-mean-zero},
stosowany także po pomnożeniu współczynnika
przez $\chi_F\chi_{(u,t]}$, $F\in\calF_u$,
daje własność martyngałową $Z^{\rho_k}$.
Ponieważ $\rho_k\uparrow\infty$ p.n.,
$Z$ jest lokalnym martyngałem.
Nieujemny lokalny martyngał jest nadmartyngałem
(twierdzenie~\ref{thm:positive-local-supermart}).

Dla $\theta>0$ zapisujemy
$Z_t=\exp\{\theta M_t-(e^\theta-1-\theta)A_t\}$.
Czas $\eta_v$ jest czasem zatrzymania dzięki
adaptacji, ciągłości $A$ i prawostronnej
ciągłości filtracji, jak w dowodzie
twierdzenia~\ref{thm:dambis-dubins-schwarz}.
Dla $t\le T\wedge\eta_v$ mamy $A_t\le v$.
Biorąc supremum, na zdarzeniu z tezy mamy
$\sup_{t\le T} Z_t\ge
e^{\theta x-(e^\theta-1-\theta)v}$
także wtedy, gdy supremum $M$ nie jest osiągnięte.
Nierówność Ville'a
(twierdzenie~\ref{thm:ville-cont}) daje
oszacowanie
\[
\exp\{-\theta x+(e^\theta-1-\theta)v\}.
\]
Wybór $\theta=\log(1+x/v)>0$ minimalizuje
wykładnik i daje $-vh(x/v)$.
Supremum procesu zatrzymanego $M^{\eta_v}$
jest mierzalne, ponieważ można je obliczyć
na przeliczalnym gęstym zbiorze chwil
wraz z końcem $T$, dzięki prawostronnej
ciągłości.
\end{proof}

\begin{uwaga}[Znaczenie ograniczenia nawiasu]
W przypadku Poissona $A_t=\lambda t$,
więc przy $v=\lambda T$ ostatnia nierówność
dotyczy maksimum skompensowanego procesu
na całym $[0,T]$, a nie tylko jego końcowej
wartości. Dla losowego $A$ oszacowanie
kontroluje przekroczenia przed wyczerpaniem
budżetu kompensatora $v$.
Nie zakładamy, że dodatni lokalny martyngał
wykładniczy jest zawsze prawdziwym martyngałem.
\end{uwaga}

\section{Martyngał ciągły i skokowy w jednym modelu}

\begin{twierdzenie}[Dodanie skończonej wariacji do procesu ciągłego]
\label{thm:finite-variation-jumps-bracket}
Niech $C$ będzie procesem o ciągłych trajektoriach, dla którego
wariacja kwadratowa $[C]$ istnieje w sensie zbieżności ucp.
Niech $A$ będzie procesem c\`adl\`ag o skończonej wariacji
całkowitej na każdym przedziale $[0,T]$, prawie na pewno.
Liczba skoków może być nieskończona. Wtedy $X=C+A$ ma wariację
kwadratową i zachodzi, z dokładnością do nierozróżnialności,
\[
[X]_t=[C]_t+\sum_{0<s\le t}(\Delta A_s)^2,
\qquad t\ge0.
\]
Nie zakładamy niezależności $C$ i $A$.
\end{twierdzenie}
\begin{proof}
Ustalmy $T>0$ i podział $\pi$. Dla każdego $t\le T$
rozwinięcie kwadratu daje
\[
V_t^\pi(C+A)=V_t^\pi(C)+V_t^\pi(A)
 +2\sum_i\Delta_i^tC\,\Delta_i^tA,
\]
gdzie $\Delta_i^tC=C_{t\wedge t_{i+1}}-C_{t\wedge t_i}$
i analogicznie definiujemy $\Delta_i^tA$. Dla każdej
ustalonej trajektorii $A$ o skończonej wariacji całkowitej
\[
\sup_{t\le T}\left|\sum_i\Delta_i^tC\,\Delta_i^tA\right|
\le \omega_C(|\pi|)\operatorname{Var}_{[0,T]}(A),
\qquad
\omega_C(\delta):=
\sup_{\substack{u,v\in[0,T]\\|u-v|\le\delta}}|C_u-C_v|.
\]
Ponieważ $C$ ma ciągłe trajektorie, $\omega_C(\delta)\to0$
przy $\delta\downarrow0$. Iloczyn po prawej stronie zbiega
więc do zera prawie na pewno, jednostajnie po $t$.
Z lematu~\ref{lem:cadlag-finite-variation-qv}
$V^\pi(A)$ zbiega jednostajnie
prawie na pewno do sumy kwadratów skoków, natomiast
$V^\pi(C)\to[C]$ w ucp z założenia.
Sumując te trzy zbieżności, otrzymujemy tezę.
Ważne jest, że ciągłość $C$ wyklucza wspólny skok;
w przeciwnym razie pojawiłyby się również iloczyny
skoków $2\Delta C_s\Delta A_s$.
\end{proof}

\begin{stwierdzenie}[Brak kowariacji niezależnych składników]
Niech $W$ będzie procesem Wienera, a $N$ niezależnym
od niego procesem Poissona o intensywności $\lambda>0$.
Dla $M_t=N_t-\lambda t$ i dowolnych deterministycznych
ciągów podziałów $\pi$ przedziału $[0,T]$
o średnicach dążących do zera zachodzi
\[
\sum_i(W_{t_{i+1}}-W_{t_i})
(M_{t_{i+1}}-M_{t_i})\longrightarrow0
\quad\text{w }L^2.
\]
W szczególności wariacja kwadratowa sumy
$X_t=W_t+M_t$ ma postać $[X]_t=t+N_t$.
Sama tożsamość wariacji kwadratowej nie wymaga
niezależności $W$ i $N$: wynika także z ciągłości $W$
i skończonej wariacji $M$ na przedziałach zwartych.
\end{stwierdzenie}
\begin{proof}
Niezależne przyrosty obu procesów oraz niezależność
całych procesów dają zanik mieszanych momentów
między różnymi komórkami siatki. Dla pojedynczej komórki
długości $h_i$ mamy
$\E(\Delta_iW)^2=h_i$ oraz
$\E(\Delta_iM)^2=\lambda h_i$.
Zatem
\[
\E\left|\sum_i\Delta_iW\Delta_iM\right|^2
=\lambda\sum_i h_i^2\le\lambda T|\pi|\longrightarrow0.
\]
Z elementarnej tożsamości
$\sum(\Delta X)^2=\sum(\Delta W)^2+
2\sum\Delta W\Delta M+\sum(\Delta M)^2$
oraz wyników wcześniejszych rozdziałów wynika
$[X]_t=t+N_t$. Dla pełnej ścisłości zbieżność ucp
wynika także bezpośrednio z twierdzenia
\ref{thm:finite-variation-jumps-bracket}: trajektoria
$M=N-\lambda t$ ma skończoną wariację i skończenie
wiele skoków na $[0,T]$, zaś $W$ jest ciągły.
Oszacowanie $L^2$ powyżej daje ponadto ilościową kontrolę
mieszanych sum w końcowej chwili $T$.
To rozdziela dwa fakty: zanik kowariacji jest tutaj
własnością regularności trajektorii, natomiast przedstawione oszacowanie $L^2$
korzysta z niezależności.
Jeżeli $W$ i $M$ są martyngałami względem tej samej
filtracji i są kwadratowo całkowalne, to
$WM$ również jest martyngałem w przypadku niezależnych
procesów w ich wspólnej uzupełnionej filtracji
naturalnej: warunkowanie względem wspólnej przeszłości
redukuje się do warunkowania w obu niezależnych
filtracjach (najpierw na prostokątach zdarzeń).
Wtedy $X^2-(1+\lambda)t$ jest martyngałem przez
rozwinięcie kwadratu i twierdzenie
\ref{thm:poisson-brackets}, a zatem
$\langle X\rangle_t=(1+\lambda)t$.
\end{proof}

\section{Co przenosimy do tomu III?}
Tom II został zbudowany wokół pytania \emph{jak zachowują
się procesy losowe}. Badaliśmy ich rozkłady, trajektorie,
filtracje, czasy zatrzymania, własność Markowa, strukturę
odnowy i rozgałęziania oraz martyngały i wariację
kwadratową. Kolejny tom zmieni pytanie na
\emph{jak przeprowadzać analizę na takich trajektoriach}.

\begin{uwaga}[Granica pomiędzy procesami a analizą stochastyczną]
W tomie III od podstaw skonstruujemy całkę względem
procesu Wienera dla procesów przewidywalnych, wzór
Itô, całkę Stratonovicha, lokalne czasy i formułę
Tanaki, a następnie całkowanie względem ogólnych
martyngałów i półmartyngałów oraz równania stochastyczne.
Później powrócimy do miary Wienera, całek funkcyjnych,
wzoru Feynmana--Kaca i relacji do formalnych całek
Feynmana po trajektoriach. Skończone sumy przewidywalne
z rozdziałów 13--14 są zapowiedzią tej konstrukcji,
ale same nie stanowią jeszcze ogólnej teorii całki Itô.
\end{uwaga}

\begin{zadanie}[Dwie różne miary fluktuacji]
Niech $W$ będzie procesem Wienera, $N$ niezależnym
procesem Poissona o intensywności $3$ oraz
$X_t=W_t+(N_t-3t)$. Obliczyć $[X]_t$ i średnią
$\E[X_t^2]$.

\emph{Rozwiązanie.} Kwadratowa zmienność ciągłej części
wynosi $t$, a skoki części poissonowskiej mają rozmiar
jeden, więc $[X]_t=t+N_t$. Z niezależności
$W_t$ oraz $N_t-3t$ mamy
$\E X_t=0$ i
$\E X_t^2=\Var(W_t)+\Var(N_t)=t+3t=4t$.
Także $\E[X]_t=t+\E N_t=4t$.
Zauważamy jednak, że $[X]_t$ jest losowy, podczas
gdy $\langle X\rangle_t=4t$ jest deterministyczny.
Tempo wzrostu przewidywalnego nawiasu wynosi $4$.
\end{zadanie}


\end{document}
